\documentclass[11pt,DIV=13,parskip=false]{scrartcl}
\usepackage{testimony}
\title{Sola gratia — Ephesians 2:8–9 and Romans 9:16}
\author{\credit{TestimonyArguments.SolaGratia}}
\date{Generated by \texttt{lake exe argtex}}
\begin{document}
\maketitle
\noindent This is one argument of the Testimony library, printed on its own. The library prints every argument together in \texttt{arguments.pdf}. This one is set as its Lean source sets it: the module docstrings as prose, in the order they are written, and under each one the package, line of reason, countermodel or result it introduces. Atoms are numbered propositional variables, numbered once per argument, so a variable means the same claim in every position within it; the legend gives each one's claim, its classification, the tradition that holds it and the source that carries it.

\noindent Nothing here asserts that the premises are true --- only that the conclusion does or does not follow from them.

\tableofcontents
\clearpage
\section*{Sola gratia — Ephesians 2:8–9 and Romans 9:16}
\addcontentsline{toc}{section}{Sola gratia — Ephesians 2:8–9 and Romans 9:16}

\subsection*{Sola gratia: what everyone grants, and what divides them}
\addcontentsline{toc}{subsection}{Sola gratia: what everyone grants, and what divides them}
“For by grace you have been saved through faith. And this is not your own doing; it is the gift of God, not a result of works, so that no one may boast” (Ephesians 2:8–9). Every party to this dispute reads the verse and agrees with it. So \emph{sola gratia}, “by grace alone”, has to be stated more sharply than the slogan if it is to separate anyone.

\subsubsection*{Everyone says “by grace”}
The Council of Trent, answering the Reformers in 1547, taught that “in adults, the beginning of the said Justification is to be derived from the prevenient grace of God”, the grace that goes before, by which they are called “without any merits existing on their parts”. It quoted Paul: “if it be a grace, it is not now by works, otherwise … grace is no more grace” (Romans 11:6). The Remonstrants, the followers of Arminius who petitioned the States of Holland in 1610, called grace “the beginning, continuance, and accomplishment of all good”. Neither is Pelagian, after Pelagius, who was taken to teach that a person can turn to God by their own power: neither thinks a person can so much as begin.

So a conclusion stated as “salvation is by grace” would follow from all three positions and separate none of them. When the Synod of Dort answered the Remonstrants in 1619, it meant something sharper. The grace that saves brings about the consent it asks for, rather than waiting on a will that could refuse it. And the new life is, in the Canons’ words, “a making alive, which God works in us without our aid”; the Westminster Confession says the sinner “is altogether passive therein, until, being quickened and renewed by the Holy Spirit, he is thereby enabled to answer this call”. The older name for this is \emph{monergism}, from the Greek for “one at work”.

The conclusion here has those three parts. Salvation is by grace; the grace that saves is \emph{effectual of itself}, that is, it brings about the consent it asks for; and the will is passive in being made alive. Each rival holds the first part. Trent’s canon, as most read it, denies both the others, and the Thomists read it otherwise, in the chapter on Aquinas; the Remonstrants deny the second.

\begin{itemize}
  \item \textbf{Trent}, on the Molinist reading, denies the third and the second. Its fourth canon on justification condemns anyone who says that the free will, “moved and excited by God”, does not cooperate “by assenting to God exciting and calling”, but “as something inanimate, … is merely passive”. The same canon condemns anyone who says the will “cannot refuse its consent, if it would”.
  \item \textbf{The Remonstrants} deny the second. Their fourth article says that grace “is not irresistible, inasmuch as it is written concerning many, that they have resisted the Holy Ghost”, citing Stephen’s charge to the council: “you always resist the Holy Spirit” (Acts 7:51).
\end{itemize}
\subsubsection*{Two routes to the conclusion}
The Reformed case has two routes, called \emph{strands} below, each from a different text. Each begins from a reading of its text and takes one or two steps from that reading to the conclusion.

\textbf{Ephesians 2.} Paul has just said that his readers “were dead in the trespasses and sins”, and that God “made us alive together with Christ” (2:1, 5). Then: “by grace you have been saved through faith. And this is not your own doing”. The Greek word for “this” (τοῦτο) is neuter, while “faith” (πίστις) and “grace” (χάρις) are both feminine, so the pronoun does not point to either word on its own. Hoehner and Robertson read it as pointing to the whole of being saved by grace through faith. The strand then takes two steps: since faith is part of that whole, faith is God’s gift; and since it is given to people who were dead until God made them alive, the turning is his work alone.

\textbf{Romans 9.} “So then it depends not on human will or exertion, but on God, who has mercy” (Romans 9:16). Read as being about God’s mercy in saving, and not only about his choice of a people for a part in history, the verse says that receiving mercy does not hang on anyone’s willing. The Romans strand does not use the Ephesians grammar at all.

\subsubsection*{What follows from each position}
Lean checks, for each position written as a list of premises, whether the conclusion follows. It does not check whether the premises are true.

The Reformed premises reach the conclusion, and neither rival’s do, though both reach its first part. A reading is called \emph{load-bearing} when removing it, with everything else kept, stops the conclusion following. Neither reading is load-bearing alone, because each strand reaches the conclusion without the other. Removed together, the conclusion no longer follows: the Reformed case needs at least one of the two readings, and either will do.

Readers who grant both readings still deny what the Reformed draw from them. Robertson reads “this” exactly as Hoehner does, and still takes the whole to be “conditioned on faith on our part”: grace is “God’s part, ‘faith’ ours”. So the step from the pronoun to faith as a gift is rated \emph{disputed}, which means that a reader the library cites grants its grounds and denies its conclusion. Calvin, himself Reformed and elsewhere clear that faith is God’s gift, still says this verse does not teach it: Paul’s “meaning is, not that faith is the gift of God, but that salvation is given to us by God”.

The next step is disputed too. Arminius acknowledged faith to be God’s “benefit and gift”, and ascribed to grace “the commencement, the continuance and the consummation of all good”. He could grant everything the Ephesians strand rests on, up to its last step. He refused that step, believing “that many persons resist the Holy Spirit and reject the grace that is offered”. Dort’s answer is its article on faith, which it calls the gift of God “not on account of its being offered by God to man, to be accepted or rejected at his pleasure, but because it is in reality conferred”. The two mean different things by \emph{gift}. The encoding writes both as one claim, so the difference shows as Arminius refusing the step: grant the texts, the reading of “this” and faith as God’s gift, leave out that step, and sola gratia does not follow (\texttt{gift\_\allowbreak{}of\_\allowbreak{}faith\_\allowbreak{}does\_\allowbreak{}not\_\allowbreak{}settle\_\allowbreak{}resistibility}).

Romans goes the same way. Arminius read Romans 9 as being about salvation, “not to him that worketh, but to him that believeth”. Abasciano, who reads its election as of a people rather than of individuals, still grants that “Paul’s argument in Romans 9 surely concerns the salvation of Israel”. The reading is widely granted, and no reader who denies it is cited here with a verified location. It is the step from it to God’s work alone that Arminius denies.

Each half of Trent’s canon meets a reader who grants its ground and denies its conclusion, and both times the reader is Aquinas. That the will “cannot refuse its consent, if it would” is condemned, which seems to deny that grace is effectual of itself. But Aquinas held both that if God moves the will, “it is incompatible with this supposition, that the will be not moved thereto. But it is not impossible simply”, and that the man God moves to grace “will infallibly attain to it”. And the canon’s “cooperates” speaks of the will “disposing and preparing itself” for justification, while Aquinas and Dort both place the will’s own activity after a first moment in which, in Aquinas’s words, it “is moved and does not move”; Dort says the will once renewed “becomes itself active”. So both of Trent’s steps are rated disputed. Not every Catholic reads the canon so: the Molinists, after Luis de Molina, take its “can refuse” to rule out a grace effectual of itself. Paul V’s decree of 1607 “allowed each party to defend its own doctrine, enjoined each from censoring or condemning the opposite opinion” (Astrain, “Congregatio de Auxiliis”). The rating stands on the Thomist reading, because one reader who grants the ground and denies the conclusion is enough.

\subsubsection*{Aquinas and the Catholic schools}
The Thomists are a position of their own, in the chapter on Aquinas. They reach the second part of sola gratia, grace effectual of itself (\texttt{thomists\_\allowbreak{}reach\_\allowbreak{}effectual\_\allowbreak{}grace}), at a step the Molinists dispute. On the third part, the will’s passivity, they say the will “does not remain passive” under operating grace, and Garrigou-Lagrange lists Calvin’s “merely passive” among the errors. But he places the will’s consent after grace “in the genus or order of formal, efficient, and final cause”, so whether that contradicts Dort depends on what “passive” means. Their own premises leave it open (\texttt{thomists\_\allowbreak{}leave\_\allowbreak{}passivity\_\allowbreak{}open}), and the dilemma over the phrase answers both readings (\texttt{what\allowbreak{}Not\allowbreak{}Merely\allowbreak{}Passive\allowbreak{}Means}). Weighed with the others, neither Dort nor the Thomists defeats the other, because the Thomist position weighed takes no side on the will’s passivity; faith alone and perseverance, on which they part, Dort’s position here does not carry.

\subsubsection*{Freedom to refuse, and two more positions}
The strongest objection comes from the other side, and it is an objection from Scripture: the Spirit is resisted (Acts 7:51), grieved (Ephesians 4:30) and quenched (1 Thessalonians 5:19), and one sin against him is never forgiven (Matthew 12:31–32). Every party grants that the Spirit’s work is resisted. What divides them is whether the grace that saves can be \emph{finally} refused, and each route from the texts to that claim needs a reading that a cited reader denies. The chapter on freedom to refuse sets them out, with two positions between Dort and the Remonstrants: the Lutherans of the Formula of Concord, for whom the will is passive in conversion and grace resistible all the same, and Wesley, for whom preventing grace restores to everyone an ability that can be used or refused.

The Lutherans are why the step from God’s work alone to grace that cannot be finally refused is rated \emph{disputed}. They hold conversion to be “altogether an operation, gift, and present, and work of the Holy Ghost alone”, and they deny that grace cannot be refused: “God does not force man to become godly”, and the perverse will “resists the Holy Ghost, who wishes to be efficacious”.

The third way, as it was put to this argument, is that God gives the ability to believe as Adam and Eve were made able to stand or to fall, and that the offer is received by faith alone. As stated, it leaves open whether grace can be finally refused (\texttt{third\_\allowbreak{}way\_\allowbreak{}leaves\_\allowbreak{}resistibility\_\allowbreak{}open}). Augustine’s distinction between the help “without which” a thing does not come to pass, which Adam had, and the help “by which” it does, which the saints have, is the dilemma it meets. Read the first way, it is the reading Wesley gives, and it cannot be held together with Dort’s position; read the second, it is Dort’s, and it cannot be held together with Trent’s canon 4, nor, by what the claims mean, with the Remonstrants. Both breaks at a premise come at steps rated \emph{disputed}, so neither shows the other side right (\texttt{what\allowbreak{}The\allowbreak{}Ability\allowbreak{}Is}).

Suppose instead that the Remonstrants are right, and the grace that saves can be finally refused. Set beside the Scriptures as the Reformed read them, with the claims read as they mean, the supposition has no model at all, and the same holds of Ephesians 2 or Romans 9:16 alone (\texttt{final\_\allowbreak{}refusal\_\allowbreak{}cannot\_\allowbreak{}be\_\allowbreak{}held\_\allowbreak{}with\_\allowbreak{}the\_\allowbreak{}scriptures}). That is sola gratia shown from its texts by the contrary, and it is exactly as strong as the readings and the disputed steps it rests on.

Weighed as one dispute, only what every party holds is settled before any choice: salvation is by grace (\texttt{only\_\allowbreak{}grace\_\allowbreak{}prevails}). Every other position, Dort’s included, can be defended and is not forced, and the third way can be defended with Dort or with Wesley.

\subsubsection*{The perseverance of the saints}
If the grace that saves brings about the consent it seeks, can the one it saves lose it? Westminster says that the “effectually called”, those called in a way that brings them to faith, “can neither totally nor finally fall away from the state of grace” (XVII.1). Augustine, Aquinas and the Formula of Concord agree that those God has predestined to life will persevere to the end. They part from Westminster over who those people are. Westminster holds that everyone God brings to true faith is among them. Augustine and Aquinas hold that God gives some people true faith, hope and love without predestining them to persevere, so that in the end they fall away. The Lutherans hold that some in whom God has begun the good work turn away by their own fault. The question is who is predestined to persevere, not whether God chose them because he foresaw their faith, which the chapter does not weigh.

So the chapter on perseverance writes the Reformed conclusion as two claims and the claim that joins them, and finds that Augustine parts from Westminster over 1 John 2:19 (“they went out from us, but they were not of us”): whether those who left never had true faith, as Calvin and Owen read it, or had it and were not elect, as Augustine does. It also weighs Romans 8:29–30; the texts Wesley reads of believers who fell (Hebrews 6 and 10, 2 Peter 2); falling for a time, read from David’s prayer in Psalm 51:11; whether a believer can know that he will persevere; and the view from eternity, Boethius’s and Aquinas’s, that God sees all of time at once. That view, as encoded, says nothing on whether the truly justified can finally fall away, so it can be added to Westminster’s position or Wesley’s without contradiction; with Owen’s reading of 1 John 2:19 it denies that they can. Weighed as a dispute of its own, Westminster is in a tie with Augustine and Aquinas, Trent, the Lutherans and Wesley, and only what the Remonstrants of 1610 granted is settled: no power from outside snatches believers out of Christ’s hand (\texttt{only\_\allowbreak{}the\_\allowbreak{}remonstrance\_\allowbreak{}prevails}). Westminster and each of the four can be defended, and none must be accepted.

\subsubsection*{What is not encoded}
The Remonstrants’ denial of the will’s passivity, and Wesley’s view of it, are not encoded: each position here concludes only what its cited sources say. Nor is the cooperation of the renewed will, which Dort, the Formula of Concord and Trent all grant in some sense, and which a dilemma over “cooperate” would split.

Molina’s own \emph{Concordia} is not encoded, because no public English text could be checked against it; the Molinist readings are cited from Pohle’s report of them.

John 6:44 (“No one can come to me unless the Father who sent me draws him”) is the third text most often cited for the Reformed view, and is not encoded either. Nor is election, the subject of the first of the Remonstrants’ articles, which is a dispute of its own.

\subsubsection*{Where things are}
\begin{longtable}{@{}>{\raggedright\arraybackslash}p{0.45\textwidth} >{\raggedright\arraybackslash}p{0.45\textwidth}@{}}
\toprule
\textbf{File} & \textbf{Contents} \\
\midrule
\endhead
\texttt{Atoms.lean} & the claims \\
\texttt{Sources.lean} & who holds each claim, and the rating of each step \\
\texttt{Lines.lean} & the steps and the lines of reason \\
\texttt{Packages.lean} & the positions, and the variants the results ask about \\
\texttt{Results.lean} & what follows, what does not, and which readings carry it \\
\texttt{Third\allowbreak{}Way.lean} & freedom to refuse, the Lutherans, Wesley, and the third way’s dilemma \\
\texttt{Thomist.lean} & Aquinas and the Catholic schools, and what “not merely passive” means \\
\texttt{Dispute.lean} & the positions weighed, and the verdicts \\
\texttt{Perseverance.lean} & the perseverance of the saints, weighed as a dispute of its own \\
\bottomrule
\end{longtable}
\subsection*{Two routes: Ephesians 2 and Romans 9}
\addcontentsline{toc}{subsection}{Two routes: Ephesians 2 and Romans 9}
Two routes reach the conclusion, one through Ephesians 2 and one through Romans 9, and each has a reading of its text that is in dispute. A third line, from Romans 11:6, delivers only the first part of the conclusion, that salvation is by grace, and every party walks it: Trent quotes the verse in the eighth chapter of its decree on justification.

Trent and the Remonstrants each deny part of what the Reformed lines deliver, so each has lines of its own here. Trent’s fourth canon on justification denies that the will is merely passive, and in the same sentence condemns anyone who says the will cannot refuse its consent: so Trent’s lines deny both the second and third parts of the conclusion. The Remonstrants deny the second part, that grace cannot be resisted, from Stephen’s speech in Acts 7.

The texts of resistance are common ground. Stephen’s hearers “always resist the Holy Spirit” (Acts 7:51), and Paul tells believers not to grieve the Spirit (Ephesians 4:30) or quench him (1 Thessalonians 5:19); every party grants that the Spirit’s work meets resistance. What is disputed is the step from there to the claim that the grace that saves can be \emph{finally} refused. The word “resist” can mean resisting for a time, as the Reformed say even the elect do before grace overcomes them, or refusing to the end. A second route to the final refusal runs through the sin that is never forgiven (Matthew 12:31–32), read with Augustine as impenitence persisted in until death.

The closing step is two steps. From God’s work alone to the will’s passivity, the Lutherans agree with Dort. From God’s work alone to grace that cannot be finally refused, they do not: the Formula of Concord holds conversion to be the Spirit’s work alone and grace resistible all the same.

\subsubsection*{The positions, and the variants the results ask about}
The rivals come first, because they were written first. Each is a position as its own documents state it: Trent from the sixth session’s decree and canons on justification, the Remonstrants from their five articles of 1610 and from Arminius. Neither is Pelagian. Both say, in so many words, that salvation is by grace, and both packages carry the line from Romans 11:6 that delivers it.

Two more positions follow, each a way of holding grace and freedom together that is neither Dort’s nor the Remonstrants’. The Lutherans of the Formula of Concord hold the will “purely passive” in conversion and grace resistible all the same. Wesley holds that God restores to everyone a measure of the freedom lost in Adam, so that each can receive the offer by faith or refuse it. And the third way as it was put to this argument, that God gives the ability to believe as Adam was made able to stand or fall, and faith alone receives, is stated without saying which kind of help that ability is, because that is what a dilemma below asks.

The variants after the Reformed package each change one thing, named, so that a result about a variant is a result about that one change.

\subsubsection*{What does and does not follow}
The results come in three groups. The first is what each position reaches: the Reformed premises reach sola gratia, neither rival’s do, and both rivals reach its first part. The second is which reading of the texts the Reformed case depends on: neither alone, and the two together. The third is where Arminius parts from the Ephesians strand, which is after the grammar of “this” and after faith is granted to be God’s gift.

\subsubsection*{Satisfiability}
A list of premises that contradict each other would “prove” anything, so each position that reaches its conclusion is shown to be holdable without contradiction, by Dort’s reading. These are claims about the health of the encoding, not about the dispute, and are not headline results.

\subsubsection*{How this is checked}
\paragraph*{The conclusion}
\leavevmode\par\nobreak
\par\medskip\noindent\textbf{\texttt{sola\allowbreak{}Gratia}}\par\nobreak\noindent
\textbf{Sola gratia, stated sharply enough that the rivals divide from it.} Salvation is by grace and not earned; the grace that saves brings about the consent it seeks; and in being made alive the will contributes nothing. The three parts are separate claims so that a rival can be shown holding the first and denying one or both of the others: Trent’s premises deny both, the Remonstrants’ the second.

\begin{derivation}
  \item[] \(\pvar{TestimonyArguments.SolaGratia}{29} \land \pvar{TestimonyArguments.SolaGratia}{30} \land \pvar{TestimonyArguments.SolaGratia}{31}\)
\end{derivation}
\paragraph*{The steps}
\leavevmode\par\nobreak
\par\medskip\noindent\textbf{\texttt{grace\allowbreak{}Step}}\par\nobreak\noindent
Romans 11:6: if salvation is by grace it is not by works, so it is by grace and not earned.

\begin{derivation}
  \item[] \(\pvar{TestimonyArguments.SolaGratia}{4} \rightarrow \pvar{TestimonyArguments.SolaGratia}{29}\)\gloss{\emph{would give:} Salvation is by grace and not earned: grace goes before every saving act}
\end{derivation}
\par\medskip\noindent\textbf{\texttt{ephesians\allowbreak{}Step}}\par\nobreak\noindent
Ephesians 2:8–9, with “this” read as pointing to the whole of being saved by grace through faith: so faith, part of that whole, is God’s gift. Rated \texttt{disputed}: Robertson reads the pronoun the same way and calls the whole “conditioned on faith on our part”, and Calvin, on the next verse, writes “not that faith is the gift of God, but that salvation is given to us by God”.

\begin{derivation}
  \item[] \((\pvar{TestimonyArguments.SolaGratia}{2} \land \pvar{TestimonyArguments.SolaGratia}{22}) \rightarrow \pvar{TestimonyArguments.SolaGratia}{24}\)\gloss{\emph{would give:} Faith itself is God’s gift}
\end{derivation}
\par\medskip\noindent\textbf{\texttt{gift\allowbreak{}Step}}\par\nobreak\noindent
\textbf{The step Arminius refuses.} If faith is God’s gift, and it is given to those who were dead until God made them alive (Ephesians 2:1–5), then conversion is God’s work alone. Rated \texttt{disputed}: Arminius grants both grounds and denies that the gift cannot be refused. Dort’s answer is its article on faith: God’s gift “not on account of its being offered by God to man, to be accepted or rejected at his pleasure, but because it is in reality conferred” (III/IV, article 14).

\begin{derivation}
  \item[] \((\pvar{TestimonyArguments.SolaGratia}{24} \land \pvar{TestimonyArguments.SolaGratia}{1}) \rightarrow \pvar{TestimonyArguments.SolaGratia}{28}\)\gloss{\emph{would give:} Conversion is God’s work alone: God regenerates, and the will does nothing to it}
\end{derivation}
\par\medskip\noindent\textbf{\texttt{romans\allowbreak{}Step}}\par\nobreak\noindent
Romans 9:16, read as being about salvation: if God’s saving mercy “depends not on human will or exertion”, conversion is God’s work alone. Rated \texttt{disputed}: Arminius grants that the verse is about salvation and reads its willing and running as seeking salvation by the works of the law, so that the mercy goes to believers and can be refused.

\begin{derivation}
  \item[] \((\pvar{TestimonyArguments.SolaGratia}{3} \land \pvar{TestimonyArguments.SolaGratia}{23}) \rightarrow \pvar{TestimonyArguments.SolaGratia}{28}\)\gloss{\emph{would give:} Conversion is God’s work alone: God regenerates, and the will does nothing to it}
\end{derivation}
\par\medskip\noindent\textbf{\texttt{passive\allowbreak{}Step}}\par\nobreak\noindent
From God’s work alone to the third part of the conclusion: the will is passive in being made alive. Rated \texttt{well\allowbreak{}Supported}: the Lutherans, who part from Dort on the next step, grant this one. It is taken as what Dort means by the phrase: the new life is “a making alive, which God works in us without our aid” (III/IV, article 12). “Passive” here means that the will is not the agent of its new life. Garrigou-Lagrange grants God’s work alone and holds the will “does not remain passive” under operating grace, but in another sense: its consent follows grace “in the genus or order of formal, efficient, and final cause”. The dilemma over the Thomists’ claim sets the two senses apart (\texttt{what\allowbreak{}Not\allowbreak{}Merely\allowbreak{}Passive\allowbreak{}Means}).

\begin{derivation}
  \item[] \(\pvar{TestimonyArguments.SolaGratia}{28} \rightarrow \pvar{TestimonyArguments.SolaGratia}{31}\)\gloss{\emph{would give:} In being made alive the will is passive, and contributes nothing to it}
\end{derivation}
\par\medskip\noindent\textbf{\texttt{operating\allowbreak{}Grace\allowbreak{}Step}}\par\nobreak\noindent
\textbf{Aquinas on operating grace.} In converting grace God is the sole mover of the will, so God alone is the one who regenerates.

\begin{derivation}
  \item[] \(\pvar{TestimonyArguments.SolaGratia}{16} \rightarrow \pvar{TestimonyArguments.SolaGratia}{28}\)\gloss{\emph{would give:} Conversion is God’s work alone: God regenerates, and the will does nothing to it}
\end{derivation}
\par\medskip\noindent\textbf{\texttt{efficacy\allowbreak{}Step}}\par\nobreak\noindent
\textbf{Aquinas on God’s intention.} If God moves the will intending that it attain grace, the grace brings about what it is given for. Rated \texttt{disputed}: the Molinists grant the intention and the infallible result, and ascribe it to the consent God foreknows, not to the grace.

\begin{derivation}
  \item[] \(\pvar{TestimonyArguments.SolaGratia}{17} \rightarrow \pvar{TestimonyArguments.SolaGratia}{30}\)\gloss{\emph{would give:} The grace that saves brings about the consent it seeks}
\end{derivation}
\par\medskip\noindent\textbf{\texttt{effectual\allowbreak{}Step}}\par\nobreak\noindent
\textbf{From God’s work alone to the second part of the conclusion: the grace that saves cannot be finally refused.} Rated \texttt{disputed}: the Formula of Concord holds conversion to be the Spirit’s work alone, with the will “purely passive”, and holds grace resistible all the same.

\begin{derivation}
  \item[] \(\pvar{TestimonyArguments.SolaGratia}{28} \rightarrow \pvar{TestimonyArguments.SolaGratia}{30}\)\gloss{\emph{would give:} The grace that saves brings about the consent it seeks}
\end{derivation}
\par\medskip\noindent\textbf{\texttt{texts\allowbreak{}Of\allowbreak{}Resistance\allowbreak{}Step}}\par\nobreak\noindent
What the texts of resistance say together: the Spirit’s work is resisted. Acts 7:51 says so of Stephen’s hearers, and Paul’s commands not to grieve the Spirit (Ephesians 4:30) or quench him (1 Thessalonians 5:19) presuppose it of believers.

\begin{derivation}
  \item[] \((\pvar{TestimonyArguments.SolaGratia}{5} \land \pvar{TestimonyArguments.SolaGratia}{6} \land \pvar{TestimonyArguments.SolaGratia}{7}) \rightarrow \pvar{TestimonyArguments.SolaGratia}{10}\)\gloss{\emph{would give:} The Holy Spirit’s work is resisted, by unbelievers and by believers}
\end{derivation}
\par\medskip\noindent\textbf{\texttt{resistance\allowbreak{}Step}}\par\nobreak\noindent
\textbf{From resistance to final refusal.} If the Spirit is resisted, and what is resisted is the grace that saves, that grace can be finally refused. The Remonstrants’ fourth article, made explicit.

\begin{derivation}
  \item[] \((\pvar{TestimonyArguments.SolaGratia}{10} \land \pvar{TestimonyArguments.SolaGratia}{25}) \rightarrow \pvar{TestimonyArguments.SolaGratia}{13}\)\gloss{\emph{would give:} The grace that saves can be finally refused}
\end{derivation}
\par\medskip\noindent\textbf{\texttt{grief\allowbreak{}Step}}\par\nobreak\noindent
\textbf{From grieving the Spirit to final refusal.} If the Spirit is resisted, and those who obstinately resist him can grieve and lose him, the grace that saves can be finally refused. The Formula of Concord’s reading of Ephesians 4:30.

\begin{derivation}
  \item[] \((\pvar{TestimonyArguments.SolaGratia}{10} \land \pvar{TestimonyArguments.SolaGratia}{11}) \rightarrow \pvar{TestimonyArguments.SolaGratia}{13}\)\gloss{\emph{would give:} The grace that saves can be finally refused}
\end{derivation}
\par\medskip\noindent\textbf{\texttt{blasphemy\allowbreak{}Step}}\par\nobreak\noindent
\textbf{From the unforgivable sin to final refusal.} If the blasphemy against the Spirit is never forgiven (Matthew 12:31–32), and it is the final refusal of the grace offered, then that grace can be finally refused.

\begin{derivation}
  \item[] \((\pvar{TestimonyArguments.SolaGratia}{8} \land \pvar{TestimonyArguments.SolaGratia}{12}) \rightarrow \pvar{TestimonyArguments.SolaGratia}{13}\)\gloss{\emph{would give:} The grace that saves can be finally refused}
\end{derivation}
\par\medskip\noindent\textbf{\texttt{cooperation\allowbreak{}Step}}\par\nobreak\noindent
\textbf{Trent’s fourth canon, first half.} The will, moved by God, cooperates by assenting, as “Return to me” calls on it to; so it is not passive. Rated \texttt{disputed}: Aquinas grants that grace cooperates with a will that moves and holds that in grace’s first effect the mind “is moved and does not move”.

\begin{derivation}
  \item[] \((\pvar{TestimonyArguments.SolaGratia}{21} \land \pvar{TestimonyArguments.SolaGratia}{26}) \rightarrow \lnot \pvar{TestimonyArguments.SolaGratia}{31}\)\gloss{\emph{would deny:} In being made alive the will is passive, and contributes nothing to it}
\end{derivation}
\par\medskip\noindent\textbf{\texttt{dissent\allowbreak{}Step}}\par\nobreak\noindent
\textbf{Trent’s fourth canon, second half.} The will, moved by God, can refuse if it wishes; so grace is not effectual of itself. Rated \texttt{disputed}: Aquinas grants that a will God moves could, taken simply, fail to be moved, and still holds that the man God moves to grace attains it infallibly.

\begin{derivation}
  \item[] \(\pvar{TestimonyArguments.SolaGratia}{27} \rightarrow \lnot \pvar{TestimonyArguments.SolaGratia}{30}\)\gloss{\emph{would deny:} The grace that saves brings about the consent it seeks}
\end{derivation}
\paragraph*{The lines}
\leavevmode\par\nobreak
\par\medskip\noindent\textbf{\texttt{grace\allowbreak{}Line}} --- Romans 11:6 (by grace, so not by works)\par\nobreak\noindent
\textbf{Romans 11:6}, the line every party walks. It delivers the first part of the conclusion and nothing more.

\begin{derivation}
  \item[(1)] \(\pvar{TestimonyArguments.SolaGratia}{4}\)\gloss{Romans 11:6 — if it is by grace, it is no longer on the basis of works}
  \item[(2)] \(\pvar{TestimonyArguments.SolaGratia}{4} \rightarrow \pvar{TestimonyArguments.SolaGratia}{29}\)
  \item[\(\vdash\)] \(\pvar{TestimonyArguments.SolaGratia}{29}\)\gloss{\emph{at issue:} Salvation is by grace and not earned: grace goes before every saving act}
\end{derivation}
\par\medskip\noindent\textbf{\texttt{ephesians\allowbreak{}Line}} --- Ephesians 2:8–9 (“this is not your own doing”)\par\nobreak\noindent
\textbf{Ephesians 2:8–9}, read with “this” pointing to the whole of being saved by grace through faith. It delivers that faith is God’s gift, which is not yet the conclusion.

\begin{derivation}
  \item[(1)] \(\pvar{TestimonyArguments.SolaGratia}{2}\)\gloss{Ephesians 2:8–9 — by grace you have been saved through faith, and this is not your own doing; it is the gift of God}
  \item[(2)] \(\pvar{TestimonyArguments.SolaGratia}{22}\)\gloss{In Ephesians 2:8, “this” (τοῦτο) points to the whole of being saved by grace through faith}
  \item[(3)] \((\pvar{TestimonyArguments.SolaGratia}{2} \land \pvar{TestimonyArguments.SolaGratia}{22}) \rightarrow \pvar{TestimonyArguments.SolaGratia}{24}\)
  \item[\(\vdash\)] \(\pvar{TestimonyArguments.SolaGratia}{24}\)\gloss{\emph{at issue:} Faith itself is God’s gift}
\end{derivation}
\par\medskip\noindent\textbf{\texttt{gift\allowbreak{}Line}} --- From faith given to the dead, to God’s work alone (Ephesians 2:1–5)\par\nobreak\noindent
\textbf{From faith given to the dead, to God’s work alone.} Its ground is Ephesians 2:1–5; that faith is God’s gift comes from the Ephesians line above.

\begin{derivation}
  \item[(1)] \(\pvar{TestimonyArguments.SolaGratia}{1}\)\gloss{Ephesians 2:1–5 — dead in trespasses, God made us alive together with Christ}
  \item[(2)] \((\pvar{TestimonyArguments.SolaGratia}{24} \land \pvar{TestimonyArguments.SolaGratia}{1}) \rightarrow \pvar{TestimonyArguments.SolaGratia}{28}\)
  \item[\(\vdash\)] \(\pvar{TestimonyArguments.SolaGratia}{28}\)\gloss{\emph{at issue:} Conversion is God’s work alone: God regenerates, and the will does nothing to it}
\end{derivation}
\par\medskip\noindent\textbf{\texttt{romans\allowbreak{}Line}} --- Romans 9:16 (“not on human will or exertion”)\par\nobreak\noindent
\textbf{Romans 9:16}, read as being about salvation. It reaches God’s work alone without the Ephesians grammar.

\begin{derivation}
  \item[(1)] \(\pvar{TestimonyArguments.SolaGratia}{3}\)\gloss{Romans 9:16 — it depends not on human will or exertion, but on God, who has mercy}
  \item[(2)] \(\pvar{TestimonyArguments.SolaGratia}{23}\)\gloss{Romans 9:16 is about God’s mercy in saving, not only in choosing a people for history}
  \item[(3)] \((\pvar{TestimonyArguments.SolaGratia}{3} \land \pvar{TestimonyArguments.SolaGratia}{23}) \rightarrow \pvar{TestimonyArguments.SolaGratia}{28}\)
  \item[\(\vdash\)] \(\pvar{TestimonyArguments.SolaGratia}{28}\)\gloss{\emph{at issue:} Conversion is God’s work alone: God regenerates, and the will does nothing to it}
\end{derivation}
\par\medskip\noindent\textbf{\texttt{cooperation\allowbreak{}Line}} --- Trent, canon 4: the will cooperates by assenting\par\nobreak\noindent
\textbf{Trent on cooperation}: Zechariah 1:3 with Lamentations 5:21, and the first half of canon 4.

\begin{derivation}
  \item[(1)] \(\pvar{TestimonyArguments.SolaGratia}{21}\)\gloss{Zechariah 1:3 (“return to me, and I will return to you”) with Lamentations 5:21 (“restore us to yourself, O Lord”)}
  \item[(2)] \(\pvar{TestimonyArguments.SolaGratia}{26}\)\gloss{The free will, moved by God, cooperates by assenting to God’s call}
  \item[(3)] \((\pvar{TestimonyArguments.SolaGratia}{21} \land \pvar{TestimonyArguments.SolaGratia}{26}) \rightarrow \lnot \pvar{TestimonyArguments.SolaGratia}{31}\)
  \item[\(\vdash\)] \(\lnot \pvar{TestimonyArguments.SolaGratia}{31}\)\gloss{\emph{at issue, the denial of:} In being made alive the will is passive, and contributes nothing to it}
\end{derivation}
\par\medskip\noindent\textbf{\texttt{dissent\allowbreak{}Line}} --- Trent, canon 4: the will can refuse if it wishes\par\nobreak\noindent
\textbf{Trent on dissent}: the second half of canon 4.

\begin{derivation}
  \item[(1)] \(\pvar{TestimonyArguments.SolaGratia}{27}\)\gloss{The free will, moved by God, can refuse its assent if it wishes}
  \item[(2)] \(\pvar{TestimonyArguments.SolaGratia}{27} \rightarrow \lnot \pvar{TestimonyArguments.SolaGratia}{30}\)
  \item[\(\vdash\)] \(\lnot \pvar{TestimonyArguments.SolaGratia}{30}\)\gloss{\emph{at issue, the denial of:} The grace that saves brings about the consent it seeks}
\end{derivation}
\par\medskip\noindent\textbf{\texttt{texts\allowbreak{}Of\allowbreak{}Resistance\allowbreak{}Line}} --- The texts of resistance (Acts 7:51; Ephesians 4:30; 1 Thessalonians 5:19)\par\nobreak\noindent
\textbf{The texts of resistance}: Acts 7:51, Ephesians 4:30 and 1 Thessalonians 5:19, delivering that the Spirit’s work is resisted. Every party walks this line.

\begin{derivation}
  \item[(1)] \(\pvar{TestimonyArguments.SolaGratia}{5}\)\gloss{Acts 7:51 — you always resist the Holy Spirit}
  \item[(2)] \(\pvar{TestimonyArguments.SolaGratia}{6}\)\gloss{Ephesians 4:30 — do not grieve the Holy Spirit of God, by whom you were sealed}
  \item[(3)] \(\pvar{TestimonyArguments.SolaGratia}{7}\)\gloss{1 Thessalonians 5:19 — do not quench the Spirit}
  \item[(4)] \((\pvar{TestimonyArguments.SolaGratia}{5} \land \pvar{TestimonyArguments.SolaGratia}{6} \land \pvar{TestimonyArguments.SolaGratia}{7}) \rightarrow \pvar{TestimonyArguments.SolaGratia}{10}\)
  \item[\(\vdash\)] \(\pvar{TestimonyArguments.SolaGratia}{10}\)\gloss{\emph{at issue:} The Holy Spirit’s work is resisted, by unbelievers and by believers}
\end{derivation}
\par\medskip\noindent\textbf{\texttt{resistance\allowbreak{}Line}} --- Remonstrance, article 4: grace is not irresistible (Acts 7:51)\par\nobreak\noindent
\textbf{The Remonstrants on resistance}: what Stephen’s hearers resisted was the grace that saves, so that grace can be finally refused.

\begin{derivation}
  \item[(1)] \(\pvar{TestimonyArguments.SolaGratia}{25}\)\gloss{What Stephen’s hearers resisted in Acts 7:51 was the grace that saves}
  \item[(2)] \((\pvar{TestimonyArguments.SolaGratia}{10} \land \pvar{TestimonyArguments.SolaGratia}{25}) \rightarrow \pvar{TestimonyArguments.SolaGratia}{13}\)
  \item[\(\vdash\)] \(\pvar{TestimonyArguments.SolaGratia}{13}\)\gloss{\emph{at issue:} The grace that saves can be finally refused}
\end{derivation}
\par\medskip\noindent\textbf{\texttt{grief\allowbreak{}Line}} --- Grieving the Spirit (Ephesians 4:30), as able to lose him\par\nobreak\noindent
\textbf{Grieving the Spirit}, read as able to end in losing him.

\begin{derivation}
  \item[(1)] \(\pvar{TestimonyArguments.SolaGratia}{11}\)\gloss{Those who obstinately resist the Spirit can grieve and lose him}
  \item[(2)] \((\pvar{TestimonyArguments.SolaGratia}{10} \land \pvar{TestimonyArguments.SolaGratia}{11}) \rightarrow \pvar{TestimonyArguments.SolaGratia}{13}\)
  \item[\(\vdash\)] \(\pvar{TestimonyArguments.SolaGratia}{13}\)\gloss{\emph{at issue:} The grace that saves can be finally refused}
\end{derivation}
\par\medskip\noindent\textbf{\texttt{blasphemy\allowbreak{}Line}} --- The sin against the Spirit (Matthew 12:31–32), as final refusal\par\nobreak\noindent
\textbf{The unforgivable sin}, read as the final refusal of grace.

\begin{derivation}
  \item[(1)] \(\pvar{TestimonyArguments.SolaGratia}{8}\)\gloss{Matthew 12:31–32 — the blasphemy against the Spirit will not be forgiven}
  \item[(2)] \(\pvar{TestimonyArguments.SolaGratia}{12}\)\gloss{The blasphemy against the Spirit is impenitence to the end: grace refused till death}
  \item[(3)] \((\pvar{TestimonyArguments.SolaGratia}{8} \land \pvar{TestimonyArguments.SolaGratia}{12}) \rightarrow \pvar{TestimonyArguments.SolaGratia}{13}\)
  \item[\(\vdash\)] \(\pvar{TestimonyArguments.SolaGratia}{13}\)\gloss{\emph{at issue:} The grace that saves can be finally refused}
\end{derivation}
\par\medskip\noindent\textbf{\texttt{operating\allowbreak{}Grace\allowbreak{}Line}} --- Operating grace: the mind is moved and does not move (Aquinas)\par\nobreak\noindent
\textbf{Operating grace} (Aquinas, I-II q.111 a.2): God alone is the one who regenerates.

\begin{derivation}
  \item[(1)] \(\pvar{TestimonyArguments.SolaGratia}{16}\)\gloss{In converting grace God is the sole mover: the mind is moved and does not move}
  \item[(2)] \(\pvar{TestimonyArguments.SolaGratia}{16} \rightarrow \pvar{TestimonyArguments.SolaGratia}{28}\)
  \item[\(\vdash\)] \(\pvar{TestimonyArguments.SolaGratia}{28}\)\gloss{\emph{at issue:} Conversion is God’s work alone: God regenerates, and the will does nothing to it}
\end{derivation}
\par\medskip\noindent\textbf{\texttt{efficacy\allowbreak{}Line}} --- God’s intention in moving: grace attained infallibly (Aquinas)\par\nobreak\noindent
\textbf{God’s intention in moving} (Aquinas, I-II q.112 a.3): the grace that saves brings about the consent it seeks.

\begin{derivation}
  \item[(1)] \(\pvar{TestimonyArguments.SolaGratia}{17}\)\gloss{God moves the will intending that it attain grace}
  \item[(2)] \(\pvar{TestimonyArguments.SolaGratia}{17} \rightarrow \pvar{TestimonyArguments.SolaGratia}{30}\)
  \item[\(\vdash\)] \(\pvar{TestimonyArguments.SolaGratia}{30}\)\gloss{\emph{at issue:} The grace that saves brings about the consent it seeks}
\end{derivation}
\paragraph*{Shared grounds}
\leavevmode\par\nobreak
\par\medskip\noindent\textbf{\texttt{reformed\allowbreak{}Shared}}\par\nobreak\noindent
What the Reformed packages rest on beyond any one line: the Reformed answer to Acts 7:51, from Dort’s account of those who refuse the call. What Stephen’s hearers resisted was the outward call of the preached word, not the grace that saves. The texts of resistance themselves are a line the Reformed walk too.

\begin{derivation}
  \item[(1)] \(\lnot \pvar{TestimonyArguments.SolaGratia}{25}\)\gloss{\emph{denied:} What Stephen’s hearers resisted in Acts 7:51 was the grace that saves}
\end{derivation}
\par\medskip\noindent\textbf{\texttt{closing\allowbreak{}Steps}}\par\nobreak\noindent
The steps that close every Reformed package: God’s work alone, unpacked into the conclusion’s last two parts, one step each.

\begin{derivation}
  \item[(1)] \(\pvar{TestimonyArguments.SolaGratia}{28} \rightarrow \pvar{TestimonyArguments.SolaGratia}{31}\)\gloss{\emph{would give:} In being made alive the will is passive, and contributes nothing to it}
  \item[(2)] \(\pvar{TestimonyArguments.SolaGratia}{28} \rightarrow \pvar{TestimonyArguments.SolaGratia}{30}\)\gloss{\emph{would give:} The grace that saves brings about the consent it seeks}
\end{derivation}
\par\medskip\noindent\textbf{\texttt{closing\allowbreak{}Inferences}}\par\nobreak\noindent
The ratings of the closing steps, which belong to no line.

{\small\begin{verbatim}
def closingInferences : List Source :=
  [passiveInference, effectualInference]
\end{verbatim}}
\paragraph*{The rivals}
\leavevmode\par\nobreak
\par\medskip\noindent\textbf{\texttt{tridentine}} --- Trent (Session VI, canon 4, on the Molinist reading)\par\nobreak\noindent
\textbf{Trent}, from the sixth session (1547). Romans 11:6, which its eighth chapter quotes; Zechariah 1:3 read with Lamentations 5:21, as its fifth chapter reads them; and canon 4, on the will that cooperates by assenting and can refuse if it wishes.

\begin{derivation}
  \item[(1)] \(\pvar{TestimonyArguments.SolaGratia}{4}\)\gloss{Romans 11:6 — if it is by grace, it is no longer on the basis of works}
  \item[(2)] \(\pvar{TestimonyArguments.SolaGratia}{21}\)\gloss{Zechariah 1:3 (“return to me, and I will return to you”) with Lamentations 5:21 (“restore us to yourself, O Lord”)}
  \item[(3)] \(\pvar{TestimonyArguments.SolaGratia}{26}\)\gloss{The free will, moved by God, cooperates by assenting to God’s call}
  \breakrow{(4)} \(\pvar{TestimonyArguments.SolaGratia}{27}\)\gloss{The free will, moved by God, can refuse its assent if it wishes}
  \breakrow{(5)} \(\lnot \pvar{TestimonyArguments.SolaGratia}{15}\)\gloss{\emph{denied:} Grace is received by faith alone}
  \breakrow{(6)} \(\pvar{TestimonyArguments.SolaGratia}{4} \rightarrow \pvar{TestimonyArguments.SolaGratia}{29}\)\gloss{\emph{would give:} Salvation is by grace and not earned: grace goes before every saving act}
  \breakrow{(7)} \((\pvar{TestimonyArguments.SolaGratia}{21} \land \pvar{TestimonyArguments.SolaGratia}{26}) \rightarrow \lnot \pvar{TestimonyArguments.SolaGratia}{31}\)\gloss{\emph{would deny:} In being made alive the will is passive, and contributes nothing to it}
  \item[(8)] \(\pvar{TestimonyArguments.SolaGratia}{27} \rightarrow \lnot \pvar{TestimonyArguments.SolaGratia}{30}\)\gloss{\emph{would deny:} The grace that saves brings about the consent it seeks}
  \item[\(\vdash\)] \(\pvar{TestimonyArguments.SolaGratia}{29} \land \pvar{TestimonyArguments.SolaGratia}{30} \land \pvar{TestimonyArguments.SolaGratia}{31}\)\gloss{\emph{at issue:} \(\pvar{TestimonyArguments.SolaGratia}{29}\) Salvation is by grace and not earned: grace goes before every saving act; \(\pvar{TestimonyArguments.SolaGratia}{30}\) The grace that saves brings about the consent it seeks; \(\pvar{TestimonyArguments.SolaGratia}{31}\) In being made alive the will is passive, and contributes nothing to it}
  \keepwithnext
\end{derivation}
\emph{Grounded in scripture alone:} Romans 11:6 — if it is by grace, it is no longer on the basis of works.
\par\medskip\noindent\textbf{\texttt{remonstrant}} --- Remonstrant (the Five Articles, 1610)\par\nobreak\noindent
\textbf{The Remonstrants}, from their articles of 1610 and from Arminius. Romans 11:6; the texts of resistance, which every party grants; the fourth article, that grace is not irresistible, read from Acts 7:51; and what Arminius grants of the Reformed case: the texts of Ephesians 2, that faith is God’s gift, and that Romans 9:16 is about salvation. What is left out is the two steps from those to God’s work alone, which are the steps he refuses.

\begin{derivation}
  \item[(1)] \(\pvar{TestimonyArguments.SolaGratia}{4}\)\gloss{Romans 11:6 — if it is by grace, it is no longer on the basis of works}
  \item[(2)] \(\pvar{TestimonyArguments.SolaGratia}{5}\)\gloss{Acts 7:51 — you always resist the Holy Spirit}
  \item[(3)] \(\pvar{TestimonyArguments.SolaGratia}{6}\)\gloss{Ephesians 4:30 — do not grieve the Holy Spirit of God, by whom you were sealed}
  \breakrow{(4)} \(\pvar{TestimonyArguments.SolaGratia}{7}\)\gloss{1 Thessalonians 5:19 — do not quench the Spirit}
  \breakrow{(5)} \(\pvar{TestimonyArguments.SolaGratia}{25}\)\gloss{What Stephen’s hearers resisted in Acts 7:51 was the grace that saves}
  \breakrow{(6)} \(\pvar{TestimonyArguments.SolaGratia}{1}\)\gloss{Ephesians 2:1–5 — dead in trespasses, God made us alive together with Christ}
  \breakrow{(7)} \(\pvar{TestimonyArguments.SolaGratia}{2}\)\gloss{Ephesians 2:8–9 — by grace you have been saved through faith, and this is not your own doing; it is the gift of God}
  \breakrow{(8)} \(\pvar{TestimonyArguments.SolaGratia}{24}\)\gloss{Faith itself is God’s gift}
  \breakrow{(9)} \(\pvar{TestimonyArguments.SolaGratia}{3}\)\gloss{Romans 9:16 — it depends not on human will or exertion, but on God, who has mercy}
  \breakrow{(10)} \(\pvar{TestimonyArguments.SolaGratia}{23}\)\gloss{Romans 9:16 is about God’s mercy in saving, not only in choosing a people for history}
  \breakrow{(11)} \(\pvar{TestimonyArguments.SolaGratia}{4} \rightarrow \pvar{TestimonyArguments.SolaGratia}{29}\)\gloss{\emph{would give:} Salvation is by grace and not earned: grace goes before every saving act}
  \breakrow{(12)} \((\pvar{TestimonyArguments.SolaGratia}{5} \land \pvar{TestimonyArguments.SolaGratia}{6} \land \pvar{TestimonyArguments.SolaGratia}{7}) \rightarrow \pvar{TestimonyArguments.SolaGratia}{10}\)\gloss{\emph{would give:} The Holy Spirit’s work is resisted, by unbelievers and by believers}
  \item[(13)] \((\pvar{TestimonyArguments.SolaGratia}{10} \land \pvar{TestimonyArguments.SolaGratia}{25}) \rightarrow \pvar{TestimonyArguments.SolaGratia}{13}\)\gloss{\emph{would give:} The grace that saves can be finally refused}
  \item[\(\vdash\)] \(\pvar{TestimonyArguments.SolaGratia}{29} \land \pvar{TestimonyArguments.SolaGratia}{30} \land \pvar{TestimonyArguments.SolaGratia}{31}\)\gloss{\emph{at issue:} \(\pvar{TestimonyArguments.SolaGratia}{29}\) Salvation is by grace and not earned: grace goes before every saving act; \(\pvar{TestimonyArguments.SolaGratia}{30}\) The grace that saves brings about the consent it seeks; \(\pvar{TestimonyArguments.SolaGratia}{31}\) In being made alive the will is passive, and contributes nothing to it}
  \keepwithnext
\end{derivation}
\emph{Grounded in scripture alone:} Romans 11:6 — if it is by grace, it is no longer on the basis of works, Acts 7:51 — you always resist the Holy Spirit, Ephesians 4:30 — do not grieve the Holy Spirit of God, by whom you were sealed, 1 Thessalonians 5:19 — do not quench the Spirit, Ephesians 2:1–5 — dead in trespasses, God made us alive together with Christ, Ephesians 2:8–9 — by grace you have been saved through faith, and this is not your own doing; it is the gift of God, Romans 9:16 — it depends not on human will or exertion, but on God, who has mercy.
\par\medskip\noindent\textbf{\texttt{lutheran}} --- Lutheran (the Formula of Concord)\par\nobreak\noindent
\textbf{The Lutherans}, from the Formula of Concord (1577) and the Augsburg Confession (1530). Conversion is the Holy Spirit’s work alone, and in it the will is “purely passive”; yet those who obstinately resist the Spirit “grieve and lose” him, so the grace that saves can be finally refused, and where it is, the fault is the person’s own. Grace is received by faith alone.

\begin{derivation}
  \item[(1)] \(\pvar{TestimonyArguments.SolaGratia}{4}\)\gloss{Romans 11:6 — if it is by grace, it is no longer on the basis of works}
  \item[(2)] \(\pvar{TestimonyArguments.SolaGratia}{5}\)\gloss{Acts 7:51 — you always resist the Holy Spirit}
  \item[(3)] \(\pvar{TestimonyArguments.SolaGratia}{6}\)\gloss{Ephesians 4:30 — do not grieve the Holy Spirit of God, by whom you were sealed}
  \breakrow{(4)} \(\pvar{TestimonyArguments.SolaGratia}{7}\)\gloss{1 Thessalonians 5:19 — do not quench the Spirit}
  \breakrow{(5)} \(\pvar{TestimonyArguments.SolaGratia}{11}\)\gloss{Those who obstinately resist the Spirit can grieve and lose him}
  \breakrow{(6)} \(\pvar{TestimonyArguments.SolaGratia}{28}\)\gloss{Conversion is God’s work alone: God regenerates, and the will does nothing to it}
  \breakrow{(7)} \(\pvar{TestimonyArguments.SolaGratia}{15}\)\gloss{Grace is received by faith alone}
  \breakrow{(8)} \(\pvar{TestimonyArguments.SolaGratia}{4} \rightarrow \pvar{TestimonyArguments.SolaGratia}{29}\)\gloss{\emph{would give:} Salvation is by grace and not earned: grace goes before every saving act}
  \breakrow{(9)} \((\pvar{TestimonyArguments.SolaGratia}{5} \land \pvar{TestimonyArguments.SolaGratia}{6} \land \pvar{TestimonyArguments.SolaGratia}{7}) \rightarrow \pvar{TestimonyArguments.SolaGratia}{10}\)\gloss{\emph{would give:} The Holy Spirit’s work is resisted, by unbelievers and by believers}
  \breakrow{(10)} \((\pvar{TestimonyArguments.SolaGratia}{10} \land \pvar{TestimonyArguments.SolaGratia}{11}) \rightarrow \pvar{TestimonyArguments.SolaGratia}{13}\)\gloss{\emph{would give:} The grace that saves can be finally refused}
  \item[(11)] \(\pvar{TestimonyArguments.SolaGratia}{28} \rightarrow \pvar{TestimonyArguments.SolaGratia}{31}\)\gloss{\emph{would give:} In being made alive the will is passive, and contributes nothing to it}
  \item[\(\vdash\)] \(\pvar{TestimonyArguments.SolaGratia}{29} \land \pvar{TestimonyArguments.SolaGratia}{30} \land \pvar{TestimonyArguments.SolaGratia}{31}\)\gloss{\emph{at issue:} \(\pvar{TestimonyArguments.SolaGratia}{29}\) Salvation is by grace and not earned: grace goes before every saving act; \(\pvar{TestimonyArguments.SolaGratia}{30}\) The grace that saves brings about the consent it seeks; \(\pvar{TestimonyArguments.SolaGratia}{31}\) In being made alive the will is passive, and contributes nothing to it}
  \keepwithnext
\end{derivation}
\emph{Grounded in scripture alone:} Romans 11:6 — if it is by grace, it is no longer on the basis of works, Acts 7:51 — you always resist the Holy Spirit, Ephesians 4:30 — do not grieve the Holy Spirit of God, by whom you were sealed, 1 Thessalonians 5:19 — do not quench the Spirit.
\par\medskip\noindent\textbf{\texttt{wesleyan}} --- Wesleyan (preventing grace)\par\nobreak\noindent
\textbf{Wesley.} God restores to every person a measure of free will, by the grace that goes before (“preventing grace”), so that each is able to believe; that grace can be resisted or yielded to; and it is received by faith alone.

\begin{derivation}
  \item[(1)] \(\pvar{TestimonyArguments.SolaGratia}{4}\)\gloss{Romans 11:6 — if it is by grace, it is no longer on the basis of works}
  \item[(2)] \(\pvar{TestimonyArguments.SolaGratia}{5}\)\gloss{Acts 7:51 — you always resist the Holy Spirit}
  \item[(3)] \(\pvar{TestimonyArguments.SolaGratia}{6}\)\gloss{Ephesians 4:30 — do not grieve the Holy Spirit of God, by whom you were sealed}
  \breakrow{(4)} \(\pvar{TestimonyArguments.SolaGratia}{7}\)\gloss{1 Thessalonians 5:19 — do not quench the Spirit}
  \breakrow{(5)} \(\pvar{TestimonyArguments.SolaGratia}{14}\)\gloss{God gives the ability to believe, as Adam was made able to stand or fall}
  \breakrow{(6)} \(\pvar{TestimonyArguments.SolaGratia}{13}\)\gloss{The grace that saves can be finally refused}
  \breakrow{(7)} \(\pvar{TestimonyArguments.SolaGratia}{15}\)\gloss{Grace is received by faith alone}
  \breakrow{(8)} \(\pvar{TestimonyArguments.SolaGratia}{4} \rightarrow \pvar{TestimonyArguments.SolaGratia}{29}\)\gloss{\emph{would give:} Salvation is by grace and not earned: grace goes before every saving act}
  \item[(9)] \((\pvar{TestimonyArguments.SolaGratia}{5} \land \pvar{TestimonyArguments.SolaGratia}{6} \land \pvar{TestimonyArguments.SolaGratia}{7}) \rightarrow \pvar{TestimonyArguments.SolaGratia}{10}\)\gloss{\emph{would give:} The Holy Spirit’s work is resisted, by unbelievers and by believers}
  \item[\(\vdash\)] \(\pvar{TestimonyArguments.SolaGratia}{29} \land \pvar{TestimonyArguments.SolaGratia}{30} \land \pvar{TestimonyArguments.SolaGratia}{31}\)\gloss{\emph{at issue:} \(\pvar{TestimonyArguments.SolaGratia}{29}\) Salvation is by grace and not earned: grace goes before every saving act; \(\pvar{TestimonyArguments.SolaGratia}{30}\) The grace that saves brings about the consent it seeks; \(\pvar{TestimonyArguments.SolaGratia}{31}\) In being made alive the will is passive, and contributes nothing to it}
  \keepwithnext
\end{derivation}
\emph{Grounded in scripture alone:} Romans 11:6 — if it is by grace, it is no longer on the basis of works, Acts 7:51 — you always resist the Holy Spirit, Ephesians 4:30 — do not grieve the Holy Spirit of God, by whom you were sealed, 1 Thessalonians 5:19 — do not quench the Spirit.
\par\medskip\noindent\textbf{\texttt{third\allowbreak{}Way}} --- The third way (an ability given, received by faith)\par\nobreak\noindent
\textbf{The third way, as it was put to this argument.} As Adam was made able to stand or to fall (Genesis 2:16–17), so God gives the one who hears the gospel the ability to believe, and grace is received by faith alone. That the offer is real is part of the third way too, and is not encoded as a premise, because no party here denies it. Stated this way it does not say what kind of help the ability is, whether one that makes believing possible or one that brings it about, and it carries nothing that would decide whether the grace that saves can be finally refused.

\begin{derivation}
  \item[(1)] \(\pvar{TestimonyArguments.SolaGratia}{4}\)\gloss{Romans 11:6 — if it is by grace, it is no longer on the basis of works}
  \item[(2)] \(\pvar{TestimonyArguments.SolaGratia}{9}\)\gloss{Genesis 2:16–17 — of the tree of the knowledge of good and evil you shall not eat}
  \item[(3)] \(\pvar{TestimonyArguments.SolaGratia}{14}\)\gloss{God gives the ability to believe, as Adam was made able to stand or fall}
  \breakrow{(4)} \(\pvar{TestimonyArguments.SolaGratia}{15}\)\gloss{Grace is received by faith alone}
  \item[(5)] \(\pvar{TestimonyArguments.SolaGratia}{4} \rightarrow \pvar{TestimonyArguments.SolaGratia}{29}\)\gloss{\emph{would give:} Salvation is by grace and not earned: grace goes before every saving act}
  \item[\(\vdash\)] \(\pvar{TestimonyArguments.SolaGratia}{29} \land \pvar{TestimonyArguments.SolaGratia}{15}\)\gloss{\emph{at issue:} \(\pvar{TestimonyArguments.SolaGratia}{29}\) Salvation is by grace and not earned: grace goes before every saving act; \(\pvar{TestimonyArguments.SolaGratia}{15}\) Grace is received by faith alone}
  \keepwithnext
\end{derivation}
\emph{Grounded in scripture alone:} Romans 11:6 — if it is by grace, it is no longer on the basis of works, Genesis 2:16–17 — of the tree of the knowledge of good and evil you shall not eat.
\par\medskip\noindent\textbf{\texttt{objection}} --- The objection from the Spirit’s resistance\par\nobreak\noindent
\textbf{The objection from the Spirit’s resistance, strengthened.} Three routes to the claim that the grace that saves can be finally refused. Two run from the texts of resistance: with the Remonstrants’ reading of what Stephen’s hearers resisted, and with the Formula of Concord’s reading of grieving the Spirit, that it can end in losing him. The third runs from the sin that is never forgiven, read with Augustine as impenitence to the end. It concludes that the grace that saves can be finally refused. That this denies the second part of sola gratia is not a premise of the objection: it is what the two claims mean, and the meanings join them (\texttt{Testimony\allowbreak{}Arguments.Meanings.Sola\allowbreak{}Gratia}).

\begin{derivation}
  \item[(1)] \(\pvar{TestimonyArguments.SolaGratia}{5}\)\gloss{Acts 7:51 — you always resist the Holy Spirit}
  \item[(2)] \(\pvar{TestimonyArguments.SolaGratia}{6}\)\gloss{Ephesians 4:30 — do not grieve the Holy Spirit of God, by whom you were sealed}
  \item[(3)] \(\pvar{TestimonyArguments.SolaGratia}{7}\)\gloss{1 Thessalonians 5:19 — do not quench the Spirit}
  \breakrow{(4)} \(\pvar{TestimonyArguments.SolaGratia}{25}\)\gloss{What Stephen’s hearers resisted in Acts 7:51 was the grace that saves}
  \breakrow{(5)} \(\pvar{TestimonyArguments.SolaGratia}{11}\)\gloss{Those who obstinately resist the Spirit can grieve and lose him}
  \breakrow{(6)} \(\pvar{TestimonyArguments.SolaGratia}{8}\)\gloss{Matthew 12:31–32 — the blasphemy against the Spirit will not be forgiven}
  \breakrow{(7)} \(\pvar{TestimonyArguments.SolaGratia}{12}\)\gloss{The blasphemy against the Spirit is impenitence to the end: grace refused till death}
  \breakrow{(8)} \((\pvar{TestimonyArguments.SolaGratia}{5} \land \pvar{TestimonyArguments.SolaGratia}{6} \land \pvar{TestimonyArguments.SolaGratia}{7}) \rightarrow \pvar{TestimonyArguments.SolaGratia}{10}\)\gloss{\emph{would give:} The Holy Spirit’s work is resisted, by unbelievers and by believers}
  \breakrow{(9)} \((\pvar{TestimonyArguments.SolaGratia}{10} \land \pvar{TestimonyArguments.SolaGratia}{25}) \rightarrow \pvar{TestimonyArguments.SolaGratia}{13}\)
  \breakrow{(10)} \((\pvar{TestimonyArguments.SolaGratia}{10} \land \pvar{TestimonyArguments.SolaGratia}{11}) \rightarrow \pvar{TestimonyArguments.SolaGratia}{13}\)
  \item[(11)] \((\pvar{TestimonyArguments.SolaGratia}{8} \land \pvar{TestimonyArguments.SolaGratia}{12}) \rightarrow \pvar{TestimonyArguments.SolaGratia}{13}\)
  \item[\(\vdash\)] \(\pvar{TestimonyArguments.SolaGratia}{13}\)\gloss{\emph{at issue:} The grace that saves can be finally refused}
  \keepwithnext
\end{derivation}
\emph{Grounded in scripture alone:} Acts 7:51 — you always resist the Holy Spirit, Ephesians 4:30 — do not grieve the Holy Spirit of God, by whom you were sealed, 1 Thessalonians 5:19 — do not quench the Spirit, Matthew 12:31–32 — the blasphemy against the Spirit will not be forgiven.
\paragraph*{The position argued for}
\leavevmode\par\nobreak
\par\medskip\noindent\textbf{\texttt{reformed}} --- Reformed (the Canons of Dort)\par\nobreak\noindent
\textbf{The Reformed position}, as the Synod of Dort (1619) and the Westminster Confession (1647) state it. Two routes to God’s work alone, one through Ephesians 2 and one through Romans 9:16, with Romans 11:6 for the first part of the conclusion, and the Reformed answer to Acts 7:51.

\begin{derivation}
  \item[(1)] \(\pvar{TestimonyArguments.SolaGratia}{4}\)\gloss{Romans 11:6 — if it is by grace, it is no longer on the basis of works}
  \item[(2)] \(\pvar{TestimonyArguments.SolaGratia}{2}\)\gloss{Ephesians 2:8–9 — by grace you have been saved through faith, and this is not your own doing; it is the gift of God}
  \item[(3)] \(\pvar{TestimonyArguments.SolaGratia}{22}\)\gloss{In Ephesians 2:8, “this” (τοῦτο) points to the whole of being saved by grace through faith}
  \breakrow{(4)} \(\pvar{TestimonyArguments.SolaGratia}{1}\)\gloss{Ephesians 2:1–5 — dead in trespasses, God made us alive together with Christ}
  \breakrow{(5)} \(\pvar{TestimonyArguments.SolaGratia}{3}\)\gloss{Romans 9:16 — it depends not on human will or exertion, but on God, who has mercy}
  \breakrow{(6)} \(\pvar{TestimonyArguments.SolaGratia}{23}\)\gloss{Romans 9:16 is about God’s mercy in saving, not only in choosing a people for history}
  \breakrow{(7)} \(\pvar{TestimonyArguments.SolaGratia}{5}\)\gloss{Acts 7:51 — you always resist the Holy Spirit}
  \breakrow{(8)} \(\pvar{TestimonyArguments.SolaGratia}{6}\)\gloss{Ephesians 4:30 — do not grieve the Holy Spirit of God, by whom you were sealed}
  \breakrow{(9)} \(\pvar{TestimonyArguments.SolaGratia}{7}\)\gloss{1 Thessalonians 5:19 — do not quench the Spirit}
  \breakrow{(10)} \(\lnot \pvar{TestimonyArguments.SolaGratia}{25}\)\gloss{\emph{denied:} What Stephen’s hearers resisted in Acts 7:51 was the grace that saves}
  \breakrow{(11)} \(\pvar{TestimonyArguments.SolaGratia}{4} \rightarrow \pvar{TestimonyArguments.SolaGratia}{29}\)\gloss{\emph{would give:} Salvation is by grace and not earned: grace goes before every saving act}
  \breakrow{(12)} \((\pvar{TestimonyArguments.SolaGratia}{2} \land \pvar{TestimonyArguments.SolaGratia}{22}) \rightarrow \pvar{TestimonyArguments.SolaGratia}{24}\)\gloss{\emph{would give:} Faith itself is God’s gift}
  \breakrow{(13)} \((\pvar{TestimonyArguments.SolaGratia}{24} \land \pvar{TestimonyArguments.SolaGratia}{1}) \rightarrow \pvar{TestimonyArguments.SolaGratia}{28}\)\gloss{\emph{would give:} Conversion is God’s work alone: God regenerates, and the will does nothing to it}
  \breakrow{(14)} \((\pvar{TestimonyArguments.SolaGratia}{3} \land \pvar{TestimonyArguments.SolaGratia}{23}) \rightarrow \pvar{TestimonyArguments.SolaGratia}{28}\)\gloss{\emph{would give:} Conversion is God’s work alone: God regenerates, and the will does nothing to it}
  \breakrow{(15)} \((\pvar{TestimonyArguments.SolaGratia}{5} \land \pvar{TestimonyArguments.SolaGratia}{6} \land \pvar{TestimonyArguments.SolaGratia}{7}) \rightarrow \pvar{TestimonyArguments.SolaGratia}{10}\)\gloss{\emph{would give:} The Holy Spirit’s work is resisted, by unbelievers and by believers}
  \breakrow{(16)} \(\pvar{TestimonyArguments.SolaGratia}{28} \rightarrow \pvar{TestimonyArguments.SolaGratia}{31}\)\gloss{\emph{would give:} In being made alive the will is passive, and contributes nothing to it}
  \item[(17)] \(\pvar{TestimonyArguments.SolaGratia}{28} \rightarrow \pvar{TestimonyArguments.SolaGratia}{30}\)\gloss{\emph{would give:} The grace that saves brings about the consent it seeks}
  \item[\(\vdash\)] \(\pvar{TestimonyArguments.SolaGratia}{29} \land \pvar{TestimonyArguments.SolaGratia}{30} \land \pvar{TestimonyArguments.SolaGratia}{31}\)\gloss{\emph{at issue:} \(\pvar{TestimonyArguments.SolaGratia}{29}\) Salvation is by grace and not earned: grace goes before every saving act; \(\pvar{TestimonyArguments.SolaGratia}{30}\) The grace that saves brings about the consent it seeks; \(\pvar{TestimonyArguments.SolaGratia}{31}\) In being made alive the will is passive, and contributes nothing to it}
  \keepwithnext
\end{derivation}
\emph{Grounded in scripture alone:} Romans 11:6 — if it is by grace, it is no longer on the basis of works, Ephesians 2:8–9 — by grace you have been saved through faith, and this is not your own doing; it is the gift of God, Ephesians 2:1–5 — dead in trespasses, God made us alive together with Christ, Romans 9:16 — it depends not on human will or exertion, but on God, who has mercy, Acts 7:51 — you always resist the Holy Spirit, Ephesians 4:30 — do not grieve the Holy Spirit of God, by whom you were sealed, 1 Thessalonians 5:19 — do not quench the Spirit.
\paragraph*{Each strand alone}
\leavevmode\par\nobreak
\par\medskip\noindent\textbf{\texttt{ephesians\allowbreak{}Strand}} --- Reformed, the Ephesians strand alone\par\nobreak\noindent
The Ephesians strand alone: Romans 11:6, Ephesians 2 with its reading of “this”, and the step from the gift to God’s work alone.

\begin{derivation}
  \item[(1)] \(\pvar{TestimonyArguments.SolaGratia}{4}\)\gloss{Romans 11:6 — if it is by grace, it is no longer on the basis of works}
  \item[(2)] \(\pvar{TestimonyArguments.SolaGratia}{2}\)\gloss{Ephesians 2:8–9 — by grace you have been saved through faith, and this is not your own doing; it is the gift of God}
  \item[(3)] \(\pvar{TestimonyArguments.SolaGratia}{22}\)\gloss{In Ephesians 2:8, “this” (τοῦτο) points to the whole of being saved by grace through faith}
  \breakrow{(4)} \(\pvar{TestimonyArguments.SolaGratia}{1}\)\gloss{Ephesians 2:1–5 — dead in trespasses, God made us alive together with Christ}
  \breakrow{(5)} \(\lnot \pvar{TestimonyArguments.SolaGratia}{25}\)\gloss{\emph{denied:} What Stephen’s hearers resisted in Acts 7:51 was the grace that saves}
  \breakrow{(6)} \(\pvar{TestimonyArguments.SolaGratia}{4} \rightarrow \pvar{TestimonyArguments.SolaGratia}{29}\)\gloss{\emph{would give:} Salvation is by grace and not earned: grace goes before every saving act}
  \breakrow{(7)} \((\pvar{TestimonyArguments.SolaGratia}{2} \land \pvar{TestimonyArguments.SolaGratia}{22}) \rightarrow \pvar{TestimonyArguments.SolaGratia}{24}\)\gloss{\emph{would give:} Faith itself is God’s gift}
  \breakrow{(8)} \((\pvar{TestimonyArguments.SolaGratia}{24} \land \pvar{TestimonyArguments.SolaGratia}{1}) \rightarrow \pvar{TestimonyArguments.SolaGratia}{28}\)\gloss{\emph{would give:} Conversion is God’s work alone: God regenerates, and the will does nothing to it}
  \breakrow{(9)} \(\pvar{TestimonyArguments.SolaGratia}{28} \rightarrow \pvar{TestimonyArguments.SolaGratia}{31}\)\gloss{\emph{would give:} In being made alive the will is passive, and contributes nothing to it}
  \item[(10)] \(\pvar{TestimonyArguments.SolaGratia}{28} \rightarrow \pvar{TestimonyArguments.SolaGratia}{30}\)\gloss{\emph{would give:} The grace that saves brings about the consent it seeks}
  \item[\(\vdash\)] \(\pvar{TestimonyArguments.SolaGratia}{29} \land \pvar{TestimonyArguments.SolaGratia}{30} \land \pvar{TestimonyArguments.SolaGratia}{31}\)\gloss{\emph{at issue:} \(\pvar{TestimonyArguments.SolaGratia}{29}\) Salvation is by grace and not earned: grace goes before every saving act; \(\pvar{TestimonyArguments.SolaGratia}{30}\) The grace that saves brings about the consent it seeks; \(\pvar{TestimonyArguments.SolaGratia}{31}\) In being made alive the will is passive, and contributes nothing to it}
  \keepwithnext
\end{derivation}
\emph{Grounded in scripture alone:} Romans 11:6 — if it is by grace, it is no longer on the basis of works, Ephesians 2:8–9 — by grace you have been saved through faith, and this is not your own doing; it is the gift of God, Ephesians 2:1–5 — dead in trespasses, God made us alive together with Christ.
\par\medskip\noindent\textbf{\texttt{romans\allowbreak{}Strand}} --- Reformed, the Romans strand alone\par\nobreak\noindent
The Romans strand alone: Romans 11:6, and Romans 9:16 read as being about salvation.

\begin{derivation}
  \item[(1)] \(\pvar{TestimonyArguments.SolaGratia}{4}\)\gloss{Romans 11:6 — if it is by grace, it is no longer on the basis of works}
  \item[(2)] \(\pvar{TestimonyArguments.SolaGratia}{3}\)\gloss{Romans 9:16 — it depends not on human will or exertion, but on God, who has mercy}
  \item[(3)] \(\pvar{TestimonyArguments.SolaGratia}{23}\)\gloss{Romans 9:16 is about God’s mercy in saving, not only in choosing a people for history}
  \breakrow{(4)} \(\lnot \pvar{TestimonyArguments.SolaGratia}{25}\)\gloss{\emph{denied:} What Stephen’s hearers resisted in Acts 7:51 was the grace that saves}
  \breakrow{(5)} \(\pvar{TestimonyArguments.SolaGratia}{4} \rightarrow \pvar{TestimonyArguments.SolaGratia}{29}\)\gloss{\emph{would give:} Salvation is by grace and not earned: grace goes before every saving act}
  \breakrow{(6)} \((\pvar{TestimonyArguments.SolaGratia}{3} \land \pvar{TestimonyArguments.SolaGratia}{23}) \rightarrow \pvar{TestimonyArguments.SolaGratia}{28}\)\gloss{\emph{would give:} Conversion is God’s work alone: God regenerates, and the will does nothing to it}
  \breakrow{(7)} \(\pvar{TestimonyArguments.SolaGratia}{28} \rightarrow \pvar{TestimonyArguments.SolaGratia}{31}\)\gloss{\emph{would give:} In being made alive the will is passive, and contributes nothing to it}
  \item[(8)] \(\pvar{TestimonyArguments.SolaGratia}{28} \rightarrow \pvar{TestimonyArguments.SolaGratia}{30}\)\gloss{\emph{would give:} The grace that saves brings about the consent it seeks}
  \item[\(\vdash\)] \(\pvar{TestimonyArguments.SolaGratia}{29} \land \pvar{TestimonyArguments.SolaGratia}{30} \land \pvar{TestimonyArguments.SolaGratia}{31}\)\gloss{\emph{at issue:} \(\pvar{TestimonyArguments.SolaGratia}{29}\) Salvation is by grace and not earned: grace goes before every saving act; \(\pvar{TestimonyArguments.SolaGratia}{30}\) The grace that saves brings about the consent it seeks; \(\pvar{TestimonyArguments.SolaGratia}{31}\) In being made alive the will is passive, and contributes nothing to it}
  \keepwithnext
\end{derivation}
\emph{Grounded in scripture alone:} Romans 11:6 — if it is by grace, it is no longer on the basis of works, Romans 9:16 — it depends not on human will or exertion, but on God, who has mercy.
\par\medskip\noindent\textbf{\texttt{ephesians\allowbreak{}Without\allowbreak{}The\allowbreak{}Gift\allowbreak{}Step}} --- Reformed, the Ephesians strand without the step from the gift\par\nobreak\noindent
The Ephesians strand without the step Arminius refuses: every text and reading it rests on, and the step to faith as God’s gift, with only the step from the gift to God’s work alone left out. The Reformed answer to Acts 7:51 is not part of the strand, and is left out with the rest of what the strand does not read.

\begin{derivation}
  \item[(1)] \(\pvar{TestimonyArguments.SolaGratia}{4}\)\gloss{Romans 11:6 — if it is by grace, it is no longer on the basis of works}
  \item[(2)] \(\pvar{TestimonyArguments.SolaGratia}{2}\)\gloss{Ephesians 2:8–9 — by grace you have been saved through faith, and this is not your own doing; it is the gift of God}
  \item[(3)] \(\pvar{TestimonyArguments.SolaGratia}{22}\)\gloss{In Ephesians 2:8, “this” (τοῦτο) points to the whole of being saved by grace through faith}
  \breakrow{(4)} \(\pvar{TestimonyArguments.SolaGratia}{1}\)\gloss{Ephesians 2:1–5 — dead in trespasses, God made us alive together with Christ}
  \breakrow{(5)} \(\pvar{TestimonyArguments.SolaGratia}{4} \rightarrow \pvar{TestimonyArguments.SolaGratia}{29}\)\gloss{\emph{would give:} Salvation is by grace and not earned: grace goes before every saving act}
  \breakrow{(6)} \((\pvar{TestimonyArguments.SolaGratia}{2} \land \pvar{TestimonyArguments.SolaGratia}{22}) \rightarrow \pvar{TestimonyArguments.SolaGratia}{24}\)\gloss{\emph{would give:} Faith itself is God’s gift}
  \breakrow{(7)} \(\pvar{TestimonyArguments.SolaGratia}{28} \rightarrow \pvar{TestimonyArguments.SolaGratia}{31}\)\gloss{\emph{would give:} In being made alive the will is passive, and contributes nothing to it}
  \item[(8)] \(\pvar{TestimonyArguments.SolaGratia}{28} \rightarrow \pvar{TestimonyArguments.SolaGratia}{30}\)\gloss{\emph{would give:} The grace that saves brings about the consent it seeks}
  \item[\(\vdash\)] \(\pvar{TestimonyArguments.SolaGratia}{29} \land \pvar{TestimonyArguments.SolaGratia}{30} \land \pvar{TestimonyArguments.SolaGratia}{31}\)\gloss{\emph{at issue:} \(\pvar{TestimonyArguments.SolaGratia}{29}\) Salvation is by grace and not earned: grace goes before every saving act; \(\pvar{TestimonyArguments.SolaGratia}{30}\) The grace that saves brings about the consent it seeks; \(\pvar{TestimonyArguments.SolaGratia}{31}\) In being made alive the will is passive, and contributes nothing to it}
  \keepwithnext
\end{derivation}
\emph{Grounded in scripture alone:} Romans 11:6 — if it is by grace, it is no longer on the basis of works, Ephesians 2:8–9 — by grace you have been saved through faith, and this is not your own doing; it is the gift of God, Ephesians 2:1–5 — dead in trespasses, God made us alive together with Christ.
\paragraph*{The readings, removed}
\leavevmode\par\nobreak
\par\medskip\noindent\textbf{\texttt{reformed\allowbreak{}Without\allowbreak{}Touto}} --- Reformed, minus its reading of “this” in Ephesians 2:8\par\nobreak\noindent
The Reformed position without its reading of “this” in Ephesians 2:8. The Romans strand is untouched.

\begin{derivation}
  \item[(1)] \(\pvar{TestimonyArguments.SolaGratia}{4}\)\gloss{Romans 11:6 — if it is by grace, it is no longer on the basis of works}
  \item[(2)] \(\pvar{TestimonyArguments.SolaGratia}{2}\)\gloss{Ephesians 2:8–9 — by grace you have been saved through faith, and this is not your own doing; it is the gift of God}
  \item[(3)] \(\pvar{TestimonyArguments.SolaGratia}{1}\)\gloss{Ephesians 2:1–5 — dead in trespasses, God made us alive together with Christ}
  \breakrow{(4)} \(\pvar{TestimonyArguments.SolaGratia}{3}\)\gloss{Romans 9:16 — it depends not on human will or exertion, but on God, who has mercy}
  \breakrow{(5)} \(\pvar{TestimonyArguments.SolaGratia}{23}\)\gloss{Romans 9:16 is about God’s mercy in saving, not only in choosing a people for history}
  \breakrow{(6)} \(\pvar{TestimonyArguments.SolaGratia}{5}\)\gloss{Acts 7:51 — you always resist the Holy Spirit}
  \breakrow{(7)} \(\pvar{TestimonyArguments.SolaGratia}{6}\)\gloss{Ephesians 4:30 — do not grieve the Holy Spirit of God, by whom you were sealed}
  \breakrow{(8)} \(\pvar{TestimonyArguments.SolaGratia}{7}\)\gloss{1 Thessalonians 5:19 — do not quench the Spirit}
  \breakrow{(9)} \(\lnot \pvar{TestimonyArguments.SolaGratia}{25}\)\gloss{\emph{denied:} What Stephen’s hearers resisted in Acts 7:51 was the grace that saves}
  \breakrow{(10)} \(\pvar{TestimonyArguments.SolaGratia}{4} \rightarrow \pvar{TestimonyArguments.SolaGratia}{29}\)\gloss{\emph{would give:} Salvation is by grace and not earned: grace goes before every saving act}
  \breakrow{(11)} \((\pvar{TestimonyArguments.SolaGratia}{2} \land \pvar{TestimonyArguments.SolaGratia}{22}) \rightarrow \pvar{TestimonyArguments.SolaGratia}{24}\)\gloss{\emph{would give:} Faith itself is God’s gift}
  \breakrow{(12)} \((\pvar{TestimonyArguments.SolaGratia}{24} \land \pvar{TestimonyArguments.SolaGratia}{1}) \rightarrow \pvar{TestimonyArguments.SolaGratia}{28}\)\gloss{\emph{would give:} Conversion is God’s work alone: God regenerates, and the will does nothing to it}
  \breakrow{(13)} \((\pvar{TestimonyArguments.SolaGratia}{3} \land \pvar{TestimonyArguments.SolaGratia}{23}) \rightarrow \pvar{TestimonyArguments.SolaGratia}{28}\)\gloss{\emph{would give:} Conversion is God’s work alone: God regenerates, and the will does nothing to it}
  \breakrow{(14)} \((\pvar{TestimonyArguments.SolaGratia}{5} \land \pvar{TestimonyArguments.SolaGratia}{6} \land \pvar{TestimonyArguments.SolaGratia}{7}) \rightarrow \pvar{TestimonyArguments.SolaGratia}{10}\)\gloss{\emph{would give:} The Holy Spirit’s work is resisted, by unbelievers and by believers}
  \breakrow{(15)} \(\pvar{TestimonyArguments.SolaGratia}{28} \rightarrow \pvar{TestimonyArguments.SolaGratia}{31}\)\gloss{\emph{would give:} In being made alive the will is passive, and contributes nothing to it}
  \item[(16)] \(\pvar{TestimonyArguments.SolaGratia}{28} \rightarrow \pvar{TestimonyArguments.SolaGratia}{30}\)\gloss{\emph{would give:} The grace that saves brings about the consent it seeks}
  \item[\(\vdash\)] \(\pvar{TestimonyArguments.SolaGratia}{29} \land \pvar{TestimonyArguments.SolaGratia}{30} \land \pvar{TestimonyArguments.SolaGratia}{31}\)\gloss{\emph{at issue:} \(\pvar{TestimonyArguments.SolaGratia}{29}\) Salvation is by grace and not earned: grace goes before every saving act; \(\pvar{TestimonyArguments.SolaGratia}{30}\) The grace that saves brings about the consent it seeks; \(\pvar{TestimonyArguments.SolaGratia}{31}\) In being made alive the will is passive, and contributes nothing to it}
  \keepwithnext
\end{derivation}
\emph{Grounded in scripture alone:} Romans 11:6 — if it is by grace, it is no longer on the basis of works, Ephesians 2:8–9 — by grace you have been saved through faith, and this is not your own doing; it is the gift of God, Ephesians 2:1–5 — dead in trespasses, God made us alive together with Christ, Romans 9:16 — it depends not on human will or exertion, but on God, who has mercy, Acts 7:51 — you always resist the Holy Spirit, Ephesians 4:30 — do not grieve the Holy Spirit of God, by whom you were sealed, 1 Thessalonians 5:19 — do not quench the Spirit.
\par\medskip\noindent\textbf{\texttt{reformed\allowbreak{}Without\allowbreak{}Romans\allowbreak{}Reading}} --- Reformed, minus Romans 9:16 as being about salvation\par\nobreak\noindent
The Reformed position without the reading of Romans 9:16 as being about salvation. The Ephesians strand is untouched.

\begin{derivation}
  \item[(1)] \(\pvar{TestimonyArguments.SolaGratia}{4}\)\gloss{Romans 11:6 — if it is by grace, it is no longer on the basis of works}
  \item[(2)] \(\pvar{TestimonyArguments.SolaGratia}{2}\)\gloss{Ephesians 2:8–9 — by grace you have been saved through faith, and this is not your own doing; it is the gift of God}
  \item[(3)] \(\pvar{TestimonyArguments.SolaGratia}{22}\)\gloss{In Ephesians 2:8, “this” (τοῦτο) points to the whole of being saved by grace through faith}
  \breakrow{(4)} \(\pvar{TestimonyArguments.SolaGratia}{1}\)\gloss{Ephesians 2:1–5 — dead in trespasses, God made us alive together with Christ}
  \breakrow{(5)} \(\pvar{TestimonyArguments.SolaGratia}{3}\)\gloss{Romans 9:16 — it depends not on human will or exertion, but on God, who has mercy}
  \breakrow{(6)} \(\pvar{TestimonyArguments.SolaGratia}{5}\)\gloss{Acts 7:51 — you always resist the Holy Spirit}
  \breakrow{(7)} \(\pvar{TestimonyArguments.SolaGratia}{6}\)\gloss{Ephesians 4:30 — do not grieve the Holy Spirit of God, by whom you were sealed}
  \breakrow{(8)} \(\pvar{TestimonyArguments.SolaGratia}{7}\)\gloss{1 Thessalonians 5:19 — do not quench the Spirit}
  \breakrow{(9)} \(\lnot \pvar{TestimonyArguments.SolaGratia}{25}\)\gloss{\emph{denied:} What Stephen’s hearers resisted in Acts 7:51 was the grace that saves}
  \breakrow{(10)} \(\pvar{TestimonyArguments.SolaGratia}{4} \rightarrow \pvar{TestimonyArguments.SolaGratia}{29}\)\gloss{\emph{would give:} Salvation is by grace and not earned: grace goes before every saving act}
  \breakrow{(11)} \((\pvar{TestimonyArguments.SolaGratia}{2} \land \pvar{TestimonyArguments.SolaGratia}{22}) \rightarrow \pvar{TestimonyArguments.SolaGratia}{24}\)\gloss{\emph{would give:} Faith itself is God’s gift}
  \breakrow{(12)} \((\pvar{TestimonyArguments.SolaGratia}{24} \land \pvar{TestimonyArguments.SolaGratia}{1}) \rightarrow \pvar{TestimonyArguments.SolaGratia}{28}\)\gloss{\emph{would give:} Conversion is God’s work alone: God regenerates, and the will does nothing to it}
  \breakrow{(13)} \((\pvar{TestimonyArguments.SolaGratia}{3} \land \pvar{TestimonyArguments.SolaGratia}{23}) \rightarrow \pvar{TestimonyArguments.SolaGratia}{28}\)\gloss{\emph{would give:} Conversion is God’s work alone: God regenerates, and the will does nothing to it}
  \breakrow{(14)} \((\pvar{TestimonyArguments.SolaGratia}{5} \land \pvar{TestimonyArguments.SolaGratia}{6} \land \pvar{TestimonyArguments.SolaGratia}{7}) \rightarrow \pvar{TestimonyArguments.SolaGratia}{10}\)\gloss{\emph{would give:} The Holy Spirit’s work is resisted, by unbelievers and by believers}
  \breakrow{(15)} \(\pvar{TestimonyArguments.SolaGratia}{28} \rightarrow \pvar{TestimonyArguments.SolaGratia}{31}\)\gloss{\emph{would give:} In being made alive the will is passive, and contributes nothing to it}
  \item[(16)] \(\pvar{TestimonyArguments.SolaGratia}{28} \rightarrow \pvar{TestimonyArguments.SolaGratia}{30}\)\gloss{\emph{would give:} The grace that saves brings about the consent it seeks}
  \item[\(\vdash\)] \(\pvar{TestimonyArguments.SolaGratia}{29} \land \pvar{TestimonyArguments.SolaGratia}{30} \land \pvar{TestimonyArguments.SolaGratia}{31}\)\gloss{\emph{at issue:} \(\pvar{TestimonyArguments.SolaGratia}{29}\) Salvation is by grace and not earned: grace goes before every saving act; \(\pvar{TestimonyArguments.SolaGratia}{30}\) The grace that saves brings about the consent it seeks; \(\pvar{TestimonyArguments.SolaGratia}{31}\) In being made alive the will is passive, and contributes nothing to it}
  \keepwithnext
\end{derivation}
\emph{Grounded in scripture alone:} Romans 11:6 — if it is by grace, it is no longer on the basis of works, Ephesians 2:8–9 — by grace you have been saved through faith, and this is not your own doing; it is the gift of God, Ephesians 2:1–5 — dead in trespasses, God made us alive together with Christ, Romans 9:16 — it depends not on human will or exertion, but on God, who has mercy, Acts 7:51 — you always resist the Holy Spirit, Ephesians 4:30 — do not grieve the Holy Spirit of God, by whom you were sealed, 1 Thessalonians 5:19 — do not quench the Spirit.
\par\medskip\noindent\textbf{\texttt{reformed\allowbreak{}Without\allowbreak{}Either\allowbreak{}Reading}} --- Reformed, minus both readings\par\nobreak\noindent
The Reformed position without either reading. Every text is kept, and every step.

\begin{derivation}
  \item[(1)] \(\pvar{TestimonyArguments.SolaGratia}{4}\)\gloss{Romans 11:6 — if it is by grace, it is no longer on the basis of works}
  \item[(2)] \(\pvar{TestimonyArguments.SolaGratia}{2}\)\gloss{Ephesians 2:8–9 — by grace you have been saved through faith, and this is not your own doing; it is the gift of God}
  \item[(3)] \(\pvar{TestimonyArguments.SolaGratia}{1}\)\gloss{Ephesians 2:1–5 — dead in trespasses, God made us alive together with Christ}
  \breakrow{(4)} \(\pvar{TestimonyArguments.SolaGratia}{3}\)\gloss{Romans 9:16 — it depends not on human will or exertion, but on God, who has mercy}
  \breakrow{(5)} \(\pvar{TestimonyArguments.SolaGratia}{5}\)\gloss{Acts 7:51 — you always resist the Holy Spirit}
  \breakrow{(6)} \(\pvar{TestimonyArguments.SolaGratia}{6}\)\gloss{Ephesians 4:30 — do not grieve the Holy Spirit of God, by whom you were sealed}
  \breakrow{(7)} \(\pvar{TestimonyArguments.SolaGratia}{7}\)\gloss{1 Thessalonians 5:19 — do not quench the Spirit}
  \breakrow{(8)} \(\lnot \pvar{TestimonyArguments.SolaGratia}{25}\)\gloss{\emph{denied:} What Stephen’s hearers resisted in Acts 7:51 was the grace that saves}
  \breakrow{(9)} \(\pvar{TestimonyArguments.SolaGratia}{4} \rightarrow \pvar{TestimonyArguments.SolaGratia}{29}\)\gloss{\emph{would give:} Salvation is by grace and not earned: grace goes before every saving act}
  \breakrow{(10)} \((\pvar{TestimonyArguments.SolaGratia}{2} \land \pvar{TestimonyArguments.SolaGratia}{22}) \rightarrow \pvar{TestimonyArguments.SolaGratia}{24}\)\gloss{\emph{would give:} Faith itself is God’s gift}
  \breakrow{(11)} \((\pvar{TestimonyArguments.SolaGratia}{24} \land \pvar{TestimonyArguments.SolaGratia}{1}) \rightarrow \pvar{TestimonyArguments.SolaGratia}{28}\)\gloss{\emph{would give:} Conversion is God’s work alone: God regenerates, and the will does nothing to it}
  \breakrow{(12)} \((\pvar{TestimonyArguments.SolaGratia}{3} \land \pvar{TestimonyArguments.SolaGratia}{23}) \rightarrow \pvar{TestimonyArguments.SolaGratia}{28}\)\gloss{\emph{would give:} Conversion is God’s work alone: God regenerates, and the will does nothing to it}
  \breakrow{(13)} \((\pvar{TestimonyArguments.SolaGratia}{5} \land \pvar{TestimonyArguments.SolaGratia}{6} \land \pvar{TestimonyArguments.SolaGratia}{7}) \rightarrow \pvar{TestimonyArguments.SolaGratia}{10}\)\gloss{\emph{would give:} The Holy Spirit’s work is resisted, by unbelievers and by believers}
  \breakrow{(14)} \(\pvar{TestimonyArguments.SolaGratia}{28} \rightarrow \pvar{TestimonyArguments.SolaGratia}{31}\)\gloss{\emph{would give:} In being made alive the will is passive, and contributes nothing to it}
  \item[(15)] \(\pvar{TestimonyArguments.SolaGratia}{28} \rightarrow \pvar{TestimonyArguments.SolaGratia}{30}\)\gloss{\emph{would give:} The grace that saves brings about the consent it seeks}
  \item[\(\vdash\)] \(\pvar{TestimonyArguments.SolaGratia}{29} \land \pvar{TestimonyArguments.SolaGratia}{30} \land \pvar{TestimonyArguments.SolaGratia}{31}\)\gloss{\emph{at issue:} \(\pvar{TestimonyArguments.SolaGratia}{29}\) Salvation is by grace and not earned: grace goes before every saving act; \(\pvar{TestimonyArguments.SolaGratia}{30}\) The grace that saves brings about the consent it seeks; \(\pvar{TestimonyArguments.SolaGratia}{31}\) In being made alive the will is passive, and contributes nothing to it}
  \keepwithnext
\end{derivation}
\emph{Grounded in scripture alone:} Romans 11:6 — if it is by grace, it is no longer on the basis of works, Ephesians 2:8–9 — by grace you have been saved through faith, and this is not your own doing; it is the gift of God, Ephesians 2:1–5 — dead in trespasses, God made us alive together with Christ, Romans 9:16 — it depends not on human will or exertion, but on God, who has mercy, Acts 7:51 — you always resist the Holy Spirit, Ephesians 4:30 — do not grieve the Holy Spirit of God, by whom you were sealed, 1 Thessalonians 5:19 — do not quench the Spirit.
\paragraph*{The positions}
\leavevmode\par\nobreak
\par\medskip\noindent\textbf{\texttt{reformed\_\allowbreak{}establishes}}\par\nobreak\noindent
\textbf{Grant the Reformed premises, and sola gratia follows.} The premises are Romans 11:6; Ephesians 2, with “this” pointing to the whole of being saved by grace through faith; Romans 9:16, read as being about salvation; the steps from each of those to God’s work alone; and the Reformed answer to Acts 7:51. The result does not say that the readings or the steps are right, only that the conclusion follows from them. Four of the steps are rated \texttt{disputed}, each on a reader who grants its grounds and denies its conclusion: the step from “this” to faith as a gift (Robertson), the steps from each strand to God’s work alone (Arminius), and the step from God’s work alone to grace that cannot be finally refused (the Formula of Concord).

{\small\begin{verbatim}
theorem reformed_establishes : Establishes reformed
-- axioms: propext, Quot.sound
\end{verbatim}}
\par\medskip\noindent\textbf{\texttt{trents\allowbreak{}Reading}}\par\nobreak\noindent
Trent’s reading, written down: salvation is by grace, the will cooperates and can refuse, so grace is not effectual of itself and the will is not passive in being made alive, and grace is not received by faith alone (canon 9). Claims Trent’s documents do not address, such as the readings of Ephesians 2 and Romans 9, are set true; nothing turns on them. On these sixteen claims Trent’s reading and Arminius’s coincide; the two positions differ in which premises they hold and whom they cite.

{\small\begin{verbatim}
def trentsReading : Valuation Claim :=
  fun a =>
    match a with
    | Claim.faithAloneReceivesGrace => False
    | Claim.conversionIsGodsWorkAlone => False
    | Claim.graceIsEffectualOfItself => False
    | Claim.willIsPassiveInRegeneration => False
    | x => True
\end{verbatim}}
\par\medskip\noindent\textbf{\texttt{tridentine\_\allowbreak{}not\_\allowbreak{}establishes}}\par\nobreak\noindent
\textbf{Trent’s premises do not lead to sola gratia.} A \emph{countermodel} shows it: a way of holding every one of Trent’s premises while denying the conclusion. On Trent’s account the will, moved by God, cooperates by assenting and can refuse its consent; so the second and third parts of the conclusion fail. Whether canon 4 really denies the second part is itself disputed, because Aquinas holds both that the will can refuse, taken simply, and that the man God moves to grace attains it infallibly; that dispute is in the rating of the canon’s step, not in this result.

{\small\begin{verbatim}
theorem tridentine_not_establishes : ¬Establishes tridentine
-- axioms: propext, Quot.sound
\end{verbatim}}
\par\medskip\noindent\textbf{\texttt{arminius\allowbreak{}Reading}}\par\nobreak\noindent
Arminius’s reading, written down: salvation is by grace, faith is God’s gift, Romans 9:16 is about salvation, and the grace that saves can be refused. It also sets the will’s passivity false. The Remonstrants denied that too, but their denial of it is not encoded, and nothing in their package turns on it. The reading of “this” in Ephesians 2:8 is set true, as Hoehner and Robertson hold it; no reading of the pronoun is cited to Arminius.

{\small\begin{verbatim}
def arminiusReading : Valuation Claim :=
  fun a =>
    match a with
    | Claim.conversionIsGodsWorkAlone => False
    | Claim.graceIsEffectualOfItself => False
    | Claim.willIsPassiveInRegeneration => False
    | x => True
\end{verbatim}}
\par\medskip\noindent\textbf{\texttt{remonstrant\_\allowbreak{}not\_\allowbreak{}establishes}}\par\nobreak\noindent
\textbf{The Remonstrants’ premises do not lead to sola gratia either}, though they grant much of what the Reformed case rests on: that faith is God’s gift, given to those dead in sin, and that Romans 9:16 is about salvation. What they add is their fourth article, that grace is not irresistible, read from Acts 7:51; what they leave out is the two steps from those grants to God’s work alone.

{\small\begin{verbatim}
theorem remonstrant_not_establishes : ¬Establishes remonstrant
-- axioms: propext, Quot.sound
\end{verbatim}}
\par\medskip\noindent\textbf{\texttt{rivals\_\allowbreak{}hold\_\allowbreak{}salvation\_\allowbreak{}by\_\allowbreak{}grace}}\par\nobreak\noindent
Both rivals reach the first part, that salvation is by grace. Trent’s premises deliver it, and so do the Remonstrants’, by the same line from Romans 11:6. A conclusion stated only as “salvation is by grace” would be established by all three positions and would divide none of them, which is why the conclusion here has two further parts.

{\small\begin{verbatim}
theorem rivals_hold_salvation_by_grace : Entails tridentine.premises (p
    Claim.salvationIsByGrace) ∧ Entails remonstrant.premises (p
    Claim.salvationIsByGrace)
-- axioms: propext, Quot.sound
\end{verbatim}}
\par\medskip\noindent\textbf{\texttt{trent\_\allowbreak{}denies\_\allowbreak{}passivity}}\par\nobreak\noindent
Trent’s premises deny the third part, that the will is passive. Canon 4 condemns anyone who says the will “is merely passive”, and the premises deliver the denial outright. The step is rated \texttt{disputed}, because the canon speaks of the will preparing itself for justification and the Reformed claim is about the moment of being made alive.

{\small\begin{verbatim}
theorem trent_denies_passivity : Entails tridentine.premises (notP
    Claim.willIsPassiveInRegeneration)
-- axioms: propext, Quot.sound
\end{verbatim}}
\par\medskip\noindent\textbf{\texttt{trent\_\allowbreak{}denies\_\allowbreak{}effectual\_\allowbreak{}grace}}\par\nobreak\noindent
Trent’s premises deny the second part too, that grace is effectual of itself. Canon 4 condemns anyone who says the will “cannot refuse its consent, if it would”, and the premises deliver the denial. The step is rated \texttt{disputed}: Aquinas grants that the will moved by God could, taken simply, refuse, and holds that the man God moves to grace “will infallibly attain to it”. So a Catholic reader in his school grants the canon and not this conclusion.

{\small\begin{verbatim}
theorem trent_denies_effectual_grace : Entails tridentine.premises (notP
    Claim.graceIsEffectualOfItself)
-- axioms: propext, Quot.sound
\end{verbatim}}
\par\medskip\noindent\textbf{\texttt{remonstrants\_\allowbreak{}deny\_\allowbreak{}effectual\_\allowbreak{}grace}}\par\nobreak\noindent
The Remonstrants deny the second part: grace can be resisted. Their fourth article reads Acts 7:51 so, and their premises deliver that the grace that saves can be finally refused. That this is the denial of grace effectual of itself is not a premise of theirs but what the two claims mean: the meanings join them, and the theorem is stated over the premises together with that join.

{\small\begin{verbatim}
theorem remonstrants_deny_effectual_grace : Entails remonstrant.held (notP
    Claim.graceIsEffectualOfItself)
-- axioms: propext, Quot.sound
\end{verbatim}}
\paragraph*{Which reading the Reformed case depends on}
\leavevmode\par\nobreak
\par\medskip\noindent\textbf{\texttt{touto\_\allowbreak{}not\_\allowbreak{}load\_\allowbreak{}bearing}}\par\nobreak\noindent
The case does not depend on the reading of “this”. Remove the reading on which “this” in Ephesians 2:8 points to the whole of being saved by grace through faith, keep everything else, and sola gratia still follows by Romans 9:16, which does not use the Ephesians grammar. The reading itself is well supported: Hoehner and Robertson both hold it.

{\small\begin{verbatim}
theorem touto_not_load_bearing : Establishes reformedWithoutTouto
-- axioms: propext, Quot.sound
\end{verbatim}}
\par\medskip\noindent\textbf{\texttt{romans\_\allowbreak{}reading\_\allowbreak{}not\_\allowbreak{}load\_\allowbreak{}bearing}}\par\nobreak\noindent
Nor does it depend on the reading of Romans 9:16. Remove the reading on which the verse is about salvation, keep everything else, and sola gratia still follows by Ephesians 2.

{\small\begin{verbatim}
theorem romans_reading_not_load_bearing : Establishes
    reformedWithoutRomansReading
-- axioms: propext, Quot.sound
\end{verbatim}}
\par\medskip\noindent\textbf{\texttt{neither\allowbreak{}Reading\allowbreak{}Granted}}\par\nobreak\noindent
The reading on which neither of the two readings holds: “this” in Ephesians 2:8 does not point to the whole, and Romans 9:16 is about God’s choice of a people for a part in history alone. No reader is cited who holds both denials; the countermodel shows what follows if both readings go, not that anyone has given them up. Everything else the Reformed position grants is kept, the Reformed answer to Acts 7:51 included.

{\small\begin{verbatim}
def neitherReadingGranted : Valuation Claim :=
  fun a =>
    match a with
    | Claim.toutoRefersToTheWhole => False
    | Claim.romans9_16ConcernsSalvation => False
    | Claim.faithIsGodsGift => False
    | Claim.conversionIsGodsWorkAlone => False
    | Claim.graceIsEffectualOfItself => False
    | Claim.willIsPassiveInRegeneration => False
    | Claim.acts7_51ResistsSavingGrace => False
    | x => True
\end{verbatim}}
\par\medskip\noindent\textbf{\texttt{readings\_\allowbreak{}jointly\_\allowbreak{}load\_\allowbreak{}bearing}}\par\nobreak\noindent
The two readings together carry it. Remove both, keep every text and every step, and sola gratia no longer follows, because nothing is left that reaches God’s work alone. So the Reformed case needs at least one of the two readings, and either will do. A reading whose removal stops the conclusion following is called \emph{load-bearing}; here only the pair is.

{\small\begin{verbatim}
theorem readings_jointly_load_bearing : ¬Establishes
    reformedWithoutEitherReading
-- axioms: propext, Quot.sound
\end{verbatim}}
\par\medskip\noindent\textbf{\texttt{ephesians\_\allowbreak{}strand\_\allowbreak{}establishes}}\par\nobreak\noindent
The Ephesians strand, argued alone, reaches sola gratia. It rests on two steps, both rated \texttt{disputed}: from “this” to faith as God’s gift, and from the gift to God’s work alone.

{\small\begin{verbatim}
theorem ephesians_strand_establishes : Establishes ephesiansStrand
-- axioms: propext, Quot.sound
\end{verbatim}}
\par\medskip\noindent\textbf{\texttt{romans\_\allowbreak{}strand\_\allowbreak{}establishes}}\par\nobreak\noindent
So does the Romans strand, on one step rated \texttt{disputed}.

{\small\begin{verbatim}
theorem romans_strand_establishes : Establishes romansStrand
-- axioms: propext, Quot.sound
\end{verbatim}}
\paragraph*{Where Arminius parts from the Ephesians strand}
\leavevmode\par\nobreak
\par\medskip\noindent\textbf{\texttt{gift\_\allowbreak{}of\_\allowbreak{}faith\_\allowbreak{}does\_\allowbreak{}not\_\allowbreak{}settle\_\allowbreak{}resistibility}}\par\nobreak\noindent
\textbf{That faith is God’s gift does not settle whether grace can be refused.} Grant everything the Ephesians strand rests on short of its last step: the texts, “this” pointing to the whole, and so faith as God’s gift, given to those who were dead in sin. Leave out only the step from that gift to God’s work alone, and sola gratia does not follow.

Arminius grants the texts and that faith is God’s gift; the reading of “this” is Hoehner’s and Robertson’s, not cited to him. On a reading that grants all of these, his, the conclusion fails. So where he parts from the strand is the step from the gift to God’s work alone. The encoding writes “faith is God’s gift” as one claim, and Dort and Arminius mean different things by it: for Dort a gift “in reality conferred”, for Arminius one that can be refused. The difference shows here as Arminius refusing the step.

{\small\begin{verbatim}
theorem gift_of_faith_does_not_settle_resistibility : ¬Establishes
    ephesiansWithoutTheGiftStep
-- axioms: propext, Quot.sound
\end{verbatim}}
\par\medskip\noindent\textbf{\texttt{dort\allowbreak{}Reading}}\par\nobreak\noindent
Dort’s reading: every claim of the Reformed position holds. What Stephen’s hearers resisted was the preached word, not the grace that saves; grieving the Spirit does not end in losing him; the unforgivable sin is not the final refusal of saving grace; and that grace cannot be finally refused.

{\small\begin{verbatim}
def dortReading : Valuation Claim :=
  fun a =>
    match a with
    | Claim.acts7_51ResistsSavingGrace => False
    | Claim.savingGraceCanBeFinallyRefused => False
    | Claim.grievingCanLoseTheSpirit => False
    | Claim.blasphemyIsFinalRefusalOfGrace => False
    | x => True
\end{verbatim}}
\par\medskip\noindent\textbf{\texttt{reformed\_\allowbreak{}is\_\allowbreak{}satisfiable}}\par\nobreak\noindent
The Reformed package has a model.

{\small\begin{verbatim}
theorem reformed_is_satisfiable : Satisfiable reformed.premises
-- axioms: propext, Quot.sound
\end{verbatim}}
\par\medskip\noindent\textbf{\texttt{reformed\allowbreak{}Without\allowbreak{}Touto\_\allowbreak{}is\_\allowbreak{}satisfiable}}\par\nobreak\noindent
The Reformed package without its reading of “this” has a model.

{\small\begin{verbatim}
theorem reformedWithoutTouto_is_satisfiable : Satisfiable
    reformedWithoutTouto.premises
-- axioms: propext, Quot.sound
\end{verbatim}}
\par\medskip\noindent\textbf{\texttt{reformed\allowbreak{}Without\allowbreak{}Romans\allowbreak{}Reading\_\allowbreak{}is\_\allowbreak{}satisfiable}}\par\nobreak\noindent
The Reformed package without its reading of Romans 9:16 has a model.

{\small\begin{verbatim}
theorem reformedWithoutRomansReading_is_satisfiable : Satisfiable
    reformedWithoutRomansReading.premises
-- axioms: propext, Quot.sound
\end{verbatim}}
\par\medskip\noindent\textbf{\texttt{ephesians\allowbreak{}Strand\_\allowbreak{}is\_\allowbreak{}satisfiable}}\par\nobreak\noindent
The Ephesians strand alone has a model.

{\small\begin{verbatim}
theorem ephesiansStrand_is_satisfiable : Satisfiable ephesiansStrand.premises
-- axioms: propext, Quot.sound
\end{verbatim}}
\par\medskip\noindent\textbf{\texttt{romans\allowbreak{}Strand\_\allowbreak{}is\_\allowbreak{}satisfiable}}\par\nobreak\noindent
The Romans strand alone has a model.

{\small\begin{verbatim}
theorem romansStrand_is_satisfiable : Satisfiable romansStrand.premises
-- axioms: propext, Quot.sound
\end{verbatim}}
\subsection*{Freedom to refuse, and the ability God gives}
\addcontentsline{toc}{subsection}{Freedom to refuse, and the ability God gives}
The strongest objection to sola gratia comes from the Arminian side, and it is an objection from Scripture. Stephen told the council, “you always resist the Holy Spirit” (Acts 7:51). Paul tells believers, “do not grieve the Holy Spirit of God, by whom you were sealed for the day of redemption” (Ephesians 4:30), and “do not quench the Spirit” (1 Thessalonians 5:19). Jesus says that “the blasphemy against the Spirit will not be forgiven … either in this age or in the age to come” (Matthew 12:31–32). If the Spirit can be resisted, grieved and quenched, and if one sin against him is never forgiven, then surely the grace he brings can be refused.

\subsubsection*{What the texts settle, and what they do not}
Every party grants what the texts say. Dort says that believers “grieve the Holy Spirit” (V, article 5); Calvin, on Ephesians 4:30, that “we grieve him when we do not follow his guidance”; the Formula of Concord, that “those who always resist the Holy Ghost … are not converted”. So the step from the texts to “the Spirit’s work is resisted” is rated \emph{consensus}, the highest rating: accepted across traditions.

The dispute is the next step, from resistance to the claim that the grace that saves can be \emph{finally} refused. “Resist” has two senses. A person can resist for a time and be overcome, which the Reformed say happens to the elect before they believe, and to believers when they grieve the Spirit and are brought back. Or a person can refuse to the end. Only the second sense denies sola gratia, and each route to it needs a reading of its own:

\begin{itemize}
  \item \textbf{Acts 7:51}, read as resistance to the grace that saves, as the Remonstrants read it. The Reformed say that what Stephen’s hearers resisted was the outward call of the preached word.
  \item \textbf{Ephesians 4:30}, read with the Formula of Concord: those who obstinately resist “do not receive, but grieve and lose, the Holy Ghost”. Dort denies the reading: God “does not wholly withdraw the Holy Spirit from his own people, even in their melancholy falls”.
  \item \textbf{Matthew 12:31–32}, read with Augustine: the sin never forgiven is impenitence to the end, “the persevering hardness of an impenitent heart”. Calvin rejects the reading outright (“Those who define it to be impenitence may be refuted without any difficulty”) and takes the sin to be knowing contempt of the Spirit, which the regenerate “cannot possibly fall into”. And the step from Augustine’s reading to final refusal of \emph{saving} grace is rated \texttt{disputed}, because Augustine himself holds both that reading and that God is the one “whom no man’s will resists when He wills to give salvation”.
\end{itemize}
Grieving the Spirit is Paul’s word, and it is addressed to believers. The unforgivable sin is Jesus’ word, in the Gospels. They are kept as two texts with two readings, because no reader cited here takes grieving to be the unforgivable sin.

\subsubsection*{Two positions between Dort and the Remonstrants}
\textbf{The Lutherans} of the Formula of Concord (1577) grant the ground of the Reformed case’s last step and deny its conclusion. They hold that conversion is “not only in part, but altogether an operation, gift, and present, and work of the Holy Ghost alone”, and that in it man “does nothing whatever towards it, but only suffers what God works in him”. That is the third part of sola gratia as Dort states it. And they deny that grace cannot be refused: “God does not force man to become godly”, and the perverse will “resists the Holy Ghost, who wishes to be efficacious”. So the step from “God’s work alone” to grace effectual of itself, grace that brings about the consent it seeks and so cannot be finally refused, is rated \emph{disputed}, and the Formula is why.

\textbf{Wesley} holds that “there is a measure of free-will supernaturally restored to every man”, so that each has “power either to resist or yield to the grace of God”; that grace, which “goes before” anything a person does, he calls “preventing grace”. Faith itself is God’s gift: “Of yourselves cometh neither your faith nor your salvation”. And it is received by faith alone: Wesley defends “salvation or justification, by faith only” against the objection that it is preaching against holiness (Sermon 1, III.1).

\subsubsection*{The third way}
The third way, as it was put to this argument: God gives the one who hears the gospel the ability to believe, as Adam and Eve were made able to stand or to fall (Genesis 2:16–17); the offer of grace is real; and it is received by faith alone. The reality of the offer is not encoded as a premise, because no party here denies it. Stated that way, the third way does not decide whether the grace that saves can be finally refused: Dort’s reading and Arminius’s both satisfy it, and disagree about exactly that (\texttt{third\_\allowbreak{}way\_\allowbreak{}leaves\_\allowbreak{}resistibility\_\allowbreak{}open}).

What decides it is what kind of help the ability is. Augustine drew the distinction. “The aid without which a thing does not come to pass is one thing, and the aid by which a thing comes to pass is another.” Adam was given the first, “not that by which it should be brought about that he should persevere, but that without which he could not of free will persevere”; the saints are given the second, “so that by means of this gift they cannot help persevering” (\emph{On Rebuke and Grace}, ch. 34). Read as Adam’s kind of help, the third way is the reading Wesley gives, and it cannot be held together with Dort’s position. Read as the saints’ kind, it is Dort’s reading, and it cannot be held together with Trent’s canon 4, nor, by what the claims mean, with the Remonstrants. Each break at a premise comes at a step rated \emph{disputed}, so neither shows the other side right (\texttt{what\allowbreak{}The\allowbreak{}Ability\allowbreak{}Is}).

\subsubsection*{The contrary, supposed}
Grant the Remonstrants their claim for a moment: the grace that saves can be finally refused. Put that beside the Reformed reading of the Scriptures, and read the claims as they mean, so that a grace which can be finally refused is not a grace that brings believing about. Then nothing can be true all at once. The same is so with Ephesians 2 alone, or Romans 9:16 alone (\texttt{final\_\allowbreak{}refusal\_\allowbreak{}cannot\_\allowbreak{}be\_\allowbreak{}held\_\allowbreak{}with\_\allowbreak{}the\_\allowbreak{}scriptures}). The supposition is not refuted by the bare texts: it is refuted by the texts as the Reformed read them, and at the steps a cited reader disputes, which is where the Remonstrants, the Lutherans and Wesley each part.

\subsubsection*{Each position, with its own conclusion}
Each position below is stated with the conclusion its own sources draw, not with sola gratia, so that it can be weighed against the others.

\subsubsection*{Final refusal, supposed}
What if the Remonstrants are right, and the grace that saves can be finally refused? Suppose it, and set the supposition beside the Scriptures as the Reformed read them.

\subsubsection*{Each position can be held without contradiction}
For each position, a reading is written down that makes every premise true.

\subsubsection*{How this is checked}
\par\medskip\noindent\textbf{\texttt{grace\allowbreak{}Case}} --- Romans 11:6 (by grace, so not by works)\par\nobreak\noindent
What every party holds: salvation is by grace, from Romans 11:6.

\begin{derivation}
  \item[(1)] \(\pvar{TestimonyArguments.SolaGratia}{4}\)\gloss{Romans 11:6 — if it is by grace, it is no longer on the basis of works}
  \item[(2)] \(\pvar{TestimonyArguments.SolaGratia}{4} \rightarrow \pvar{TestimonyArguments.SolaGratia}{29}\)
  \item[\(\vdash\)] \(\pvar{TestimonyArguments.SolaGratia}{29}\)\gloss{\emph{at issue:} Salvation is by grace and not earned: grace goes before every saving act}
  \keepwithnext
\end{derivation}
\emph{Grounded in scripture alone:} Romans 11:6 — if it is by grace, it is no longer on the basis of works.
\par\medskip\noindent\textbf{\texttt{tridentine\allowbreak{}Position}} --- Trent (Session VI, canon 4, on the Molinist reading)\par\nobreak\noindent
Trent’s position on the Molinist reading of canon 4: salvation by grace, the will not passive, grace not effectual of itself, and not by faith alone. The Thomists, who read the canon as compatible with grace effectual of itself, are cited against its steps and are not a position here.

\begin{derivation}
  \item[(1)] \(\pvar{TestimonyArguments.SolaGratia}{4}\)\gloss{Romans 11:6 — if it is by grace, it is no longer on the basis of works}
  \item[(2)] \(\pvar{TestimonyArguments.SolaGratia}{21}\)\gloss{Zechariah 1:3 (“return to me, and I will return to you”) with Lamentations 5:21 (“restore us to yourself, O Lord”)}
  \item[(3)] \(\pvar{TestimonyArguments.SolaGratia}{26}\)\gloss{The free will, moved by God, cooperates by assenting to God’s call}
  \breakrow{(4)} \(\pvar{TestimonyArguments.SolaGratia}{27}\)\gloss{The free will, moved by God, can refuse its assent if it wishes}
  \breakrow{(5)} \(\lnot \pvar{TestimonyArguments.SolaGratia}{15}\)\gloss{\emph{denied:} Grace is received by faith alone}
  \breakrow{(6)} \(\pvar{TestimonyArguments.SolaGratia}{4} \rightarrow \pvar{TestimonyArguments.SolaGratia}{29}\)\gloss{\emph{would give:} Salvation is by grace and not earned: grace goes before every saving act}
  \breakrow{(7)} \((\pvar{TestimonyArguments.SolaGratia}{21} \land \pvar{TestimonyArguments.SolaGratia}{26}) \rightarrow \lnot \pvar{TestimonyArguments.SolaGratia}{31}\)\gloss{\emph{would deny:} In being made alive the will is passive, and contributes nothing to it}
  \item[(8)] \(\pvar{TestimonyArguments.SolaGratia}{27} \rightarrow \lnot \pvar{TestimonyArguments.SolaGratia}{30}\)\gloss{\emph{would deny:} The grace that saves brings about the consent it seeks}
  \item[\(\vdash\)] \(\pvar{TestimonyArguments.SolaGratia}{29} \land \lnot \pvar{TestimonyArguments.SolaGratia}{31} \land \lnot \pvar{TestimonyArguments.SolaGratia}{30} \land \lnot \pvar{TestimonyArguments.SolaGratia}{15}\)\gloss{\emph{at issue:} \(\pvar{TestimonyArguments.SolaGratia}{29}\) Salvation is by grace and not earned: grace goes before every saving act; \(\pvar{TestimonyArguments.SolaGratia}{31}\) \emph{denied:} In being made alive the will is passive, and contributes nothing to it; \(\pvar{TestimonyArguments.SolaGratia}{30}\) \emph{denied:} The grace that saves brings about the consent it seeks; \(\pvar{TestimonyArguments.SolaGratia}{15}\) \emph{denied:} Grace is received by faith alone}
  \keepwithnext
\end{derivation}
\emph{Grounded in scripture alone:} Romans 11:6 — if it is by grace, it is no longer on the basis of works.
\par\medskip\noindent\textbf{\texttt{remonstrant\allowbreak{}Position}} --- Remonstrant (the Five Articles, 1610)\par\nobreak\noindent
The Remonstrants’ position: salvation by grace, and a grace that can be resisted to the end.

\begin{derivation}
  \item[(1)] \(\pvar{TestimonyArguments.SolaGratia}{4}\)\gloss{Romans 11:6 — if it is by grace, it is no longer on the basis of works}
  \item[(2)] \(\pvar{TestimonyArguments.SolaGratia}{5}\)\gloss{Acts 7:51 — you always resist the Holy Spirit}
  \item[(3)] \(\pvar{TestimonyArguments.SolaGratia}{6}\)\gloss{Ephesians 4:30 — do not grieve the Holy Spirit of God, by whom you were sealed}
  \breakrow{(4)} \(\pvar{TestimonyArguments.SolaGratia}{7}\)\gloss{1 Thessalonians 5:19 — do not quench the Spirit}
  \breakrow{(5)} \(\pvar{TestimonyArguments.SolaGratia}{25}\)\gloss{What Stephen’s hearers resisted in Acts 7:51 was the grace that saves}
  \breakrow{(6)} \(\pvar{TestimonyArguments.SolaGratia}{1}\)\gloss{Ephesians 2:1–5 — dead in trespasses, God made us alive together with Christ}
  \breakrow{(7)} \(\pvar{TestimonyArguments.SolaGratia}{2}\)\gloss{Ephesians 2:8–9 — by grace you have been saved through faith, and this is not your own doing; it is the gift of God}
  \breakrow{(8)} \(\pvar{TestimonyArguments.SolaGratia}{24}\)\gloss{Faith itself is God’s gift}
  \breakrow{(9)} \(\pvar{TestimonyArguments.SolaGratia}{3}\)\gloss{Romans 9:16 — it depends not on human will or exertion, but on God, who has mercy}
  \breakrow{(10)} \(\pvar{TestimonyArguments.SolaGratia}{23}\)\gloss{Romans 9:16 is about God’s mercy in saving, not only in choosing a people for history}
  \breakrow{(11)} \(\pvar{TestimonyArguments.SolaGratia}{4} \rightarrow \pvar{TestimonyArguments.SolaGratia}{29}\)\gloss{\emph{would give:} Salvation is by grace and not earned: grace goes before every saving act}
  \breakrow{(12)} \((\pvar{TestimonyArguments.SolaGratia}{5} \land \pvar{TestimonyArguments.SolaGratia}{6} \land \pvar{TestimonyArguments.SolaGratia}{7}) \rightarrow \pvar{TestimonyArguments.SolaGratia}{10}\)\gloss{\emph{would give:} The Holy Spirit’s work is resisted, by unbelievers and by believers}
  \item[(13)] \((\pvar{TestimonyArguments.SolaGratia}{10} \land \pvar{TestimonyArguments.SolaGratia}{25}) \rightarrow \pvar{TestimonyArguments.SolaGratia}{13}\)\gloss{\emph{would give:} The grace that saves can be finally refused}
  \item[\(\vdash\)] \(\pvar{TestimonyArguments.SolaGratia}{29} \land \pvar{TestimonyArguments.SolaGratia}{13}\)\gloss{\emph{at issue:} \(\pvar{TestimonyArguments.SolaGratia}{29}\) Salvation is by grace and not earned: grace goes before every saving act; \(\pvar{TestimonyArguments.SolaGratia}{13}\) The grace that saves can be finally refused}
  \keepwithnext
\end{derivation}
\emph{Grounded in scripture alone:} Romans 11:6 — if it is by grace, it is no longer on the basis of works, Acts 7:51 — you always resist the Holy Spirit, Ephesians 4:30 — do not grieve the Holy Spirit of God, by whom you were sealed, 1 Thessalonians 5:19 — do not quench the Spirit, Ephesians 2:1–5 — dead in trespasses, God made us alive together with Christ, Ephesians 2:8–9 — by grace you have been saved through faith, and this is not your own doing; it is the gift of God, Romans 9:16 — it depends not on human will or exertion, but on God, who has mercy.
\par\medskip\noindent\textbf{\texttt{lutheran\allowbreak{}Position}} --- Lutheran (the Formula of Concord)\par\nobreak\noindent
The Lutherans’ position: salvation by grace, the will passive in conversion, the grace that converts resistible, and faith alone.

\begin{derivation}
  \item[(1)] \(\pvar{TestimonyArguments.SolaGratia}{4}\)\gloss{Romans 11:6 — if it is by grace, it is no longer on the basis of works}
  \item[(2)] \(\pvar{TestimonyArguments.SolaGratia}{5}\)\gloss{Acts 7:51 — you always resist the Holy Spirit}
  \item[(3)] \(\pvar{TestimonyArguments.SolaGratia}{6}\)\gloss{Ephesians 4:30 — do not grieve the Holy Spirit of God, by whom you were sealed}
  \breakrow{(4)} \(\pvar{TestimonyArguments.SolaGratia}{7}\)\gloss{1 Thessalonians 5:19 — do not quench the Spirit}
  \breakrow{(5)} \(\pvar{TestimonyArguments.SolaGratia}{11}\)\gloss{Those who obstinately resist the Spirit can grieve and lose him}
  \breakrow{(6)} \(\pvar{TestimonyArguments.SolaGratia}{28}\)\gloss{Conversion is God’s work alone: God regenerates, and the will does nothing to it}
  \breakrow{(7)} \(\pvar{TestimonyArguments.SolaGratia}{15}\)\gloss{Grace is received by faith alone}
  \breakrow{(8)} \(\pvar{TestimonyArguments.SolaGratia}{4} \rightarrow \pvar{TestimonyArguments.SolaGratia}{29}\)\gloss{\emph{would give:} Salvation is by grace and not earned: grace goes before every saving act}
  \breakrow{(9)} \((\pvar{TestimonyArguments.SolaGratia}{5} \land \pvar{TestimonyArguments.SolaGratia}{6} \land \pvar{TestimonyArguments.SolaGratia}{7}) \rightarrow \pvar{TestimonyArguments.SolaGratia}{10}\)\gloss{\emph{would give:} The Holy Spirit’s work is resisted, by unbelievers and by believers}
  \breakrow{(10)} \((\pvar{TestimonyArguments.SolaGratia}{10} \land \pvar{TestimonyArguments.SolaGratia}{11}) \rightarrow \pvar{TestimonyArguments.SolaGratia}{13}\)\gloss{\emph{would give:} The grace that saves can be finally refused}
  \item[(11)] \(\pvar{TestimonyArguments.SolaGratia}{28} \rightarrow \pvar{TestimonyArguments.SolaGratia}{31}\)\gloss{\emph{would give:} In being made alive the will is passive, and contributes nothing to it}
  \item[\(\vdash\)] \(\pvar{TestimonyArguments.SolaGratia}{29} \land \pvar{TestimonyArguments.SolaGratia}{31} \land \pvar{TestimonyArguments.SolaGratia}{13} \land \pvar{TestimonyArguments.SolaGratia}{15}\)\gloss{\emph{at issue:} \(\pvar{TestimonyArguments.SolaGratia}{29}\) Salvation is by grace and not earned: grace goes before every saving act; \(\pvar{TestimonyArguments.SolaGratia}{31}\) In being made alive the will is passive, and contributes nothing to it; \(\pvar{TestimonyArguments.SolaGratia}{13}\) The grace that saves can be finally refused; \(\pvar{TestimonyArguments.SolaGratia}{15}\) Grace is received by faith alone}
  \keepwithnext
\end{derivation}
\emph{Grounded in scripture alone:} Romans 11:6 — if it is by grace, it is no longer on the basis of works, Acts 7:51 — you always resist the Holy Spirit, Ephesians 4:30 — do not grieve the Holy Spirit of God, by whom you were sealed, 1 Thessalonians 5:19 — do not quench the Spirit.
\par\medskip\noindent\textbf{\texttt{wesleyan\allowbreak{}Position}} --- Wesleyan (preventing grace)\par\nobreak\noindent
Wesley’s position: salvation by grace, a grace that goes before and can be refused, and faith alone.

\begin{derivation}
  \item[(1)] \(\pvar{TestimonyArguments.SolaGratia}{4}\)\gloss{Romans 11:6 — if it is by grace, it is no longer on the basis of works}
  \item[(2)] \(\pvar{TestimonyArguments.SolaGratia}{5}\)\gloss{Acts 7:51 — you always resist the Holy Spirit}
  \item[(3)] \(\pvar{TestimonyArguments.SolaGratia}{6}\)\gloss{Ephesians 4:30 — do not grieve the Holy Spirit of God, by whom you were sealed}
  \breakrow{(4)} \(\pvar{TestimonyArguments.SolaGratia}{7}\)\gloss{1 Thessalonians 5:19 — do not quench the Spirit}
  \breakrow{(5)} \(\pvar{TestimonyArguments.SolaGratia}{14}\)\gloss{God gives the ability to believe, as Adam was made able to stand or fall}
  \breakrow{(6)} \(\pvar{TestimonyArguments.SolaGratia}{13}\)\gloss{The grace that saves can be finally refused}
  \breakrow{(7)} \(\pvar{TestimonyArguments.SolaGratia}{15}\)\gloss{Grace is received by faith alone}
  \breakrow{(8)} \(\pvar{TestimonyArguments.SolaGratia}{4} \rightarrow \pvar{TestimonyArguments.SolaGratia}{29}\)\gloss{\emph{would give:} Salvation is by grace and not earned: grace goes before every saving act}
  \item[(9)] \((\pvar{TestimonyArguments.SolaGratia}{5} \land \pvar{TestimonyArguments.SolaGratia}{6} \land \pvar{TestimonyArguments.SolaGratia}{7}) \rightarrow \pvar{TestimonyArguments.SolaGratia}{10}\)\gloss{\emph{would give:} The Holy Spirit’s work is resisted, by unbelievers and by believers}
  \item[\(\vdash\)] \(\pvar{TestimonyArguments.SolaGratia}{29} \land \pvar{TestimonyArguments.SolaGratia}{13} \land \pvar{TestimonyArguments.SolaGratia}{15}\)\gloss{\emph{at issue:} \(\pvar{TestimonyArguments.SolaGratia}{29}\) Salvation is by grace and not earned: grace goes before every saving act; \(\pvar{TestimonyArguments.SolaGratia}{13}\) The grace that saves can be finally refused; \(\pvar{TestimonyArguments.SolaGratia}{15}\) Grace is received by faith alone}
  \keepwithnext
\end{derivation}
\emph{Grounded in scripture alone:} Romans 11:6 — if it is by grace, it is no longer on the basis of works, Acts 7:51 — you always resist the Holy Spirit, Ephesians 4:30 — do not grieve the Holy Spirit of God, by whom you were sealed, 1 Thessalonians 5:19 — do not quench the Spirit.
\paragraph*{What each position reaches}
\leavevmode\par\nobreak
\par\medskip\noindent\textbf{\texttt{grace\allowbreak{}Case\_\allowbreak{}establishes}}\par\nobreak\noindent
What every party holds follows from Romans 11:6.

{\small\begin{verbatim}
theorem graceCase_establishes : Establishes graceCase
-- axioms: propext, Quot.sound
\end{verbatim}}
\par\medskip\noindent\textbf{\texttt{tridentine\allowbreak{}Position\_\allowbreak{}establishes}}\par\nobreak\noindent
Trent’s premises, on the Molinist reading of canon 4, deliver that position.

{\small\begin{verbatim}
theorem tridentinePosition_establishes : Establishes tridentinePosition
-- axioms: propext, Quot.sound
\end{verbatim}}
\par\medskip\noindent\textbf{\texttt{remonstrant\allowbreak{}Position\_\allowbreak{}establishes}}\par\nobreak\noindent
The Remonstrants’ fourth article delivers grace that can be refused.

{\small\begin{verbatim}
theorem remonstrantPosition_establishes : Establishes remonstrantPosition
-- axioms: propext, Quot.sound
\end{verbatim}}
\par\medskip\noindent\textbf{\texttt{lutheran\allowbreak{}Position\_\allowbreak{}establishes}}\par\nobreak\noindent
\textbf{The Lutheran premises reach the Lutheran position}: the will passive in conversion, and the grace that converts resistible. The Reformed case joins the two; the Formula of Concord holds them apart. This does not show the Lutherans right, only that their premises deliver their position; that it can be held without contradiction is \texttt{lutheran\allowbreak{}Position\_\allowbreak{}is\_\allowbreak{}satisfiable}.

{\small\begin{verbatim}
theorem lutheranPosition_establishes : Establishes lutheranPosition
-- axioms: propext, Quot.sound
\end{verbatim}}
\par\medskip\noindent\textbf{\texttt{wesleyan\allowbreak{}Position\_\allowbreak{}establishes}}\par\nobreak\noindent
Wesley’s premises reach grace that goes before, can be refused, and is received by faith alone.

{\small\begin{verbatim}
theorem wesleyanPosition_establishes : Establishes wesleyanPosition
-- axioms: propext, Quot.sound
\end{verbatim}}
\par\medskip\noindent\textbf{\texttt{third\allowbreak{}Way\_\allowbreak{}establishes}}\par\nobreak\noindent
The third way’s premises reach salvation by grace, received by faith alone.

{\small\begin{verbatim}
theorem thirdWay_establishes : Establishes thirdWay
-- axioms: propext, Quot.sound
\end{verbatim}}
\par\medskip\noindent\textbf{\texttt{objection\_\allowbreak{}establishes}}\par\nobreak\noindent
\textbf{The objection from the Spirit’s resistance is valid on its own grounds.} Grant the texts of resistance and any one of its three readings, and the grace that saves is not effectual of itself.

{\small\begin{verbatim}
theorem objection_establishes : Establishes objection
-- axioms: propext, Quot.sound
\end{verbatim}}
\par\medskip\noindent\textbf{\texttt{third\_\allowbreak{}way\_\allowbreak{}leaves\_\allowbreak{}resistibility\_\allowbreak{}open}}\par\nobreak\noindent
\textbf{The third way leaves resistibility open.} Grant everything it says: Adam made able to stand or fall, the ability to believe given by God, faith alone, and that grace which can be finally refused is not effectual of itself. Whether the grace that saves is effectual of itself neither follows nor fails. Arminius’s reading satisfies every premise and denies it; Dort’s satisfies every premise and holds it. So the third way, as stated, neither blocks sola gratia nor establishes it.

{\small\begin{verbatim}
theorem third_way_leaves_resistibility_open : Independent thirdWay.premises (p
    Claim.graceIsEffectualOfItself)
-- axioms: propext, Quot.sound
\end{verbatim}}
\par\medskip\noindent\textbf{\texttt{final\_\allowbreak{}refusal\_\allowbreak{}cannot\_\allowbreak{}be\_\allowbreak{}held\_\allowbreak{}with\_\allowbreak{}the\_\allowbreak{}scriptures}}\par\nobreak\noindent
\textbf{Granted the Scriptures as the Reformed read them, final refusal has no place.} Suppose the grace that saves can be finally refused. Add that supposition to the Reformed premises: Romans 11:6; Ephesians 2, with “this” pointing to the whole of being saved by grace through faith; Romans 9:16, read as being about salvation; the texts of resistance; and the steps from each to God’s work alone. Read the claims as they mean, so that grace which can be finally refused is not grace that brings believing about. Then no reading makes everything true at once. The same holds of either strand alone, Ephesians 2 or Romans 9:16, so either text carries it.

What this claims is conditional, as every result here is: on the Reformed readings and on the steps from them, four of which are rated \emph{disputed}, each on a reader who grants its grounds and denies its conclusion. A reader who holds that the grace of God can be finally refused denies one of those steps or readings, and the chapters say which: Robertson and Calvin at the step from Ephesians 2:8 to faith as a gift, Arminius at the steps to God’s work alone, and the Lutherans at the step from God’s work alone to grace effectual of itself. It does not claim the supposition is refuted by the texts alone, without those readings.

{\small\begin{verbatim}
theorem final_refusal_cannot_be_held_with_the_scriptures : ¬Satisfiable
    (reformed.held ++ [p Claim.savingGraceCanBeFinallyRefused]) ∧ ¬Satisfiable
    (ephesiansStrand.held ++ [p Claim.savingGraceCanBeFinallyRefused]) ∧
    ¬Satisfiable (romansStrand.held ++ [p
    Claim.savingGraceCanBeFinallyRefused])
-- axioms: propext, Classical.choice, Quot.sound
\end{verbatim}}
\par\medskip\noindent\textbf{\texttt{third\allowbreak{}Way\_\allowbreak{}is\_\allowbreak{}satisfiable}}\par\nobreak\noindent
The third way has a model: Dort’s reading.

{\small\begin{verbatim}
theorem thirdWay_is_satisfiable : Satisfiable thirdWay.premises
-- axioms: propext, Quot.sound
\end{verbatim}}
\par\medskip\noindent\textbf{\texttt{objection\_\allowbreak{}is\_\allowbreak{}satisfiable}}\par\nobreak\noindent
The objection has a model: Arminius’s reading.

{\small\begin{verbatim}
theorem objection_is_satisfiable : Satisfiable objection.premises
-- axioms: propext, Quot.sound
\end{verbatim}}
\par\medskip\noindent\textbf{\texttt{lutheran\allowbreak{}Reading}}\par\nobreak\noindent
The Lutheran reading: conversion is God’s work alone and the will passive, and the grace that does it can be resisted and lost.

{\small\begin{verbatim}
def lutheranReading : Valuation Claim :=
  fun a =>
    match a with
    | Claim.graceIsEffectualOfItself => False
    | x => True
\end{verbatim}}
\par\medskip\noindent\textbf{\texttt{lutheran\allowbreak{}Position\_\allowbreak{}is\_\allowbreak{}satisfiable}}\par\nobreak\noindent
The Lutheran position has a model.

{\small\begin{verbatim}
theorem lutheranPosition_is_satisfiable : Satisfiable
    lutheranPosition.premises
-- axioms: propext, Quot.sound
\end{verbatim}}
\par\medskip\noindent\textbf{\texttt{tridentine\allowbreak{}Position\_\allowbreak{}is\_\allowbreak{}satisfiable}}\par\nobreak\noindent
Trent’s position has a model: Trent’s reading.

{\small\begin{verbatim}
theorem tridentinePosition_is_satisfiable : Satisfiable
    tridentinePosition.premises
-- axioms: propext, Quot.sound
\end{verbatim}}
\par\medskip\noindent\textbf{\texttt{remonstrant\allowbreak{}Position\_\allowbreak{}is\_\allowbreak{}satisfiable}}\par\nobreak\noindent
The Remonstrants’ position has a model: Arminius’s reading.

{\small\begin{verbatim}
theorem remonstrantPosition_is_satisfiable : Satisfiable
    remonstrantPosition.premises
-- axioms: propext, Quot.sound
\end{verbatim}}
\par\medskip\noindent\textbf{\texttt{wesleyan\allowbreak{}Position\_\allowbreak{}is\_\allowbreak{}satisfiable}}\par\nobreak\noindent
Wesley’s position has a model.

{\small\begin{verbatim}
theorem wesleyanPosition_is_satisfiable : Satisfiable
    wesleyanPosition.premises
-- axioms: propext, Quot.sound
\end{verbatim}}
\par\medskip\noindent\textbf{\texttt{grace\allowbreak{}Case\_\allowbreak{}is\_\allowbreak{}satisfiable}}\par\nobreak\noindent
What every party holds has a model.

{\small\begin{verbatim}
theorem graceCase_is_satisfiable : Satisfiable graceCase.premises
-- axioms: propext, Quot.sound
\end{verbatim}}
\paragraph*{What kind of help the ability is}
\leavevmode\par\nobreak
\par\medskip\noindent\textbf{\texttt{as\allowbreak{}Aid\allowbreak{}Without\allowbreak{}Which}}\par\nobreak\noindent
The ability read as Adam’s kind of help, “not that by which it should be brought about” that a person believes, “but that without which” they could not: so the grace that gives it is not effectual of itself. Wesley’s reading.

{\small\begin{verbatim}
def asAidWithoutWhich : Reading Claim :=
  { name := "as the help without which a person cannot believe, Adam's kind",
    commits := notP Claim.graceIsEffectualOfItself,
    source := asAidWithoutWhichSource }
\end{verbatim}}
\par\medskip\noindent\textbf{\texttt{as\allowbreak{}Aid\allowbreak{}By\allowbreak{}Which}}\par\nobreak\noindent
The ability read as the saints’ kind of help: the grace that gives it brings believing about. Augustine’s reading of the grace given to the saints, and Dort’s.

{\small\begin{verbatim}
def asAidByWhich : Reading Claim :=
  { name := "as the help by which a person believes, the saints' kind",
    commits := p Claim.graceIsEffectualOfItself,
    source := asAidByWhichSource }
\end{verbatim}}
\par\medskip\noindent\textbf{\texttt{third\allowbreak{}Way\allowbreak{}As\allowbreak{}Aid\allowbreak{}Without\allowbreak{}Which}} --- The third way (an ability given, received by faith), read as the help without which a person cannot believe, Adam’s kind\par\nobreak\noindent
The third way, with the ability read as Adam’s kind of help.

\begin{derivation}
  \item[(1)] \(\pvar{TestimonyArguments.SolaGratia}{4}\)\gloss{Romans 11:6 — if it is by grace, it is no longer on the basis of works}
  \item[(2)] \(\pvar{TestimonyArguments.SolaGratia}{9}\)\gloss{Genesis 2:16–17 — of the tree of the knowledge of good and evil you shall not eat}
  \item[(3)] \(\pvar{TestimonyArguments.SolaGratia}{14}\)\gloss{God gives the ability to believe, as Adam was made able to stand or fall}
  \breakrow{(4)} \(\pvar{TestimonyArguments.SolaGratia}{15}\)\gloss{Grace is received by faith alone}
  \breakrow{(5)} \(\pvar{TestimonyArguments.SolaGratia}{4} \rightarrow \pvar{TestimonyArguments.SolaGratia}{29}\)\gloss{\emph{would give:} Salvation is by grace and not earned: grace goes before every saving act}
  \item[(6)] \(\pvar{TestimonyArguments.SolaGratia}{14} \rightarrow \lnot \pvar{TestimonyArguments.SolaGratia}{30}\)\gloss{\emph{would deny:} The grace that saves brings about the consent it seeks}
  \item[\(\vdash\)] \(\pvar{TestimonyArguments.SolaGratia}{29} \land \pvar{TestimonyArguments.SolaGratia}{15}\)\gloss{\emph{at issue:} \(\pvar{TestimonyArguments.SolaGratia}{29}\) Salvation is by grace and not earned: grace goes before every saving act; \(\pvar{TestimonyArguments.SolaGratia}{15}\) Grace is received by faith alone}
  \keepwithnext
\end{derivation}
\emph{Grounded in scripture alone:} Romans 11:6 — if it is by grace, it is no longer on the basis of works, Genesis 2:16–17 — of the tree of the knowledge of good and evil you shall not eat.
\par\medskip\noindent\textbf{\texttt{third\allowbreak{}Way\allowbreak{}As\allowbreak{}Aid\allowbreak{}By\allowbreak{}Which}} --- The third way (an ability given, received by faith), read as the help by which a person believes, the saints’ kind\par\nobreak\noindent
The third way, with the ability read as the saints’ kind of help.

\begin{derivation}
  \item[(1)] \(\pvar{TestimonyArguments.SolaGratia}{4}\)\gloss{Romans 11:6 — if it is by grace, it is no longer on the basis of works}
  \item[(2)] \(\pvar{TestimonyArguments.SolaGratia}{9}\)\gloss{Genesis 2:16–17 — of the tree of the knowledge of good and evil you shall not eat}
  \item[(3)] \(\pvar{TestimonyArguments.SolaGratia}{14}\)\gloss{God gives the ability to believe, as Adam was made able to stand or fall}
  \breakrow{(4)} \(\pvar{TestimonyArguments.SolaGratia}{15}\)\gloss{Grace is received by faith alone}
  \breakrow{(5)} \(\pvar{TestimonyArguments.SolaGratia}{4} \rightarrow \pvar{TestimonyArguments.SolaGratia}{29}\)\gloss{\emph{would give:} Salvation is by grace and not earned: grace goes before every saving act}
  \item[(6)] \(\pvar{TestimonyArguments.SolaGratia}{14} \rightarrow \pvar{TestimonyArguments.SolaGratia}{30}\)\gloss{\emph{would give:} The grace that saves brings about the consent it seeks}
  \item[\(\vdash\)] \(\pvar{TestimonyArguments.SolaGratia}{29} \land \pvar{TestimonyArguments.SolaGratia}{15}\)\gloss{\emph{at issue:} \(\pvar{TestimonyArguments.SolaGratia}{29}\) Salvation is by grace and not earned: grace goes before every saving act; \(\pvar{TestimonyArguments.SolaGratia}{15}\) Grace is received by faith alone}
  \keepwithnext
\end{derivation}
\emph{Grounded in scripture alone:} Romans 11:6 — if it is by grace, it is no longer on the basis of works, Genesis 2:16–17 — of the tree of the knowledge of good and evil you shall not eat.
\par\medskip\noindent\textbf{\texttt{third\allowbreak{}Way\allowbreak{}As\allowbreak{}Aid\allowbreak{}Without\allowbreak{}Which\_\allowbreak{}is\_\allowbreak{}satisfiable}}\par\nobreak\noindent
The third way, read as giving Adam’s kind of help, can be held without contradiction: Arminius’s reading satisfies it, and it is the reading Wesley gives.

{\small\begin{verbatim}
theorem thirdWayAsAidWithoutWhich_is_satisfiable : Satisfiable
    thirdWayAsAidWithoutWhich.premises
-- axioms: propext, Quot.sound
\end{verbatim}}
\par\medskip\noindent\textbf{\texttt{third\allowbreak{}Way\allowbreak{}As\allowbreak{}Aid\allowbreak{}By\allowbreak{}Which\_\allowbreak{}is\_\allowbreak{}satisfiable}}\par\nobreak\noindent
The third way, read as giving the saints’ kind of help, can be held without contradiction: Dort’s reading satisfies it, and it is the reading Dort gives.

{\small\begin{verbatim}
theorem thirdWayAsAidByWhich_is_satisfiable : Satisfiable
    thirdWayAsAidByWhich.premises
-- axioms: propext, Quot.sound
\end{verbatim}}
\par\medskip\noindent\textbf{\texttt{remonstrants\_\allowbreak{}do\_\allowbreak{}not\_\allowbreak{}reach\_\allowbreak{}adams\_\allowbreak{}kind}}\par\nobreak\noindent
The Remonstrants hold together with the first reading: Arminius’s reading satisfies both.

{\small\begin{verbatim}
theorem remonstrants_do_not_reach_adams_kind : Satisfiable
    (remonstrantPosition.premises ++ thirdWayAsAidWithoutWhich.premises)
-- axioms: propext, Quot.sound
\end{verbatim}}
\par\medskip\noindent\textbf{\texttt{dort\_\allowbreak{}does\_\allowbreak{}not\_\allowbreak{}reach\_\allowbreak{}the\_\allowbreak{}saints\_\allowbreak{}kind}}\par\nobreak\noindent
Dort holds together with the second reading: Dort’s reading satisfies both.

{\small\begin{verbatim}
theorem dort_does_not_reach_the_saints_kind : Satisfiable (reformed.premises
    ++ thirdWayAsAidByWhich.premises)
-- axioms: propext, Quot.sound
\end{verbatim}}
\par\medskip\noindent\textbf{\texttt{where\allowbreak{}Adams\allowbreak{}Kind\allowbreak{}Falls}}\par\nobreak\noindent
\textbf{Why Dort stands against the third way, read as Adam’s kind of help.} The crux is Dort’s step from God’s work alone to grace effectual of itself, that is, grace that brings about the consent it seeks and so cannot be finally refused; the Lutherans dispute it. With the Romans strand granted, the reading of Romans 9:16 as being about God’s mercy in saving, the crux breaks the third way at two of its premises: the ability given, and the reading of it as a help that does not bring believing about.

\paragraph{Why \emph{Reformed (the Canons of Dort)} stands against \emph{The third way (an ability given, received by faith), read as the help without which a person cannot believe, Adam’s kind}.}
\begin{itemize}
\item \textbf{The crux:} \(\pvar{TestimonyArguments.SolaGratia}{28} \rightarrow \pvar{TestimonyArguments.SolaGratia}{30}\), a premise of \emph{Reformed (the Canons of Dort)}.
\item \textbf{What it does:} its conclusion does not follow without it.
\item \textbf{Granted:} \(\pvar{TestimonyArguments.SolaGratia}{3}\), \(\pvar{TestimonyArguments.SolaGratia}{23}\), \((\pvar{TestimonyArguments.SolaGratia}{3} \land \pvar{TestimonyArguments.SolaGratia}{23}) \rightarrow \pvar{TestimonyArguments.SolaGratia}{28}\), which the break also needs.
\item \textbf{Where the rival breaks:} \(\pvar{TestimonyArguments.SolaGratia}{14}\), \(\pvar{TestimonyArguments.SolaGratia}{14} \rightarrow \lnot \pvar{TestimonyArguments.SolaGratia}{30}\) cannot be held together with \(\pvar{TestimonyArguments.SolaGratia}{28} \rightarrow \pvar{TestimonyArguments.SolaGratia}{30}\) and \(\pvar{TestimonyArguments.SolaGratia}{3}\), \(\pvar{TestimonyArguments.SolaGratia}{23}\), \((\pvar{TestimonyArguments.SolaGratia}{3} \land \pvar{TestimonyArguments.SolaGratia}{23}) \rightarrow \pvar{TestimonyArguments.SolaGratia}{28}\); each is needed for the break, and without the crux they stand.
\item \textbf{What this rests on:}
\begin{itemize}
\item \(\pvar{TestimonyArguments.SolaGratia}{3}\) Romans 9:16 — it depends not on human will or exertion, but on God, who has mercy --- \emph{consensus}: Rom 9:16
\item \(\pvar{TestimonyArguments.SolaGratia}{23}\) Romans 9:16 is about God’s mercy in saving, not only in choosing a people for history --- \emph{well supported}: \cite[527--565]{arminius-works-3-1853}; \cite[351]{abasciano-corporate-election-romans-9-2006}
\item \(\pvar{TestimonyArguments.SolaGratia}{28}\) Conversion is God’s work alone: God regenerates, and the will does nothing to it --- \emph{disputed}: \cite[Canons of Dort, III/IV, arts. 10–12]{schaff-creeds-christendom-3-1977}
\item the step itself --- \emph{consensus}: \cite[Trent, Session VI (1547), ch. 8]{tanner-decrees-1990}; \cite[36]{waterworth-trent-1848}; \cite[Five Arminian Articles (1610), art. 4]{schaff-creeds-christendom-3-1977}; \cite[Canons of Dort, III/IV, art. 10]{schaff-creeds-christendom-3-1977}; \emph{disputed}: \cite[ad loc. Eph 2:8]{robertson-word-pictures-4-1931}; \cite[228--229]{calvin-galatians-ephesians-1854}; \emph{disputed}: \cite[600]{arminius-works-1-1825}; \cite[681]{arminius-works-1-1825}; \cite[Five Arminian Articles (1610), art. 4]{schaff-creeds-christendom-3-1977}; \emph{disputed}: \cite[527--565]{arminius-works-3-1853}; \emph{consensus}: \cite[593]{schaff-creeds-christendom-3-1977}; \cite[301]{calvin-galatians-ephesians-1854}; \cite[Solid Declaration II.60]{concordia-triglotta-1921}; \emph{well supported}: \cite[Solid Declaration II.89]{concordia-triglotta-1921}; \cite[Canons of Dort, III/IV, art. 12]{schaff-creeds-christendom-3-1977}; \emph{disputed}: \cite[Solid Declaration II.89]{concordia-triglotta-1921}; \cite[Solid Declaration II.60]{concordia-triglotta-1921}; \cite[Solid Declaration XI.41]{concordia-triglotta-1921}
\end{itemize}
\end{itemize}
\par\medskip\noindent\textbf{\texttt{where\allowbreak{}The\allowbreak{}Saints\allowbreak{}Kind\allowbreak{}Falls}}\par\nobreak\noindent
\textbf{Why Trent stands against the third way, read as the saints’ kind of help.} The crux is the second half of Trent’s canon 4: the will, moved by God, can refuse its consent if it would, so grace is not effectual of itself. With that power to refuse granted, it breaks the third way at two of its premises: the ability given, and the reading of it as a help that brings believing about. The crux is rated \emph{disputed}, because Aquinas grants that the will moved by God could, taken simply, refuse, and holds that the man God moves to grace “will infallibly attain to it”.

The Remonstrants stand against this reading too, but not at a premise either side states: their premises deliver that the grace that saves can be finally refused, and that excludes grace which brings believing about by what the two claims mean. The meanings hold that join, and every dispute over the argument holds it; a break at a premise, which is what this records, is Trent’s.

\paragraph{Why \emph{Trent (Session VI, canon 4, on the Molinist reading)} stands against \emph{The third way (an ability given, received by faith), read as the help by which a person believes, the saints’ kind}.}
\begin{itemize}
\item \textbf{The crux:} \(\pvar{TestimonyArguments.SolaGratia}{27} \rightarrow \lnot \pvar{TestimonyArguments.SolaGratia}{30}\), a premise of \emph{Trent (Session VI, canon 4, on the Molinist reading)}.
\item \textbf{What it does:} its conclusion does not follow without it.
\item \textbf{Granted:} \(\pvar{TestimonyArguments.SolaGratia}{27}\), which the break also needs.
\item \textbf{Where the rival breaks:} \(\pvar{TestimonyArguments.SolaGratia}{14}\), \(\pvar{TestimonyArguments.SolaGratia}{14} \rightarrow \pvar{TestimonyArguments.SolaGratia}{30}\) cannot be held together with \(\pvar{TestimonyArguments.SolaGratia}{27} \rightarrow \lnot \pvar{TestimonyArguments.SolaGratia}{30}\) and \(\pvar{TestimonyArguments.SolaGratia}{27}\); each is needed for the break, and without the crux they stand.
\item \textbf{What this rests on:}
\begin{itemize}
\item \(\pvar{TestimonyArguments.SolaGratia}{27}\) The free will, moved by God, can refuse its assent if it wishes --- \emph{disputed}: \cite[Trent, Session VI (1547), canon 4]{tanner-decrees-1990}; \cite[45]{waterworth-trent-1848}
\item the step itself --- \emph{consensus}: \cite[Trent, Session VI (1547), ch. 8]{tanner-decrees-1990}; \cite[36]{waterworth-trent-1848}; \cite[Five Arminian Articles (1610), art. 4]{schaff-creeds-christendom-3-1977}; \cite[Canons of Dort, III/IV, art. 10]{schaff-creeds-christendom-3-1977}; \emph{disputed}: \cite[I-II q.111 a.2]{aquinas-summa-1920}; \cite[Canons of Dort, III/IV, art. 12]{schaff-creeds-christendom-3-1977}; \cite[Trent, Session VI (1547), canon 4]{tanner-decrees-1990}; \cite[45]{waterworth-trent-1848}; \emph{disputed}: \cite[I-II q.10 a.4 ad 3]{aquinas-summa-1920}; \cite[I-II q.112 a.3]{aquinas-summa-1920}; \cite[Trent, Session VI (1547), canon 4]{tanner-decrees-1990}; \cite[45]{waterworth-trent-1848}
\end{itemize}
\end{itemize}
\par\medskip\noindent\textbf{\texttt{what\allowbreak{}The\allowbreak{}Ability\allowbreak{}Is}}\par\nobreak\noindent
\textbf{What the ability God gives is, read both ways.} The claim is the third way’s: God gives the one who hears the gospel the ability to believe, as Adam was made able to stand or to fall.

Read as Adam’s kind of help, the help “without which” a person cannot believe, it is the reading Wesley gives, and the Remonstrants can hold it without contradiction. It cannot be held together with Dort’s position. The break comes at Dort’s step from God’s work alone to grace effectual of itself, and that step is rated \emph{disputed}, because the Lutherans grant its ground and deny its conclusion; so the break does not show Dort right.

Read as the saints’ kind, the help “by which” a person believes, it is Dort’s reading, and Dort can hold it without contradiction. It cannot be held together with Trent’s canon 4, at the step from the will’s power to refuse, which is rated \emph{disputed} too: Aquinas grants the power and holds grace effectual all the same. Nor can it be held with the Remonstrants, though there the break is in what the claims mean rather than at a premise either side states.

Each reading is fair: the third way read either way can be held without contradiction. Which reading its holder means is not something the dilemma decides. What it shows is that the third way takes a side the moment it says which kind of help God gives, and which side depends on the answer.

\paragraph{Every reading of \(P_{14}\), answered.} \emph{The third way (an ability given, received by faith)} holds \(\pvar{TestimonyArguments.SolaGratia}{14}\). It is read 2 ways here, and each reading is checked; none can be left out.

\paragraph{1. Read as the help without which a person cannot believe, Adam’s kind.} The claim commits \emph{The third way (an ability given, received by faith)} to \(\lnot \pvar{TestimonyArguments.SolaGratia}{30}\) --- so read by \cite[397]{wesley-works-9-1827}; \cite[485]{npnf1-05-1887} (\emph{disputed}).

Against \emph{Reformed (the Canons of Dort)}, it cannot be held. Where it breaks:
\paragraph{Why \emph{Reformed (the Canons of Dort)} stands against \emph{The third way (an ability given, received by faith), read as the help without which a person cannot believe, Adam’s kind}.}
\begin{itemize}
\item \textbf{The crux:} \(\pvar{TestimonyArguments.SolaGratia}{28} \rightarrow \pvar{TestimonyArguments.SolaGratia}{30}\), a premise of \emph{Reformed (the Canons of Dort)}.
\item \textbf{What it does:} its conclusion does not follow without it.
\item \textbf{Granted:} \(\pvar{TestimonyArguments.SolaGratia}{3}\), \(\pvar{TestimonyArguments.SolaGratia}{23}\), \((\pvar{TestimonyArguments.SolaGratia}{3} \land \pvar{TestimonyArguments.SolaGratia}{23}) \rightarrow \pvar{TestimonyArguments.SolaGratia}{28}\), which the break also needs.
\item \textbf{Where the rival breaks:} \(\pvar{TestimonyArguments.SolaGratia}{14}\), \(\pvar{TestimonyArguments.SolaGratia}{14} \rightarrow \lnot \pvar{TestimonyArguments.SolaGratia}{30}\) cannot be held together with \(\pvar{TestimonyArguments.SolaGratia}{28} \rightarrow \pvar{TestimonyArguments.SolaGratia}{30}\) and \(\pvar{TestimonyArguments.SolaGratia}{3}\), \(\pvar{TestimonyArguments.SolaGratia}{23}\), \((\pvar{TestimonyArguments.SolaGratia}{3} \land \pvar{TestimonyArguments.SolaGratia}{23}) \rightarrow \pvar{TestimonyArguments.SolaGratia}{28}\); each is needed for the break, and without the crux they stand.
\item \textbf{What this rests on:}
\begin{itemize}
\item \(\pvar{TestimonyArguments.SolaGratia}{3}\) Romans 9:16 — it depends not on human will or exertion, but on God, who has mercy --- \emph{consensus}: Rom 9:16
\item \(\pvar{TestimonyArguments.SolaGratia}{23}\) Romans 9:16 is about God’s mercy in saving, not only in choosing a people for history --- \emph{well supported}: \cite[527--565]{arminius-works-3-1853}; \cite[351]{abasciano-corporate-election-romans-9-2006}
\item \(\pvar{TestimonyArguments.SolaGratia}{28}\) Conversion is God’s work alone: God regenerates, and the will does nothing to it --- \emph{disputed}: \cite[Canons of Dort, III/IV, arts. 10–12]{schaff-creeds-christendom-3-1977}
\item the step itself --- \emph{consensus}: \cite[Trent, Session VI (1547), ch. 8]{tanner-decrees-1990}; \cite[36]{waterworth-trent-1848}; \cite[Five Arminian Articles (1610), art. 4]{schaff-creeds-christendom-3-1977}; \cite[Canons of Dort, III/IV, art. 10]{schaff-creeds-christendom-3-1977}; \emph{disputed}: \cite[ad loc. Eph 2:8]{robertson-word-pictures-4-1931}; \cite[228--229]{calvin-galatians-ephesians-1854}; \emph{disputed}: \cite[600]{arminius-works-1-1825}; \cite[681]{arminius-works-1-1825}; \cite[Five Arminian Articles (1610), art. 4]{schaff-creeds-christendom-3-1977}; \emph{disputed}: \cite[527--565]{arminius-works-3-1853}; \emph{consensus}: \cite[593]{schaff-creeds-christendom-3-1977}; \cite[301]{calvin-galatians-ephesians-1854}; \cite[Solid Declaration II.60]{concordia-triglotta-1921}; \emph{well supported}: \cite[Solid Declaration II.89]{concordia-triglotta-1921}; \cite[Canons of Dort, III/IV, art. 12]{schaff-creeds-christendom-3-1977}; \emph{disputed}: \cite[Solid Declaration II.89]{concordia-triglotta-1921}; \cite[Solid Declaration II.60]{concordia-triglotta-1921}; \cite[Solid Declaration XI.41]{concordia-triglotta-1921}
\end{itemize}
\end{itemize}
Against \emph{Remonstrant (the Five Articles, 1610)}, it is not reached: \emph{Remonstrant (the Five Articles, 1610)} holds, and can be held together with it.

\paragraph{2. Read as the help by which a person believes, the saints’ kind.} The claim commits \emph{The third way (an ability given, received by faith)} to \(\pvar{TestimonyArguments.SolaGratia}{30}\) --- so read by \cite[485]{npnf1-05-1887}; \cite[591]{schaff-creeds-christendom-3-1977} (\emph{disputed}).

Against \emph{Trent (Session VI, canon 4, on the Molinist reading)}, it cannot be held. Where it breaks:
\paragraph{Why \emph{Trent (Session VI, canon 4, on the Molinist reading)} stands against \emph{The third way (an ability given, received by faith), read as the help by which a person believes, the saints’ kind}.}
\begin{itemize}
\item \textbf{The crux:} \(\pvar{TestimonyArguments.SolaGratia}{27} \rightarrow \lnot \pvar{TestimonyArguments.SolaGratia}{30}\), a premise of \emph{Trent (Session VI, canon 4, on the Molinist reading)}.
\item \textbf{What it does:} its conclusion does not follow without it.
\item \textbf{Granted:} \(\pvar{TestimonyArguments.SolaGratia}{27}\), which the break also needs.
\item \textbf{Where the rival breaks:} \(\pvar{TestimonyArguments.SolaGratia}{14}\), \(\pvar{TestimonyArguments.SolaGratia}{14} \rightarrow \pvar{TestimonyArguments.SolaGratia}{30}\) cannot be held together with \(\pvar{TestimonyArguments.SolaGratia}{27} \rightarrow \lnot \pvar{TestimonyArguments.SolaGratia}{30}\) and \(\pvar{TestimonyArguments.SolaGratia}{27}\); each is needed for the break, and without the crux they stand.
\item \textbf{What this rests on:}
\begin{itemize}
\item \(\pvar{TestimonyArguments.SolaGratia}{27}\) The free will, moved by God, can refuse its assent if it wishes --- \emph{disputed}: \cite[Trent, Session VI (1547), canon 4]{tanner-decrees-1990}; \cite[45]{waterworth-trent-1848}
\item the step itself --- \emph{consensus}: \cite[Trent, Session VI (1547), ch. 8]{tanner-decrees-1990}; \cite[36]{waterworth-trent-1848}; \cite[Five Arminian Articles (1610), art. 4]{schaff-creeds-christendom-3-1977}; \cite[Canons of Dort, III/IV, art. 10]{schaff-creeds-christendom-3-1977}; \emph{disputed}: \cite[I-II q.111 a.2]{aquinas-summa-1920}; \cite[Canons of Dort, III/IV, art. 12]{schaff-creeds-christendom-3-1977}; \cite[Trent, Session VI (1547), canon 4]{tanner-decrees-1990}; \cite[45]{waterworth-trent-1848}; \emph{disputed}: \cite[I-II q.10 a.4 ad 3]{aquinas-summa-1920}; \cite[I-II q.112 a.3]{aquinas-summa-1920}; \cite[Trent, Session VI (1547), canon 4]{tanner-decrees-1990}; \cite[45]{waterworth-trent-1848}
\end{itemize}
\end{itemize}
Against \emph{Reformed (the Canons of Dort)}, it is not reached: \emph{Reformed (the Canons of Dort)} holds, and can be held together with it.

\subsection*{Aquinas and the Catholic schools}
\addcontentsline{toc}{subsection}{Aquinas and the Catholic schools}
Is sola gratia, stated sharply, a Protestant doctrine? Thomas Aquinas wrote, three centuries before Dort, that “if God intends, while moving, that the one whose heart He moves should attain to grace, he will infallibly attain to it” (\emph{Summa Theologiae}, First Part of the Second Part, question 112, article 3; cited below as I-II q.112 a.3). And in the grace by which the will is converted, “our mind is moved and does not move, but … God is the sole mover” (I-II q.111 a.2). Aquinas calls that \emph{operating grace}: grace in which God alone moves the will. In \emph{cooperating grace} the will, once moved, moves as well.

The Dominican school that followed him, Domingo Báñez in the sixteenth century and Réginald Garrigou-Lagrange in the twentieth, held with Garrigou-Lagrange that grace is “efficacious of itself, producing of itself the virtuous act” (\emph{Grace}, ch. V): when it is given, the grace itself, by its own power, brings about the will’s consent. This argument calls that grace \emph{effectual of itself}, the second part of sola gratia.

The Jesuit school of Luis de Molina read Trent otherwise. For the Molinists, as Joseph Pohle reports them, “the infallible success of efficacious grace is not to be ascribed to its own intrinsic nature, but to” God’s “middle knowledge” (\emph{scientia media}), his knowledge of what a person would freely do in any circumstance; “it is rather the will itself which by its consent” makes grace efficacious. So for them efficacious grace is infallibly followed by consent, but because God foresaw the consent, not because the grace produces it. The two schools argued before the popes for twenty years. Paul V’s decree of 5 September 1607 “allowed each party to defend its own doctrine, enjoined each from censoring or condemning the opposite opinion” (Astrain, “Congregatio de Auxiliis”), and Rome has not decided between them since. The Trent of this argument’s other chapters is Trent on the Molinist reading of its fourth canon.

\subsubsection*{What the Thomists reach}
The Thomists accept Trent’s fourth canon: the will, moved by God, cooperates by assenting and can dissent if it would. But they read it so that it denies neither part of sola gratia that Trent seems to deny. Garrigou-Lagrange explains the canon so, in his own translation of it: “The Council does not maintain that man does, in fact, sometimes dissent, but that ‘he can dissent if he so wills.’ In other words, the contrary power remains, but under efficacious grace man never wills to resist, nor does he.” The schoolmen call this the \emph{divided sense}: taken apart from grace, the will can dissent; in the \emph{composite sense}, taken together with efficacious grace, it does not. So the step from the canon to “grace is not effectual of itself”, rated \emph{disputed} because Aquinas grants the canon’s ground and denies that it follows, is a step the Thomists refuse.

Grant the Thomists their premises, and the second part of sola gratia follows: the grace that saves brings about the consent it seeks (\texttt{thomists\_\allowbreak{}reach\_\allowbreak{}effectual\_\allowbreak{}grace}). The step that delivers it is rated \emph{disputed}, because the Molinists grant God’s intention and the infallible result and ascribe the efficacy to the consent God foresees.

Two Thomist claims are about resisting grace, and neither is the claim that the grace that saves can be finally refused. Grace “really sufficient in its own order” (Garrigou-Lagrange, after Lemos) is resisted, and through that resistance the efficacious grace offered with it is withheld. And justifying grace, once given, can be lost: Aquinas holds that the charity of one still on the way (in this life) can be lost, and Trent condemns anyone who says otherwise (canon 23). The grace resisted in the first is not the grace that converts, and the grace lost in the second is grace once given. So they are encoded as claims of their own. Dort parts from the Thomists on the second, the perseverance of the saints, which the chapter on perseverance weighs, with Aquinas beside Augustine.

\subsubsection*{The word “passive”}
The Thomists’ operating grace gives them that conversion is God’s work alone. What they do not take is Dort’s next step, from God’s work alone to the will’s passivity. Instead they say the will is “not merely passive”. Garrigou-Lagrange: under operating grace “the will elicits its act vitally, in fact, it freely consents to the divine motion or inspiration, but it does not strictly move itself”, and “in these acts the will does not remain passive”. He lists Calvin’s “merely passive” among the errors. Taken at their face, these words deny the third part of sola gratia.

But he is exact about the sense. The will’s consent, “although as a disposition it precedes it [grace] in the genus or order of material cause, it nevertheless follows it in the genus or order of formal, efficient, and final cause”. These are Aristotle’s four kinds of cause: the material, what receives a thing; the efficient, what brings it about; the formal, what makes it what it is; the final, what it is for. So the consent comes before grace only as the readiness that receives it, and grace is what brings the consent about. Aquinas says the same: the movement of the free will “is not the cause of grace, but the effect”. The third part of sola gratia, as this argument states it, says that the will is not the agent of its own new life. On Garrigou-Lagrange’s own explanation the Thomist grants it in that sense; but that is a reading of him, rated \emph{disputed}, not something his premises settle.

If “not merely passive” makes the will’s consent a cause of the new life, as the Molinists hold of the will’s part, the Thomist breaks with Dort at Dort’s step from God’s work alone to the will’s passivity. If it means the will acts while God alone gives the new life, as Garrigou-Lagrange explains it, the Thomist can stand with Dort, and breaks instead with Trent on the Molinist reading, at the step from the will’s cooperation, which Aquinas already disputes (\texttt{what\allowbreak{}Not\allowbreak{}Merely\allowbreak{}Passive\allowbreak{}Means}). Granted the Thomists’ premises and no more, the third part neither follows nor fails (\texttt{thomists\_\allowbreak{}leave\_\allowbreak{}passivity\_\allowbreak{}open}).

\subsubsection*{So is sola gratia a Protestant doctrine?}
The Thomists’ premises reach the first two parts of sola gratia and leave the third open. Nothing here shows that they hold sola gratia as stated, and nothing shows that they deny it. Weighed with the other positions, the Thomists can be defended with Dort and are not forced (\texttt{thomist\_\allowbreak{}defensible}, \texttt{thomist\_\allowbreak{}not\_\allowbreak{}forced}). Where the Thomists plainly part from the Reformers is elsewhere: on faith formed by charity, which Aquinas calls “the form of faith”, and so not faith alone, which Dort’s position here does not carry; and on perseverance. Neither is part of sola gratia as this argument states it.

\subsubsection*{How this is checked}
\paragraph*{The Thomist position}
\leavevmode\par\nobreak
\par\medskip\noindent\textbf{\texttt{thomist\allowbreak{}Cite}}\par\nobreak\noindent
The Thomist citations: Aquinas, with Trent’s ninth canon, for the denial of faith alone, faith being formed by charity (II-II q.4 a.3).

{\small\begin{verbatim}
def thomistCite : Claim → AtomMeta
\end{verbatim}}
\par\medskip\noindent\textbf{\texttt{thomist}} --- Thomist (Aquinas; Garrigou-Lagrange)\par\nobreak\noindent
\textbf{The Thomists}, from Aquinas (\emph{Summa} I-II qq.109–113) and Garrigou-Lagrange. They hold Romans 11:6, and operating grace, in which God is the sole mover. They hold that God moves the will intending it to attain grace, so that it attains it infallibly. They accept Trent’s canon 4, that the will cooperates and can dissent, but draw neither of the steps that Trent’s Molinist reading draws from it. They hold that the will under operating grace is not merely passive, that grace truly sufficient is resisted, that justifying grace can be lost, and that grace is not received by faith alone, faith being formed by charity.

The conclusion set against these premises is sola gratia as this argument states it, to test how far they reach. It is not the Thomists’ own conclusion, which is \texttt{thomist\allowbreak{}Position} below.

\begin{derivation}
  \item[(1)] \(\pvar{TestimonyArguments.SolaGratia}{4}\)\gloss{Romans 11:6 — if it is by grace, it is no longer on the basis of works}
  \item[(2)] \(\pvar{TestimonyArguments.SolaGratia}{16}\)\gloss{In converting grace God is the sole mover: the mind is moved and does not move}
  \item[(3)] \(\pvar{TestimonyArguments.SolaGratia}{17}\)\gloss{God moves the will intending that it attain grace}
  \breakrow{(4)} \(\pvar{TestimonyArguments.SolaGratia}{26}\)\gloss{The free will, moved by God, cooperates by assenting to God’s call}
  \breakrow{(5)} \(\pvar{TestimonyArguments.SolaGratia}{27}\)\gloss{The free will, moved by God, can refuse its assent if it wishes}
  \breakrow{(6)} \(\pvar{TestimonyArguments.SolaGratia}{18}\)\gloss{Under operating grace the will is not merely passive: it freely consents}
  \breakrow{(7)} \(\pvar{TestimonyArguments.SolaGratia}{19}\)\gloss{Grace truly sufficient, giving the power to do good, is resisted}
  \breakrow{(8)} \(\pvar{TestimonyArguments.SolaGratia}{20}\)\gloss{Justifying grace, once given, can be lost}
  \breakrow{(9)} \(\lnot \pvar{TestimonyArguments.SolaGratia}{15}\)\gloss{\emph{denied:} Grace is received by faith alone}
  \breakrow{(10)} \(\pvar{TestimonyArguments.SolaGratia}{4} \rightarrow \pvar{TestimonyArguments.SolaGratia}{29}\)\gloss{\emph{would give:} Salvation is by grace and not earned: grace goes before every saving act}
  \breakrow{(11)} \(\pvar{TestimonyArguments.SolaGratia}{16} \rightarrow \pvar{TestimonyArguments.SolaGratia}{28}\)\gloss{\emph{would give:} Conversion is God’s work alone: God regenerates, and the will does nothing to it}
  \item[(12)] \(\pvar{TestimonyArguments.SolaGratia}{17} \rightarrow \pvar{TestimonyArguments.SolaGratia}{30}\)\gloss{\emph{would give:} The grace that saves brings about the consent it seeks}
  \item[\(\vdash\)] \(\pvar{TestimonyArguments.SolaGratia}{29} \land \pvar{TestimonyArguments.SolaGratia}{30} \land \pvar{TestimonyArguments.SolaGratia}{31}\)\gloss{\emph{at issue:} \(\pvar{TestimonyArguments.SolaGratia}{29}\) Salvation is by grace and not earned: grace goes before every saving act; \(\pvar{TestimonyArguments.SolaGratia}{30}\) The grace that saves brings about the consent it seeks; \(\pvar{TestimonyArguments.SolaGratia}{31}\) In being made alive the will is passive, and contributes nothing to it}
  \keepwithnext
\end{derivation}
\emph{Grounded in scripture alone:} Romans 11:6 — if it is by grace, it is no longer on the basis of works.
\par\medskip\noindent\textbf{\texttt{thomist\allowbreak{}Position}} --- Thomist (Aquinas; Garrigou-Lagrange)\par\nobreak\noindent
The Thomists’ position, concluding what their sources draw: salvation by grace, and grace effectual of itself.

\begin{derivation}
  \item[(1)] \(\pvar{TestimonyArguments.SolaGratia}{4}\)\gloss{Romans 11:6 — if it is by grace, it is no longer on the basis of works}
  \item[(2)] \(\pvar{TestimonyArguments.SolaGratia}{16}\)\gloss{In converting grace God is the sole mover: the mind is moved and does not move}
  \item[(3)] \(\pvar{TestimonyArguments.SolaGratia}{17}\)\gloss{God moves the will intending that it attain grace}
  \breakrow{(4)} \(\pvar{TestimonyArguments.SolaGratia}{26}\)\gloss{The free will, moved by God, cooperates by assenting to God’s call}
  \breakrow{(5)} \(\pvar{TestimonyArguments.SolaGratia}{27}\)\gloss{The free will, moved by God, can refuse its assent if it wishes}
  \breakrow{(6)} \(\pvar{TestimonyArguments.SolaGratia}{18}\)\gloss{Under operating grace the will is not merely passive: it freely consents}
  \breakrow{(7)} \(\pvar{TestimonyArguments.SolaGratia}{19}\)\gloss{Grace truly sufficient, giving the power to do good, is resisted}
  \breakrow{(8)} \(\pvar{TestimonyArguments.SolaGratia}{20}\)\gloss{Justifying grace, once given, can be lost}
  \breakrow{(9)} \(\lnot \pvar{TestimonyArguments.SolaGratia}{15}\)\gloss{\emph{denied:} Grace is received by faith alone}
  \breakrow{(10)} \(\pvar{TestimonyArguments.SolaGratia}{4} \rightarrow \pvar{TestimonyArguments.SolaGratia}{29}\)\gloss{\emph{would give:} Salvation is by grace and not earned: grace goes before every saving act}
  \breakrow{(11)} \(\pvar{TestimonyArguments.SolaGratia}{16} \rightarrow \pvar{TestimonyArguments.SolaGratia}{28}\)\gloss{\emph{would give:} Conversion is God’s work alone: God regenerates, and the will does nothing to it}
  \item[(12)] \(\pvar{TestimonyArguments.SolaGratia}{17} \rightarrow \pvar{TestimonyArguments.SolaGratia}{30}\)\gloss{\emph{would give:} The grace that saves brings about the consent it seeks}
  \item[\(\vdash\)] \(\pvar{TestimonyArguments.SolaGratia}{29} \land \pvar{TestimonyArguments.SolaGratia}{30}\)\gloss{\emph{at issue:} \(\pvar{TestimonyArguments.SolaGratia}{29}\) Salvation is by grace and not earned: grace goes before every saving act; \(\pvar{TestimonyArguments.SolaGratia}{30}\) The grace that saves brings about the consent it seeks}
  \keepwithnext
\end{derivation}
\emph{Grounded in scripture alone:} Romans 11:6 — if it is by grace, it is no longer on the basis of works.
\par\medskip\noindent\textbf{\texttt{thomist\allowbreak{}Reading}}\par\nobreak\noindent
The Thomists’ reading, with “not merely passive” taken as the Molinists take it: grace effectual of itself, God alone the one who regenerates, the will not passive in it, and not faith alone.

{\small\begin{verbatim}
def thomistReading : Valuation Claim :=
  fun a =>
    match a with
    | Claim.willIsPassiveInRegeneration => False
    | Claim.faithAloneReceivesGrace => False
    | Claim.savingGraceCanBeFinallyRefused => False
    | x => True
\end{verbatim}}
\par\medskip\noindent\textbf{\texttt{thomist\allowbreak{}With\allowbreak{}Dort\allowbreak{}Reading}}\par\nobreak\noindent
The Thomists’ reading, with “not merely passive” taken as Garrigou-Lagrange explains it, so that the will is not the agent of its new life; and Dort’s answer to Acts 7:51 granted besides, so that Dort’s premises hold with it.

{\small\begin{verbatim}
def thomistWithDortReading : Valuation Claim :=
  fun a =>
    match a with
    | Claim.faithAloneReceivesGrace => False
    | Claim.savingGraceCanBeFinallyRefused => False
    | Claim.acts7_51ResistsSavingGrace => False
    | x => True
\end{verbatim}}
\paragraph*{What the Thomist reaches}
\leavevmode\par\nobreak
\par\medskip\noindent\textbf{\texttt{thomists\_\allowbreak{}reach\_\allowbreak{}effectual\_\allowbreak{}grace}}\par\nobreak\noindent
\textbf{The Thomists’ premises reach grace effectual of itself}, the second part of sola gratia. Aquinas’s “if God intends, while moving … he will infallibly attain to it”, read as Garrigou-Lagrange reads it, delivers that the grace that saves brings about the consent it seeks. The step is rated \texttt{disputed}: the Molinists grant the intention and the infallible result, and ascribe the efficacy to the consent God foresees.

{\small\begin{verbatim}
theorem thomists_reach_effectual_grace : Entails thomist.premises (p
    Claim.graceIsEffectualOfItself)
-- axioms: propext, Quot.sound
\end{verbatim}}
\par\medskip\noindent\textbf{\texttt{thomist\allowbreak{}Position\_\allowbreak{}establishes}}\par\nobreak\noindent
The Thomists’ premises deliver their position.

{\small\begin{verbatim}
theorem thomistPosition_establishes : Establishes thomistPosition
-- axioms: propext, Quot.sound
\end{verbatim}}
\par\medskip\noindent\textbf{\texttt{thomist\allowbreak{}Position\_\allowbreak{}is\_\allowbreak{}satisfiable}}\par\nobreak\noindent
The Thomists’ position can be held without contradiction.

{\small\begin{verbatim}
theorem thomistPosition_is_satisfiable : Satisfiable thomistPosition.premises
-- axioms: propext, Quot.sound
\end{verbatim}}
\par\medskip\noindent\textbf{\texttt{thomists\_\allowbreak{}leave\_\allowbreak{}passivity\_\allowbreak{}open}}\par\nobreak\noindent
\textbf{Granted the Thomists’ premises, the will’s passivity neither follows nor fails.} Their operating grace gives them that conversion is God’s work alone. What they do not take is Dort’s next step, from God’s work alone to the will’s passivity, which is rated \emph{well supported}. Instead they hold that the will is “not merely passive”, which can be read either way: as denying the third part of sola gratia, or as consistent with it. Their premises alone do not choose: a reading that takes the will as passive in the sense of not being the agent of its new life satisfies them, and so does one that does not. So the Thomists do not reach sola gratia as stated, and they do not refute it either.

{\small\begin{verbatim}
theorem thomists_leave_passivity_open : Independent thomist.premises (p
    Claim.willIsPassiveInRegeneration)
-- axioms: propext, Quot.sound
\end{verbatim}}
\paragraph*{What “not merely passive” means}
\leavevmode\par\nobreak
\par\medskip\noindent\textbf{\texttt{as\allowbreak{}Agent\allowbreak{}Of\allowbreak{}New\allowbreak{}Life}}\par\nobreak\noindent
“Not merely passive” read as making the will’s consent a cause of the new life, as the Molinists hold of the will’s part, for whom “it is rather the will itself which by its consent” makes grace efficacious (Pohle). No Thomist reads his own claim so; the horn asks what follows if the phrase meant this.

{\small\begin{verbatim}
def asAgentOfNewLife : Reading Claim :=
  {
    name :=
      "as making the will's consent a cause of the new life, in the Molinists' sense",
    commits := notP Claim.willIsPassiveInRegeneration,
    source := asAgentOfNewLifeSource }
\end{verbatim}}
\par\medskip\noindent\textbf{\texttt{as\allowbreak{}Disposition}}\par\nobreak\noindent
“Not merely passive” read as the will acting while God alone is the agent of the new life: the will is passive in the sense the third part of sola gratia means. Garrigou-Lagrange’s reading.

{\small\begin{verbatim}
def asDisposition : Reading Claim :=
  {
    name :=
      "as the will consenting while God alone is the agent, Garrigou-Lagrange's sense",
    commits := p Claim.willIsPassiveInRegeneration,
    source := asDispositionSource }
\end{verbatim}}
\par\medskip\noindent\textbf{\texttt{thomist\allowbreak{}As\allowbreak{}Agent}} --- Thomist (Aquinas; Garrigou-Lagrange), read as making the will’s consent a cause of the new life, in the Molinists’ sense\par\nobreak\noindent
The Thomists, with “not merely passive” read in the Molinists’ sense.

\begin{derivation}
  \item[(1)] \(\pvar{TestimonyArguments.SolaGratia}{4}\)\gloss{Romans 11:6 — if it is by grace, it is no longer on the basis of works}
  \item[(2)] \(\pvar{TestimonyArguments.SolaGratia}{16}\)\gloss{In converting grace God is the sole mover: the mind is moved and does not move}
  \item[(3)] \(\pvar{TestimonyArguments.SolaGratia}{17}\)\gloss{God moves the will intending that it attain grace}
  \breakrow{(4)} \(\pvar{TestimonyArguments.SolaGratia}{26}\)\gloss{The free will, moved by God, cooperates by assenting to God’s call}
  \breakrow{(5)} \(\pvar{TestimonyArguments.SolaGratia}{27}\)\gloss{The free will, moved by God, can refuse its assent if it wishes}
  \breakrow{(6)} \(\pvar{TestimonyArguments.SolaGratia}{18}\)\gloss{Under operating grace the will is not merely passive: it freely consents}
  \breakrow{(7)} \(\pvar{TestimonyArguments.SolaGratia}{19}\)\gloss{Grace truly sufficient, giving the power to do good, is resisted}
  \breakrow{(8)} \(\pvar{TestimonyArguments.SolaGratia}{20}\)\gloss{Justifying grace, once given, can be lost}
  \breakrow{(9)} \(\lnot \pvar{TestimonyArguments.SolaGratia}{15}\)\gloss{\emph{denied:} Grace is received by faith alone}
  \breakrow{(10)} \(\pvar{TestimonyArguments.SolaGratia}{4} \rightarrow \pvar{TestimonyArguments.SolaGratia}{29}\)\gloss{\emph{would give:} Salvation is by grace and not earned: grace goes before every saving act}
  \breakrow{(11)} \(\pvar{TestimonyArguments.SolaGratia}{16} \rightarrow \pvar{TestimonyArguments.SolaGratia}{28}\)\gloss{\emph{would give:} Conversion is God’s work alone: God regenerates, and the will does nothing to it}
  \breakrow{(12)} \(\pvar{TestimonyArguments.SolaGratia}{17} \rightarrow \pvar{TestimonyArguments.SolaGratia}{30}\)\gloss{\emph{would give:} The grace that saves brings about the consent it seeks}
  \item[(13)] \(\pvar{TestimonyArguments.SolaGratia}{18} \rightarrow \lnot \pvar{TestimonyArguments.SolaGratia}{31}\)\gloss{\emph{would deny:} In being made alive the will is passive, and contributes nothing to it}
  \item[\(\vdash\)] \(\pvar{TestimonyArguments.SolaGratia}{29} \land \pvar{TestimonyArguments.SolaGratia}{30}\)\gloss{\emph{at issue:} \(\pvar{TestimonyArguments.SolaGratia}{29}\) Salvation is by grace and not earned: grace goes before every saving act; \(\pvar{TestimonyArguments.SolaGratia}{30}\) The grace that saves brings about the consent it seeks}
  \keepwithnext
\end{derivation}
\emph{Grounded in scripture alone:} Romans 11:6 — if it is by grace, it is no longer on the basis of works.
\par\medskip\noindent\textbf{\texttt{thomist\allowbreak{}As\allowbreak{}Disposition}} --- Thomist (Aquinas; Garrigou-Lagrange), read as the will consenting while God alone is the agent, Garrigou-Lagrange’s sense\par\nobreak\noindent
The Thomists, with “not merely passive” read in Garrigou-Lagrange’s sense.

\begin{derivation}
  \item[(1)] \(\pvar{TestimonyArguments.SolaGratia}{4}\)\gloss{Romans 11:6 — if it is by grace, it is no longer on the basis of works}
  \item[(2)] \(\pvar{TestimonyArguments.SolaGratia}{16}\)\gloss{In converting grace God is the sole mover: the mind is moved and does not move}
  \item[(3)] \(\pvar{TestimonyArguments.SolaGratia}{17}\)\gloss{God moves the will intending that it attain grace}
  \breakrow{(4)} \(\pvar{TestimonyArguments.SolaGratia}{26}\)\gloss{The free will, moved by God, cooperates by assenting to God’s call}
  \breakrow{(5)} \(\pvar{TestimonyArguments.SolaGratia}{27}\)\gloss{The free will, moved by God, can refuse its assent if it wishes}
  \breakrow{(6)} \(\pvar{TestimonyArguments.SolaGratia}{18}\)\gloss{Under operating grace the will is not merely passive: it freely consents}
  \breakrow{(7)} \(\pvar{TestimonyArguments.SolaGratia}{19}\)\gloss{Grace truly sufficient, giving the power to do good, is resisted}
  \breakrow{(8)} \(\pvar{TestimonyArguments.SolaGratia}{20}\)\gloss{Justifying grace, once given, can be lost}
  \breakrow{(9)} \(\lnot \pvar{TestimonyArguments.SolaGratia}{15}\)\gloss{\emph{denied:} Grace is received by faith alone}
  \breakrow{(10)} \(\pvar{TestimonyArguments.SolaGratia}{4} \rightarrow \pvar{TestimonyArguments.SolaGratia}{29}\)\gloss{\emph{would give:} Salvation is by grace and not earned: grace goes before every saving act}
  \breakrow{(11)} \(\pvar{TestimonyArguments.SolaGratia}{16} \rightarrow \pvar{TestimonyArguments.SolaGratia}{28}\)\gloss{\emph{would give:} Conversion is God’s work alone: God regenerates, and the will does nothing to it}
  \breakrow{(12)} \(\pvar{TestimonyArguments.SolaGratia}{17} \rightarrow \pvar{TestimonyArguments.SolaGratia}{30}\)\gloss{\emph{would give:} The grace that saves brings about the consent it seeks}
  \item[(13)] \(\pvar{TestimonyArguments.SolaGratia}{18} \rightarrow \pvar{TestimonyArguments.SolaGratia}{31}\)\gloss{\emph{would give:} In being made alive the will is passive, and contributes nothing to it}
  \item[\(\vdash\)] \(\pvar{TestimonyArguments.SolaGratia}{29} \land \pvar{TestimonyArguments.SolaGratia}{30}\)\gloss{\emph{at issue:} \(\pvar{TestimonyArguments.SolaGratia}{29}\) Salvation is by grace and not earned: grace goes before every saving act; \(\pvar{TestimonyArguments.SolaGratia}{30}\) The grace that saves brings about the consent it seeks}
  \keepwithnext
\end{derivation}
\emph{Grounded in scripture alone:} Romans 11:6 — if it is by grace, it is no longer on the basis of works.
\par\medskip\noindent\textbf{\texttt{thomist\allowbreak{}As\allowbreak{}Agent\_\allowbreak{}is\_\allowbreak{}satisfiable}}\par\nobreak\noindent
The Thomists read the first way can be held without contradiction.

{\small\begin{verbatim}
theorem thomistAsAgent_is_satisfiable : Satisfiable thomistAsAgent.premises
-- axioms: propext, Quot.sound
\end{verbatim}}
\par\medskip\noindent\textbf{\texttt{thomist\allowbreak{}As\allowbreak{}Disposition\_\allowbreak{}is\_\allowbreak{}satisfiable}}\par\nobreak\noindent
The Thomists read the second way can be held without contradiction.

{\small\begin{verbatim}
theorem thomistAsDisposition_is_satisfiable : Satisfiable
    thomistAsDisposition.premises
-- axioms: propext, Quot.sound
\end{verbatim}}
\par\medskip\noindent\textbf{\texttt{dort\_\allowbreak{}holds\_\allowbreak{}with\_\allowbreak{}garrigou\_\allowbreak{}lagrange}}\par\nobreak\noindent
Dort holds together with the Thomists read the second way.

{\small\begin{verbatim}
theorem dort_holds_with_garrigou_lagrange : Satisfiable (reformed.premises ++
    thomistAsDisposition.premises)
-- axioms: propext, Quot.sound
\end{verbatim}}
\par\medskip\noindent\textbf{\texttt{where\allowbreak{}The\allowbreak{}Agent\allowbreak{}Reading\allowbreak{}Falls}}\par\nobreak\noindent
\textbf{Why Dort stands against the Thomists, with “not merely passive” in the Molinists’ sense.} The Thomists’ own operating grace gives them that God alone does the work. Dort’s step says that if God alone does it, the will is passive. Read in the Molinists’ sense, the Thomists’ “not merely passive” says it is not. They cannot hold all three, and the crux is Dort’s step.

\paragraph{Why \emph{Reformed (the Canons of Dort)} stands against \emph{Thomist (Aquinas; Garrigou-Lagrange), read as making the will’s consent a cause of the new life, in the Molinists’ sense}.}
\begin{itemize}
\item \textbf{The crux:} \(\pvar{TestimonyArguments.SolaGratia}{28} \rightarrow \pvar{TestimonyArguments.SolaGratia}{31}\), a premise of \emph{Reformed (the Canons of Dort)}.
\item \textbf{What it does:} its conclusion does not follow without it.
\item \textbf{Where the rival breaks:} \(\pvar{TestimonyArguments.SolaGratia}{16}\), \(\pvar{TestimonyArguments.SolaGratia}{16} \rightarrow \pvar{TestimonyArguments.SolaGratia}{28}\), \(\pvar{TestimonyArguments.SolaGratia}{18}\), \(\pvar{TestimonyArguments.SolaGratia}{18} \rightarrow \lnot \pvar{TestimonyArguments.SolaGratia}{31}\) cannot be held together with \(\pvar{TestimonyArguments.SolaGratia}{28} \rightarrow \pvar{TestimonyArguments.SolaGratia}{31}\); each is needed for the break, and without the crux they stand.
\item \textbf{What this rests on:}
\begin{itemize}
\item the step itself --- \emph{consensus}: \cite[Trent, Session VI (1547), ch. 8]{tanner-decrees-1990}; \cite[36]{waterworth-trent-1848}; \cite[Five Arminian Articles (1610), art. 4]{schaff-creeds-christendom-3-1977}; \cite[Canons of Dort, III/IV, art. 10]{schaff-creeds-christendom-3-1977}; \emph{disputed}: \cite[ad loc. Eph 2:8]{robertson-word-pictures-4-1931}; \cite[228--229]{calvin-galatians-ephesians-1854}; \emph{disputed}: \cite[600]{arminius-works-1-1825}; \cite[681]{arminius-works-1-1825}; \cite[Five Arminian Articles (1610), art. 4]{schaff-creeds-christendom-3-1977}; \emph{disputed}: \cite[527--565]{arminius-works-3-1853}; \emph{consensus}: \cite[593]{schaff-creeds-christendom-3-1977}; \cite[301]{calvin-galatians-ephesians-1854}; \cite[Solid Declaration II.60]{concordia-triglotta-1921}; \emph{well supported}: \cite[Solid Declaration II.89]{concordia-triglotta-1921}; \cite[Canons of Dort, III/IV, art. 12]{schaff-creeds-christendom-3-1977}; \emph{disputed}: \cite[Solid Declaration II.89]{concordia-triglotta-1921}; \cite[Solid Declaration II.60]{concordia-triglotta-1921}; \cite[Solid Declaration XI.41]{concordia-triglotta-1921}
\end{itemize}
\end{itemize}
\par\medskip\noindent\textbf{\texttt{where\allowbreak{}The\allowbreak{}Disposition\allowbreak{}Reading\allowbreak{}Falls}}\par\nobreak\noindent
\textbf{Why Trent, on the Molinist reading, stands against the Thomists, with “not merely passive” in Garrigou-Lagrange’s sense.} The crux is the first half of Trent’s canon 4 as the Molinist reading takes it: the will cooperates by assenting, so it is not passive. That step is rated \emph{disputed}, because Aquinas grants that grace cooperates with a will that moves and holds that in grace’s first effect the mind “is moved and does not move”.

\paragraph{Why \emph{Trent (Session VI, canon 4, on the Molinist reading)} stands against \emph{Thomist (Aquinas; Garrigou-Lagrange), read as the will consenting while God alone is the agent, Garrigou-Lagrange’s sense}.}
\begin{itemize}
\item \textbf{The crux:} \((\pvar{TestimonyArguments.SolaGratia}{21} \land \pvar{TestimonyArguments.SolaGratia}{26}) \rightarrow \lnot \pvar{TestimonyArguments.SolaGratia}{31}\), a premise of \emph{Trent (Session VI, canon 4, on the Molinist reading)}.
\item \textbf{What it does:} its conclusion does not follow without it.
\item \textbf{Granted:} \(\pvar{TestimonyArguments.SolaGratia}{21}\), \(\pvar{TestimonyArguments.SolaGratia}{26}\), which the break also needs.
\item \textbf{Where the rival breaks:} \(\pvar{TestimonyArguments.SolaGratia}{18}\), \(\pvar{TestimonyArguments.SolaGratia}{18} \rightarrow \pvar{TestimonyArguments.SolaGratia}{31}\) cannot be held together with \((\pvar{TestimonyArguments.SolaGratia}{21} \land \pvar{TestimonyArguments.SolaGratia}{26}) \rightarrow \lnot \pvar{TestimonyArguments.SolaGratia}{31}\) and \(\pvar{TestimonyArguments.SolaGratia}{21}\), \(\pvar{TestimonyArguments.SolaGratia}{26}\); each is needed for the break, and without the crux they stand.
\item \textbf{What this rests on:}
\begin{itemize}
\item \(\pvar{TestimonyArguments.SolaGratia}{21}\) Zechariah 1:3 (“return to me, and I will return to you”) with Lamentations 5:21 (“restore us to yourself, O Lord”) --- \emph{consensus}: Zech 1:3; Lam 5:21; \cite[Trent, Session VI (1547), ch. 5]{tanner-decrees-1990}; \cite[33]{waterworth-trent-1848}
\item \(\pvar{TestimonyArguments.SolaGratia}{26}\) The free will, moved by God, cooperates by assenting to God’s call --- \emph{disputed}: \cite[Trent, Session VI (1547), canon 4]{tanner-decrees-1990}; \cite[45]{waterworth-trent-1848}
\item the step itself --- \emph{consensus}: \cite[Trent, Session VI (1547), ch. 8]{tanner-decrees-1990}; \cite[36]{waterworth-trent-1848}; \cite[Five Arminian Articles (1610), art. 4]{schaff-creeds-christendom-3-1977}; \cite[Canons of Dort, III/IV, art. 10]{schaff-creeds-christendom-3-1977}; \emph{disputed}: \cite[I-II q.111 a.2]{aquinas-summa-1920}; \cite[Canons of Dort, III/IV, art. 12]{schaff-creeds-christendom-3-1977}; \cite[Trent, Session VI (1547), canon 4]{tanner-decrees-1990}; \cite[45]{waterworth-trent-1848}; \emph{disputed}: \cite[I-II q.10 a.4 ad 3]{aquinas-summa-1920}; \cite[I-II q.112 a.3]{aquinas-summa-1920}; \cite[Trent, Session VI (1547), canon 4]{tanner-decrees-1990}; \cite[45]{waterworth-trent-1848}
\end{itemize}
\end{itemize}
\par\medskip\noindent\textbf{\texttt{what\allowbreak{}Not\allowbreak{}Merely\allowbreak{}Passive\allowbreak{}Means}}\par\nobreak\noindent
\textbf{What “not merely passive” means, read both ways.} The claim is the Thomists’: under operating grace the will is not merely passive, but freely consents.

Read as making the will’s consent a cause of the new life, in the Molinists’ sense, it cannot be held with Dort, at Dort’s step from God’s work alone to the will’s passivity.

Read as the will consenting while God alone is the agent of the new life, Garrigou-Lagrange’s sense (“it nevertheless follows it in the genus or order of formal, efficient, and final cause”), Dort can hold it without contradiction. It cannot be held with Trent on the Molinist reading, at the step from the will’s cooperation, which is rated \texttt{disputed} because Aquinas grants that grace cooperates with a will that moves and holds that in grace’s first effect the mind “is moved and does not move”.

Each reading is consistent: the Thomists read either way can be held without contradiction. Which one the Thomists mean is not something the dilemma decides, but Garrigou-Lagrange’s own explanation is the second.

\paragraph{Every reading of \(P_{18}\), answered.} \emph{Thomist (Aquinas; Garrigou-Lagrange)} holds \(\pvar{TestimonyArguments.SolaGratia}{18}\). It is read 2 ways here, and each reading is checked; none can be left out.

\paragraph{1. Read as making the will’s consent a cause of the new life, in the Molinists’ sense.} The claim commits \emph{Thomist (Aquinas; Garrigou-Lagrange)} to \(\lnot \pvar{TestimonyArguments.SolaGratia}{31}\) --- so read by \cite{pohle-molinism-1911} (\emph{disputed}).

Against \emph{Reformed (the Canons of Dort)}, it cannot be held. Where it breaks:
\paragraph{Why \emph{Reformed (the Canons of Dort)} stands against \emph{Thomist (Aquinas; Garrigou-Lagrange), read as making the will’s consent a cause of the new life, in the Molinists’ sense}.}
\begin{itemize}
\item \textbf{The crux:} \(\pvar{TestimonyArguments.SolaGratia}{28} \rightarrow \pvar{TestimonyArguments.SolaGratia}{31}\), a premise of \emph{Reformed (the Canons of Dort)}.
\item \textbf{What it does:} its conclusion does not follow without it.
\item \textbf{Where the rival breaks:} \(\pvar{TestimonyArguments.SolaGratia}{16}\), \(\pvar{TestimonyArguments.SolaGratia}{16} \rightarrow \pvar{TestimonyArguments.SolaGratia}{28}\), \(\pvar{TestimonyArguments.SolaGratia}{18}\), \(\pvar{TestimonyArguments.SolaGratia}{18} \rightarrow \lnot \pvar{TestimonyArguments.SolaGratia}{31}\) cannot be held together with \(\pvar{TestimonyArguments.SolaGratia}{28} \rightarrow \pvar{TestimonyArguments.SolaGratia}{31}\); each is needed for the break, and without the crux they stand.
\item \textbf{What this rests on:}
\begin{itemize}
\item the step itself --- \emph{consensus}: \cite[Trent, Session VI (1547), ch. 8]{tanner-decrees-1990}; \cite[36]{waterworth-trent-1848}; \cite[Five Arminian Articles (1610), art. 4]{schaff-creeds-christendom-3-1977}; \cite[Canons of Dort, III/IV, art. 10]{schaff-creeds-christendom-3-1977}; \emph{disputed}: \cite[ad loc. Eph 2:8]{robertson-word-pictures-4-1931}; \cite[228--229]{calvin-galatians-ephesians-1854}; \emph{disputed}: \cite[600]{arminius-works-1-1825}; \cite[681]{arminius-works-1-1825}; \cite[Five Arminian Articles (1610), art. 4]{schaff-creeds-christendom-3-1977}; \emph{disputed}: \cite[527--565]{arminius-works-3-1853}; \emph{consensus}: \cite[593]{schaff-creeds-christendom-3-1977}; \cite[301]{calvin-galatians-ephesians-1854}; \cite[Solid Declaration II.60]{concordia-triglotta-1921}; \emph{well supported}: \cite[Solid Declaration II.89]{concordia-triglotta-1921}; \cite[Canons of Dort, III/IV, art. 12]{schaff-creeds-christendom-3-1977}; \emph{disputed}: \cite[Solid Declaration II.89]{concordia-triglotta-1921}; \cite[Solid Declaration II.60]{concordia-triglotta-1921}; \cite[Solid Declaration XI.41]{concordia-triglotta-1921}
\end{itemize}
\end{itemize}
\paragraph{2. Read as the will consenting while God alone is the agent, Garrigou-Lagrange’s sense.} The claim commits \emph{Thomist (Aquinas; Garrigou-Lagrange)} to \(\pvar{TestimonyArguments.SolaGratia}{31}\) --- so read by \cite[ch. X]{garrigou-lagrange-grace-1952}; \cite[I-II q.111 a.2 ad 2]{aquinas-summa-1920} (\emph{disputed}).

Against \emph{Trent (Session VI, canon 4, on the Molinist reading)}, it cannot be held. Where it breaks:
\paragraph{Why \emph{Trent (Session VI, canon 4, on the Molinist reading)} stands against \emph{Thomist (Aquinas; Garrigou-Lagrange), read as the will consenting while God alone is the agent, Garrigou-Lagrange’s sense}.}
\begin{itemize}
\item \textbf{The crux:} \((\pvar{TestimonyArguments.SolaGratia}{21} \land \pvar{TestimonyArguments.SolaGratia}{26}) \rightarrow \lnot \pvar{TestimonyArguments.SolaGratia}{31}\), a premise of \emph{Trent (Session VI, canon 4, on the Molinist reading)}.
\item \textbf{What it does:} its conclusion does not follow without it.
\item \textbf{Granted:} \(\pvar{TestimonyArguments.SolaGratia}{21}\), \(\pvar{TestimonyArguments.SolaGratia}{26}\), which the break also needs.
\item \textbf{Where the rival breaks:} \(\pvar{TestimonyArguments.SolaGratia}{18}\), \(\pvar{TestimonyArguments.SolaGratia}{18} \rightarrow \pvar{TestimonyArguments.SolaGratia}{31}\) cannot be held together with \((\pvar{TestimonyArguments.SolaGratia}{21} \land \pvar{TestimonyArguments.SolaGratia}{26}) \rightarrow \lnot \pvar{TestimonyArguments.SolaGratia}{31}\) and \(\pvar{TestimonyArguments.SolaGratia}{21}\), \(\pvar{TestimonyArguments.SolaGratia}{26}\); each is needed for the break, and without the crux they stand.
\item \textbf{What this rests on:}
\begin{itemize}
\item \(\pvar{TestimonyArguments.SolaGratia}{21}\) Zechariah 1:3 (“return to me, and I will return to you”) with Lamentations 5:21 (“restore us to yourself, O Lord”) --- \emph{consensus}: Zech 1:3; Lam 5:21; \cite[Trent, Session VI (1547), ch. 5]{tanner-decrees-1990}; \cite[33]{waterworth-trent-1848}
\item \(\pvar{TestimonyArguments.SolaGratia}{26}\) The free will, moved by God, cooperates by assenting to God’s call --- \emph{disputed}: \cite[Trent, Session VI (1547), canon 4]{tanner-decrees-1990}; \cite[45]{waterworth-trent-1848}
\item the step itself --- \emph{consensus}: \cite[Trent, Session VI (1547), ch. 8]{tanner-decrees-1990}; \cite[36]{waterworth-trent-1848}; \cite[Five Arminian Articles (1610), art. 4]{schaff-creeds-christendom-3-1977}; \cite[Canons of Dort, III/IV, art. 10]{schaff-creeds-christendom-3-1977}; \emph{disputed}: \cite[I-II q.111 a.2]{aquinas-summa-1920}; \cite[Canons of Dort, III/IV, art. 12]{schaff-creeds-christendom-3-1977}; \cite[Trent, Session VI (1547), canon 4]{tanner-decrees-1990}; \cite[45]{waterworth-trent-1848}; \emph{disputed}: \cite[I-II q.10 a.4 ad 3]{aquinas-summa-1920}; \cite[I-II q.112 a.3]{aquinas-summa-1920}; \cite[Trent, Session VI (1547), canon 4]{tanner-decrees-1990}; \cite[45]{waterworth-trent-1848}
\end{itemize}
\end{itemize}
Against \emph{Reformed (the Canons of Dort)}, it is not reached: \emph{Reformed (the Canons of Dort)} holds, and can be held together with it.

\subsection*{Weighing the positions}
\addcontentsline{toc}{subsection}{Weighing the positions}
Nine parties are heard. One holds only what all of them hold: salvation is by grace, from Romans 11:6. Then Dort; Trent, on the Molinist reading of its fourth canon; the Thomists, on theirs; the Remonstrants; the Lutherans of the Formula of Concord; Wesley; the third way as it was put to this argument; and the objection from the Spirit’s resistance, with its three routes to the claim that the grace that saves can be finally refused. Who defeats whom is computed from their premises and checked by the kernel.

\subsubsection*{How to read the verdicts}
\begin{itemize}
  \item \textbf{Weakest link.} A position is only as strong as its least-supported premise or step. What every party holds rests on Romans 11:6 and a step rated \emph{consensus}, accepted across traditions. Every other position rests on at least one reading or step rated \emph{disputed}, the lowest rating, so no rating breaks a standoff between them.
  \item \textbf{Defeats.} One position contradicts another (it denies one of the other’s premises, its conclusion, or a claim it derives on the way), and its weakest link is not rated below the other’s. A defeat that runs both ways is a \emph{standoff}.
  \item \textbf{Settled before any choice.} What prevails once every unanswered attack is taken into account, before any standoff is decided either way.
  \item \textbf{Can be defended.} Held by some consistent group of positions that answers every attack on its members, by itself or through an ally. Such a group that no further party can join is a \emph{maximal defensible position}.
  \item \textbf{Forced.} Held in every maximal defensible position, however the standoffs are resolved. A party is shown not forced by naming a maximal defensible position that does not hold it.
\end{itemize}
Dort’s position here does not carry faith alone: the Reformed case is stated from Ephesians 2 and Romans 9, and sola fide is an argument of its own. So Dort and Trent are never set against each other over it; they stand off over the will’s passivity and grace effectual of itself.

\subsubsection*{What the weighing found}
Only what every party holds is settled before any choice: salvation is by grace (\texttt{only\_\allowbreak{}grace\_\allowbreak{}prevails}). Everything else is in a standoff with something.

\begin{itemize}
  \item \textbf{Dort stands off against every party that denies part of sola gratia:} Trent, the Remonstrants, the Lutherans, Wesley and the objection. Each of them defeats Dort, and Dort defeats each back.
  \item \textbf{Neither Dort nor the Thomists defeats the other} (\texttt{dort\_\allowbreak{}stands\_\allowbreak{}with\_\allowbreak{}the\_\allowbreak{}thomists}), because the Thomist position weighed here takes no side on the will’s passivity: it holds that the will is “not merely passive” and draws nothing from it. So this does not show that the Thomists hold the third part of sola gratia. In the Molinists’ sense the phrase cannot be held with Dort; in Garrigou-Lagrange’s it can (the chapter on Aquinas). On faith alone and perseverance, on which they part, Dort’s position here carries nothing. So the Thomists can be defended with Dort (\texttt{thomist\_\allowbreak{}defensible}).
  \item \textbf{The Thomists stand off against every party but Dort and what every party holds.} Against the Remonstrants, the Lutherans, Wesley and the objection, the standoff is over grace effectual of itself, which the claim that the grace that saves can be finally refused excludes. Against the third way, the Lutherans and Wesley, it is also over faith alone. Against Trent on the Molinist reading, it is the question Rome left open in 1607.
  \item \textbf{Trent stands off against Dort, the Thomists, the Lutherans, Wesley and the third way.} With the Thomists the standoff is over grace effectual of itself, the question of the two Catholic schools that Rome left open in 1607 (\texttt{thomist\_\allowbreak{}not\_\allowbreak{}forced}); with the Lutherans, Wesley and the third way it is over faith alone, which Trent’s ninth canon condemns; with the Lutherans also over the will’s passivity.
  \item \textbf{The Remonstrants, the Lutherans, Wesley and the objection do not defeat one another.} They agree that the grace that saves can be finally refused. The Lutherans also hold the will passive; the Remonstrants and Wesley are not encoded with a view on that, so nothing here sets them against each other.
  \item \textbf{The third way is attacked only over faith alone}, by Trent and by the Thomists. It takes no side between Dort and Wesley, and so it can be defended with either: with Dort, against Trent (\texttt{third\_\allowbreak{}way\_\allowbreak{}defensible\_\allowbreak{}with\_\allowbreak{}dort}); and with the Remonstrants, the Lutherans and Wesley, against both Dort and Trent (\texttt{third\_\allowbreak{}way\_\allowbreak{}defensible\_\allowbreak{}with\_\allowbreak{}wesley}). The dilemma found the same thing: until the third way says which kind of help God gives, it has taken neither side.
  \item \textbf{Trent can be defended with the Remonstrants and the objection}, which agree with it that grace can be resisted, against Dort.
\end{itemize}
Every party can be defended, and only the first is forced: each of the others is left out of some maximal defensible position. Each standoff rests on readings and steps rated \emph{disputed}, so the verdict moves only if a new, cited argument moves one of those ratings.

\subsubsection*{How this is checked}
\paragraph*{The dispute}
\leavevmode\par\nobreak
\par\medskip\noindent\textbf{\texttt{party\allowbreak{}Node}}\par\nobreak\noindent
The package each party argues from.

{\small\begin{verbatim}
def partyNode : Party → ArgumentPackage Claim :=
  fun x =>
    match x with
    | Party.grace => graceCase
    | Party.reformed => reformed
    | Party.trent => tridentinePosition
    | Party.remonstrant => remonstrantPosition
    | Party.lutheran => lutheranPosition
    | Party.wesleyan => wesleyanPosition
    | Party.thirdWay => thirdWay
    | Party.objection => objection
    | Party.thomist => thomistPosition
\end{verbatim}}
\par\medskip\noindent\textbf{\texttt{sola\allowbreak{}Gratia\allowbreak{}Dispute}}\par\nobreak\noindent
The dispute over sola gratia: every party can be held without contradiction, delivers its conclusion, and has its steps rated.

{\small\begin{verbatim}
def solaGratiaDispute : Dispute Claim Party :=
  { node := partyNode,
    consistent := solaGratiaDispute._proof_1,
    sound := solaGratiaDispute._proof_2,
    rated := solaGratiaDispute._proof_12 }
\end{verbatim}}
\par\medskip\noindent\textbf{\texttt{sola\allowbreak{}Gratia\allowbreak{}Dispute\_\allowbreak{}defeats}}\par\nobreak\noindent
\textbf{Who defeats whom}, all 81 pairs: twenty-eight defeats, every one of them half of a standoff. Dort stands off against Trent, the Remonstrants, the Lutherans, Wesley and the objection; Trent against Dort, the Lutherans, Wesley, the third way and the Thomists; the Thomists against every party but Dort and what every party holds. Dort and the Thomists defeat neither each other. Nothing attacks what every party holds. Every cell is checked by the kernel.

{\small\begin{verbatim}
theorem solaGratiaDispute_defeats : ∀ (i j : Party), solaGratiaDispute.defeats
    i j ↔ (i, j) ∈ [(Party.reformed, Party.trent), (Party.reformed,
    Party.remonstrant), (Party.reformed, Party.lutheran), (Party.reformed,
    Party.wesleyan), (Party.reformed, Party.objection), (Party.trent,
    Party.reformed), (Party.trent, Party.lutheran), (Party.trent,
    Party.wesleyan), (Party.trent, Party.thirdWay), (Party.trent,
    Party.thomist), (Party.remonstrant, Party.reformed), (Party.remonstrant,
    Party.thomist), (Party.lutheran, Party.reformed), (Party.lutheran,
    Party.trent), (Party.lutheran, Party.thomist), (Party.wesleyan,
    Party.reformed), (Party.wesleyan, Party.trent), (Party.wesleyan,
    Party.thomist), (Party.thirdWay, Party.trent), (Party.thirdWay,
    Party.thomist), (Party.objection, Party.reformed), (Party.objection,
    Party.thomist), (Party.thomist, Party.trent), (Party.thomist,
    Party.remonstrant), (Party.thomist, Party.lutheran), (Party.thomist,
    Party.wesleyan), (Party.thomist, Party.thirdWay), (Party.thomist,
    Party.objection)]
-- axioms: propext, Classical.choice, Quot.sound
\end{verbatim}}
\paragraph*{What the dispute decides}
\leavevmode\par\nobreak
\par\medskip\noindent\textbf{\texttt{grace\allowbreak{}Unanswered}}\par\nobreak\noindent
How far the dispute is settled before any choice. Nothing defeats what every party holds, so it is settled first; and nothing is settled with it, because every other party has a defeater that what every party holds does not answer.

\paragraph{What the dispute forces is exactly \emph{Romans 11:6 (by grace, so not by works)}.}
\begin{itemize}
\item Stage 1: \emph{Romans 11:6 (by grace, so not by works)}.
\begin{itemize}
\item \emph{Romans 11:6 (by grace, so not by works)} is defeated by nothing.
\end{itemize}
\item Nothing else is forced:
\begin{itemize}
\item \emph{Trent (Session VI, canon 4, on the Molinist reading)} defeats \emph{Reformed (the Canons of Dort)}, and nothing forced defeats \emph{Trent (Session VI, canon 4, on the Molinist reading)}.
\item \emph{Reformed (the Canons of Dort)} defeats \emph{Trent (Session VI, canon 4, on the Molinist reading)}, and nothing forced defeats \emph{Reformed (the Canons of Dort)}.
\item \emph{Reformed (the Canons of Dort)} defeats \emph{Remonstrant (the Five Articles, 1610)}, and nothing forced defeats \emph{Reformed (the Canons of Dort)}.
\item \emph{Reformed (the Canons of Dort)} defeats \emph{Lutheran (the Formula of Concord)}, and nothing forced defeats \emph{Reformed (the Canons of Dort)}.
\item \emph{Reformed (the Canons of Dort)} defeats \emph{Wesleyan (preventing grace)}, and nothing forced defeats \emph{Reformed (the Canons of Dort)}.
\item \emph{Trent (Session VI, canon 4, on the Molinist reading)} defeats \emph{The third way (an ability given, received by faith)}, and nothing forced defeats \emph{Trent (Session VI, canon 4, on the Molinist reading)}.
\item \emph{Reformed (the Canons of Dort)} defeats \emph{The objection from the Spirit’s resistance}, and nothing forced defeats \emph{Reformed (the Canons of Dort)}.
\item \emph{Trent (Session VI, canon 4, on the Molinist reading)} defeats \emph{Thomist (Aquinas; Garrigou-Lagrange)}, and nothing forced defeats \emph{Trent (Session VI, canon 4, on the Molinist reading)}.
\end{itemize}
\end{itemize}
\par\noindent\textbf{What this rests on.} The reasons state 8 defeats and 11 absences of defeat, each a cell of the defeat table, \texttt{sola\allowbreak{}Gratia\allowbreak{}Dispute\_\allowbreak{}defeats}, computed from the two parties' premises and checked by the kernel. Whether an attack survives turns on the attacker's weakest link and on the rating of what it attacks, so a changed rating can change the verdict. The weakest links:
\begin{itemize}
\item \emph{Romans 11:6 (by grace, so not by works)}, weakest at \emph{consensus}:
\begin{itemize}
\item Romans 11:6 — if it is by grace, it is no longer on the basis of works --- Rom 11:6
\item an inference step --- \cite[Trent, Session VI (1547), ch. 8]{tanner-decrees-1990}; \cite[36]{waterworth-trent-1848}; \cite[Five Arminian Articles (1610), art. 4]{schaff-creeds-christendom-3-1977}; \cite[Canons of Dort, III/IV, art. 10]{schaff-creeds-christendom-3-1977}
\end{itemize}
\item \emph{Reformed (the Canons of Dort)}, weakest at \emph{disputed}:
\begin{itemize}
\item What Stephen’s hearers resisted in Acts 7:51 was the grace that saves (encoded as a denial, so it ranks disputed) --- \cite[Canons of Dort, III/IV, art. 9]{schaff-creeds-christendom-3-1977}
\item an inference step --- \cite[ad loc. Eph 2:8]{robertson-word-pictures-4-1931}; \cite[228--229]{calvin-galatians-ephesians-1854}
\item an inference step --- \cite[600]{arminius-works-1-1825}; \cite[681]{arminius-works-1-1825}; \cite[Five Arminian Articles (1610), art. 4]{schaff-creeds-christendom-3-1977}
\item an inference step --- \cite[527--565]{arminius-works-3-1853}
\item an inference step --- \cite[Solid Declaration II.89]{concordia-triglotta-1921}; \cite[Solid Declaration II.60]{concordia-triglotta-1921}; \cite[Solid Declaration XI.41]{concordia-triglotta-1921}
\end{itemize}
\item \emph{Trent (Session VI, canon 4, on the Molinist reading)}, weakest at \emph{disputed}:
\begin{itemize}
\item The free will, moved by God, cooperates by assenting to God’s call --- \cite[Trent, Session VI (1547), canon 4]{tanner-decrees-1990}; \cite[45]{waterworth-trent-1848}
\item The free will, moved by God, can refuse its assent if it wishes --- \cite[Trent, Session VI (1547), canon 4]{tanner-decrees-1990}; \cite[45]{waterworth-trent-1848}
\item Grace is received by faith alone (encoded as a denial, so it ranks disputed) --- \cite[Augsburg Confession IV.1–3]{concordia-triglotta-1921}; \cite[Sermon 1, III.1]{wesley-sermons-1872}; \cite[Trent, Session VI (1547), canon 9]{tanner-decrees-1990}
\item an inference step --- \cite[I-II q.111 a.2]{aquinas-summa-1920}; \cite[Canons of Dort, III/IV, art. 12]{schaff-creeds-christendom-3-1977}; \cite[Trent, Session VI (1547), canon 4]{tanner-decrees-1990}; \cite[45]{waterworth-trent-1848}
\item an inference step --- \cite[I-II q.10 a.4 ad 3]{aquinas-summa-1920}; \cite[I-II q.112 a.3]{aquinas-summa-1920}; \cite[Trent, Session VI (1547), canon 4]{tanner-decrees-1990}; \cite[45]{waterworth-trent-1848}
\end{itemize}
\item \emph{Remonstrant (the Five Articles, 1610)}, weakest at \emph{disputed}:
\begin{itemize}
\item What Stephen’s hearers resisted in Acts 7:51 was the grace that saves --- \cite[Five Arminian Articles (1610), art. 4]{schaff-creeds-christendom-3-1977}
\end{itemize}
\item \emph{Lutheran (the Formula of Concord)}, weakest at \emph{disputed}:
\begin{itemize}
\item Those who obstinately resist the Spirit can grieve and lose him --- \cite[Solid Declaration II.83]{concordia-triglotta-1921}; \cite[Solid Declaration XI.42]{concordia-triglotta-1921}
\item Conversion is God’s work alone: God regenerates, and the will does nothing to it --- \cite[Solid Declaration II.89]{concordia-triglotta-1921}; \cite[Solid Declaration II.90]{concordia-triglotta-1921}
\item Grace is received by faith alone --- \cite[Augsburg Confession IV.1–3]{concordia-triglotta-1921}; \cite[Sermon 1, III.1]{wesley-sermons-1872}; \cite[Trent, Session VI (1547), canon 9]{tanner-decrees-1990}
\end{itemize}
\item \emph{Wesleyan (preventing grace)}, weakest at \emph{disputed}:
\begin{itemize}
\item The grace that saves can be finally refused --- \cite[397]{wesley-works-9-1827}
\item Grace is received by faith alone --- \cite[Augsburg Confession IV.1–3]{concordia-triglotta-1921}; \cite[Sermon 1, III.1]{wesley-sermons-1872}; \cite[Trent, Session VI (1547), canon 9]{tanner-decrees-1990}
\end{itemize}
\item \emph{The third way (an ability given, received by faith)}, weakest at \emph{disputed}:
\begin{itemize}
\item Grace is received by faith alone --- \cite[Augsburg Confession IV.1–3]{concordia-triglotta-1921}; \cite[Sermon 1, III.1]{wesley-sermons-1872}; \cite[Trent, Session VI (1547), canon 9]{tanner-decrees-1990}
\end{itemize}
\item \emph{The objection from the Spirit’s resistance}, weakest at \emph{disputed}:
\begin{itemize}
\item What Stephen’s hearers resisted in Acts 7:51 was the grace that saves --- \cite[Five Arminian Articles (1610), art. 4]{schaff-creeds-christendom-3-1977}
\item Those who obstinately resist the Spirit can grieve and lose him --- \cite[Solid Declaration II.83]{concordia-triglotta-1921}; \cite[Solid Declaration XI.42]{concordia-triglotta-1921}
\item The blasphemy against the Spirit is impenitence to the end: grace refused till death --- \cite[325]{npnf1-06-1888}; \cite[II-II q.14 a.1]{aquinas-summa-1920}; \cite[74]{calvin-harmony-evangelists-2-1845}
\item an inference step --- \cite[325]{npnf1-06-1888}; \cite[On Rebuke and Grace, ch. 43]{npnf1-05-1887}
\end{itemize}
\item \emph{Thomist (Aquinas; Garrigou-Lagrange)}, weakest at \emph{disputed}:
\begin{itemize}
\item In converting grace God is the sole mover: the mind is moved and does not move --- \cite[I-II q.111 a.2]{aquinas-summa-1920}; \cite[ch. IV]{garrigou-lagrange-grace-1952}
\item The free will, moved by God, cooperates by assenting to God’s call --- \cite[Trent, Session VI (1547), canon 4]{tanner-decrees-1990}; \cite[45]{waterworth-trent-1848}
\item The free will, moved by God, can refuse its assent if it wishes --- \cite[Trent, Session VI (1547), canon 4]{tanner-decrees-1990}; \cite[45]{waterworth-trent-1848}
\item Grace truly sufficient, giving the power to do good, is resisted --- \cite[ch. VI]{garrigou-lagrange-grace-1952}
\item Justifying grace, once given, can be lost --- \cite[II-II q.24 a.11]{aquinas-summa-1920}; \cite[Trent, Session VI (1547), canon 23]{tanner-decrees-1990}
\item Grace is received by faith alone (encoded as a denial, so it ranks disputed) --- \cite[II-II q.4 a.3]{aquinas-summa-1920}; \cite[I-II q.113 a.4 ad 1]{aquinas-summa-1920}; \cite[Trent, Session VI (1547), canon 9]{tanner-decrees-1990}
\item an inference step --- \cite{pohle-molinism-1911}; \cite[I-II q.112 a.3]{aquinas-summa-1920}; \cite[ch. VII]{garrigou-lagrange-grace-1952}
\end{itemize}
\end{itemize}
\par\medskip\noindent\textbf{\texttt{only\_\allowbreak{}grace\_\allowbreak{}prevails}}\par\nobreak\noindent
\textbf{Only what every party holds prevails: salvation is by grace.} What the dispute settles before any choice between rivals is that, and nothing more. Nothing defeats it, because every party holds every premise it holds.

What this does not claim: that sola gratia is false, or that any rival is right. Every other position is in a standoff, and none is forced.

{\small\begin{verbatim}
theorem only_grace_prevails : Framework.grounded solaGratiaDispute.defeats =
    Solver.toSet [Party.grace]
-- axioms: propext, Classical.choice, Quot.sound
\end{verbatim}}
\par\medskip\noindent\textbf{\texttt{dort\allowbreak{}Defended}}\par\nobreak\noindent
Why Dort can be defended: with what every party holds and the third way, it answers every attacker itself.

\paragraph{\emph{Reformed (the Canons of Dort)} is accepted on some resolution.}
\begin{itemize}
\item A position holding \emph{Romans 11:6 (by grace, so not by works)}, \emph{Reformed (the Canons of Dort)}, \emph{The third way (an ability given, received by faith)} can be held.
\item None of \emph{Romans 11:6 (by grace, so not by works)}, \emph{Reformed (the Canons of Dort)}, \emph{The third way (an ability given, received by faith)} defeats another, and each attack on them is answered from among them:
\begin{itemize}
\item \emph{Trent (Session VI, canon 4, on the Molinist reading)} defeats \emph{Reformed (the Canons of Dort)}, and \emph{Reformed (the Canons of Dort)} defeats \emph{Trent (Session VI, canon 4, on the Molinist reading)}.
\item \emph{Remonstrant (the Five Articles, 1610)} defeats \emph{Reformed (the Canons of Dort)}, and \emph{Reformed (the Canons of Dort)} defeats \emph{Remonstrant (the Five Articles, 1610)}.
\item \emph{Lutheran (the Formula of Concord)} defeats \emph{Reformed (the Canons of Dort)}, and \emph{Reformed (the Canons of Dort)} defeats \emph{Lutheran (the Formula of Concord)}.
\item \emph{Wesleyan (preventing grace)} defeats \emph{Reformed (the Canons of Dort)}, and \emph{Reformed (the Canons of Dort)} defeats \emph{Wesleyan (preventing grace)}.
\item \emph{The objection from the Spirit’s resistance} defeats \emph{Reformed (the Canons of Dort)}, and \emph{Reformed (the Canons of Dort)} defeats \emph{The objection from the Spirit’s resistance}.
\item \emph{Trent (Session VI, canon 4, on the Molinist reading)} defeats \emph{The third way (an ability given, received by faith)}, and \emph{Reformed (the Canons of Dort)} defeats \emph{Trent (Session VI, canon 4, on the Molinist reading)}.
\item \emph{Thomist (Aquinas; Garrigou-Lagrange)} defeats \emph{The third way (an ability given, received by faith)}, and \emph{The third way (an ability given, received by faith)} defeats \emph{Thomist (Aquinas; Garrigou-Lagrange)}.
\end{itemize}
\end{itemize}
\par\noindent\textbf{What this rests on.} The reasons state 13 defeats and 20 absences of defeat, each a cell of the defeat table, \texttt{sola\allowbreak{}Gratia\allowbreak{}Dispute\_\allowbreak{}defeats}, computed from the two parties' premises and checked by the kernel. Whether an attack survives turns on the attacker's weakest link and on the rating of what it attacks, so a changed rating can change the verdict. The weakest links:
\begin{itemize}
\item \emph{Romans 11:6 (by grace, so not by works)}, weakest at \emph{consensus}:
\begin{itemize}
\item Romans 11:6 — if it is by grace, it is no longer on the basis of works --- Rom 11:6
\item an inference step --- \cite[Trent, Session VI (1547), ch. 8]{tanner-decrees-1990}; \cite[36]{waterworth-trent-1848}; \cite[Five Arminian Articles (1610), art. 4]{schaff-creeds-christendom-3-1977}; \cite[Canons of Dort, III/IV, art. 10]{schaff-creeds-christendom-3-1977}
\end{itemize}
\item \emph{Reformed (the Canons of Dort)}, weakest at \emph{disputed}:
\begin{itemize}
\item What Stephen’s hearers resisted in Acts 7:51 was the grace that saves (encoded as a denial, so it ranks disputed) --- \cite[Canons of Dort, III/IV, art. 9]{schaff-creeds-christendom-3-1977}
\item an inference step --- \cite[ad loc. Eph 2:8]{robertson-word-pictures-4-1931}; \cite[228--229]{calvin-galatians-ephesians-1854}
\item an inference step --- \cite[600]{arminius-works-1-1825}; \cite[681]{arminius-works-1-1825}; \cite[Five Arminian Articles (1610), art. 4]{schaff-creeds-christendom-3-1977}
\item an inference step --- \cite[527--565]{arminius-works-3-1853}
\item an inference step --- \cite[Solid Declaration II.89]{concordia-triglotta-1921}; \cite[Solid Declaration II.60]{concordia-triglotta-1921}; \cite[Solid Declaration XI.41]{concordia-triglotta-1921}
\end{itemize}
\item \emph{The third way (an ability given, received by faith)}, weakest at \emph{disputed}:
\begin{itemize}
\item Grace is received by faith alone --- \cite[Augsburg Confession IV.1–3]{concordia-triglotta-1921}; \cite[Sermon 1, III.1]{wesley-sermons-1872}; \cite[Trent, Session VI (1547), canon 9]{tanner-decrees-1990}
\end{itemize}
\item \emph{Trent (Session VI, canon 4, on the Molinist reading)}, weakest at \emph{disputed}:
\begin{itemize}
\item The free will, moved by God, cooperates by assenting to God’s call --- \cite[Trent, Session VI (1547), canon 4]{tanner-decrees-1990}; \cite[45]{waterworth-trent-1848}
\item The free will, moved by God, can refuse its assent if it wishes --- \cite[Trent, Session VI (1547), canon 4]{tanner-decrees-1990}; \cite[45]{waterworth-trent-1848}
\item Grace is received by faith alone (encoded as a denial, so it ranks disputed) --- \cite[Augsburg Confession IV.1–3]{concordia-triglotta-1921}; \cite[Sermon 1, III.1]{wesley-sermons-1872}; \cite[Trent, Session VI (1547), canon 9]{tanner-decrees-1990}
\item an inference step --- \cite[I-II q.111 a.2]{aquinas-summa-1920}; \cite[Canons of Dort, III/IV, art. 12]{schaff-creeds-christendom-3-1977}; \cite[Trent, Session VI (1547), canon 4]{tanner-decrees-1990}; \cite[45]{waterworth-trent-1848}
\item an inference step --- \cite[I-II q.10 a.4 ad 3]{aquinas-summa-1920}; \cite[I-II q.112 a.3]{aquinas-summa-1920}; \cite[Trent, Session VI (1547), canon 4]{tanner-decrees-1990}; \cite[45]{waterworth-trent-1848}
\end{itemize}
\item \emph{Remonstrant (the Five Articles, 1610)}, weakest at \emph{disputed}:
\begin{itemize}
\item What Stephen’s hearers resisted in Acts 7:51 was the grace that saves --- \cite[Five Arminian Articles (1610), art. 4]{schaff-creeds-christendom-3-1977}
\end{itemize}
\item \emph{Lutheran (the Formula of Concord)}, weakest at \emph{disputed}:
\begin{itemize}
\item Those who obstinately resist the Spirit can grieve and lose him --- \cite[Solid Declaration II.83]{concordia-triglotta-1921}; \cite[Solid Declaration XI.42]{concordia-triglotta-1921}
\item Conversion is God’s work alone: God regenerates, and the will does nothing to it --- \cite[Solid Declaration II.89]{concordia-triglotta-1921}; \cite[Solid Declaration II.90]{concordia-triglotta-1921}
\item Grace is received by faith alone --- \cite[Augsburg Confession IV.1–3]{concordia-triglotta-1921}; \cite[Sermon 1, III.1]{wesley-sermons-1872}; \cite[Trent, Session VI (1547), canon 9]{tanner-decrees-1990}
\end{itemize}
\item \emph{Wesleyan (preventing grace)}, weakest at \emph{disputed}:
\begin{itemize}
\item The grace that saves can be finally refused --- \cite[397]{wesley-works-9-1827}
\item Grace is received by faith alone --- \cite[Augsburg Confession IV.1–3]{concordia-triglotta-1921}; \cite[Sermon 1, III.1]{wesley-sermons-1872}; \cite[Trent, Session VI (1547), canon 9]{tanner-decrees-1990}
\end{itemize}
\item \emph{The objection from the Spirit’s resistance}, weakest at \emph{disputed}:
\begin{itemize}
\item What Stephen’s hearers resisted in Acts 7:51 was the grace that saves --- \cite[Five Arminian Articles (1610), art. 4]{schaff-creeds-christendom-3-1977}
\item Those who obstinately resist the Spirit can grieve and lose him --- \cite[Solid Declaration II.83]{concordia-triglotta-1921}; \cite[Solid Declaration XI.42]{concordia-triglotta-1921}
\item The blasphemy against the Spirit is impenitence to the end: grace refused till death --- \cite[325]{npnf1-06-1888}; \cite[II-II q.14 a.1]{aquinas-summa-1920}; \cite[74]{calvin-harmony-evangelists-2-1845}
\item an inference step --- \cite[325]{npnf1-06-1888}; \cite[On Rebuke and Grace, ch. 43]{npnf1-05-1887}
\end{itemize}
\item \emph{Thomist (Aquinas; Garrigou-Lagrange)}, weakest at \emph{disputed}:
\begin{itemize}
\item In converting grace God is the sole mover: the mind is moved and does not move --- \cite[I-II q.111 a.2]{aquinas-summa-1920}; \cite[ch. IV]{garrigou-lagrange-grace-1952}
\item The free will, moved by God, cooperates by assenting to God’s call --- \cite[Trent, Session VI (1547), canon 4]{tanner-decrees-1990}; \cite[45]{waterworth-trent-1848}
\item The free will, moved by God, can refuse its assent if it wishes --- \cite[Trent, Session VI (1547), canon 4]{tanner-decrees-1990}; \cite[45]{waterworth-trent-1848}
\item Grace truly sufficient, giving the power to do good, is resisted --- \cite[ch. VI]{garrigou-lagrange-grace-1952}
\item Justifying grace, once given, can be lost --- \cite[II-II q.24 a.11]{aquinas-summa-1920}; \cite[Trent, Session VI (1547), canon 23]{tanner-decrees-1990}
\item Grace is received by faith alone (encoded as a denial, so it ranks disputed) --- \cite[II-II q.4 a.3]{aquinas-summa-1920}; \cite[I-II q.113 a.4 ad 1]{aquinas-summa-1920}; \cite[Trent, Session VI (1547), canon 9]{tanner-decrees-1990}
\item an inference step --- \cite{pohle-molinism-1911}; \cite[I-II q.112 a.3]{aquinas-summa-1920}; \cite[ch. VII]{garrigou-lagrange-grace-1952}
\end{itemize}
\end{itemize}
\par\medskip\noindent\textbf{\texttt{reformed\_\allowbreak{}defensible}}\par\nobreak\noindent
\textbf{Dort can be defended}, together with what every party holds and the third way: some maximal defensible position, one no other party can join, holds all three.

{\small\begin{verbatim}
theorem reformed_defensible : Framework.CredulouslyAccepted
    solaGratiaDispute.defeats Party.reformed
-- axioms: propext, Classical.choice, Quot.sound
\end{verbatim}}
\par\medskip\noindent\textbf{\texttt{dort\allowbreak{}Answered\allowbreak{}By\allowbreak{}The\allowbreak{}Remonstrants}}\par\nobreak\noindent
Why Dort is not forced: a defensible position holds the Remonstrants, and they defeat Dort.

\paragraph{\emph{Reformed (the Canons of Dort)} is not accepted on every resolution.}
\begin{itemize}
\item A position holding \emph{Romans 11:6 (by grace, so not by works)}, \emph{Remonstrant (the Five Articles, 1610)}, \emph{Lutheran (the Formula of Concord)}, \emph{Wesleyan (preventing grace)}, \emph{The third way (an ability given, received by faith)}, \emph{The objection from the Spirit’s resistance} can be held, and it holds \emph{Remonstrant (the Five Articles, 1610)}; \emph{Remonstrant (the Five Articles, 1610)} conflicts with \emph{Reformed (the Canons of Dort)}: \emph{Remonstrant (the Five Articles, 1610)} defeats \emph{Reformed (the Canons of Dort)}, and \emph{Reformed (the Canons of Dort)} defeats \emph{Remonstrant (the Five Articles, 1610)}.
\item None of \emph{Romans 11:6 (by grace, so not by works)}, \emph{Remonstrant (the Five Articles, 1610)}, \emph{Lutheran (the Formula of Concord)}, \emph{Wesleyan (preventing grace)}, \emph{The third way (an ability given, received by faith)}, \emph{The objection from the Spirit’s resistance} defeats another, and each attack on them is answered from among them:
\begin{itemize}
\item \emph{Reformed (the Canons of Dort)} defeats \emph{Remonstrant (the Five Articles, 1610)}, and \emph{Remonstrant (the Five Articles, 1610)} defeats \emph{Reformed (the Canons of Dort)}.
\item \emph{Thomist (Aquinas; Garrigou-Lagrange)} defeats \emph{Remonstrant (the Five Articles, 1610)}, and \emph{Remonstrant (the Five Articles, 1610)} defeats \emph{Thomist (Aquinas; Garrigou-Lagrange)}.
\item \emph{Reformed (the Canons of Dort)} defeats \emph{Lutheran (the Formula of Concord)}, and \emph{Remonstrant (the Five Articles, 1610)} defeats \emph{Reformed (the Canons of Dort)}.
\item \emph{Trent (Session VI, canon 4, on the Molinist reading)} defeats \emph{Lutheran (the Formula of Concord)}, and \emph{Lutheran (the Formula of Concord)} defeats \emph{Trent (Session VI, canon 4, on the Molinist reading)}.
\item \emph{Thomist (Aquinas; Garrigou-Lagrange)} defeats \emph{Lutheran (the Formula of Concord)}, and \emph{Remonstrant (the Five Articles, 1610)} defeats \emph{Thomist (Aquinas; Garrigou-Lagrange)}.
\item \emph{Reformed (the Canons of Dort)} defeats \emph{Wesleyan (preventing grace)}, and \emph{Remonstrant (the Five Articles, 1610)} defeats \emph{Reformed (the Canons of Dort)}.
\item \emph{Trent (Session VI, canon 4, on the Molinist reading)} defeats \emph{Wesleyan (preventing grace)}, and \emph{Lutheran (the Formula of Concord)} defeats \emph{Trent (Session VI, canon 4, on the Molinist reading)}.
\item \emph{Thomist (Aquinas; Garrigou-Lagrange)} defeats \emph{Wesleyan (preventing grace)}, and \emph{Remonstrant (the Five Articles, 1610)} defeats \emph{Thomist (Aquinas; Garrigou-Lagrange)}.
\item \emph{Trent (Session VI, canon 4, on the Molinist reading)} defeats \emph{The third way (an ability given, received by faith)}, and \emph{Lutheran (the Formula of Concord)} defeats \emph{Trent (Session VI, canon 4, on the Molinist reading)}.
\item \emph{Thomist (Aquinas; Garrigou-Lagrange)} defeats \emph{The third way (an ability given, received by faith)}, and \emph{Remonstrant (the Five Articles, 1610)} defeats \emph{Thomist (Aquinas; Garrigou-Lagrange)}.
\item \emph{Reformed (the Canons of Dort)} defeats \emph{The objection from the Spirit’s resistance}, and \emph{Remonstrant (the Five Articles, 1610)} defeats \emph{Reformed (the Canons of Dort)}.
\item \emph{Thomist (Aquinas; Garrigou-Lagrange)} defeats \emph{The objection from the Spirit’s resistance}, and \emph{Remonstrant (the Five Articles, 1610)} defeats \emph{Thomist (Aquinas; Garrigou-Lagrange)}.
\end{itemize}
\end{itemize}
\par\noindent\textbf{What this rests on.} The reasons state 15 defeats and 42 absences of defeat, each a cell of the defeat table, \texttt{sola\allowbreak{}Gratia\allowbreak{}Dispute\_\allowbreak{}defeats}, computed from the two parties' premises and checked by the kernel. Whether an attack survives turns on the attacker's weakest link and on the rating of what it attacks, so a changed rating can change the verdict. The weakest links:
\begin{itemize}
\item \emph{Remonstrant (the Five Articles, 1610)}, weakest at \emph{disputed}:
\begin{itemize}
\item What Stephen’s hearers resisted in Acts 7:51 was the grace that saves --- \cite[Five Arminian Articles (1610), art. 4]{schaff-creeds-christendom-3-1977}
\end{itemize}
\item \emph{Reformed (the Canons of Dort)}, weakest at \emph{disputed}:
\begin{itemize}
\item What Stephen’s hearers resisted in Acts 7:51 was the grace that saves (encoded as a denial, so it ranks disputed) --- \cite[Canons of Dort, III/IV, art. 9]{schaff-creeds-christendom-3-1977}
\item an inference step --- \cite[ad loc. Eph 2:8]{robertson-word-pictures-4-1931}; \cite[228--229]{calvin-galatians-ephesians-1854}
\item an inference step --- \cite[600]{arminius-works-1-1825}; \cite[681]{arminius-works-1-1825}; \cite[Five Arminian Articles (1610), art. 4]{schaff-creeds-christendom-3-1977}
\item an inference step --- \cite[527--565]{arminius-works-3-1853}
\item an inference step --- \cite[Solid Declaration II.89]{concordia-triglotta-1921}; \cite[Solid Declaration II.60]{concordia-triglotta-1921}; \cite[Solid Declaration XI.41]{concordia-triglotta-1921}
\end{itemize}
\item \emph{Romans 11:6 (by grace, so not by works)}, weakest at \emph{consensus}:
\begin{itemize}
\item Romans 11:6 — if it is by grace, it is no longer on the basis of works --- Rom 11:6
\item an inference step --- \cite[Trent, Session VI (1547), ch. 8]{tanner-decrees-1990}; \cite[36]{waterworth-trent-1848}; \cite[Five Arminian Articles (1610), art. 4]{schaff-creeds-christendom-3-1977}; \cite[Canons of Dort, III/IV, art. 10]{schaff-creeds-christendom-3-1977}
\end{itemize}
\item \emph{Lutheran (the Formula of Concord)}, weakest at \emph{disputed}:
\begin{itemize}
\item Those who obstinately resist the Spirit can grieve and lose him --- \cite[Solid Declaration II.83]{concordia-triglotta-1921}; \cite[Solid Declaration XI.42]{concordia-triglotta-1921}
\item Conversion is God’s work alone: God regenerates, and the will does nothing to it --- \cite[Solid Declaration II.89]{concordia-triglotta-1921}; \cite[Solid Declaration II.90]{concordia-triglotta-1921}
\item Grace is received by faith alone --- \cite[Augsburg Confession IV.1–3]{concordia-triglotta-1921}; \cite[Sermon 1, III.1]{wesley-sermons-1872}; \cite[Trent, Session VI (1547), canon 9]{tanner-decrees-1990}
\end{itemize}
\item \emph{Wesleyan (preventing grace)}, weakest at \emph{disputed}:
\begin{itemize}
\item The grace that saves can be finally refused --- \cite[397]{wesley-works-9-1827}
\item Grace is received by faith alone --- \cite[Augsburg Confession IV.1–3]{concordia-triglotta-1921}; \cite[Sermon 1, III.1]{wesley-sermons-1872}; \cite[Trent, Session VI (1547), canon 9]{tanner-decrees-1990}
\end{itemize}
\item \emph{The third way (an ability given, received by faith)}, weakest at \emph{disputed}:
\begin{itemize}
\item Grace is received by faith alone --- \cite[Augsburg Confession IV.1–3]{concordia-triglotta-1921}; \cite[Sermon 1, III.1]{wesley-sermons-1872}; \cite[Trent, Session VI (1547), canon 9]{tanner-decrees-1990}
\end{itemize}
\item \emph{The objection from the Spirit’s resistance}, weakest at \emph{disputed}:
\begin{itemize}
\item What Stephen’s hearers resisted in Acts 7:51 was the grace that saves --- \cite[Five Arminian Articles (1610), art. 4]{schaff-creeds-christendom-3-1977}
\item Those who obstinately resist the Spirit can grieve and lose him --- \cite[Solid Declaration II.83]{concordia-triglotta-1921}; \cite[Solid Declaration XI.42]{concordia-triglotta-1921}
\item The blasphemy against the Spirit is impenitence to the end: grace refused till death --- \cite[325]{npnf1-06-1888}; \cite[II-II q.14 a.1]{aquinas-summa-1920}; \cite[74]{calvin-harmony-evangelists-2-1845}
\item an inference step --- \cite[325]{npnf1-06-1888}; \cite[On Rebuke and Grace, ch. 43]{npnf1-05-1887}
\end{itemize}
\item \emph{Trent (Session VI, canon 4, on the Molinist reading)}, weakest at \emph{disputed}:
\begin{itemize}
\item The free will, moved by God, cooperates by assenting to God’s call --- \cite[Trent, Session VI (1547), canon 4]{tanner-decrees-1990}; \cite[45]{waterworth-trent-1848}
\item The free will, moved by God, can refuse its assent if it wishes --- \cite[Trent, Session VI (1547), canon 4]{tanner-decrees-1990}; \cite[45]{waterworth-trent-1848}
\item Grace is received by faith alone (encoded as a denial, so it ranks disputed) --- \cite[Augsburg Confession IV.1–3]{concordia-triglotta-1921}; \cite[Sermon 1, III.1]{wesley-sermons-1872}; \cite[Trent, Session VI (1547), canon 9]{tanner-decrees-1990}
\item an inference step --- \cite[I-II q.111 a.2]{aquinas-summa-1920}; \cite[Canons of Dort, III/IV, art. 12]{schaff-creeds-christendom-3-1977}; \cite[Trent, Session VI (1547), canon 4]{tanner-decrees-1990}; \cite[45]{waterworth-trent-1848}
\item an inference step --- \cite[I-II q.10 a.4 ad 3]{aquinas-summa-1920}; \cite[I-II q.112 a.3]{aquinas-summa-1920}; \cite[Trent, Session VI (1547), canon 4]{tanner-decrees-1990}; \cite[45]{waterworth-trent-1848}
\end{itemize}
\item \emph{Thomist (Aquinas; Garrigou-Lagrange)}, weakest at \emph{disputed}:
\begin{itemize}
\item In converting grace God is the sole mover: the mind is moved and does not move --- \cite[I-II q.111 a.2]{aquinas-summa-1920}; \cite[ch. IV]{garrigou-lagrange-grace-1952}
\item The free will, moved by God, cooperates by assenting to God’s call --- \cite[Trent, Session VI (1547), canon 4]{tanner-decrees-1990}; \cite[45]{waterworth-trent-1848}
\item The free will, moved by God, can refuse its assent if it wishes --- \cite[Trent, Session VI (1547), canon 4]{tanner-decrees-1990}; \cite[45]{waterworth-trent-1848}
\item Grace truly sufficient, giving the power to do good, is resisted --- \cite[ch. VI]{garrigou-lagrange-grace-1952}
\item Justifying grace, once given, can be lost --- \cite[II-II q.24 a.11]{aquinas-summa-1920}; \cite[Trent, Session VI (1547), canon 23]{tanner-decrees-1990}
\item Grace is received by faith alone (encoded as a denial, so it ranks disputed) --- \cite[II-II q.4 a.3]{aquinas-summa-1920}; \cite[I-II q.113 a.4 ad 1]{aquinas-summa-1920}; \cite[Trent, Session VI (1547), canon 9]{tanner-decrees-1990}
\item an inference step --- \cite{pohle-molinism-1911}; \cite[I-II q.112 a.3]{aquinas-summa-1920}; \cite[ch. VII]{garrigou-lagrange-grace-1952}
\end{itemize}
\end{itemize}
\par\medskip\noindent\textbf{\texttt{reformed\_\allowbreak{}not\_\allowbreak{}forced}}\par\nobreak\noindent
Nor is Dort forced. A maximal defensible position holds the Remonstrants, the Lutherans, Wesley, the third way and the objection, and cannot hold Dort with them.

{\small\begin{verbatim}
theorem reformed_not_forced : ¬Framework.SkepticallyAccepted
    solaGratiaDispute.defeats Party.reformed
-- axioms: propext, Classical.choice, Quot.sound
\end{verbatim}}
\par\medskip\noindent\textbf{\texttt{lutherans\allowbreak{}Defended}}\par\nobreak\noindent
Why the Lutherans can be defended: with the Remonstrants and Wesley they answer both Dort and Trent.

\paragraph{\emph{Lutheran (the Formula of Concord)} is accepted on some resolution.}
\begin{itemize}
\item A position holding \emph{Romans 11:6 (by grace, so not by works)}, \emph{Remonstrant (the Five Articles, 1610)}, \emph{Lutheran (the Formula of Concord)}, \emph{Wesleyan (preventing grace)}, \emph{The third way (an ability given, received by faith)}, \emph{The objection from the Spirit’s resistance} can be held.
\item None of \emph{Romans 11:6 (by grace, so not by works)}, \emph{Remonstrant (the Five Articles, 1610)}, \emph{Lutheran (the Formula of Concord)}, \emph{Wesleyan (preventing grace)}, \emph{The third way (an ability given, received by faith)}, \emph{The objection from the Spirit’s resistance} defeats another, and each attack on them is answered from among them:
\begin{itemize}
\item \emph{Reformed (the Canons of Dort)} defeats \emph{Remonstrant (the Five Articles, 1610)}, and \emph{Remonstrant (the Five Articles, 1610)} defeats \emph{Reformed (the Canons of Dort)}.
\item \emph{Thomist (Aquinas; Garrigou-Lagrange)} defeats \emph{Remonstrant (the Five Articles, 1610)}, and \emph{Remonstrant (the Five Articles, 1610)} defeats \emph{Thomist (Aquinas; Garrigou-Lagrange)}.
\item \emph{Reformed (the Canons of Dort)} defeats \emph{Lutheran (the Formula of Concord)}, and \emph{Remonstrant (the Five Articles, 1610)} defeats \emph{Reformed (the Canons of Dort)}.
\item \emph{Trent (Session VI, canon 4, on the Molinist reading)} defeats \emph{Lutheran (the Formula of Concord)}, and \emph{Lutheran (the Formula of Concord)} defeats \emph{Trent (Session VI, canon 4, on the Molinist reading)}.
\item \emph{Thomist (Aquinas; Garrigou-Lagrange)} defeats \emph{Lutheran (the Formula of Concord)}, and \emph{Remonstrant (the Five Articles, 1610)} defeats \emph{Thomist (Aquinas; Garrigou-Lagrange)}.
\item \emph{Reformed (the Canons of Dort)} defeats \emph{Wesleyan (preventing grace)}, and \emph{Remonstrant (the Five Articles, 1610)} defeats \emph{Reformed (the Canons of Dort)}.
\item \emph{Trent (Session VI, canon 4, on the Molinist reading)} defeats \emph{Wesleyan (preventing grace)}, and \emph{Lutheran (the Formula of Concord)} defeats \emph{Trent (Session VI, canon 4, on the Molinist reading)}.
\item \emph{Thomist (Aquinas; Garrigou-Lagrange)} defeats \emph{Wesleyan (preventing grace)}, and \emph{Remonstrant (the Five Articles, 1610)} defeats \emph{Thomist (Aquinas; Garrigou-Lagrange)}.
\item \emph{Trent (Session VI, canon 4, on the Molinist reading)} defeats \emph{The third way (an ability given, received by faith)}, and \emph{Lutheran (the Formula of Concord)} defeats \emph{Trent (Session VI, canon 4, on the Molinist reading)}.
\item \emph{Thomist (Aquinas; Garrigou-Lagrange)} defeats \emph{The third way (an ability given, received by faith)}, and \emph{Remonstrant (the Five Articles, 1610)} defeats \emph{Thomist (Aquinas; Garrigou-Lagrange)}.
\item \emph{Reformed (the Canons of Dort)} defeats \emph{The objection from the Spirit’s resistance}, and \emph{Remonstrant (the Five Articles, 1610)} defeats \emph{Reformed (the Canons of Dort)}.
\item \emph{Thomist (Aquinas; Garrigou-Lagrange)} defeats \emph{The objection from the Spirit’s resistance}, and \emph{Remonstrant (the Five Articles, 1610)} defeats \emph{Thomist (Aquinas; Garrigou-Lagrange)}.
\end{itemize}
\end{itemize}
\par\noindent\textbf{What this rests on.} The reasons state 15 defeats and 42 absences of defeat, each a cell of the defeat table, \texttt{sola\allowbreak{}Gratia\allowbreak{}Dispute\_\allowbreak{}defeats}, computed from the two parties' premises and checked by the kernel. Whether an attack survives turns on the attacker's weakest link and on the rating of what it attacks, so a changed rating can change the verdict. The weakest links:
\begin{itemize}
\item \emph{Romans 11:6 (by grace, so not by works)}, weakest at \emph{consensus}:
\begin{itemize}
\item Romans 11:6 — if it is by grace, it is no longer on the basis of works --- Rom 11:6
\item an inference step --- \cite[Trent, Session VI (1547), ch. 8]{tanner-decrees-1990}; \cite[36]{waterworth-trent-1848}; \cite[Five Arminian Articles (1610), art. 4]{schaff-creeds-christendom-3-1977}; \cite[Canons of Dort, III/IV, art. 10]{schaff-creeds-christendom-3-1977}
\end{itemize}
\item \emph{Remonstrant (the Five Articles, 1610)}, weakest at \emph{disputed}:
\begin{itemize}
\item What Stephen’s hearers resisted in Acts 7:51 was the grace that saves --- \cite[Five Arminian Articles (1610), art. 4]{schaff-creeds-christendom-3-1977}
\end{itemize}
\item \emph{Lutheran (the Formula of Concord)}, weakest at \emph{disputed}:
\begin{itemize}
\item Those who obstinately resist the Spirit can grieve and lose him --- \cite[Solid Declaration II.83]{concordia-triglotta-1921}; \cite[Solid Declaration XI.42]{concordia-triglotta-1921}
\item Conversion is God’s work alone: God regenerates, and the will does nothing to it --- \cite[Solid Declaration II.89]{concordia-triglotta-1921}; \cite[Solid Declaration II.90]{concordia-triglotta-1921}
\item Grace is received by faith alone --- \cite[Augsburg Confession IV.1–3]{concordia-triglotta-1921}; \cite[Sermon 1, III.1]{wesley-sermons-1872}; \cite[Trent, Session VI (1547), canon 9]{tanner-decrees-1990}
\end{itemize}
\item \emph{Wesleyan (preventing grace)}, weakest at \emph{disputed}:
\begin{itemize}
\item The grace that saves can be finally refused --- \cite[397]{wesley-works-9-1827}
\item Grace is received by faith alone --- \cite[Augsburg Confession IV.1–3]{concordia-triglotta-1921}; \cite[Sermon 1, III.1]{wesley-sermons-1872}; \cite[Trent, Session VI (1547), canon 9]{tanner-decrees-1990}
\end{itemize}
\item \emph{The third way (an ability given, received by faith)}, weakest at \emph{disputed}:
\begin{itemize}
\item Grace is received by faith alone --- \cite[Augsburg Confession IV.1–3]{concordia-triglotta-1921}; \cite[Sermon 1, III.1]{wesley-sermons-1872}; \cite[Trent, Session VI (1547), canon 9]{tanner-decrees-1990}
\end{itemize}
\item \emph{The objection from the Spirit’s resistance}, weakest at \emph{disputed}:
\begin{itemize}
\item What Stephen’s hearers resisted in Acts 7:51 was the grace that saves --- \cite[Five Arminian Articles (1610), art. 4]{schaff-creeds-christendom-3-1977}
\item Those who obstinately resist the Spirit can grieve and lose him --- \cite[Solid Declaration II.83]{concordia-triglotta-1921}; \cite[Solid Declaration XI.42]{concordia-triglotta-1921}
\item The blasphemy against the Spirit is impenitence to the end: grace refused till death --- \cite[325]{npnf1-06-1888}; \cite[II-II q.14 a.1]{aquinas-summa-1920}; \cite[74]{calvin-harmony-evangelists-2-1845}
\item an inference step --- \cite[325]{npnf1-06-1888}; \cite[On Rebuke and Grace, ch. 43]{npnf1-05-1887}
\end{itemize}
\item \emph{Reformed (the Canons of Dort)}, weakest at \emph{disputed}:
\begin{itemize}
\item What Stephen’s hearers resisted in Acts 7:51 was the grace that saves (encoded as a denial, so it ranks disputed) --- \cite[Canons of Dort, III/IV, art. 9]{schaff-creeds-christendom-3-1977}
\item an inference step --- \cite[ad loc. Eph 2:8]{robertson-word-pictures-4-1931}; \cite[228--229]{calvin-galatians-ephesians-1854}
\item an inference step --- \cite[600]{arminius-works-1-1825}; \cite[681]{arminius-works-1-1825}; \cite[Five Arminian Articles (1610), art. 4]{schaff-creeds-christendom-3-1977}
\item an inference step --- \cite[527--565]{arminius-works-3-1853}
\item an inference step --- \cite[Solid Declaration II.89]{concordia-triglotta-1921}; \cite[Solid Declaration II.60]{concordia-triglotta-1921}; \cite[Solid Declaration XI.41]{concordia-triglotta-1921}
\end{itemize}
\item \emph{Trent (Session VI, canon 4, on the Molinist reading)}, weakest at \emph{disputed}:
\begin{itemize}
\item The free will, moved by God, cooperates by assenting to God’s call --- \cite[Trent, Session VI (1547), canon 4]{tanner-decrees-1990}; \cite[45]{waterworth-trent-1848}
\item The free will, moved by God, can refuse its assent if it wishes --- \cite[Trent, Session VI (1547), canon 4]{tanner-decrees-1990}; \cite[45]{waterworth-trent-1848}
\item Grace is received by faith alone (encoded as a denial, so it ranks disputed) --- \cite[Augsburg Confession IV.1–3]{concordia-triglotta-1921}; \cite[Sermon 1, III.1]{wesley-sermons-1872}; \cite[Trent, Session VI (1547), canon 9]{tanner-decrees-1990}
\item an inference step --- \cite[I-II q.111 a.2]{aquinas-summa-1920}; \cite[Canons of Dort, III/IV, art. 12]{schaff-creeds-christendom-3-1977}; \cite[Trent, Session VI (1547), canon 4]{tanner-decrees-1990}; \cite[45]{waterworth-trent-1848}
\item an inference step --- \cite[I-II q.10 a.4 ad 3]{aquinas-summa-1920}; \cite[I-II q.112 a.3]{aquinas-summa-1920}; \cite[Trent, Session VI (1547), canon 4]{tanner-decrees-1990}; \cite[45]{waterworth-trent-1848}
\end{itemize}
\item \emph{Thomist (Aquinas; Garrigou-Lagrange)}, weakest at \emph{disputed}:
\begin{itemize}
\item In converting grace God is the sole mover: the mind is moved and does not move --- \cite[I-II q.111 a.2]{aquinas-summa-1920}; \cite[ch. IV]{garrigou-lagrange-grace-1952}
\item The free will, moved by God, cooperates by assenting to God’s call --- \cite[Trent, Session VI (1547), canon 4]{tanner-decrees-1990}; \cite[45]{waterworth-trent-1848}
\item The free will, moved by God, can refuse its assent if it wishes --- \cite[Trent, Session VI (1547), canon 4]{tanner-decrees-1990}; \cite[45]{waterworth-trent-1848}
\item Grace truly sufficient, giving the power to do good, is resisted --- \cite[ch. VI]{garrigou-lagrange-grace-1952}
\item Justifying grace, once given, can be lost --- \cite[II-II q.24 a.11]{aquinas-summa-1920}; \cite[Trent, Session VI (1547), canon 23]{tanner-decrees-1990}
\item Grace is received by faith alone (encoded as a denial, so it ranks disputed) --- \cite[II-II q.4 a.3]{aquinas-summa-1920}; \cite[I-II q.113 a.4 ad 1]{aquinas-summa-1920}; \cite[Trent, Session VI (1547), canon 9]{tanner-decrees-1990}
\item an inference step --- \cite{pohle-molinism-1911}; \cite[I-II q.112 a.3]{aquinas-summa-1920}; \cite[ch. VII]{garrigou-lagrange-grace-1952}
\end{itemize}
\end{itemize}
\par\medskip\noindent\textbf{\texttt{lutheran\_\allowbreak{}defensible}}\par\nobreak\noindent
\textbf{The Lutheran position can be defended}, with the Remonstrants, Wesley, the third way and the objection.

{\small\begin{verbatim}
theorem lutheran_defensible : Framework.CredulouslyAccepted
    solaGratiaDispute.defeats Party.lutheran
-- axioms: propext, Classical.choice, Quot.sound
\end{verbatim}}
\par\medskip\noindent\textbf{\texttt{trent\allowbreak{}Defended}}\par\nobreak\noindent
Why Trent can be defended: it answers every attacker of its group itself.

\paragraph{\emph{Trent (Session VI, canon 4, on the Molinist reading)} is accepted on some resolution.}
\begin{itemize}
\item A position holding \emph{Romans 11:6 (by grace, so not by works)}, \emph{Trent (Session VI, canon 4, on the Molinist reading)}, \emph{Remonstrant (the Five Articles, 1610)}, \emph{The objection from the Spirit’s resistance} can be held.
\item None of \emph{Romans 11:6 (by grace, so not by works)}, \emph{Trent (Session VI, canon 4, on the Molinist reading)}, \emph{Remonstrant (the Five Articles, 1610)}, \emph{The objection from the Spirit’s resistance} defeats another, and each attack on them is answered from among them:
\begin{itemize}
\item \emph{Reformed (the Canons of Dort)} defeats \emph{Trent (Session VI, canon 4, on the Molinist reading)}, and \emph{Trent (Session VI, canon 4, on the Molinist reading)} defeats \emph{Reformed (the Canons of Dort)}.
\item \emph{Lutheran (the Formula of Concord)} defeats \emph{Trent (Session VI, canon 4, on the Molinist reading)}, and \emph{Trent (Session VI, canon 4, on the Molinist reading)} defeats \emph{Lutheran (the Formula of Concord)}.
\item \emph{Wesleyan (preventing grace)} defeats \emph{Trent (Session VI, canon 4, on the Molinist reading)}, and \emph{Trent (Session VI, canon 4, on the Molinist reading)} defeats \emph{Wesleyan (preventing grace)}.
\item \emph{The third way (an ability given, received by faith)} defeats \emph{Trent (Session VI, canon 4, on the Molinist reading)}, and \emph{Trent (Session VI, canon 4, on the Molinist reading)} defeats \emph{The third way (an ability given, received by faith)}.
\item \emph{Thomist (Aquinas; Garrigou-Lagrange)} defeats \emph{Trent (Session VI, canon 4, on the Molinist reading)}, and \emph{Trent (Session VI, canon 4, on the Molinist reading)} defeats \emph{Thomist (Aquinas; Garrigou-Lagrange)}.
\item \emph{Reformed (the Canons of Dort)} defeats \emph{Remonstrant (the Five Articles, 1610)}, and \emph{Trent (Session VI, canon 4, on the Molinist reading)} defeats \emph{Reformed (the Canons of Dort)}.
\item \emph{Thomist (Aquinas; Garrigou-Lagrange)} defeats \emph{Remonstrant (the Five Articles, 1610)}, and \emph{Trent (Session VI, canon 4, on the Molinist reading)} defeats \emph{Thomist (Aquinas; Garrigou-Lagrange)}.
\item \emph{Reformed (the Canons of Dort)} defeats \emph{The objection from the Spirit’s resistance}, and \emph{Trent (Session VI, canon 4, on the Molinist reading)} defeats \emph{Reformed (the Canons of Dort)}.
\item \emph{Thomist (Aquinas; Garrigou-Lagrange)} defeats \emph{The objection from the Spirit’s resistance}, and \emph{Trent (Session VI, canon 4, on the Molinist reading)} defeats \emph{Thomist (Aquinas; Garrigou-Lagrange)}.
\end{itemize}
\end{itemize}
\par\noindent\textbf{What this rests on.} The reasons state 14 defeats and 27 absences of defeat, each a cell of the defeat table, \texttt{sola\allowbreak{}Gratia\allowbreak{}Dispute\_\allowbreak{}defeats}, computed from the two parties' premises and checked by the kernel. Whether an attack survives turns on the attacker's weakest link and on the rating of what it attacks, so a changed rating can change the verdict. The weakest links:
\begin{itemize}
\item \emph{Romans 11:6 (by grace, so not by works)}, weakest at \emph{consensus}:
\begin{itemize}
\item Romans 11:6 — if it is by grace, it is no longer on the basis of works --- Rom 11:6
\item an inference step --- \cite[Trent, Session VI (1547), ch. 8]{tanner-decrees-1990}; \cite[36]{waterworth-trent-1848}; \cite[Five Arminian Articles (1610), art. 4]{schaff-creeds-christendom-3-1977}; \cite[Canons of Dort, III/IV, art. 10]{schaff-creeds-christendom-3-1977}
\end{itemize}
\item \emph{Trent (Session VI, canon 4, on the Molinist reading)}, weakest at \emph{disputed}:
\begin{itemize}
\item The free will, moved by God, cooperates by assenting to God’s call --- \cite[Trent, Session VI (1547), canon 4]{tanner-decrees-1990}; \cite[45]{waterworth-trent-1848}
\item The free will, moved by God, can refuse its assent if it wishes --- \cite[Trent, Session VI (1547), canon 4]{tanner-decrees-1990}; \cite[45]{waterworth-trent-1848}
\item Grace is received by faith alone (encoded as a denial, so it ranks disputed) --- \cite[Augsburg Confession IV.1–3]{concordia-triglotta-1921}; \cite[Sermon 1, III.1]{wesley-sermons-1872}; \cite[Trent, Session VI (1547), canon 9]{tanner-decrees-1990}
\item an inference step --- \cite[I-II q.111 a.2]{aquinas-summa-1920}; \cite[Canons of Dort, III/IV, art. 12]{schaff-creeds-christendom-3-1977}; \cite[Trent, Session VI (1547), canon 4]{tanner-decrees-1990}; \cite[45]{waterworth-trent-1848}
\item an inference step --- \cite[I-II q.10 a.4 ad 3]{aquinas-summa-1920}; \cite[I-II q.112 a.3]{aquinas-summa-1920}; \cite[Trent, Session VI (1547), canon 4]{tanner-decrees-1990}; \cite[45]{waterworth-trent-1848}
\end{itemize}
\item \emph{Remonstrant (the Five Articles, 1610)}, weakest at \emph{disputed}:
\begin{itemize}
\item What Stephen’s hearers resisted in Acts 7:51 was the grace that saves --- \cite[Five Arminian Articles (1610), art. 4]{schaff-creeds-christendom-3-1977}
\end{itemize}
\item \emph{The objection from the Spirit’s resistance}, weakest at \emph{disputed}:
\begin{itemize}
\item What Stephen’s hearers resisted in Acts 7:51 was the grace that saves --- \cite[Five Arminian Articles (1610), art. 4]{schaff-creeds-christendom-3-1977}
\item Those who obstinately resist the Spirit can grieve and lose him --- \cite[Solid Declaration II.83]{concordia-triglotta-1921}; \cite[Solid Declaration XI.42]{concordia-triglotta-1921}
\item The blasphemy against the Spirit is impenitence to the end: grace refused till death --- \cite[325]{npnf1-06-1888}; \cite[II-II q.14 a.1]{aquinas-summa-1920}; \cite[74]{calvin-harmony-evangelists-2-1845}
\item an inference step --- \cite[325]{npnf1-06-1888}; \cite[On Rebuke and Grace, ch. 43]{npnf1-05-1887}
\end{itemize}
\item \emph{Reformed (the Canons of Dort)}, weakest at \emph{disputed}:
\begin{itemize}
\item What Stephen’s hearers resisted in Acts 7:51 was the grace that saves (encoded as a denial, so it ranks disputed) --- \cite[Canons of Dort, III/IV, art. 9]{schaff-creeds-christendom-3-1977}
\item an inference step --- \cite[ad loc. Eph 2:8]{robertson-word-pictures-4-1931}; \cite[228--229]{calvin-galatians-ephesians-1854}
\item an inference step --- \cite[600]{arminius-works-1-1825}; \cite[681]{arminius-works-1-1825}; \cite[Five Arminian Articles (1610), art. 4]{schaff-creeds-christendom-3-1977}
\item an inference step --- \cite[527--565]{arminius-works-3-1853}
\item an inference step --- \cite[Solid Declaration II.89]{concordia-triglotta-1921}; \cite[Solid Declaration II.60]{concordia-triglotta-1921}; \cite[Solid Declaration XI.41]{concordia-triglotta-1921}
\end{itemize}
\item \emph{Lutheran (the Formula of Concord)}, weakest at \emph{disputed}:
\begin{itemize}
\item Those who obstinately resist the Spirit can grieve and lose him --- \cite[Solid Declaration II.83]{concordia-triglotta-1921}; \cite[Solid Declaration XI.42]{concordia-triglotta-1921}
\item Conversion is God’s work alone: God regenerates, and the will does nothing to it --- \cite[Solid Declaration II.89]{concordia-triglotta-1921}; \cite[Solid Declaration II.90]{concordia-triglotta-1921}
\item Grace is received by faith alone --- \cite[Augsburg Confession IV.1–3]{concordia-triglotta-1921}; \cite[Sermon 1, III.1]{wesley-sermons-1872}; \cite[Trent, Session VI (1547), canon 9]{tanner-decrees-1990}
\end{itemize}
\item \emph{Wesleyan (preventing grace)}, weakest at \emph{disputed}:
\begin{itemize}
\item The grace that saves can be finally refused --- \cite[397]{wesley-works-9-1827}
\item Grace is received by faith alone --- \cite[Augsburg Confession IV.1–3]{concordia-triglotta-1921}; \cite[Sermon 1, III.1]{wesley-sermons-1872}; \cite[Trent, Session VI (1547), canon 9]{tanner-decrees-1990}
\end{itemize}
\item \emph{The third way (an ability given, received by faith)}, weakest at \emph{disputed}:
\begin{itemize}
\item Grace is received by faith alone --- \cite[Augsburg Confession IV.1–3]{concordia-triglotta-1921}; \cite[Sermon 1, III.1]{wesley-sermons-1872}; \cite[Trent, Session VI (1547), canon 9]{tanner-decrees-1990}
\end{itemize}
\item \emph{Thomist (Aquinas; Garrigou-Lagrange)}, weakest at \emph{disputed}:
\begin{itemize}
\item In converting grace God is the sole mover: the mind is moved and does not move --- \cite[I-II q.111 a.2]{aquinas-summa-1920}; \cite[ch. IV]{garrigou-lagrange-grace-1952}
\item The free will, moved by God, cooperates by assenting to God’s call --- \cite[Trent, Session VI (1547), canon 4]{tanner-decrees-1990}; \cite[45]{waterworth-trent-1848}
\item The free will, moved by God, can refuse its assent if it wishes --- \cite[Trent, Session VI (1547), canon 4]{tanner-decrees-1990}; \cite[45]{waterworth-trent-1848}
\item Grace truly sufficient, giving the power to do good, is resisted --- \cite[ch. VI]{garrigou-lagrange-grace-1952}
\item Justifying grace, once given, can be lost --- \cite[II-II q.24 a.11]{aquinas-summa-1920}; \cite[Trent, Session VI (1547), canon 23]{tanner-decrees-1990}
\item Grace is received by faith alone (encoded as a denial, so it ranks disputed) --- \cite[II-II q.4 a.3]{aquinas-summa-1920}; \cite[I-II q.113 a.4 ad 1]{aquinas-summa-1920}; \cite[Trent, Session VI (1547), canon 9]{tanner-decrees-1990}
\item an inference step --- \cite{pohle-molinism-1911}; \cite[I-II q.112 a.3]{aquinas-summa-1920}; \cite[ch. VII]{garrigou-lagrange-grace-1952}
\end{itemize}
\end{itemize}
\par\medskip\noindent\textbf{\texttt{tridentine\_\allowbreak{}defensible}}\par\nobreak\noindent
\textbf{Trent can be defended}, with what every party holds, the Remonstrants and the objection, which agree with it that grace can be resisted. Only Dort attacks the Remonstrants and the objection, and Trent answers Dort.

{\small\begin{verbatim}
theorem tridentine_defensible : Framework.CredulouslyAccepted
    solaGratiaDispute.defeats Party.trent
-- axioms: propext, Classical.choice, Quot.sound
\end{verbatim}}
\par\medskip\noindent\textbf{\texttt{trent\allowbreak{}Answered\allowbreak{}By\allowbreak{}Dort}}\par\nobreak\noindent
Why Trent is not forced: a maximal defensible position holds Dort, and Dort defeats Trent.

\paragraph{\emph{Trent (Session VI, canon 4, on the Molinist reading)} is not accepted on every resolution.}
\begin{itemize}
\item A position holding \emph{Romans 11:6 (by grace, so not by works)}, \emph{Reformed (the Canons of Dort)}, \emph{The third way (an ability given, received by faith)} can be held, and it holds \emph{Reformed (the Canons of Dort)}; \emph{Reformed (the Canons of Dort)} conflicts with \emph{Trent (Session VI, canon 4, on the Molinist reading)}: \emph{Reformed (the Canons of Dort)} defeats \emph{Trent (Session VI, canon 4, on the Molinist reading)}, and \emph{Trent (Session VI, canon 4, on the Molinist reading)} defeats \emph{Reformed (the Canons of Dort)}.
\item None of \emph{Romans 11:6 (by grace, so not by works)}, \emph{Reformed (the Canons of Dort)}, \emph{The third way (an ability given, received by faith)} defeats another, and each attack on them is answered from among them:
\begin{itemize}
\item \emph{Trent (Session VI, canon 4, on the Molinist reading)} defeats \emph{Reformed (the Canons of Dort)}, and \emph{Reformed (the Canons of Dort)} defeats \emph{Trent (Session VI, canon 4, on the Molinist reading)}.
\item \emph{Remonstrant (the Five Articles, 1610)} defeats \emph{Reformed (the Canons of Dort)}, and \emph{Reformed (the Canons of Dort)} defeats \emph{Remonstrant (the Five Articles, 1610)}.
\item \emph{Lutheran (the Formula of Concord)} defeats \emph{Reformed (the Canons of Dort)}, and \emph{Reformed (the Canons of Dort)} defeats \emph{Lutheran (the Formula of Concord)}.
\item \emph{Wesleyan (preventing grace)} defeats \emph{Reformed (the Canons of Dort)}, and \emph{Reformed (the Canons of Dort)} defeats \emph{Wesleyan (preventing grace)}.
\item \emph{The objection from the Spirit’s resistance} defeats \emph{Reformed (the Canons of Dort)}, and \emph{Reformed (the Canons of Dort)} defeats \emph{The objection from the Spirit’s resistance}.
\item \emph{Trent (Session VI, canon 4, on the Molinist reading)} defeats \emph{The third way (an ability given, received by faith)}, and \emph{Reformed (the Canons of Dort)} defeats \emph{Trent (Session VI, canon 4, on the Molinist reading)}.
\item \emph{Thomist (Aquinas; Garrigou-Lagrange)} defeats \emph{The third way (an ability given, received by faith)}, and \emph{The third way (an ability given, received by faith)} defeats \emph{Thomist (Aquinas; Garrigou-Lagrange)}.
\end{itemize}
\end{itemize}
\par\noindent\textbf{What this rests on.} The reasons state 13 defeats and 20 absences of defeat, each a cell of the defeat table, \texttt{sola\allowbreak{}Gratia\allowbreak{}Dispute\_\allowbreak{}defeats}, computed from the two parties' premises and checked by the kernel. Whether an attack survives turns on the attacker's weakest link and on the rating of what it attacks, so a changed rating can change the verdict. The weakest links:
\begin{itemize}
\item \emph{Reformed (the Canons of Dort)}, weakest at \emph{disputed}:
\begin{itemize}
\item What Stephen’s hearers resisted in Acts 7:51 was the grace that saves (encoded as a denial, so it ranks disputed) --- \cite[Canons of Dort, III/IV, art. 9]{schaff-creeds-christendom-3-1977}
\item an inference step --- \cite[ad loc. Eph 2:8]{robertson-word-pictures-4-1931}; \cite[228--229]{calvin-galatians-ephesians-1854}
\item an inference step --- \cite[600]{arminius-works-1-1825}; \cite[681]{arminius-works-1-1825}; \cite[Five Arminian Articles (1610), art. 4]{schaff-creeds-christendom-3-1977}
\item an inference step --- \cite[527--565]{arminius-works-3-1853}
\item an inference step --- \cite[Solid Declaration II.89]{concordia-triglotta-1921}; \cite[Solid Declaration II.60]{concordia-triglotta-1921}; \cite[Solid Declaration XI.41]{concordia-triglotta-1921}
\end{itemize}
\item \emph{Trent (Session VI, canon 4, on the Molinist reading)}, weakest at \emph{disputed}:
\begin{itemize}
\item The free will, moved by God, cooperates by assenting to God’s call --- \cite[Trent, Session VI (1547), canon 4]{tanner-decrees-1990}; \cite[45]{waterworth-trent-1848}
\item The free will, moved by God, can refuse its assent if it wishes --- \cite[Trent, Session VI (1547), canon 4]{tanner-decrees-1990}; \cite[45]{waterworth-trent-1848}
\item Grace is received by faith alone (encoded as a denial, so it ranks disputed) --- \cite[Augsburg Confession IV.1–3]{concordia-triglotta-1921}; \cite[Sermon 1, III.1]{wesley-sermons-1872}; \cite[Trent, Session VI (1547), canon 9]{tanner-decrees-1990}
\item an inference step --- \cite[I-II q.111 a.2]{aquinas-summa-1920}; \cite[Canons of Dort, III/IV, art. 12]{schaff-creeds-christendom-3-1977}; \cite[Trent, Session VI (1547), canon 4]{tanner-decrees-1990}; \cite[45]{waterworth-trent-1848}
\item an inference step --- \cite[I-II q.10 a.4 ad 3]{aquinas-summa-1920}; \cite[I-II q.112 a.3]{aquinas-summa-1920}; \cite[Trent, Session VI (1547), canon 4]{tanner-decrees-1990}; \cite[45]{waterworth-trent-1848}
\end{itemize}
\item \emph{Romans 11:6 (by grace, so not by works)}, weakest at \emph{consensus}:
\begin{itemize}
\item Romans 11:6 — if it is by grace, it is no longer on the basis of works --- Rom 11:6
\item an inference step --- \cite[Trent, Session VI (1547), ch. 8]{tanner-decrees-1990}; \cite[36]{waterworth-trent-1848}; \cite[Five Arminian Articles (1610), art. 4]{schaff-creeds-christendom-3-1977}; \cite[Canons of Dort, III/IV, art. 10]{schaff-creeds-christendom-3-1977}
\end{itemize}
\item \emph{The third way (an ability given, received by faith)}, weakest at \emph{disputed}:
\begin{itemize}
\item Grace is received by faith alone --- \cite[Augsburg Confession IV.1–3]{concordia-triglotta-1921}; \cite[Sermon 1, III.1]{wesley-sermons-1872}; \cite[Trent, Session VI (1547), canon 9]{tanner-decrees-1990}
\end{itemize}
\item \emph{Remonstrant (the Five Articles, 1610)}, weakest at \emph{disputed}:
\begin{itemize}
\item What Stephen’s hearers resisted in Acts 7:51 was the grace that saves --- \cite[Five Arminian Articles (1610), art. 4]{schaff-creeds-christendom-3-1977}
\end{itemize}
\item \emph{Lutheran (the Formula of Concord)}, weakest at \emph{disputed}:
\begin{itemize}
\item Those who obstinately resist the Spirit can grieve and lose him --- \cite[Solid Declaration II.83]{concordia-triglotta-1921}; \cite[Solid Declaration XI.42]{concordia-triglotta-1921}
\item Conversion is God’s work alone: God regenerates, and the will does nothing to it --- \cite[Solid Declaration II.89]{concordia-triglotta-1921}; \cite[Solid Declaration II.90]{concordia-triglotta-1921}
\item Grace is received by faith alone --- \cite[Augsburg Confession IV.1–3]{concordia-triglotta-1921}; \cite[Sermon 1, III.1]{wesley-sermons-1872}; \cite[Trent, Session VI (1547), canon 9]{tanner-decrees-1990}
\end{itemize}
\item \emph{Wesleyan (preventing grace)}, weakest at \emph{disputed}:
\begin{itemize}
\item The grace that saves can be finally refused --- \cite[397]{wesley-works-9-1827}
\item Grace is received by faith alone --- \cite[Augsburg Confession IV.1–3]{concordia-triglotta-1921}; \cite[Sermon 1, III.1]{wesley-sermons-1872}; \cite[Trent, Session VI (1547), canon 9]{tanner-decrees-1990}
\end{itemize}
\item \emph{The objection from the Spirit’s resistance}, weakest at \emph{disputed}:
\begin{itemize}
\item What Stephen’s hearers resisted in Acts 7:51 was the grace that saves --- \cite[Five Arminian Articles (1610), art. 4]{schaff-creeds-christendom-3-1977}
\item Those who obstinately resist the Spirit can grieve and lose him --- \cite[Solid Declaration II.83]{concordia-triglotta-1921}; \cite[Solid Declaration XI.42]{concordia-triglotta-1921}
\item The blasphemy against the Spirit is impenitence to the end: grace refused till death --- \cite[325]{npnf1-06-1888}; \cite[II-II q.14 a.1]{aquinas-summa-1920}; \cite[74]{calvin-harmony-evangelists-2-1845}
\item an inference step --- \cite[325]{npnf1-06-1888}; \cite[On Rebuke and Grace, ch. 43]{npnf1-05-1887}
\end{itemize}
\item \emph{Thomist (Aquinas; Garrigou-Lagrange)}, weakest at \emph{disputed}:
\begin{itemize}
\item In converting grace God is the sole mover: the mind is moved and does not move --- \cite[I-II q.111 a.2]{aquinas-summa-1920}; \cite[ch. IV]{garrigou-lagrange-grace-1952}
\item The free will, moved by God, cooperates by assenting to God’s call --- \cite[Trent, Session VI (1547), canon 4]{tanner-decrees-1990}; \cite[45]{waterworth-trent-1848}
\item The free will, moved by God, can refuse its assent if it wishes --- \cite[Trent, Session VI (1547), canon 4]{tanner-decrees-1990}; \cite[45]{waterworth-trent-1848}
\item Grace truly sufficient, giving the power to do good, is resisted --- \cite[ch. VI]{garrigou-lagrange-grace-1952}
\item Justifying grace, once given, can be lost --- \cite[II-II q.24 a.11]{aquinas-summa-1920}; \cite[Trent, Session VI (1547), canon 23]{tanner-decrees-1990}
\item Grace is received by faith alone (encoded as a denial, so it ranks disputed) --- \cite[II-II q.4 a.3]{aquinas-summa-1920}; \cite[I-II q.113 a.4 ad 1]{aquinas-summa-1920}; \cite[Trent, Session VI (1547), canon 9]{tanner-decrees-1990}
\item an inference step --- \cite{pohle-molinism-1911}; \cite[I-II q.112 a.3]{aquinas-summa-1920}; \cite[ch. VII]{garrigou-lagrange-grace-1952}
\end{itemize}
\end{itemize}
\par\medskip\noindent\textbf{\texttt{tridentine\_\allowbreak{}not\_\allowbreak{}forced}}\par\nobreak\noindent
Nor is Trent forced. The maximal defensible position that holds Dort and the third way cannot hold Trent with them.

{\small\begin{verbatim}
theorem tridentine_not_forced : ¬Framework.SkepticallyAccepted
    solaGratiaDispute.defeats Party.trent
-- axioms: propext, Classical.choice, Quot.sound
\end{verbatim}}
\par\medskip\noindent\textbf{\texttt{remonstrants\allowbreak{}Answered\allowbreak{}By\allowbreak{}Dort}}\par\nobreak\noindent
Why the Remonstrants are not forced: Dort’s position defeats them.

\paragraph{\emph{Remonstrant (the Five Articles, 1610)} is not accepted on every resolution.}
\begin{itemize}
\item A position holding \emph{Romans 11:6 (by grace, so not by works)}, \emph{Reformed (the Canons of Dort)}, \emph{The third way (an ability given, received by faith)} can be held, and it holds \emph{Reformed (the Canons of Dort)}; \emph{Reformed (the Canons of Dort)} conflicts with \emph{Remonstrant (the Five Articles, 1610)}: \emph{Reformed (the Canons of Dort)} defeats \emph{Remonstrant (the Five Articles, 1610)}, and \emph{Remonstrant (the Five Articles, 1610)} defeats \emph{Reformed (the Canons of Dort)}.
\item None of \emph{Romans 11:6 (by grace, so not by works)}, \emph{Reformed (the Canons of Dort)}, \emph{The third way (an ability given, received by faith)} defeats another, and each attack on them is answered from among them:
\begin{itemize}
\item \emph{Trent (Session VI, canon 4, on the Molinist reading)} defeats \emph{Reformed (the Canons of Dort)}, and \emph{Reformed (the Canons of Dort)} defeats \emph{Trent (Session VI, canon 4, on the Molinist reading)}.
\item \emph{Remonstrant (the Five Articles, 1610)} defeats \emph{Reformed (the Canons of Dort)}, and \emph{Reformed (the Canons of Dort)} defeats \emph{Remonstrant (the Five Articles, 1610)}.
\item \emph{Lutheran (the Formula of Concord)} defeats \emph{Reformed (the Canons of Dort)}, and \emph{Reformed (the Canons of Dort)} defeats \emph{Lutheran (the Formula of Concord)}.
\item \emph{Wesleyan (preventing grace)} defeats \emph{Reformed (the Canons of Dort)}, and \emph{Reformed (the Canons of Dort)} defeats \emph{Wesleyan (preventing grace)}.
\item \emph{The objection from the Spirit’s resistance} defeats \emph{Reformed (the Canons of Dort)}, and \emph{Reformed (the Canons of Dort)} defeats \emph{The objection from the Spirit’s resistance}.
\item \emph{Trent (Session VI, canon 4, on the Molinist reading)} defeats \emph{The third way (an ability given, received by faith)}, and \emph{Reformed (the Canons of Dort)} defeats \emph{Trent (Session VI, canon 4, on the Molinist reading)}.
\item \emph{Thomist (Aquinas; Garrigou-Lagrange)} defeats \emph{The third way (an ability given, received by faith)}, and \emph{The third way (an ability given, received by faith)} defeats \emph{Thomist (Aquinas; Garrigou-Lagrange)}.
\end{itemize}
\end{itemize}
\par\noindent\textbf{What this rests on.} The reasons state 13 defeats and 20 absences of defeat, each a cell of the defeat table, \texttt{sola\allowbreak{}Gratia\allowbreak{}Dispute\_\allowbreak{}defeats}, computed from the two parties' premises and checked by the kernel. Whether an attack survives turns on the attacker's weakest link and on the rating of what it attacks, so a changed rating can change the verdict. The weakest links:
\begin{itemize}
\item \emph{Reformed (the Canons of Dort)}, weakest at \emph{disputed}:
\begin{itemize}
\item What Stephen’s hearers resisted in Acts 7:51 was the grace that saves (encoded as a denial, so it ranks disputed) --- \cite[Canons of Dort, III/IV, art. 9]{schaff-creeds-christendom-3-1977}
\item an inference step --- \cite[ad loc. Eph 2:8]{robertson-word-pictures-4-1931}; \cite[228--229]{calvin-galatians-ephesians-1854}
\item an inference step --- \cite[600]{arminius-works-1-1825}; \cite[681]{arminius-works-1-1825}; \cite[Five Arminian Articles (1610), art. 4]{schaff-creeds-christendom-3-1977}
\item an inference step --- \cite[527--565]{arminius-works-3-1853}
\item an inference step --- \cite[Solid Declaration II.89]{concordia-triglotta-1921}; \cite[Solid Declaration II.60]{concordia-triglotta-1921}; \cite[Solid Declaration XI.41]{concordia-triglotta-1921}
\end{itemize}
\item \emph{Remonstrant (the Five Articles, 1610)}, weakest at \emph{disputed}:
\begin{itemize}
\item What Stephen’s hearers resisted in Acts 7:51 was the grace that saves --- \cite[Five Arminian Articles (1610), art. 4]{schaff-creeds-christendom-3-1977}
\end{itemize}
\item \emph{Romans 11:6 (by grace, so not by works)}, weakest at \emph{consensus}:
\begin{itemize}
\item Romans 11:6 — if it is by grace, it is no longer on the basis of works --- Rom 11:6
\item an inference step --- \cite[Trent, Session VI (1547), ch. 8]{tanner-decrees-1990}; \cite[36]{waterworth-trent-1848}; \cite[Five Arminian Articles (1610), art. 4]{schaff-creeds-christendom-3-1977}; \cite[Canons of Dort, III/IV, art. 10]{schaff-creeds-christendom-3-1977}
\end{itemize}
\item \emph{The third way (an ability given, received by faith)}, weakest at \emph{disputed}:
\begin{itemize}
\item Grace is received by faith alone --- \cite[Augsburg Confession IV.1–3]{concordia-triglotta-1921}; \cite[Sermon 1, III.1]{wesley-sermons-1872}; \cite[Trent, Session VI (1547), canon 9]{tanner-decrees-1990}
\end{itemize}
\item \emph{Trent (Session VI, canon 4, on the Molinist reading)}, weakest at \emph{disputed}:
\begin{itemize}
\item The free will, moved by God, cooperates by assenting to God’s call --- \cite[Trent, Session VI (1547), canon 4]{tanner-decrees-1990}; \cite[45]{waterworth-trent-1848}
\item The free will, moved by God, can refuse its assent if it wishes --- \cite[Trent, Session VI (1547), canon 4]{tanner-decrees-1990}; \cite[45]{waterworth-trent-1848}
\item Grace is received by faith alone (encoded as a denial, so it ranks disputed) --- \cite[Augsburg Confession IV.1–3]{concordia-triglotta-1921}; \cite[Sermon 1, III.1]{wesley-sermons-1872}; \cite[Trent, Session VI (1547), canon 9]{tanner-decrees-1990}
\item an inference step --- \cite[I-II q.111 a.2]{aquinas-summa-1920}; \cite[Canons of Dort, III/IV, art. 12]{schaff-creeds-christendom-3-1977}; \cite[Trent, Session VI (1547), canon 4]{tanner-decrees-1990}; \cite[45]{waterworth-trent-1848}
\item an inference step --- \cite[I-II q.10 a.4 ad 3]{aquinas-summa-1920}; \cite[I-II q.112 a.3]{aquinas-summa-1920}; \cite[Trent, Session VI (1547), canon 4]{tanner-decrees-1990}; \cite[45]{waterworth-trent-1848}
\end{itemize}
\item \emph{Lutheran (the Formula of Concord)}, weakest at \emph{disputed}:
\begin{itemize}
\item Those who obstinately resist the Spirit can grieve and lose him --- \cite[Solid Declaration II.83]{concordia-triglotta-1921}; \cite[Solid Declaration XI.42]{concordia-triglotta-1921}
\item Conversion is God’s work alone: God regenerates, and the will does nothing to it --- \cite[Solid Declaration II.89]{concordia-triglotta-1921}; \cite[Solid Declaration II.90]{concordia-triglotta-1921}
\item Grace is received by faith alone --- \cite[Augsburg Confession IV.1–3]{concordia-triglotta-1921}; \cite[Sermon 1, III.1]{wesley-sermons-1872}; \cite[Trent, Session VI (1547), canon 9]{tanner-decrees-1990}
\end{itemize}
\item \emph{Wesleyan (preventing grace)}, weakest at \emph{disputed}:
\begin{itemize}
\item The grace that saves can be finally refused --- \cite[397]{wesley-works-9-1827}
\item Grace is received by faith alone --- \cite[Augsburg Confession IV.1–3]{concordia-triglotta-1921}; \cite[Sermon 1, III.1]{wesley-sermons-1872}; \cite[Trent, Session VI (1547), canon 9]{tanner-decrees-1990}
\end{itemize}
\item \emph{The objection from the Spirit’s resistance}, weakest at \emph{disputed}:
\begin{itemize}
\item What Stephen’s hearers resisted in Acts 7:51 was the grace that saves --- \cite[Five Arminian Articles (1610), art. 4]{schaff-creeds-christendom-3-1977}
\item Those who obstinately resist the Spirit can grieve and lose him --- \cite[Solid Declaration II.83]{concordia-triglotta-1921}; \cite[Solid Declaration XI.42]{concordia-triglotta-1921}
\item The blasphemy against the Spirit is impenitence to the end: grace refused till death --- \cite[325]{npnf1-06-1888}; \cite[II-II q.14 a.1]{aquinas-summa-1920}; \cite[74]{calvin-harmony-evangelists-2-1845}
\item an inference step --- \cite[325]{npnf1-06-1888}; \cite[On Rebuke and Grace, ch. 43]{npnf1-05-1887}
\end{itemize}
\item \emph{Thomist (Aquinas; Garrigou-Lagrange)}, weakest at \emph{disputed}:
\begin{itemize}
\item In converting grace God is the sole mover: the mind is moved and does not move --- \cite[I-II q.111 a.2]{aquinas-summa-1920}; \cite[ch. IV]{garrigou-lagrange-grace-1952}
\item The free will, moved by God, cooperates by assenting to God’s call --- \cite[Trent, Session VI (1547), canon 4]{tanner-decrees-1990}; \cite[45]{waterworth-trent-1848}
\item The free will, moved by God, can refuse its assent if it wishes --- \cite[Trent, Session VI (1547), canon 4]{tanner-decrees-1990}; \cite[45]{waterworth-trent-1848}
\item Grace truly sufficient, giving the power to do good, is resisted --- \cite[ch. VI]{garrigou-lagrange-grace-1952}
\item Justifying grace, once given, can be lost --- \cite[II-II q.24 a.11]{aquinas-summa-1920}; \cite[Trent, Session VI (1547), canon 23]{tanner-decrees-1990}
\item Grace is received by faith alone (encoded as a denial, so it ranks disputed) --- \cite[II-II q.4 a.3]{aquinas-summa-1920}; \cite[I-II q.113 a.4 ad 1]{aquinas-summa-1920}; \cite[Trent, Session VI (1547), canon 9]{tanner-decrees-1990}
\item an inference step --- \cite{pohle-molinism-1911}; \cite[I-II q.112 a.3]{aquinas-summa-1920}; \cite[ch. VII]{garrigou-lagrange-grace-1952}
\end{itemize}
\end{itemize}
\par\medskip\noindent\textbf{\texttt{remonstrant\_\allowbreak{}not\_\allowbreak{}forced}}\par\nobreak\noindent
The Remonstrants can be defended (with the Lutherans and Wesley, as \texttt{lutheran\_\allowbreak{}defensible} shows), and are not forced: the maximal defensible position that holds Dort cannot hold them.

{\small\begin{verbatim}
theorem remonstrant_not_forced : ¬Framework.SkepticallyAccepted
    solaGratiaDispute.defeats Party.remonstrant
-- axioms: propext, Classical.choice, Quot.sound
\end{verbatim}}
\par\medskip\noindent\textbf{\texttt{lutherans\allowbreak{}Answered\allowbreak{}By\allowbreak{}Dort}}\par\nobreak\noindent
Why the Lutherans are not forced: Dort’s position defeats them.

\paragraph{\emph{Lutheran (the Formula of Concord)} is not accepted on every resolution.}
\begin{itemize}
\item A position holding \emph{Romans 11:6 (by grace, so not by works)}, \emph{Reformed (the Canons of Dort)}, \emph{The third way (an ability given, received by faith)} can be held, and it holds \emph{Reformed (the Canons of Dort)}; \emph{Reformed (the Canons of Dort)} conflicts with \emph{Lutheran (the Formula of Concord)}: \emph{Reformed (the Canons of Dort)} defeats \emph{Lutheran (the Formula of Concord)}, and \emph{Lutheran (the Formula of Concord)} defeats \emph{Reformed (the Canons of Dort)}.
\item None of \emph{Romans 11:6 (by grace, so not by works)}, \emph{Reformed (the Canons of Dort)}, \emph{The third way (an ability given, received by faith)} defeats another, and each attack on them is answered from among them:
\begin{itemize}
\item \emph{Trent (Session VI, canon 4, on the Molinist reading)} defeats \emph{Reformed (the Canons of Dort)}, and \emph{Reformed (the Canons of Dort)} defeats \emph{Trent (Session VI, canon 4, on the Molinist reading)}.
\item \emph{Remonstrant (the Five Articles, 1610)} defeats \emph{Reformed (the Canons of Dort)}, and \emph{Reformed (the Canons of Dort)} defeats \emph{Remonstrant (the Five Articles, 1610)}.
\item \emph{Lutheran (the Formula of Concord)} defeats \emph{Reformed (the Canons of Dort)}, and \emph{Reformed (the Canons of Dort)} defeats \emph{Lutheran (the Formula of Concord)}.
\item \emph{Wesleyan (preventing grace)} defeats \emph{Reformed (the Canons of Dort)}, and \emph{Reformed (the Canons of Dort)} defeats \emph{Wesleyan (preventing grace)}.
\item \emph{The objection from the Spirit’s resistance} defeats \emph{Reformed (the Canons of Dort)}, and \emph{Reformed (the Canons of Dort)} defeats \emph{The objection from the Spirit’s resistance}.
\item \emph{Trent (Session VI, canon 4, on the Molinist reading)} defeats \emph{The third way (an ability given, received by faith)}, and \emph{Reformed (the Canons of Dort)} defeats \emph{Trent (Session VI, canon 4, on the Molinist reading)}.
\item \emph{Thomist (Aquinas; Garrigou-Lagrange)} defeats \emph{The third way (an ability given, received by faith)}, and \emph{The third way (an ability given, received by faith)} defeats \emph{Thomist (Aquinas; Garrigou-Lagrange)}.
\end{itemize}
\end{itemize}
\par\noindent\textbf{What this rests on.} The reasons state 13 defeats and 20 absences of defeat, each a cell of the defeat table, \texttt{sola\allowbreak{}Gratia\allowbreak{}Dispute\_\allowbreak{}defeats}, computed from the two parties' premises and checked by the kernel. Whether an attack survives turns on the attacker's weakest link and on the rating of what it attacks, so a changed rating can change the verdict. The weakest links:
\begin{itemize}
\item \emph{Reformed (the Canons of Dort)}, weakest at \emph{disputed}:
\begin{itemize}
\item What Stephen’s hearers resisted in Acts 7:51 was the grace that saves (encoded as a denial, so it ranks disputed) --- \cite[Canons of Dort, III/IV, art. 9]{schaff-creeds-christendom-3-1977}
\item an inference step --- \cite[ad loc. Eph 2:8]{robertson-word-pictures-4-1931}; \cite[228--229]{calvin-galatians-ephesians-1854}
\item an inference step --- \cite[600]{arminius-works-1-1825}; \cite[681]{arminius-works-1-1825}; \cite[Five Arminian Articles (1610), art. 4]{schaff-creeds-christendom-3-1977}
\item an inference step --- \cite[527--565]{arminius-works-3-1853}
\item an inference step --- \cite[Solid Declaration II.89]{concordia-triglotta-1921}; \cite[Solid Declaration II.60]{concordia-triglotta-1921}; \cite[Solid Declaration XI.41]{concordia-triglotta-1921}
\end{itemize}
\item \emph{Lutheran (the Formula of Concord)}, weakest at \emph{disputed}:
\begin{itemize}
\item Those who obstinately resist the Spirit can grieve and lose him --- \cite[Solid Declaration II.83]{concordia-triglotta-1921}; \cite[Solid Declaration XI.42]{concordia-triglotta-1921}
\item Conversion is God’s work alone: God regenerates, and the will does nothing to it --- \cite[Solid Declaration II.89]{concordia-triglotta-1921}; \cite[Solid Declaration II.90]{concordia-triglotta-1921}
\item Grace is received by faith alone --- \cite[Augsburg Confession IV.1–3]{concordia-triglotta-1921}; \cite[Sermon 1, III.1]{wesley-sermons-1872}; \cite[Trent, Session VI (1547), canon 9]{tanner-decrees-1990}
\end{itemize}
\item \emph{Romans 11:6 (by grace, so not by works)}, weakest at \emph{consensus}:
\begin{itemize}
\item Romans 11:6 — if it is by grace, it is no longer on the basis of works --- Rom 11:6
\item an inference step --- \cite[Trent, Session VI (1547), ch. 8]{tanner-decrees-1990}; \cite[36]{waterworth-trent-1848}; \cite[Five Arminian Articles (1610), art. 4]{schaff-creeds-christendom-3-1977}; \cite[Canons of Dort, III/IV, art. 10]{schaff-creeds-christendom-3-1977}
\end{itemize}
\item \emph{The third way (an ability given, received by faith)}, weakest at \emph{disputed}:
\begin{itemize}
\item Grace is received by faith alone --- \cite[Augsburg Confession IV.1–3]{concordia-triglotta-1921}; \cite[Sermon 1, III.1]{wesley-sermons-1872}; \cite[Trent, Session VI (1547), canon 9]{tanner-decrees-1990}
\end{itemize}
\item \emph{Trent (Session VI, canon 4, on the Molinist reading)}, weakest at \emph{disputed}:
\begin{itemize}
\item The free will, moved by God, cooperates by assenting to God’s call --- \cite[Trent, Session VI (1547), canon 4]{tanner-decrees-1990}; \cite[45]{waterworth-trent-1848}
\item The free will, moved by God, can refuse its assent if it wishes --- \cite[Trent, Session VI (1547), canon 4]{tanner-decrees-1990}; \cite[45]{waterworth-trent-1848}
\item Grace is received by faith alone (encoded as a denial, so it ranks disputed) --- \cite[Augsburg Confession IV.1–3]{concordia-triglotta-1921}; \cite[Sermon 1, III.1]{wesley-sermons-1872}; \cite[Trent, Session VI (1547), canon 9]{tanner-decrees-1990}
\item an inference step --- \cite[I-II q.111 a.2]{aquinas-summa-1920}; \cite[Canons of Dort, III/IV, art. 12]{schaff-creeds-christendom-3-1977}; \cite[Trent, Session VI (1547), canon 4]{tanner-decrees-1990}; \cite[45]{waterworth-trent-1848}
\item an inference step --- \cite[I-II q.10 a.4 ad 3]{aquinas-summa-1920}; \cite[I-II q.112 a.3]{aquinas-summa-1920}; \cite[Trent, Session VI (1547), canon 4]{tanner-decrees-1990}; \cite[45]{waterworth-trent-1848}
\end{itemize}
\item \emph{Remonstrant (the Five Articles, 1610)}, weakest at \emph{disputed}:
\begin{itemize}
\item What Stephen’s hearers resisted in Acts 7:51 was the grace that saves --- \cite[Five Arminian Articles (1610), art. 4]{schaff-creeds-christendom-3-1977}
\end{itemize}
\item \emph{Wesleyan (preventing grace)}, weakest at \emph{disputed}:
\begin{itemize}
\item The grace that saves can be finally refused --- \cite[397]{wesley-works-9-1827}
\item Grace is received by faith alone --- \cite[Augsburg Confession IV.1–3]{concordia-triglotta-1921}; \cite[Sermon 1, III.1]{wesley-sermons-1872}; \cite[Trent, Session VI (1547), canon 9]{tanner-decrees-1990}
\end{itemize}
\item \emph{The objection from the Spirit’s resistance}, weakest at \emph{disputed}:
\begin{itemize}
\item What Stephen’s hearers resisted in Acts 7:51 was the grace that saves --- \cite[Five Arminian Articles (1610), art. 4]{schaff-creeds-christendom-3-1977}
\item Those who obstinately resist the Spirit can grieve and lose him --- \cite[Solid Declaration II.83]{concordia-triglotta-1921}; \cite[Solid Declaration XI.42]{concordia-triglotta-1921}
\item The blasphemy against the Spirit is impenitence to the end: grace refused till death --- \cite[325]{npnf1-06-1888}; \cite[II-II q.14 a.1]{aquinas-summa-1920}; \cite[74]{calvin-harmony-evangelists-2-1845}
\item an inference step --- \cite[325]{npnf1-06-1888}; \cite[On Rebuke and Grace, ch. 43]{npnf1-05-1887}
\end{itemize}
\item \emph{Thomist (Aquinas; Garrigou-Lagrange)}, weakest at \emph{disputed}:
\begin{itemize}
\item In converting grace God is the sole mover: the mind is moved and does not move --- \cite[I-II q.111 a.2]{aquinas-summa-1920}; \cite[ch. IV]{garrigou-lagrange-grace-1952}
\item The free will, moved by God, cooperates by assenting to God’s call --- \cite[Trent, Session VI (1547), canon 4]{tanner-decrees-1990}; \cite[45]{waterworth-trent-1848}
\item The free will, moved by God, can refuse its assent if it wishes --- \cite[Trent, Session VI (1547), canon 4]{tanner-decrees-1990}; \cite[45]{waterworth-trent-1848}
\item Grace truly sufficient, giving the power to do good, is resisted --- \cite[ch. VI]{garrigou-lagrange-grace-1952}
\item Justifying grace, once given, can be lost --- \cite[II-II q.24 a.11]{aquinas-summa-1920}; \cite[Trent, Session VI (1547), canon 23]{tanner-decrees-1990}
\item Grace is received by faith alone (encoded as a denial, so it ranks disputed) --- \cite[II-II q.4 a.3]{aquinas-summa-1920}; \cite[I-II q.113 a.4 ad 1]{aquinas-summa-1920}; \cite[Trent, Session VI (1547), canon 9]{tanner-decrees-1990}
\item an inference step --- \cite{pohle-molinism-1911}; \cite[I-II q.112 a.3]{aquinas-summa-1920}; \cite[ch. VII]{garrigou-lagrange-grace-1952}
\end{itemize}
\end{itemize}
\par\medskip\noindent\textbf{\texttt{lutheran\_\allowbreak{}not\_\allowbreak{}forced}}\par\nobreak\noindent
Nor are the Lutherans forced: the maximal defensible position that holds Dort cannot hold them.

{\small\begin{verbatim}
theorem lutheran_not_forced : ¬Framework.SkepticallyAccepted
    solaGratiaDispute.defeats Party.lutheran
-- axioms: propext, Classical.choice, Quot.sound
\end{verbatim}}
\par\medskip\noindent\textbf{\texttt{wesley\allowbreak{}Answered\allowbreak{}By\allowbreak{}Dort}}\par\nobreak\noindent
Why Wesley is not forced: Dort’s position defeats him.

\paragraph{\emph{Wesleyan (preventing grace)} is not accepted on every resolution.}
\begin{itemize}
\item A position holding \emph{Romans 11:6 (by grace, so not by works)}, \emph{Reformed (the Canons of Dort)}, \emph{The third way (an ability given, received by faith)} can be held, and it holds \emph{Reformed (the Canons of Dort)}; \emph{Reformed (the Canons of Dort)} conflicts with \emph{Wesleyan (preventing grace)}: \emph{Reformed (the Canons of Dort)} defeats \emph{Wesleyan (preventing grace)}, and \emph{Wesleyan (preventing grace)} defeats \emph{Reformed (the Canons of Dort)}.
\item None of \emph{Romans 11:6 (by grace, so not by works)}, \emph{Reformed (the Canons of Dort)}, \emph{The third way (an ability given, received by faith)} defeats another, and each attack on them is answered from among them:
\begin{itemize}
\item \emph{Trent (Session VI, canon 4, on the Molinist reading)} defeats \emph{Reformed (the Canons of Dort)}, and \emph{Reformed (the Canons of Dort)} defeats \emph{Trent (Session VI, canon 4, on the Molinist reading)}.
\item \emph{Remonstrant (the Five Articles, 1610)} defeats \emph{Reformed (the Canons of Dort)}, and \emph{Reformed (the Canons of Dort)} defeats \emph{Remonstrant (the Five Articles, 1610)}.
\item \emph{Lutheran (the Formula of Concord)} defeats \emph{Reformed (the Canons of Dort)}, and \emph{Reformed (the Canons of Dort)} defeats \emph{Lutheran (the Formula of Concord)}.
\item \emph{Wesleyan (preventing grace)} defeats \emph{Reformed (the Canons of Dort)}, and \emph{Reformed (the Canons of Dort)} defeats \emph{Wesleyan (preventing grace)}.
\item \emph{The objection from the Spirit’s resistance} defeats \emph{Reformed (the Canons of Dort)}, and \emph{Reformed (the Canons of Dort)} defeats \emph{The objection from the Spirit’s resistance}.
\item \emph{Trent (Session VI, canon 4, on the Molinist reading)} defeats \emph{The third way (an ability given, received by faith)}, and \emph{Reformed (the Canons of Dort)} defeats \emph{Trent (Session VI, canon 4, on the Molinist reading)}.
\item \emph{Thomist (Aquinas; Garrigou-Lagrange)} defeats \emph{The third way (an ability given, received by faith)}, and \emph{The third way (an ability given, received by faith)} defeats \emph{Thomist (Aquinas; Garrigou-Lagrange)}.
\end{itemize}
\end{itemize}
\par\noindent\textbf{What this rests on.} The reasons state 13 defeats and 20 absences of defeat, each a cell of the defeat table, \texttt{sola\allowbreak{}Gratia\allowbreak{}Dispute\_\allowbreak{}defeats}, computed from the two parties' premises and checked by the kernel. Whether an attack survives turns on the attacker's weakest link and on the rating of what it attacks, so a changed rating can change the verdict. The weakest links:
\begin{itemize}
\item \emph{Reformed (the Canons of Dort)}, weakest at \emph{disputed}:
\begin{itemize}
\item What Stephen’s hearers resisted in Acts 7:51 was the grace that saves (encoded as a denial, so it ranks disputed) --- \cite[Canons of Dort, III/IV, art. 9]{schaff-creeds-christendom-3-1977}
\item an inference step --- \cite[ad loc. Eph 2:8]{robertson-word-pictures-4-1931}; \cite[228--229]{calvin-galatians-ephesians-1854}
\item an inference step --- \cite[600]{arminius-works-1-1825}; \cite[681]{arminius-works-1-1825}; \cite[Five Arminian Articles (1610), art. 4]{schaff-creeds-christendom-3-1977}
\item an inference step --- \cite[527--565]{arminius-works-3-1853}
\item an inference step --- \cite[Solid Declaration II.89]{concordia-triglotta-1921}; \cite[Solid Declaration II.60]{concordia-triglotta-1921}; \cite[Solid Declaration XI.41]{concordia-triglotta-1921}
\end{itemize}
\item \emph{Wesleyan (preventing grace)}, weakest at \emph{disputed}:
\begin{itemize}
\item The grace that saves can be finally refused --- \cite[397]{wesley-works-9-1827}
\item Grace is received by faith alone --- \cite[Augsburg Confession IV.1–3]{concordia-triglotta-1921}; \cite[Sermon 1, III.1]{wesley-sermons-1872}; \cite[Trent, Session VI (1547), canon 9]{tanner-decrees-1990}
\end{itemize}
\item \emph{Romans 11:6 (by grace, so not by works)}, weakest at \emph{consensus}:
\begin{itemize}
\item Romans 11:6 — if it is by grace, it is no longer on the basis of works --- Rom 11:6
\item an inference step --- \cite[Trent, Session VI (1547), ch. 8]{tanner-decrees-1990}; \cite[36]{waterworth-trent-1848}; \cite[Five Arminian Articles (1610), art. 4]{schaff-creeds-christendom-3-1977}; \cite[Canons of Dort, III/IV, art. 10]{schaff-creeds-christendom-3-1977}
\end{itemize}
\item \emph{The third way (an ability given, received by faith)}, weakest at \emph{disputed}:
\begin{itemize}
\item Grace is received by faith alone --- \cite[Augsburg Confession IV.1–3]{concordia-triglotta-1921}; \cite[Sermon 1, III.1]{wesley-sermons-1872}; \cite[Trent, Session VI (1547), canon 9]{tanner-decrees-1990}
\end{itemize}
\item \emph{Trent (Session VI, canon 4, on the Molinist reading)}, weakest at \emph{disputed}:
\begin{itemize}
\item The free will, moved by God, cooperates by assenting to God’s call --- \cite[Trent, Session VI (1547), canon 4]{tanner-decrees-1990}; \cite[45]{waterworth-trent-1848}
\item The free will, moved by God, can refuse its assent if it wishes --- \cite[Trent, Session VI (1547), canon 4]{tanner-decrees-1990}; \cite[45]{waterworth-trent-1848}
\item Grace is received by faith alone (encoded as a denial, so it ranks disputed) --- \cite[Augsburg Confession IV.1–3]{concordia-triglotta-1921}; \cite[Sermon 1, III.1]{wesley-sermons-1872}; \cite[Trent, Session VI (1547), canon 9]{tanner-decrees-1990}
\item an inference step --- \cite[I-II q.111 a.2]{aquinas-summa-1920}; \cite[Canons of Dort, III/IV, art. 12]{schaff-creeds-christendom-3-1977}; \cite[Trent, Session VI (1547), canon 4]{tanner-decrees-1990}; \cite[45]{waterworth-trent-1848}
\item an inference step --- \cite[I-II q.10 a.4 ad 3]{aquinas-summa-1920}; \cite[I-II q.112 a.3]{aquinas-summa-1920}; \cite[Trent, Session VI (1547), canon 4]{tanner-decrees-1990}; \cite[45]{waterworth-trent-1848}
\end{itemize}
\item \emph{Remonstrant (the Five Articles, 1610)}, weakest at \emph{disputed}:
\begin{itemize}
\item What Stephen’s hearers resisted in Acts 7:51 was the grace that saves --- \cite[Five Arminian Articles (1610), art. 4]{schaff-creeds-christendom-3-1977}
\end{itemize}
\item \emph{Lutheran (the Formula of Concord)}, weakest at \emph{disputed}:
\begin{itemize}
\item Those who obstinately resist the Spirit can grieve and lose him --- \cite[Solid Declaration II.83]{concordia-triglotta-1921}; \cite[Solid Declaration XI.42]{concordia-triglotta-1921}
\item Conversion is God’s work alone: God regenerates, and the will does nothing to it --- \cite[Solid Declaration II.89]{concordia-triglotta-1921}; \cite[Solid Declaration II.90]{concordia-triglotta-1921}
\item Grace is received by faith alone --- \cite[Augsburg Confession IV.1–3]{concordia-triglotta-1921}; \cite[Sermon 1, III.1]{wesley-sermons-1872}; \cite[Trent, Session VI (1547), canon 9]{tanner-decrees-1990}
\end{itemize}
\item \emph{The objection from the Spirit’s resistance}, weakest at \emph{disputed}:
\begin{itemize}
\item What Stephen’s hearers resisted in Acts 7:51 was the grace that saves --- \cite[Five Arminian Articles (1610), art. 4]{schaff-creeds-christendom-3-1977}
\item Those who obstinately resist the Spirit can grieve and lose him --- \cite[Solid Declaration II.83]{concordia-triglotta-1921}; \cite[Solid Declaration XI.42]{concordia-triglotta-1921}
\item The blasphemy against the Spirit is impenitence to the end: grace refused till death --- \cite[325]{npnf1-06-1888}; \cite[II-II q.14 a.1]{aquinas-summa-1920}; \cite[74]{calvin-harmony-evangelists-2-1845}
\item an inference step --- \cite[325]{npnf1-06-1888}; \cite[On Rebuke and Grace, ch. 43]{npnf1-05-1887}
\end{itemize}
\item \emph{Thomist (Aquinas; Garrigou-Lagrange)}, weakest at \emph{disputed}:
\begin{itemize}
\item In converting grace God is the sole mover: the mind is moved and does not move --- \cite[I-II q.111 a.2]{aquinas-summa-1920}; \cite[ch. IV]{garrigou-lagrange-grace-1952}
\item The free will, moved by God, cooperates by assenting to God’s call --- \cite[Trent, Session VI (1547), canon 4]{tanner-decrees-1990}; \cite[45]{waterworth-trent-1848}
\item The free will, moved by God, can refuse its assent if it wishes --- \cite[Trent, Session VI (1547), canon 4]{tanner-decrees-1990}; \cite[45]{waterworth-trent-1848}
\item Grace truly sufficient, giving the power to do good, is resisted --- \cite[ch. VI]{garrigou-lagrange-grace-1952}
\item Justifying grace, once given, can be lost --- \cite[II-II q.24 a.11]{aquinas-summa-1920}; \cite[Trent, Session VI (1547), canon 23]{tanner-decrees-1990}
\item Grace is received by faith alone (encoded as a denial, so it ranks disputed) --- \cite[II-II q.4 a.3]{aquinas-summa-1920}; \cite[I-II q.113 a.4 ad 1]{aquinas-summa-1920}; \cite[Trent, Session VI (1547), canon 9]{tanner-decrees-1990}
\item an inference step --- \cite{pohle-molinism-1911}; \cite[I-II q.112 a.3]{aquinas-summa-1920}; \cite[ch. VII]{garrigou-lagrange-grace-1952}
\end{itemize}
\end{itemize}
\par\medskip\noindent\textbf{\texttt{wesleyan\_\allowbreak{}not\_\allowbreak{}forced}}\par\nobreak\noindent
Nor is Wesley’s position forced: the maximal defensible position that holds Dort cannot hold it.

{\small\begin{verbatim}
theorem wesleyan_not_forced : ¬Framework.SkepticallyAccepted
    solaGratiaDispute.defeats Party.wesleyan
-- axioms: propext, Classical.choice, Quot.sound
\end{verbatim}}
\par\medskip\noindent\textbf{\texttt{objection\allowbreak{}Answered\allowbreak{}By\allowbreak{}Dort}}\par\nobreak\noindent
Why the objection is not forced: Dort’s position defeats it.

\paragraph{\emph{The objection from the Spirit’s resistance} is not accepted on every resolution.}
\begin{itemize}
\item A position holding \emph{Romans 11:6 (by grace, so not by works)}, \emph{Reformed (the Canons of Dort)}, \emph{The third way (an ability given, received by faith)} can be held, and it holds \emph{Reformed (the Canons of Dort)}; \emph{Reformed (the Canons of Dort)} conflicts with \emph{The objection from the Spirit’s resistance}: \emph{Reformed (the Canons of Dort)} defeats \emph{The objection from the Spirit’s resistance}, and \emph{The objection from the Spirit’s resistance} defeats \emph{Reformed (the Canons of Dort)}.
\item None of \emph{Romans 11:6 (by grace, so not by works)}, \emph{Reformed (the Canons of Dort)}, \emph{The third way (an ability given, received by faith)} defeats another, and each attack on them is answered from among them:
\begin{itemize}
\item \emph{Trent (Session VI, canon 4, on the Molinist reading)} defeats \emph{Reformed (the Canons of Dort)}, and \emph{Reformed (the Canons of Dort)} defeats \emph{Trent (Session VI, canon 4, on the Molinist reading)}.
\item \emph{Remonstrant (the Five Articles, 1610)} defeats \emph{Reformed (the Canons of Dort)}, and \emph{Reformed (the Canons of Dort)} defeats \emph{Remonstrant (the Five Articles, 1610)}.
\item \emph{Lutheran (the Formula of Concord)} defeats \emph{Reformed (the Canons of Dort)}, and \emph{Reformed (the Canons of Dort)} defeats \emph{Lutheran (the Formula of Concord)}.
\item \emph{Wesleyan (preventing grace)} defeats \emph{Reformed (the Canons of Dort)}, and \emph{Reformed (the Canons of Dort)} defeats \emph{Wesleyan (preventing grace)}.
\item \emph{The objection from the Spirit’s resistance} defeats \emph{Reformed (the Canons of Dort)}, and \emph{Reformed (the Canons of Dort)} defeats \emph{The objection from the Spirit’s resistance}.
\item \emph{Trent (Session VI, canon 4, on the Molinist reading)} defeats \emph{The third way (an ability given, received by faith)}, and \emph{Reformed (the Canons of Dort)} defeats \emph{Trent (Session VI, canon 4, on the Molinist reading)}.
\item \emph{Thomist (Aquinas; Garrigou-Lagrange)} defeats \emph{The third way (an ability given, received by faith)}, and \emph{The third way (an ability given, received by faith)} defeats \emph{Thomist (Aquinas; Garrigou-Lagrange)}.
\end{itemize}
\end{itemize}
\par\noindent\textbf{What this rests on.} The reasons state 13 defeats and 20 absences of defeat, each a cell of the defeat table, \texttt{sola\allowbreak{}Gratia\allowbreak{}Dispute\_\allowbreak{}defeats}, computed from the two parties' premises and checked by the kernel. Whether an attack survives turns on the attacker's weakest link and on the rating of what it attacks, so a changed rating can change the verdict. The weakest links:
\begin{itemize}
\item \emph{Reformed (the Canons of Dort)}, weakest at \emph{disputed}:
\begin{itemize}
\item What Stephen’s hearers resisted in Acts 7:51 was the grace that saves (encoded as a denial, so it ranks disputed) --- \cite[Canons of Dort, III/IV, art. 9]{schaff-creeds-christendom-3-1977}
\item an inference step --- \cite[ad loc. Eph 2:8]{robertson-word-pictures-4-1931}; \cite[228--229]{calvin-galatians-ephesians-1854}
\item an inference step --- \cite[600]{arminius-works-1-1825}; \cite[681]{arminius-works-1-1825}; \cite[Five Arminian Articles (1610), art. 4]{schaff-creeds-christendom-3-1977}
\item an inference step --- \cite[527--565]{arminius-works-3-1853}
\item an inference step --- \cite[Solid Declaration II.89]{concordia-triglotta-1921}; \cite[Solid Declaration II.60]{concordia-triglotta-1921}; \cite[Solid Declaration XI.41]{concordia-triglotta-1921}
\end{itemize}
\item \emph{The objection from the Spirit’s resistance}, weakest at \emph{disputed}:
\begin{itemize}
\item What Stephen’s hearers resisted in Acts 7:51 was the grace that saves --- \cite[Five Arminian Articles (1610), art. 4]{schaff-creeds-christendom-3-1977}
\item Those who obstinately resist the Spirit can grieve and lose him --- \cite[Solid Declaration II.83]{concordia-triglotta-1921}; \cite[Solid Declaration XI.42]{concordia-triglotta-1921}
\item The blasphemy against the Spirit is impenitence to the end: grace refused till death --- \cite[325]{npnf1-06-1888}; \cite[II-II q.14 a.1]{aquinas-summa-1920}; \cite[74]{calvin-harmony-evangelists-2-1845}
\item an inference step --- \cite[325]{npnf1-06-1888}; \cite[On Rebuke and Grace, ch. 43]{npnf1-05-1887}
\end{itemize}
\item \emph{Romans 11:6 (by grace, so not by works)}, weakest at \emph{consensus}:
\begin{itemize}
\item Romans 11:6 — if it is by grace, it is no longer on the basis of works --- Rom 11:6
\item an inference step --- \cite[Trent, Session VI (1547), ch. 8]{tanner-decrees-1990}; \cite[36]{waterworth-trent-1848}; \cite[Five Arminian Articles (1610), art. 4]{schaff-creeds-christendom-3-1977}; \cite[Canons of Dort, III/IV, art. 10]{schaff-creeds-christendom-3-1977}
\end{itemize}
\item \emph{The third way (an ability given, received by faith)}, weakest at \emph{disputed}:
\begin{itemize}
\item Grace is received by faith alone --- \cite[Augsburg Confession IV.1–3]{concordia-triglotta-1921}; \cite[Sermon 1, III.1]{wesley-sermons-1872}; \cite[Trent, Session VI (1547), canon 9]{tanner-decrees-1990}
\end{itemize}
\item \emph{Trent (Session VI, canon 4, on the Molinist reading)}, weakest at \emph{disputed}:
\begin{itemize}
\item The free will, moved by God, cooperates by assenting to God’s call --- \cite[Trent, Session VI (1547), canon 4]{tanner-decrees-1990}; \cite[45]{waterworth-trent-1848}
\item The free will, moved by God, can refuse its assent if it wishes --- \cite[Trent, Session VI (1547), canon 4]{tanner-decrees-1990}; \cite[45]{waterworth-trent-1848}
\item Grace is received by faith alone (encoded as a denial, so it ranks disputed) --- \cite[Augsburg Confession IV.1–3]{concordia-triglotta-1921}; \cite[Sermon 1, III.1]{wesley-sermons-1872}; \cite[Trent, Session VI (1547), canon 9]{tanner-decrees-1990}
\item an inference step --- \cite[I-II q.111 a.2]{aquinas-summa-1920}; \cite[Canons of Dort, III/IV, art. 12]{schaff-creeds-christendom-3-1977}; \cite[Trent, Session VI (1547), canon 4]{tanner-decrees-1990}; \cite[45]{waterworth-trent-1848}
\item an inference step --- \cite[I-II q.10 a.4 ad 3]{aquinas-summa-1920}; \cite[I-II q.112 a.3]{aquinas-summa-1920}; \cite[Trent, Session VI (1547), canon 4]{tanner-decrees-1990}; \cite[45]{waterworth-trent-1848}
\end{itemize}
\item \emph{Remonstrant (the Five Articles, 1610)}, weakest at \emph{disputed}:
\begin{itemize}
\item What Stephen’s hearers resisted in Acts 7:51 was the grace that saves --- \cite[Five Arminian Articles (1610), art. 4]{schaff-creeds-christendom-3-1977}
\end{itemize}
\item \emph{Lutheran (the Formula of Concord)}, weakest at \emph{disputed}:
\begin{itemize}
\item Those who obstinately resist the Spirit can grieve and lose him --- \cite[Solid Declaration II.83]{concordia-triglotta-1921}; \cite[Solid Declaration XI.42]{concordia-triglotta-1921}
\item Conversion is God’s work alone: God regenerates, and the will does nothing to it --- \cite[Solid Declaration II.89]{concordia-triglotta-1921}; \cite[Solid Declaration II.90]{concordia-triglotta-1921}
\item Grace is received by faith alone --- \cite[Augsburg Confession IV.1–3]{concordia-triglotta-1921}; \cite[Sermon 1, III.1]{wesley-sermons-1872}; \cite[Trent, Session VI (1547), canon 9]{tanner-decrees-1990}
\end{itemize}
\item \emph{Wesleyan (preventing grace)}, weakest at \emph{disputed}:
\begin{itemize}
\item The grace that saves can be finally refused --- \cite[397]{wesley-works-9-1827}
\item Grace is received by faith alone --- \cite[Augsburg Confession IV.1–3]{concordia-triglotta-1921}; \cite[Sermon 1, III.1]{wesley-sermons-1872}; \cite[Trent, Session VI (1547), canon 9]{tanner-decrees-1990}
\end{itemize}
\item \emph{Thomist (Aquinas; Garrigou-Lagrange)}, weakest at \emph{disputed}:
\begin{itemize}
\item In converting grace God is the sole mover: the mind is moved and does not move --- \cite[I-II q.111 a.2]{aquinas-summa-1920}; \cite[ch. IV]{garrigou-lagrange-grace-1952}
\item The free will, moved by God, cooperates by assenting to God’s call --- \cite[Trent, Session VI (1547), canon 4]{tanner-decrees-1990}; \cite[45]{waterworth-trent-1848}
\item The free will, moved by God, can refuse its assent if it wishes --- \cite[Trent, Session VI (1547), canon 4]{tanner-decrees-1990}; \cite[45]{waterworth-trent-1848}
\item Grace truly sufficient, giving the power to do good, is resisted --- \cite[ch. VI]{garrigou-lagrange-grace-1952}
\item Justifying grace, once given, can be lost --- \cite[II-II q.24 a.11]{aquinas-summa-1920}; \cite[Trent, Session VI (1547), canon 23]{tanner-decrees-1990}
\item Grace is received by faith alone (encoded as a denial, so it ranks disputed) --- \cite[II-II q.4 a.3]{aquinas-summa-1920}; \cite[I-II q.113 a.4 ad 1]{aquinas-summa-1920}; \cite[Trent, Session VI (1547), canon 9]{tanner-decrees-1990}
\item an inference step --- \cite{pohle-molinism-1911}; \cite[I-II q.112 a.3]{aquinas-summa-1920}; \cite[ch. VII]{garrigou-lagrange-grace-1952}
\end{itemize}
\end{itemize}
\par\medskip\noindent\textbf{\texttt{objection\_\allowbreak{}not\_\allowbreak{}forced}}\par\nobreak\noindent
Nor is the objection from the Spirit’s resistance forced: the maximal defensible position that holds Dort cannot hold it.

{\small\begin{verbatim}
theorem objection_not_forced : ¬Framework.SkepticallyAccepted
    solaGratiaDispute.defeats Party.objection
-- axioms: propext, Classical.choice, Quot.sound
\end{verbatim}}
\par\medskip\noindent\textbf{\texttt{third\allowbreak{}Way\allowbreak{}With\allowbreak{}Dort}}\par\nobreak\noindent
The third way, defended with Dort.

\paragraph{\emph{The third way (an ability given, received by faith)} is accepted on some resolution.}
\begin{itemize}
\item A position holding \emph{Romans 11:6 (by grace, so not by works)}, \emph{Reformed (the Canons of Dort)}, \emph{The third way (an ability given, received by faith)} can be held.
\item None of \emph{Romans 11:6 (by grace, so not by works)}, \emph{Reformed (the Canons of Dort)}, \emph{The third way (an ability given, received by faith)} defeats another, and each attack on them is answered from among them:
\begin{itemize}
\item \emph{Trent (Session VI, canon 4, on the Molinist reading)} defeats \emph{Reformed (the Canons of Dort)}, and \emph{Reformed (the Canons of Dort)} defeats \emph{Trent (Session VI, canon 4, on the Molinist reading)}.
\item \emph{Remonstrant (the Five Articles, 1610)} defeats \emph{Reformed (the Canons of Dort)}, and \emph{Reformed (the Canons of Dort)} defeats \emph{Remonstrant (the Five Articles, 1610)}.
\item \emph{Lutheran (the Formula of Concord)} defeats \emph{Reformed (the Canons of Dort)}, and \emph{Reformed (the Canons of Dort)} defeats \emph{Lutheran (the Formula of Concord)}.
\item \emph{Wesleyan (preventing grace)} defeats \emph{Reformed (the Canons of Dort)}, and \emph{Reformed (the Canons of Dort)} defeats \emph{Wesleyan (preventing grace)}.
\item \emph{The objection from the Spirit’s resistance} defeats \emph{Reformed (the Canons of Dort)}, and \emph{Reformed (the Canons of Dort)} defeats \emph{The objection from the Spirit’s resistance}.
\item \emph{Trent (Session VI, canon 4, on the Molinist reading)} defeats \emph{The third way (an ability given, received by faith)}, and \emph{Reformed (the Canons of Dort)} defeats \emph{Trent (Session VI, canon 4, on the Molinist reading)}.
\item \emph{Thomist (Aquinas; Garrigou-Lagrange)} defeats \emph{The third way (an ability given, received by faith)}, and \emph{The third way (an ability given, received by faith)} defeats \emph{Thomist (Aquinas; Garrigou-Lagrange)}.
\end{itemize}
\end{itemize}
\par\noindent\textbf{What this rests on.} The reasons state 13 defeats and 20 absences of defeat, each a cell of the defeat table, \texttt{sola\allowbreak{}Gratia\allowbreak{}Dispute\_\allowbreak{}defeats}, computed from the two parties' premises and checked by the kernel. Whether an attack survives turns on the attacker's weakest link and on the rating of what it attacks, so a changed rating can change the verdict. The weakest links:
\begin{itemize}
\item \emph{Romans 11:6 (by grace, so not by works)}, weakest at \emph{consensus}:
\begin{itemize}
\item Romans 11:6 — if it is by grace, it is no longer on the basis of works --- Rom 11:6
\item an inference step --- \cite[Trent, Session VI (1547), ch. 8]{tanner-decrees-1990}; \cite[36]{waterworth-trent-1848}; \cite[Five Arminian Articles (1610), art. 4]{schaff-creeds-christendom-3-1977}; \cite[Canons of Dort, III/IV, art. 10]{schaff-creeds-christendom-3-1977}
\end{itemize}
\item \emph{Reformed (the Canons of Dort)}, weakest at \emph{disputed}:
\begin{itemize}
\item What Stephen’s hearers resisted in Acts 7:51 was the grace that saves (encoded as a denial, so it ranks disputed) --- \cite[Canons of Dort, III/IV, art. 9]{schaff-creeds-christendom-3-1977}
\item an inference step --- \cite[ad loc. Eph 2:8]{robertson-word-pictures-4-1931}; \cite[228--229]{calvin-galatians-ephesians-1854}
\item an inference step --- \cite[600]{arminius-works-1-1825}; \cite[681]{arminius-works-1-1825}; \cite[Five Arminian Articles (1610), art. 4]{schaff-creeds-christendom-3-1977}
\item an inference step --- \cite[527--565]{arminius-works-3-1853}
\item an inference step --- \cite[Solid Declaration II.89]{concordia-triglotta-1921}; \cite[Solid Declaration II.60]{concordia-triglotta-1921}; \cite[Solid Declaration XI.41]{concordia-triglotta-1921}
\end{itemize}
\item \emph{The third way (an ability given, received by faith)}, weakest at \emph{disputed}:
\begin{itemize}
\item Grace is received by faith alone --- \cite[Augsburg Confession IV.1–3]{concordia-triglotta-1921}; \cite[Sermon 1, III.1]{wesley-sermons-1872}; \cite[Trent, Session VI (1547), canon 9]{tanner-decrees-1990}
\end{itemize}
\item \emph{Trent (Session VI, canon 4, on the Molinist reading)}, weakest at \emph{disputed}:
\begin{itemize}
\item The free will, moved by God, cooperates by assenting to God’s call --- \cite[Trent, Session VI (1547), canon 4]{tanner-decrees-1990}; \cite[45]{waterworth-trent-1848}
\item The free will, moved by God, can refuse its assent if it wishes --- \cite[Trent, Session VI (1547), canon 4]{tanner-decrees-1990}; \cite[45]{waterworth-trent-1848}
\item Grace is received by faith alone (encoded as a denial, so it ranks disputed) --- \cite[Augsburg Confession IV.1–3]{concordia-triglotta-1921}; \cite[Sermon 1, III.1]{wesley-sermons-1872}; \cite[Trent, Session VI (1547), canon 9]{tanner-decrees-1990}
\item an inference step --- \cite[I-II q.111 a.2]{aquinas-summa-1920}; \cite[Canons of Dort, III/IV, art. 12]{schaff-creeds-christendom-3-1977}; \cite[Trent, Session VI (1547), canon 4]{tanner-decrees-1990}; \cite[45]{waterworth-trent-1848}
\item an inference step --- \cite[I-II q.10 a.4 ad 3]{aquinas-summa-1920}; \cite[I-II q.112 a.3]{aquinas-summa-1920}; \cite[Trent, Session VI (1547), canon 4]{tanner-decrees-1990}; \cite[45]{waterworth-trent-1848}
\end{itemize}
\item \emph{Remonstrant (the Five Articles, 1610)}, weakest at \emph{disputed}:
\begin{itemize}
\item What Stephen’s hearers resisted in Acts 7:51 was the grace that saves --- \cite[Five Arminian Articles (1610), art. 4]{schaff-creeds-christendom-3-1977}
\end{itemize}
\item \emph{Lutheran (the Formula of Concord)}, weakest at \emph{disputed}:
\begin{itemize}
\item Those who obstinately resist the Spirit can grieve and lose him --- \cite[Solid Declaration II.83]{concordia-triglotta-1921}; \cite[Solid Declaration XI.42]{concordia-triglotta-1921}
\item Conversion is God’s work alone: God regenerates, and the will does nothing to it --- \cite[Solid Declaration II.89]{concordia-triglotta-1921}; \cite[Solid Declaration II.90]{concordia-triglotta-1921}
\item Grace is received by faith alone --- \cite[Augsburg Confession IV.1–3]{concordia-triglotta-1921}; \cite[Sermon 1, III.1]{wesley-sermons-1872}; \cite[Trent, Session VI (1547), canon 9]{tanner-decrees-1990}
\end{itemize}
\item \emph{Wesleyan (preventing grace)}, weakest at \emph{disputed}:
\begin{itemize}
\item The grace that saves can be finally refused --- \cite[397]{wesley-works-9-1827}
\item Grace is received by faith alone --- \cite[Augsburg Confession IV.1–3]{concordia-triglotta-1921}; \cite[Sermon 1, III.1]{wesley-sermons-1872}; \cite[Trent, Session VI (1547), canon 9]{tanner-decrees-1990}
\end{itemize}
\item \emph{The objection from the Spirit’s resistance}, weakest at \emph{disputed}:
\begin{itemize}
\item What Stephen’s hearers resisted in Acts 7:51 was the grace that saves --- \cite[Five Arminian Articles (1610), art. 4]{schaff-creeds-christendom-3-1977}
\item Those who obstinately resist the Spirit can grieve and lose him --- \cite[Solid Declaration II.83]{concordia-triglotta-1921}; \cite[Solid Declaration XI.42]{concordia-triglotta-1921}
\item The blasphemy against the Spirit is impenitence to the end: grace refused till death --- \cite[325]{npnf1-06-1888}; \cite[II-II q.14 a.1]{aquinas-summa-1920}; \cite[74]{calvin-harmony-evangelists-2-1845}
\item an inference step --- \cite[325]{npnf1-06-1888}; \cite[On Rebuke and Grace, ch. 43]{npnf1-05-1887}
\end{itemize}
\item \emph{Thomist (Aquinas; Garrigou-Lagrange)}, weakest at \emph{disputed}:
\begin{itemize}
\item In converting grace God is the sole mover: the mind is moved and does not move --- \cite[I-II q.111 a.2]{aquinas-summa-1920}; \cite[ch. IV]{garrigou-lagrange-grace-1952}
\item The free will, moved by God, cooperates by assenting to God’s call --- \cite[Trent, Session VI (1547), canon 4]{tanner-decrees-1990}; \cite[45]{waterworth-trent-1848}
\item The free will, moved by God, can refuse its assent if it wishes --- \cite[Trent, Session VI (1547), canon 4]{tanner-decrees-1990}; \cite[45]{waterworth-trent-1848}
\item Grace truly sufficient, giving the power to do good, is resisted --- \cite[ch. VI]{garrigou-lagrange-grace-1952}
\item Justifying grace, once given, can be lost --- \cite[II-II q.24 a.11]{aquinas-summa-1920}; \cite[Trent, Session VI (1547), canon 23]{tanner-decrees-1990}
\item Grace is received by faith alone (encoded as a denial, so it ranks disputed) --- \cite[II-II q.4 a.3]{aquinas-summa-1920}; \cite[I-II q.113 a.4 ad 1]{aquinas-summa-1920}; \cite[Trent, Session VI (1547), canon 9]{tanner-decrees-1990}
\item an inference step --- \cite{pohle-molinism-1911}; \cite[I-II q.112 a.3]{aquinas-summa-1920}; \cite[ch. VII]{garrigou-lagrange-grace-1952}
\end{itemize}
\end{itemize}
\par\medskip\noindent\textbf{\texttt{third\_\allowbreak{}way\_\allowbreak{}defensible\_\allowbreak{}with\_\allowbreak{}dort}}\par\nobreak\noindent
\textbf{The third way can be defended with Dort}, against Trent, by the position shown above. Trent attacks the third way over faith alone; Dort, which defeats Trent over the will and grace, answers that attack; and nothing in the third way contradicts Dort.

{\small\begin{verbatim}
theorem third_way_defensible_with_dort : Framework.CredulouslyAccepted
    solaGratiaDispute.defeats Party.thirdWay
-- axioms: propext, Classical.choice, Quot.sound
\end{verbatim}}
\par\medskip\noindent\textbf{\texttt{third\allowbreak{}Way\allowbreak{}With\allowbreak{}Wesley}}\par\nobreak\noindent
The third way, defended with Wesley.

\paragraph{\emph{The third way (an ability given, received by faith)} is accepted on some resolution.}
\begin{itemize}
\item A position holding \emph{Romans 11:6 (by grace, so not by works)}, \emph{Remonstrant (the Five Articles, 1610)}, \emph{Lutheran (the Formula of Concord)}, \emph{Wesleyan (preventing grace)}, \emph{The third way (an ability given, received by faith)}, \emph{The objection from the Spirit’s resistance} can be held.
\item None of \emph{Romans 11:6 (by grace, so not by works)}, \emph{Remonstrant (the Five Articles, 1610)}, \emph{Lutheran (the Formula of Concord)}, \emph{Wesleyan (preventing grace)}, \emph{The third way (an ability given, received by faith)}, \emph{The objection from the Spirit’s resistance} defeats another, and each attack on them is answered from among them:
\begin{itemize}
\item \emph{Reformed (the Canons of Dort)} defeats \emph{Remonstrant (the Five Articles, 1610)}, and \emph{Remonstrant (the Five Articles, 1610)} defeats \emph{Reformed (the Canons of Dort)}.
\item \emph{Thomist (Aquinas; Garrigou-Lagrange)} defeats \emph{Remonstrant (the Five Articles, 1610)}, and \emph{Remonstrant (the Five Articles, 1610)} defeats \emph{Thomist (Aquinas; Garrigou-Lagrange)}.
\item \emph{Reformed (the Canons of Dort)} defeats \emph{Lutheran (the Formula of Concord)}, and \emph{Remonstrant (the Five Articles, 1610)} defeats \emph{Reformed (the Canons of Dort)}.
\item \emph{Trent (Session VI, canon 4, on the Molinist reading)} defeats \emph{Lutheran (the Formula of Concord)}, and \emph{Lutheran (the Formula of Concord)} defeats \emph{Trent (Session VI, canon 4, on the Molinist reading)}.
\item \emph{Thomist (Aquinas; Garrigou-Lagrange)} defeats \emph{Lutheran (the Formula of Concord)}, and \emph{Remonstrant (the Five Articles, 1610)} defeats \emph{Thomist (Aquinas; Garrigou-Lagrange)}.
\item \emph{Reformed (the Canons of Dort)} defeats \emph{Wesleyan (preventing grace)}, and \emph{Remonstrant (the Five Articles, 1610)} defeats \emph{Reformed (the Canons of Dort)}.
\item \emph{Trent (Session VI, canon 4, on the Molinist reading)} defeats \emph{Wesleyan (preventing grace)}, and \emph{Lutheran (the Formula of Concord)} defeats \emph{Trent (Session VI, canon 4, on the Molinist reading)}.
\item \emph{Thomist (Aquinas; Garrigou-Lagrange)} defeats \emph{Wesleyan (preventing grace)}, and \emph{Remonstrant (the Five Articles, 1610)} defeats \emph{Thomist (Aquinas; Garrigou-Lagrange)}.
\item \emph{Trent (Session VI, canon 4, on the Molinist reading)} defeats \emph{The third way (an ability given, received by faith)}, and \emph{Lutheran (the Formula of Concord)} defeats \emph{Trent (Session VI, canon 4, on the Molinist reading)}.
\item \emph{Thomist (Aquinas; Garrigou-Lagrange)} defeats \emph{The third way (an ability given, received by faith)}, and \emph{Remonstrant (the Five Articles, 1610)} defeats \emph{Thomist (Aquinas; Garrigou-Lagrange)}.
\item \emph{Reformed (the Canons of Dort)} defeats \emph{The objection from the Spirit’s resistance}, and \emph{Remonstrant (the Five Articles, 1610)} defeats \emph{Reformed (the Canons of Dort)}.
\item \emph{Thomist (Aquinas; Garrigou-Lagrange)} defeats \emph{The objection from the Spirit’s resistance}, and \emph{Remonstrant (the Five Articles, 1610)} defeats \emph{Thomist (Aquinas; Garrigou-Lagrange)}.
\end{itemize}
\end{itemize}
\par\noindent\textbf{What this rests on.} The reasons state 15 defeats and 42 absences of defeat, each a cell of the defeat table, \texttt{sola\allowbreak{}Gratia\allowbreak{}Dispute\_\allowbreak{}defeats}, computed from the two parties' premises and checked by the kernel. Whether an attack survives turns on the attacker's weakest link and on the rating of what it attacks, so a changed rating can change the verdict. The weakest links:
\begin{itemize}
\item \emph{Romans 11:6 (by grace, so not by works)}, weakest at \emph{consensus}:
\begin{itemize}
\item Romans 11:6 — if it is by grace, it is no longer on the basis of works --- Rom 11:6
\item an inference step --- \cite[Trent, Session VI (1547), ch. 8]{tanner-decrees-1990}; \cite[36]{waterworth-trent-1848}; \cite[Five Arminian Articles (1610), art. 4]{schaff-creeds-christendom-3-1977}; \cite[Canons of Dort, III/IV, art. 10]{schaff-creeds-christendom-3-1977}
\end{itemize}
\item \emph{Remonstrant (the Five Articles, 1610)}, weakest at \emph{disputed}:
\begin{itemize}
\item What Stephen’s hearers resisted in Acts 7:51 was the grace that saves --- \cite[Five Arminian Articles (1610), art. 4]{schaff-creeds-christendom-3-1977}
\end{itemize}
\item \emph{Lutheran (the Formula of Concord)}, weakest at \emph{disputed}:
\begin{itemize}
\item Those who obstinately resist the Spirit can grieve and lose him --- \cite[Solid Declaration II.83]{concordia-triglotta-1921}; \cite[Solid Declaration XI.42]{concordia-triglotta-1921}
\item Conversion is God’s work alone: God regenerates, and the will does nothing to it --- \cite[Solid Declaration II.89]{concordia-triglotta-1921}; \cite[Solid Declaration II.90]{concordia-triglotta-1921}
\item Grace is received by faith alone --- \cite[Augsburg Confession IV.1–3]{concordia-triglotta-1921}; \cite[Sermon 1, III.1]{wesley-sermons-1872}; \cite[Trent, Session VI (1547), canon 9]{tanner-decrees-1990}
\end{itemize}
\item \emph{Wesleyan (preventing grace)}, weakest at \emph{disputed}:
\begin{itemize}
\item The grace that saves can be finally refused --- \cite[397]{wesley-works-9-1827}
\item Grace is received by faith alone --- \cite[Augsburg Confession IV.1–3]{concordia-triglotta-1921}; \cite[Sermon 1, III.1]{wesley-sermons-1872}; \cite[Trent, Session VI (1547), canon 9]{tanner-decrees-1990}
\end{itemize}
\item \emph{The third way (an ability given, received by faith)}, weakest at \emph{disputed}:
\begin{itemize}
\item Grace is received by faith alone --- \cite[Augsburg Confession IV.1–3]{concordia-triglotta-1921}; \cite[Sermon 1, III.1]{wesley-sermons-1872}; \cite[Trent, Session VI (1547), canon 9]{tanner-decrees-1990}
\end{itemize}
\item \emph{The objection from the Spirit’s resistance}, weakest at \emph{disputed}:
\begin{itemize}
\item What Stephen’s hearers resisted in Acts 7:51 was the grace that saves --- \cite[Five Arminian Articles (1610), art. 4]{schaff-creeds-christendom-3-1977}
\item Those who obstinately resist the Spirit can grieve and lose him --- \cite[Solid Declaration II.83]{concordia-triglotta-1921}; \cite[Solid Declaration XI.42]{concordia-triglotta-1921}
\item The blasphemy against the Spirit is impenitence to the end: grace refused till death --- \cite[325]{npnf1-06-1888}; \cite[II-II q.14 a.1]{aquinas-summa-1920}; \cite[74]{calvin-harmony-evangelists-2-1845}
\item an inference step --- \cite[325]{npnf1-06-1888}; \cite[On Rebuke and Grace, ch. 43]{npnf1-05-1887}
\end{itemize}
\item \emph{Reformed (the Canons of Dort)}, weakest at \emph{disputed}:
\begin{itemize}
\item What Stephen’s hearers resisted in Acts 7:51 was the grace that saves (encoded as a denial, so it ranks disputed) --- \cite[Canons of Dort, III/IV, art. 9]{schaff-creeds-christendom-3-1977}
\item an inference step --- \cite[ad loc. Eph 2:8]{robertson-word-pictures-4-1931}; \cite[228--229]{calvin-galatians-ephesians-1854}
\item an inference step --- \cite[600]{arminius-works-1-1825}; \cite[681]{arminius-works-1-1825}; \cite[Five Arminian Articles (1610), art. 4]{schaff-creeds-christendom-3-1977}
\item an inference step --- \cite[527--565]{arminius-works-3-1853}
\item an inference step --- \cite[Solid Declaration II.89]{concordia-triglotta-1921}; \cite[Solid Declaration II.60]{concordia-triglotta-1921}; \cite[Solid Declaration XI.41]{concordia-triglotta-1921}
\end{itemize}
\item \emph{Trent (Session VI, canon 4, on the Molinist reading)}, weakest at \emph{disputed}:
\begin{itemize}
\item The free will, moved by God, cooperates by assenting to God’s call --- \cite[Trent, Session VI (1547), canon 4]{tanner-decrees-1990}; \cite[45]{waterworth-trent-1848}
\item The free will, moved by God, can refuse its assent if it wishes --- \cite[Trent, Session VI (1547), canon 4]{tanner-decrees-1990}; \cite[45]{waterworth-trent-1848}
\item Grace is received by faith alone (encoded as a denial, so it ranks disputed) --- \cite[Augsburg Confession IV.1–3]{concordia-triglotta-1921}; \cite[Sermon 1, III.1]{wesley-sermons-1872}; \cite[Trent, Session VI (1547), canon 9]{tanner-decrees-1990}
\item an inference step --- \cite[I-II q.111 a.2]{aquinas-summa-1920}; \cite[Canons of Dort, III/IV, art. 12]{schaff-creeds-christendom-3-1977}; \cite[Trent, Session VI (1547), canon 4]{tanner-decrees-1990}; \cite[45]{waterworth-trent-1848}
\item an inference step --- \cite[I-II q.10 a.4 ad 3]{aquinas-summa-1920}; \cite[I-II q.112 a.3]{aquinas-summa-1920}; \cite[Trent, Session VI (1547), canon 4]{tanner-decrees-1990}; \cite[45]{waterworth-trent-1848}
\end{itemize}
\item \emph{Thomist (Aquinas; Garrigou-Lagrange)}, weakest at \emph{disputed}:
\begin{itemize}
\item In converting grace God is the sole mover: the mind is moved and does not move --- \cite[I-II q.111 a.2]{aquinas-summa-1920}; \cite[ch. IV]{garrigou-lagrange-grace-1952}
\item The free will, moved by God, cooperates by assenting to God’s call --- \cite[Trent, Session VI (1547), canon 4]{tanner-decrees-1990}; \cite[45]{waterworth-trent-1848}
\item The free will, moved by God, can refuse its assent if it wishes --- \cite[Trent, Session VI (1547), canon 4]{tanner-decrees-1990}; \cite[45]{waterworth-trent-1848}
\item Grace truly sufficient, giving the power to do good, is resisted --- \cite[ch. VI]{garrigou-lagrange-grace-1952}
\item Justifying grace, once given, can be lost --- \cite[II-II q.24 a.11]{aquinas-summa-1920}; \cite[Trent, Session VI (1547), canon 23]{tanner-decrees-1990}
\item Grace is received by faith alone (encoded as a denial, so it ranks disputed) --- \cite[II-II q.4 a.3]{aquinas-summa-1920}; \cite[I-II q.113 a.4 ad 1]{aquinas-summa-1920}; \cite[Trent, Session VI (1547), canon 9]{tanner-decrees-1990}
\item an inference step --- \cite{pohle-molinism-1911}; \cite[I-II q.112 a.3]{aquinas-summa-1920}; \cite[ch. VII]{garrigou-lagrange-grace-1952}
\end{itemize}
\end{itemize}
\par\medskip\noindent\textbf{\texttt{third\_\allowbreak{}way\_\allowbreak{}defensible\_\allowbreak{}with\_\allowbreak{}wesley}}\par\nobreak\noindent
It can also be defended with Wesley, the Lutherans and the Remonstrants, against both Dort and Trent, by the position shown above. The third way, as stated, takes no side between Dort and Wesley.

{\small\begin{verbatim}
theorem third_way_defensible_with_wesley : Framework.CredulouslyAccepted
    solaGratiaDispute.defeats Party.thirdWay
-- axioms: propext, Classical.choice, Quot.sound
\end{verbatim}}
\par\medskip\noindent\textbf{\texttt{third\allowbreak{}Way\allowbreak{}Answered\allowbreak{}By\allowbreak{}Trent}}\par\nobreak\noindent
Why the third way is not forced: a defensible position holds Trent, and Trent defeats it.

\paragraph{\emph{The third way (an ability given, received by faith)} is not accepted on every resolution.}
\begin{itemize}
\item A position holding \emph{Romans 11:6 (by grace, so not by works)}, \emph{Trent (Session VI, canon 4, on the Molinist reading)}, \emph{Remonstrant (the Five Articles, 1610)}, \emph{The objection from the Spirit’s resistance} can be held, and it holds \emph{Trent (Session VI, canon 4, on the Molinist reading)}; \emph{Trent (Session VI, canon 4, on the Molinist reading)} conflicts with \emph{The third way (an ability given, received by faith)}: \emph{Trent (Session VI, canon 4, on the Molinist reading)} defeats \emph{The third way (an ability given, received by faith)}, and \emph{The third way (an ability given, received by faith)} defeats \emph{Trent (Session VI, canon 4, on the Molinist reading)}.
\item None of \emph{Romans 11:6 (by grace, so not by works)}, \emph{Trent (Session VI, canon 4, on the Molinist reading)}, \emph{Remonstrant (the Five Articles, 1610)}, \emph{The objection from the Spirit’s resistance} defeats another, and each attack on them is answered from among them:
\begin{itemize}
\item \emph{Reformed (the Canons of Dort)} defeats \emph{Trent (Session VI, canon 4, on the Molinist reading)}, and \emph{Trent (Session VI, canon 4, on the Molinist reading)} defeats \emph{Reformed (the Canons of Dort)}.
\item \emph{Lutheran (the Formula of Concord)} defeats \emph{Trent (Session VI, canon 4, on the Molinist reading)}, and \emph{Trent (Session VI, canon 4, on the Molinist reading)} defeats \emph{Lutheran (the Formula of Concord)}.
\item \emph{Wesleyan (preventing grace)} defeats \emph{Trent (Session VI, canon 4, on the Molinist reading)}, and \emph{Trent (Session VI, canon 4, on the Molinist reading)} defeats \emph{Wesleyan (preventing grace)}.
\item \emph{The third way (an ability given, received by faith)} defeats \emph{Trent (Session VI, canon 4, on the Molinist reading)}, and \emph{Trent (Session VI, canon 4, on the Molinist reading)} defeats \emph{The third way (an ability given, received by faith)}.
\item \emph{Thomist (Aquinas; Garrigou-Lagrange)} defeats \emph{Trent (Session VI, canon 4, on the Molinist reading)}, and \emph{Trent (Session VI, canon 4, on the Molinist reading)} defeats \emph{Thomist (Aquinas; Garrigou-Lagrange)}.
\item \emph{Reformed (the Canons of Dort)} defeats \emph{Remonstrant (the Five Articles, 1610)}, and \emph{Trent (Session VI, canon 4, on the Molinist reading)} defeats \emph{Reformed (the Canons of Dort)}.
\item \emph{Thomist (Aquinas; Garrigou-Lagrange)} defeats \emph{Remonstrant (the Five Articles, 1610)}, and \emph{Trent (Session VI, canon 4, on the Molinist reading)} defeats \emph{Thomist (Aquinas; Garrigou-Lagrange)}.
\item \emph{Reformed (the Canons of Dort)} defeats \emph{The objection from the Spirit’s resistance}, and \emph{Trent (Session VI, canon 4, on the Molinist reading)} defeats \emph{Reformed (the Canons of Dort)}.
\item \emph{Thomist (Aquinas; Garrigou-Lagrange)} defeats \emph{The objection from the Spirit’s resistance}, and \emph{Trent (Session VI, canon 4, on the Molinist reading)} defeats \emph{Thomist (Aquinas; Garrigou-Lagrange)}.
\end{itemize}
\end{itemize}
\par\noindent\textbf{What this rests on.} The reasons state 14 defeats and 27 absences of defeat, each a cell of the defeat table, \texttt{sola\allowbreak{}Gratia\allowbreak{}Dispute\_\allowbreak{}defeats}, computed from the two parties' premises and checked by the kernel. Whether an attack survives turns on the attacker's weakest link and on the rating of what it attacks, so a changed rating can change the verdict. The weakest links:
\begin{itemize}
\item \emph{Trent (Session VI, canon 4, on the Molinist reading)}, weakest at \emph{disputed}:
\begin{itemize}
\item The free will, moved by God, cooperates by assenting to God’s call --- \cite[Trent, Session VI (1547), canon 4]{tanner-decrees-1990}; \cite[45]{waterworth-trent-1848}
\item The free will, moved by God, can refuse its assent if it wishes --- \cite[Trent, Session VI (1547), canon 4]{tanner-decrees-1990}; \cite[45]{waterworth-trent-1848}
\item Grace is received by faith alone (encoded as a denial, so it ranks disputed) --- \cite[Augsburg Confession IV.1–3]{concordia-triglotta-1921}; \cite[Sermon 1, III.1]{wesley-sermons-1872}; \cite[Trent, Session VI (1547), canon 9]{tanner-decrees-1990}
\item an inference step --- \cite[I-II q.111 a.2]{aquinas-summa-1920}; \cite[Canons of Dort, III/IV, art. 12]{schaff-creeds-christendom-3-1977}; \cite[Trent, Session VI (1547), canon 4]{tanner-decrees-1990}; \cite[45]{waterworth-trent-1848}
\item an inference step --- \cite[I-II q.10 a.4 ad 3]{aquinas-summa-1920}; \cite[I-II q.112 a.3]{aquinas-summa-1920}; \cite[Trent, Session VI (1547), canon 4]{tanner-decrees-1990}; \cite[45]{waterworth-trent-1848}
\end{itemize}
\item \emph{The third way (an ability given, received by faith)}, weakest at \emph{disputed}:
\begin{itemize}
\item Grace is received by faith alone --- \cite[Augsburg Confession IV.1–3]{concordia-triglotta-1921}; \cite[Sermon 1, III.1]{wesley-sermons-1872}; \cite[Trent, Session VI (1547), canon 9]{tanner-decrees-1990}
\end{itemize}
\item \emph{Romans 11:6 (by grace, so not by works)}, weakest at \emph{consensus}:
\begin{itemize}
\item Romans 11:6 — if it is by grace, it is no longer on the basis of works --- Rom 11:6
\item an inference step --- \cite[Trent, Session VI (1547), ch. 8]{tanner-decrees-1990}; \cite[36]{waterworth-trent-1848}; \cite[Five Arminian Articles (1610), art. 4]{schaff-creeds-christendom-3-1977}; \cite[Canons of Dort, III/IV, art. 10]{schaff-creeds-christendom-3-1977}
\end{itemize}
\item \emph{Remonstrant (the Five Articles, 1610)}, weakest at \emph{disputed}:
\begin{itemize}
\item What Stephen’s hearers resisted in Acts 7:51 was the grace that saves --- \cite[Five Arminian Articles (1610), art. 4]{schaff-creeds-christendom-3-1977}
\end{itemize}
\item \emph{The objection from the Spirit’s resistance}, weakest at \emph{disputed}:
\begin{itemize}
\item What Stephen’s hearers resisted in Acts 7:51 was the grace that saves --- \cite[Five Arminian Articles (1610), art. 4]{schaff-creeds-christendom-3-1977}
\item Those who obstinately resist the Spirit can grieve and lose him --- \cite[Solid Declaration II.83]{concordia-triglotta-1921}; \cite[Solid Declaration XI.42]{concordia-triglotta-1921}
\item The blasphemy against the Spirit is impenitence to the end: grace refused till death --- \cite[325]{npnf1-06-1888}; \cite[II-II q.14 a.1]{aquinas-summa-1920}; \cite[74]{calvin-harmony-evangelists-2-1845}
\item an inference step --- \cite[325]{npnf1-06-1888}; \cite[On Rebuke and Grace, ch. 43]{npnf1-05-1887}
\end{itemize}
\item \emph{Reformed (the Canons of Dort)}, weakest at \emph{disputed}:
\begin{itemize}
\item What Stephen’s hearers resisted in Acts 7:51 was the grace that saves (encoded as a denial, so it ranks disputed) --- \cite[Canons of Dort, III/IV, art. 9]{schaff-creeds-christendom-3-1977}
\item an inference step --- \cite[ad loc. Eph 2:8]{robertson-word-pictures-4-1931}; \cite[228--229]{calvin-galatians-ephesians-1854}
\item an inference step --- \cite[600]{arminius-works-1-1825}; \cite[681]{arminius-works-1-1825}; \cite[Five Arminian Articles (1610), art. 4]{schaff-creeds-christendom-3-1977}
\item an inference step --- \cite[527--565]{arminius-works-3-1853}
\item an inference step --- \cite[Solid Declaration II.89]{concordia-triglotta-1921}; \cite[Solid Declaration II.60]{concordia-triglotta-1921}; \cite[Solid Declaration XI.41]{concordia-triglotta-1921}
\end{itemize}
\item \emph{Lutheran (the Formula of Concord)}, weakest at \emph{disputed}:
\begin{itemize}
\item Those who obstinately resist the Spirit can grieve and lose him --- \cite[Solid Declaration II.83]{concordia-triglotta-1921}; \cite[Solid Declaration XI.42]{concordia-triglotta-1921}
\item Conversion is God’s work alone: God regenerates, and the will does nothing to it --- \cite[Solid Declaration II.89]{concordia-triglotta-1921}; \cite[Solid Declaration II.90]{concordia-triglotta-1921}
\item Grace is received by faith alone --- \cite[Augsburg Confession IV.1–3]{concordia-triglotta-1921}; \cite[Sermon 1, III.1]{wesley-sermons-1872}; \cite[Trent, Session VI (1547), canon 9]{tanner-decrees-1990}
\end{itemize}
\item \emph{Wesleyan (preventing grace)}, weakest at \emph{disputed}:
\begin{itemize}
\item The grace that saves can be finally refused --- \cite[397]{wesley-works-9-1827}
\item Grace is received by faith alone --- \cite[Augsburg Confession IV.1–3]{concordia-triglotta-1921}; \cite[Sermon 1, III.1]{wesley-sermons-1872}; \cite[Trent, Session VI (1547), canon 9]{tanner-decrees-1990}
\end{itemize}
\item \emph{Thomist (Aquinas; Garrigou-Lagrange)}, weakest at \emph{disputed}:
\begin{itemize}
\item In converting grace God is the sole mover: the mind is moved and does not move --- \cite[I-II q.111 a.2]{aquinas-summa-1920}; \cite[ch. IV]{garrigou-lagrange-grace-1952}
\item The free will, moved by God, cooperates by assenting to God’s call --- \cite[Trent, Session VI (1547), canon 4]{tanner-decrees-1990}; \cite[45]{waterworth-trent-1848}
\item The free will, moved by God, can refuse its assent if it wishes --- \cite[Trent, Session VI (1547), canon 4]{tanner-decrees-1990}; \cite[45]{waterworth-trent-1848}
\item Grace truly sufficient, giving the power to do good, is resisted --- \cite[ch. VI]{garrigou-lagrange-grace-1952}
\item Justifying grace, once given, can be lost --- \cite[II-II q.24 a.11]{aquinas-summa-1920}; \cite[Trent, Session VI (1547), canon 23]{tanner-decrees-1990}
\item Grace is received by faith alone (encoded as a denial, so it ranks disputed) --- \cite[II-II q.4 a.3]{aquinas-summa-1920}; \cite[I-II q.113 a.4 ad 1]{aquinas-summa-1920}; \cite[Trent, Session VI (1547), canon 9]{tanner-decrees-1990}
\item an inference step --- \cite{pohle-molinism-1911}; \cite[I-II q.112 a.3]{aquinas-summa-1920}; \cite[ch. VII]{garrigou-lagrange-grace-1952}
\end{itemize}
\end{itemize}
\par\medskip\noindent\textbf{\texttt{third\_\allowbreak{}way\_\allowbreak{}not\_\allowbreak{}forced}}\par\nobreak\noindent
Nor is the third way forced. A maximal defensible position holds Trent, which condemns faith alone, and cannot hold the third way with it.

{\small\begin{verbatim}
theorem third_way_not_forced : ¬Framework.SkepticallyAccepted
    solaGratiaDispute.defeats Party.thirdWay
-- axioms: propext, Classical.choice, Quot.sound
\end{verbatim}}
\par\medskip\noindent\textbf{\texttt{dort\_\allowbreak{}stands\_\allowbreak{}with\_\allowbreak{}the\_\allowbreak{}thomists}}\par\nobreak\noindent
\textbf{Dort and the Thomists stand together}: one reading satisfies every premise and every conclusion of both, so neither defeats the other. The Thomist position weighed here takes no side on the will’s passivity. It holds that the will is “not merely passive” and draws nothing from it, so this result does not show that the Thomists hold the third part of sola gratia. Read in the Molinists’ sense, the phrase cannot be held with Dort (\texttt{where\allowbreak{}The\allowbreak{}Agent\allowbreak{}Reading\allowbreak{}Falls}); read in Garrigou-Lagrange’s, it can (\texttt{dort\_\allowbreak{}holds\_\allowbreak{}with\_\allowbreak{}garrigou\_\allowbreak{}lagrange}). They differ on faith alone and on perseverance, which this dispute does not weigh between them.

{\small\begin{verbatim}
theorem dort_stands_with_the_thomists : solaGratiaDispute.StandTogether
    [Party.reformed, Party.thomist]
-- axioms: propext, Classical.choice, Quot.sound
\end{verbatim}}
\par\medskip\noindent\textbf{\texttt{thomists\allowbreak{}With\allowbreak{}Dort}}\par\nobreak\noindent
The Thomists, defended with Dort.

\paragraph{\emph{Thomist (Aquinas; Garrigou-Lagrange)} is accepted on some resolution.}
\begin{itemize}
\item A position holding \emph{Romans 11:6 (by grace, so not by works)}, \emph{Reformed (the Canons of Dort)}, \emph{Thomist (Aquinas; Garrigou-Lagrange)} can be held.
\item None of \emph{Romans 11:6 (by grace, so not by works)}, \emph{Reformed (the Canons of Dort)}, \emph{Thomist (Aquinas; Garrigou-Lagrange)} defeats another, and each attack on them is answered from among them:
\begin{itemize}
\item \emph{Trent (Session VI, canon 4, on the Molinist reading)} defeats \emph{Reformed (the Canons of Dort)}, and \emph{Reformed (the Canons of Dort)} defeats \emph{Trent (Session VI, canon 4, on the Molinist reading)}.
\item \emph{Remonstrant (the Five Articles, 1610)} defeats \emph{Reformed (the Canons of Dort)}, and \emph{Reformed (the Canons of Dort)} defeats \emph{Remonstrant (the Five Articles, 1610)}.
\item \emph{Lutheran (the Formula of Concord)} defeats \emph{Reformed (the Canons of Dort)}, and \emph{Reformed (the Canons of Dort)} defeats \emph{Lutheran (the Formula of Concord)}.
\item \emph{Wesleyan (preventing grace)} defeats \emph{Reformed (the Canons of Dort)}, and \emph{Reformed (the Canons of Dort)} defeats \emph{Wesleyan (preventing grace)}.
\item \emph{The objection from the Spirit’s resistance} defeats \emph{Reformed (the Canons of Dort)}, and \emph{Reformed (the Canons of Dort)} defeats \emph{The objection from the Spirit’s resistance}.
\item \emph{Trent (Session VI, canon 4, on the Molinist reading)} defeats \emph{Thomist (Aquinas; Garrigou-Lagrange)}, and \emph{Reformed (the Canons of Dort)} defeats \emph{Trent (Session VI, canon 4, on the Molinist reading)}.
\item \emph{Remonstrant (the Five Articles, 1610)} defeats \emph{Thomist (Aquinas; Garrigou-Lagrange)}, and \emph{Reformed (the Canons of Dort)} defeats \emph{Remonstrant (the Five Articles, 1610)}.
\item \emph{Lutheran (the Formula of Concord)} defeats \emph{Thomist (Aquinas; Garrigou-Lagrange)}, and \emph{Reformed (the Canons of Dort)} defeats \emph{Lutheran (the Formula of Concord)}.
\item \emph{Wesleyan (preventing grace)} defeats \emph{Thomist (Aquinas; Garrigou-Lagrange)}, and \emph{Reformed (the Canons of Dort)} defeats \emph{Wesleyan (preventing grace)}.
\item \emph{The third way (an ability given, received by faith)} defeats \emph{Thomist (Aquinas; Garrigou-Lagrange)}, and \emph{Thomist (Aquinas; Garrigou-Lagrange)} defeats \emph{The third way (an ability given, received by faith)}.
\item \emph{The objection from the Spirit’s resistance} defeats \emph{Thomist (Aquinas; Garrigou-Lagrange)}, and \emph{Reformed (the Canons of Dort)} defeats \emph{The objection from the Spirit’s resistance}.
\end{itemize}
\end{itemize}
\par\noindent\textbf{What this rests on.} The reasons state 17 defeats and 16 absences of defeat, each a cell of the defeat table, \texttt{sola\allowbreak{}Gratia\allowbreak{}Dispute\_\allowbreak{}defeats}, computed from the two parties' premises and checked by the kernel. Whether an attack survives turns on the attacker's weakest link and on the rating of what it attacks, so a changed rating can change the verdict. The weakest links:
\begin{itemize}
\item \emph{Romans 11:6 (by grace, so not by works)}, weakest at \emph{consensus}:
\begin{itemize}
\item Romans 11:6 — if it is by grace, it is no longer on the basis of works --- Rom 11:6
\item an inference step --- \cite[Trent, Session VI (1547), ch. 8]{tanner-decrees-1990}; \cite[36]{waterworth-trent-1848}; \cite[Five Arminian Articles (1610), art. 4]{schaff-creeds-christendom-3-1977}; \cite[Canons of Dort, III/IV, art. 10]{schaff-creeds-christendom-3-1977}
\end{itemize}
\item \emph{Reformed (the Canons of Dort)}, weakest at \emph{disputed}:
\begin{itemize}
\item What Stephen’s hearers resisted in Acts 7:51 was the grace that saves (encoded as a denial, so it ranks disputed) --- \cite[Canons of Dort, III/IV, art. 9]{schaff-creeds-christendom-3-1977}
\item an inference step --- \cite[ad loc. Eph 2:8]{robertson-word-pictures-4-1931}; \cite[228--229]{calvin-galatians-ephesians-1854}
\item an inference step --- \cite[600]{arminius-works-1-1825}; \cite[681]{arminius-works-1-1825}; \cite[Five Arminian Articles (1610), art. 4]{schaff-creeds-christendom-3-1977}
\item an inference step --- \cite[527--565]{arminius-works-3-1853}
\item an inference step --- \cite[Solid Declaration II.89]{concordia-triglotta-1921}; \cite[Solid Declaration II.60]{concordia-triglotta-1921}; \cite[Solid Declaration XI.41]{concordia-triglotta-1921}
\end{itemize}
\item \emph{Thomist (Aquinas; Garrigou-Lagrange)}, weakest at \emph{disputed}:
\begin{itemize}
\item In converting grace God is the sole mover: the mind is moved and does not move --- \cite[I-II q.111 a.2]{aquinas-summa-1920}; \cite[ch. IV]{garrigou-lagrange-grace-1952}
\item The free will, moved by God, cooperates by assenting to God’s call --- \cite[Trent, Session VI (1547), canon 4]{tanner-decrees-1990}; \cite[45]{waterworth-trent-1848}
\item The free will, moved by God, can refuse its assent if it wishes --- \cite[Trent, Session VI (1547), canon 4]{tanner-decrees-1990}; \cite[45]{waterworth-trent-1848}
\item Grace truly sufficient, giving the power to do good, is resisted --- \cite[ch. VI]{garrigou-lagrange-grace-1952}
\item Justifying grace, once given, can be lost --- \cite[II-II q.24 a.11]{aquinas-summa-1920}; \cite[Trent, Session VI (1547), canon 23]{tanner-decrees-1990}
\item Grace is received by faith alone (encoded as a denial, so it ranks disputed) --- \cite[II-II q.4 a.3]{aquinas-summa-1920}; \cite[I-II q.113 a.4 ad 1]{aquinas-summa-1920}; \cite[Trent, Session VI (1547), canon 9]{tanner-decrees-1990}
\item an inference step --- \cite{pohle-molinism-1911}; \cite[I-II q.112 a.3]{aquinas-summa-1920}; \cite[ch. VII]{garrigou-lagrange-grace-1952}
\end{itemize}
\item \emph{Trent (Session VI, canon 4, on the Molinist reading)}, weakest at \emph{disputed}:
\begin{itemize}
\item The free will, moved by God, cooperates by assenting to God’s call --- \cite[Trent, Session VI (1547), canon 4]{tanner-decrees-1990}; \cite[45]{waterworth-trent-1848}
\item The free will, moved by God, can refuse its assent if it wishes --- \cite[Trent, Session VI (1547), canon 4]{tanner-decrees-1990}; \cite[45]{waterworth-trent-1848}
\item Grace is received by faith alone (encoded as a denial, so it ranks disputed) --- \cite[Augsburg Confession IV.1–3]{concordia-triglotta-1921}; \cite[Sermon 1, III.1]{wesley-sermons-1872}; \cite[Trent, Session VI (1547), canon 9]{tanner-decrees-1990}
\item an inference step --- \cite[I-II q.111 a.2]{aquinas-summa-1920}; \cite[Canons of Dort, III/IV, art. 12]{schaff-creeds-christendom-3-1977}; \cite[Trent, Session VI (1547), canon 4]{tanner-decrees-1990}; \cite[45]{waterworth-trent-1848}
\item an inference step --- \cite[I-II q.10 a.4 ad 3]{aquinas-summa-1920}; \cite[I-II q.112 a.3]{aquinas-summa-1920}; \cite[Trent, Session VI (1547), canon 4]{tanner-decrees-1990}; \cite[45]{waterworth-trent-1848}
\end{itemize}
\item \emph{Remonstrant (the Five Articles, 1610)}, weakest at \emph{disputed}:
\begin{itemize}
\item What Stephen’s hearers resisted in Acts 7:51 was the grace that saves --- \cite[Five Arminian Articles (1610), art. 4]{schaff-creeds-christendom-3-1977}
\end{itemize}
\item \emph{Lutheran (the Formula of Concord)}, weakest at \emph{disputed}:
\begin{itemize}
\item Those who obstinately resist the Spirit can grieve and lose him --- \cite[Solid Declaration II.83]{concordia-triglotta-1921}; \cite[Solid Declaration XI.42]{concordia-triglotta-1921}
\item Conversion is God’s work alone: God regenerates, and the will does nothing to it --- \cite[Solid Declaration II.89]{concordia-triglotta-1921}; \cite[Solid Declaration II.90]{concordia-triglotta-1921}
\item Grace is received by faith alone --- \cite[Augsburg Confession IV.1–3]{concordia-triglotta-1921}; \cite[Sermon 1, III.1]{wesley-sermons-1872}; \cite[Trent, Session VI (1547), canon 9]{tanner-decrees-1990}
\end{itemize}
\item \emph{Wesleyan (preventing grace)}, weakest at \emph{disputed}:
\begin{itemize}
\item The grace that saves can be finally refused --- \cite[397]{wesley-works-9-1827}
\item Grace is received by faith alone --- \cite[Augsburg Confession IV.1–3]{concordia-triglotta-1921}; \cite[Sermon 1, III.1]{wesley-sermons-1872}; \cite[Trent, Session VI (1547), canon 9]{tanner-decrees-1990}
\end{itemize}
\item \emph{The third way (an ability given, received by faith)}, weakest at \emph{disputed}:
\begin{itemize}
\item Grace is received by faith alone --- \cite[Augsburg Confession IV.1–3]{concordia-triglotta-1921}; \cite[Sermon 1, III.1]{wesley-sermons-1872}; \cite[Trent, Session VI (1547), canon 9]{tanner-decrees-1990}
\end{itemize}
\item \emph{The objection from the Spirit’s resistance}, weakest at \emph{disputed}:
\begin{itemize}
\item What Stephen’s hearers resisted in Acts 7:51 was the grace that saves --- \cite[Five Arminian Articles (1610), art. 4]{schaff-creeds-christendom-3-1977}
\item Those who obstinately resist the Spirit can grieve and lose him --- \cite[Solid Declaration II.83]{concordia-triglotta-1921}; \cite[Solid Declaration XI.42]{concordia-triglotta-1921}
\item The blasphemy against the Spirit is impenitence to the end: grace refused till death --- \cite[325]{npnf1-06-1888}; \cite[II-II q.14 a.1]{aquinas-summa-1920}; \cite[74]{calvin-harmony-evangelists-2-1845}
\item an inference step --- \cite[325]{npnf1-06-1888}; \cite[On Rebuke and Grace, ch. 43]{npnf1-05-1887}
\end{itemize}
\end{itemize}
\par\medskip\noindent\textbf{\texttt{thomist\_\allowbreak{}defensible}}\par\nobreak\noindent
\textbf{The Thomists can be defended, with Dort}: some maximal defensible position holds both, with what every party holds. Between them they answer every attack on either.

{\small\begin{verbatim}
theorem thomist_defensible : Framework.CredulouslyAccepted
    solaGratiaDispute.defeats Party.thomist
-- axioms: propext, Classical.choice, Quot.sound
\end{verbatim}}
\par\medskip\noindent\textbf{\texttt{thomists\allowbreak{}Answered\allowbreak{}By\allowbreak{}Trent}}\par\nobreak\noindent
Why the Thomists are not forced: a maximal defensible position holds Trent on the Molinist reading, and Trent defeats them.

\paragraph{\emph{Thomist (Aquinas; Garrigou-Lagrange)} is not accepted on every resolution.}
\begin{itemize}
\item A position holding \emph{Romans 11:6 (by grace, so not by works)}, \emph{Trent (Session VI, canon 4, on the Molinist reading)}, \emph{Remonstrant (the Five Articles, 1610)}, \emph{The objection from the Spirit’s resistance} can be held, and it holds \emph{Trent (Session VI, canon 4, on the Molinist reading)}; \emph{Trent (Session VI, canon 4, on the Molinist reading)} conflicts with \emph{Thomist (Aquinas; Garrigou-Lagrange)}: \emph{Trent (Session VI, canon 4, on the Molinist reading)} defeats \emph{Thomist (Aquinas; Garrigou-Lagrange)}, and \emph{Thomist (Aquinas; Garrigou-Lagrange)} defeats \emph{Trent (Session VI, canon 4, on the Molinist reading)}.
\item None of \emph{Romans 11:6 (by grace, so not by works)}, \emph{Trent (Session VI, canon 4, on the Molinist reading)}, \emph{Remonstrant (the Five Articles, 1610)}, \emph{The objection from the Spirit’s resistance} defeats another, and each attack on them is answered from among them:
\begin{itemize}
\item \emph{Reformed (the Canons of Dort)} defeats \emph{Trent (Session VI, canon 4, on the Molinist reading)}, and \emph{Trent (Session VI, canon 4, on the Molinist reading)} defeats \emph{Reformed (the Canons of Dort)}.
\item \emph{Lutheran (the Formula of Concord)} defeats \emph{Trent (Session VI, canon 4, on the Molinist reading)}, and \emph{Trent (Session VI, canon 4, on the Molinist reading)} defeats \emph{Lutheran (the Formula of Concord)}.
\item \emph{Wesleyan (preventing grace)} defeats \emph{Trent (Session VI, canon 4, on the Molinist reading)}, and \emph{Trent (Session VI, canon 4, on the Molinist reading)} defeats \emph{Wesleyan (preventing grace)}.
\item \emph{The third way (an ability given, received by faith)} defeats \emph{Trent (Session VI, canon 4, on the Molinist reading)}, and \emph{Trent (Session VI, canon 4, on the Molinist reading)} defeats \emph{The third way (an ability given, received by faith)}.
\item \emph{Thomist (Aquinas; Garrigou-Lagrange)} defeats \emph{Trent (Session VI, canon 4, on the Molinist reading)}, and \emph{Trent (Session VI, canon 4, on the Molinist reading)} defeats \emph{Thomist (Aquinas; Garrigou-Lagrange)}.
\item \emph{Reformed (the Canons of Dort)} defeats \emph{Remonstrant (the Five Articles, 1610)}, and \emph{Trent (Session VI, canon 4, on the Molinist reading)} defeats \emph{Reformed (the Canons of Dort)}.
\item \emph{Thomist (Aquinas; Garrigou-Lagrange)} defeats \emph{Remonstrant (the Five Articles, 1610)}, and \emph{Trent (Session VI, canon 4, on the Molinist reading)} defeats \emph{Thomist (Aquinas; Garrigou-Lagrange)}.
\item \emph{Reformed (the Canons of Dort)} defeats \emph{The objection from the Spirit’s resistance}, and \emph{Trent (Session VI, canon 4, on the Molinist reading)} defeats \emph{Reformed (the Canons of Dort)}.
\item \emph{Thomist (Aquinas; Garrigou-Lagrange)} defeats \emph{The objection from the Spirit’s resistance}, and \emph{Trent (Session VI, canon 4, on the Molinist reading)} defeats \emph{Thomist (Aquinas; Garrigou-Lagrange)}.
\end{itemize}
\end{itemize}
\par\noindent\textbf{What this rests on.} The reasons state 14 defeats and 27 absences of defeat, each a cell of the defeat table, \texttt{sola\allowbreak{}Gratia\allowbreak{}Dispute\_\allowbreak{}defeats}, computed from the two parties' premises and checked by the kernel. Whether an attack survives turns on the attacker's weakest link and on the rating of what it attacks, so a changed rating can change the verdict. The weakest links:
\begin{itemize}
\item \emph{Trent (Session VI, canon 4, on the Molinist reading)}, weakest at \emph{disputed}:
\begin{itemize}
\item The free will, moved by God, cooperates by assenting to God’s call --- \cite[Trent, Session VI (1547), canon 4]{tanner-decrees-1990}; \cite[45]{waterworth-trent-1848}
\item The free will, moved by God, can refuse its assent if it wishes --- \cite[Trent, Session VI (1547), canon 4]{tanner-decrees-1990}; \cite[45]{waterworth-trent-1848}
\item Grace is received by faith alone (encoded as a denial, so it ranks disputed) --- \cite[Augsburg Confession IV.1–3]{concordia-triglotta-1921}; \cite[Sermon 1, III.1]{wesley-sermons-1872}; \cite[Trent, Session VI (1547), canon 9]{tanner-decrees-1990}
\item an inference step --- \cite[I-II q.111 a.2]{aquinas-summa-1920}; \cite[Canons of Dort, III/IV, art. 12]{schaff-creeds-christendom-3-1977}; \cite[Trent, Session VI (1547), canon 4]{tanner-decrees-1990}; \cite[45]{waterworth-trent-1848}
\item an inference step --- \cite[I-II q.10 a.4 ad 3]{aquinas-summa-1920}; \cite[I-II q.112 a.3]{aquinas-summa-1920}; \cite[Trent, Session VI (1547), canon 4]{tanner-decrees-1990}; \cite[45]{waterworth-trent-1848}
\end{itemize}
\item \emph{Thomist (Aquinas; Garrigou-Lagrange)}, weakest at \emph{disputed}:
\begin{itemize}
\item In converting grace God is the sole mover: the mind is moved and does not move --- \cite[I-II q.111 a.2]{aquinas-summa-1920}; \cite[ch. IV]{garrigou-lagrange-grace-1952}
\item The free will, moved by God, cooperates by assenting to God’s call --- \cite[Trent, Session VI (1547), canon 4]{tanner-decrees-1990}; \cite[45]{waterworth-trent-1848}
\item The free will, moved by God, can refuse its assent if it wishes --- \cite[Trent, Session VI (1547), canon 4]{tanner-decrees-1990}; \cite[45]{waterworth-trent-1848}
\item Grace truly sufficient, giving the power to do good, is resisted --- \cite[ch. VI]{garrigou-lagrange-grace-1952}
\item Justifying grace, once given, can be lost --- \cite[II-II q.24 a.11]{aquinas-summa-1920}; \cite[Trent, Session VI (1547), canon 23]{tanner-decrees-1990}
\item Grace is received by faith alone (encoded as a denial, so it ranks disputed) --- \cite[II-II q.4 a.3]{aquinas-summa-1920}; \cite[I-II q.113 a.4 ad 1]{aquinas-summa-1920}; \cite[Trent, Session VI (1547), canon 9]{tanner-decrees-1990}
\item an inference step --- \cite{pohle-molinism-1911}; \cite[I-II q.112 a.3]{aquinas-summa-1920}; \cite[ch. VII]{garrigou-lagrange-grace-1952}
\end{itemize}
\item \emph{Romans 11:6 (by grace, so not by works)}, weakest at \emph{consensus}:
\begin{itemize}
\item Romans 11:6 — if it is by grace, it is no longer on the basis of works --- Rom 11:6
\item an inference step --- \cite[Trent, Session VI (1547), ch. 8]{tanner-decrees-1990}; \cite[36]{waterworth-trent-1848}; \cite[Five Arminian Articles (1610), art. 4]{schaff-creeds-christendom-3-1977}; \cite[Canons of Dort, III/IV, art. 10]{schaff-creeds-christendom-3-1977}
\end{itemize}
\item \emph{Remonstrant (the Five Articles, 1610)}, weakest at \emph{disputed}:
\begin{itemize}
\item What Stephen’s hearers resisted in Acts 7:51 was the grace that saves --- \cite[Five Arminian Articles (1610), art. 4]{schaff-creeds-christendom-3-1977}
\end{itemize}
\item \emph{The objection from the Spirit’s resistance}, weakest at \emph{disputed}:
\begin{itemize}
\item What Stephen’s hearers resisted in Acts 7:51 was the grace that saves --- \cite[Five Arminian Articles (1610), art. 4]{schaff-creeds-christendom-3-1977}
\item Those who obstinately resist the Spirit can grieve and lose him --- \cite[Solid Declaration II.83]{concordia-triglotta-1921}; \cite[Solid Declaration XI.42]{concordia-triglotta-1921}
\item The blasphemy against the Spirit is impenitence to the end: grace refused till death --- \cite[325]{npnf1-06-1888}; \cite[II-II q.14 a.1]{aquinas-summa-1920}; \cite[74]{calvin-harmony-evangelists-2-1845}
\item an inference step --- \cite[325]{npnf1-06-1888}; \cite[On Rebuke and Grace, ch. 43]{npnf1-05-1887}
\end{itemize}
\item \emph{Reformed (the Canons of Dort)}, weakest at \emph{disputed}:
\begin{itemize}
\item What Stephen’s hearers resisted in Acts 7:51 was the grace that saves (encoded as a denial, so it ranks disputed) --- \cite[Canons of Dort, III/IV, art. 9]{schaff-creeds-christendom-3-1977}
\item an inference step --- \cite[ad loc. Eph 2:8]{robertson-word-pictures-4-1931}; \cite[228--229]{calvin-galatians-ephesians-1854}
\item an inference step --- \cite[600]{arminius-works-1-1825}; \cite[681]{arminius-works-1-1825}; \cite[Five Arminian Articles (1610), art. 4]{schaff-creeds-christendom-3-1977}
\item an inference step --- \cite[527--565]{arminius-works-3-1853}
\item an inference step --- \cite[Solid Declaration II.89]{concordia-triglotta-1921}; \cite[Solid Declaration II.60]{concordia-triglotta-1921}; \cite[Solid Declaration XI.41]{concordia-triglotta-1921}
\end{itemize}
\item \emph{Lutheran (the Formula of Concord)}, weakest at \emph{disputed}:
\begin{itemize}
\item Those who obstinately resist the Spirit can grieve and lose him --- \cite[Solid Declaration II.83]{concordia-triglotta-1921}; \cite[Solid Declaration XI.42]{concordia-triglotta-1921}
\item Conversion is God’s work alone: God regenerates, and the will does nothing to it --- \cite[Solid Declaration II.89]{concordia-triglotta-1921}; \cite[Solid Declaration II.90]{concordia-triglotta-1921}
\item Grace is received by faith alone --- \cite[Augsburg Confession IV.1–3]{concordia-triglotta-1921}; \cite[Sermon 1, III.1]{wesley-sermons-1872}; \cite[Trent, Session VI (1547), canon 9]{tanner-decrees-1990}
\end{itemize}
\item \emph{Wesleyan (preventing grace)}, weakest at \emph{disputed}:
\begin{itemize}
\item The grace that saves can be finally refused --- \cite[397]{wesley-works-9-1827}
\item Grace is received by faith alone --- \cite[Augsburg Confession IV.1–3]{concordia-triglotta-1921}; \cite[Sermon 1, III.1]{wesley-sermons-1872}; \cite[Trent, Session VI (1547), canon 9]{tanner-decrees-1990}
\end{itemize}
\item \emph{The third way (an ability given, received by faith)}, weakest at \emph{disputed}:
\begin{itemize}
\item Grace is received by faith alone --- \cite[Augsburg Confession IV.1–3]{concordia-triglotta-1921}; \cite[Sermon 1, III.1]{wesley-sermons-1872}; \cite[Trent, Session VI (1547), canon 9]{tanner-decrees-1990}
\end{itemize}
\end{itemize}
\par\medskip\noindent\textbf{\texttt{thomist\_\allowbreak{}not\_\allowbreak{}forced}}\par\nobreak\noindent
Nor are the Thomists forced. A maximal defensible position holds Trent on the Molinist reading, with the Remonstrants and the objection, and cannot hold the Thomists with them: the dispute between the two Catholic schools, which Rome left open in 1607, is a standoff here too.

{\small\begin{verbatim}
theorem thomist_not_forced : ¬Framework.SkepticallyAccepted
    solaGratiaDispute.defeats Party.thomist
-- axioms: propext, Classical.choice, Quot.sound
\end{verbatim}}
\subsection*{The perseverance of the saints}
\addcontentsline{toc}{subsection}{The perseverance of the saints}
“My sheep hear my voice, and I know them, and they follow me. I give them eternal life, and they will never perish, and no one will snatch them out of my hand” (John 10:27–28). If the grace that saves brings a person to believe, can the one it saves lose it again? The Westminster Confession of 1647 answers that those “whom God hath accepted in his Beloved, effectually called and sanctified by his Spirit, can neither totally nor finally fall away from the state of grace; but shall certainly persevere therein to the end, and be eternally saved” (XVII.1). To be \emph{effectually called} is to be called in a way that brings one to faith. The Synod of Dort had said the same in its Fifth Head of Doctrine in 1619: God “does not wholly withdraw the Holy Spirit from his own people, even in their melancholy falls” (V, article 6). This is the \emph{perseverance of the saints}.

\subsubsection*{Two claims, and the claim that joins them}
Westminster’s sentence is about the effectually called, and its tenth chapter makes them the elect, those God chose to save. Of the elect, Augustine had said the same about twelve centuries earlier. Those “called according to God’s purpose” are elected, and those who are not to persevere “are not elected” (\emph{On Rebuke and Grace}, ch. 16). Aquinas says that predestination “most certainly and infallibly takes effect” (\emph{Summa} I q.23 a.6), and the Formula of Concord that God placed the elect’s salvation “in the almighty hand of our Savior Jesus Christ, from which no one can pluck us, John 10:28” (Solid Declaration XI.45). So Augustine, Aquinas and the Lutherans agree with Westminster that the elect persevere, and divide from it over who the elect are: whether everyone God brings to true faith is among those predestined to persevere. That is not the question of whether God chose them because he foresaw their faith, which this chapter does not weigh. Trent and Wesley are not encoded as taking a side on the elect.

Augustine held that some who are truly \emph{regenerate}, given new life by the Spirit, are not elect: “to some of His own children — whom He has regenerated in Christ — to whom He has given faith, hope, and love, God does not give perseverance also” (ch. 18). The Reformed hold the contrary. Everyone truly justified, that is, put right with God by faith, is among the elect, and those who fall away never had justifying faith. Among the texts they cite is 1 John 2:19: “They went out from us, but they were not of us; for if they had been of us, they would have continued with us.” Calvin puts it so on Hebrews 6: “God indeed favours none but the elect alone with the Spirit of regeneration”.

So the Reformed conclusion is written as two claims: the elect persevere, and no one truly justified can finally fall away. A third claim joins them, that the truly justified are elect. If “the perseverance of the saints” were one claim, “saints” would have to mean both the elect and everyone truly justified. Augustine agrees about the first and denies the second, so as one claim he would seem to contradict Westminster even where he agrees with it.

\subsubsection*{The three routes, and where each is disputed}
\textbf{John 10 and 1 John 2.} From John 10:27–29, the Reformed conclude that the elect persevere. Calvin reads the verse so: “the salvation of all the elect is not less certain than the power of God is invincible”. Wesley grants the text and reads its promise as conditional, quoting 2 Peter 1:10: “They are \emph{my sheep}, that \emph{hear my voice}, that \emph{follow me} in all holiness. And ‘if ye do these things, ye shall never fall’”. So the step is rated \emph{disputed}, which means that a reader the library cites grants its grounds and denies its conclusion, and here gives his reason, the condition he reads in the promise.

From 1 John 2:19, “They went out from us, but they were not of us”, the Reformed conclude that the truly justified are elect. They get there through a reading of the verse. Calvin: “He plainly declares that those who fell away had never been members of the Church”, having had “only a light and a transient taste” of the knowledge of Christ. John Owen, the Puritan, makes it a rule for every apostate: “If all they who fall away totally … were never true believers … then those who are true believers cannot fall away; but the first is true, therefore the latter.” Here too the reading is the doctrine found in the verse, so the step from it to “the truly justified are elect” is rated \emph{well supported} and the whole weight falls on the reading. Augustine reads the same verse and holds that those who left “were in goodness”, “in the faith of children”, and still were not elect (ch. 20). So the reading is rated \emph{disputed}, and the question it turns on is whether one who leaves the faith never had it, or had it and gave it up.

\textbf{Philippians 1.} “He who began a good work in you will bring it to completion at the day of Jesus Christ” (Philippians 1:6). Trent’s chapter on perseverance first says that “all ought to place and repose a most firm hope in God’s help”, and then quotes the verse with a condition: “For God, unless men be themselves wanting to His grace, as he has begun the good work, so will he perfect it” (Session VI, ch. 13). So the step from the verse to “no one truly justified can finally fall away” is rated \emph{disputed}.

\textbf{Romans 8.} “And those whom he predestined he also called, and those whom he called he also justified, and those whom he justified he also glorified” (Romans 8:29–30). Dort reads it as a chain from which no link drops anyone, so that every one justified is glorified: God decreed to give the elect “true faith, justification, and sanctification; and having powerfully preserved them … finally to glorify them” (I, article 7). Augustine could say as much of the elect. What is Dort’s own is that the chain holds every one justified: its Rejection of Errors condemns those who teach that God chooses some to justifying faith whom he does not choose to salvation, “dissolving this golden chain of salvation” (I.2, in Latin, \emph{auream hanc salutis catenam dissolvens}; the translation is this page’s). Dort’s reading is the doctrine itself found in the verse, so everything turns on whether the verse says it, and the step from it is rated \emph{well supported}: no cited reader grants the reading and denies what follows. Wesley’s quarrel is with the reading, and he gives his reason. He takes the verse as the order in which God brings to glory those who are glorified, not as a promise of glory to every one justified: “This proves no more than that all who are glorified, are pardoned and sanctified first” (\emph{Predestination Calmly Considered}, §76), and Paul “does not affirm … that precisely the same number of men are called, justified, and glorified” (\emph{Notes}, on the verse). So the reading is rated \emph{disputed}, and it is where Wesley parts from the Romans route (\texttt{where\allowbreak{}Wesley\allowbreak{}Parts\allowbreak{}Over\allowbreak{}Romans8}). Augustine reads the chain as being about those “called according to His purpose”, the elect, and still holds that others who are truly justified fall away.

The John route reaches the conclusion without Philippians or Romans, and the Philippians route reaches it without 1 John 2:19, though it still needs John 10 for the elect (\texttt{john\_\allowbreak{}strand\_\allowbreak{}establishes}, \texttt{philippians\_\allowbreak{}strand\_\allowbreak{}establishes}). The Romans route, like the Philippians route, needs John 10 for the elect, and is not checked alone.

\subsubsection*{Augustine, Trent, the Lutherans and Wesley}
Four positions hold that one truly justified can finally fall away.

\begin{itemize}
  \item \textbf{Augustine and Aquinas}, on Augustine’s ground that some regenerate are not given perseverance. They hold that the elect persevere as firmly as Westminster does. They part from it over 1 John 2:19: whether those who left never had true faith, or had it and were not elect (\texttt{where\allowbreak{}Augustine\allowbreak{}Parts\allowbreak{}From\allowbreak{}Johns\allowbreak{}Strand}).
  \item \textbf{Trent}, from the condition in its chapter 13. Its canon 23 condemns anyone who says “that a man once justified can sin no more, nor lose grace, and that therefore he that falls and sins was never truly justified”. The canon is about losing grace; that it can be lost finally is what the condition says.
  \item \textbf{The Lutherans}, from the Formula of Concord. They hold that the elect persevere, citing John 10:28, and that some in whom God began the good work fall away by their own fault: it “is not as though God were unwilling to grant grace for perseverance to those in whom He has begun the good work” (Solid Declaration XI.42).
  \item \textbf{Wesley}, from the condition he reads in John 10 and from three texts he takes of true believers who fell. “The sum of all is this. If the Scriptures are true, those who are holy or righteous in the judgment of God himself … may nevertheless so fall from God, as to perish everlastingly” (\emph{Serious Thoughts upon the Perseverance of the Saints}, §30).
  \item Hebrews 6:4–6: “it is impossible, in the case of those who have once been enlightened, who have tasted the heavenly gift, and have shared in the Holy Spirit, and have tasted the goodness of the word of God and the powers of the age to come, and then have fallen away, to restore them again to repentance”. Calvin takes “tasted” at its word: God grants “the reprobate also some taste of his grace”, a knowledge “which afterwards vanishes away”, while “God indeed favours none but the elect alone with the Spirit of regeneration”.
  \item Hebrews 10:29: the one “who has trampled underfoot the Son of God, and has profaned the blood of the covenant by which he was sanctified, and has outraged the Spirit of grace”. Wesley: “It is undeniably plain, 1. That the person mentioned here, was once sanctified by the blood of the covenant” (§27). John Owen, the Puritan, gives two answers. If “he” is the apostate, “It is not real or internal sanctification that is here intended, but it is a separation and dedication unto God”, the outward setting apart of those who “by baptism, and confession of faith in the church of Christ, were separated from all others”. But he holds that “the design of the apostle in the context leads plainly to another application”: “It is Christ himself that is spoken of, who was sanctified and dedicated unto God to be an eternal high priest”. Either way the verse does not say the apostate was inwardly made holy.
  \item 2 Peter 2:20–22, those who “escaped the defilements of the world through the knowledge of our Lord” and returned, as “the dog returns to its own vomit”. Wesley calls their knowledge “an inward, experimental knowledge” (§21), that is, known by experience. Calvin holds that “The faithful also do indeed sin; but as they allow not dominion to sin, they do not fall away from the grace of God”, so those who return to the mire were not among the faithful.
\end{itemize}
  If Hebrews 6, Hebrews 10 or 2 Peter 2 does describe true believers who were   lost, it follows that a true believer can be lost, and no one cited here grants   the first and denies the second; those steps are rated \emph{well supported}. The   disagreement is over what the passages describe, and each of Wesley’s readings   is rated \emph{disputed}, against Calvin on Hebrews 6 and 2 Peter 2 and Owen on   Hebrews 10.

Trent, Wesley and the Lutherans each make the promise hold unless the believer turns away: Trent “unless men be themselves wanting to His grace”, Wesley “if ye do these things, ye shall never fall”, and the Formula of Concord that those who fall away do so by their own fault. Taken as a condition that can fail, that simply is the claim that the justified can finally fall away. Taken in the sense of Dort’s Fifth Head, as a condition God keeps in his own sheep (“God is faithful, who having conferred grace, mercifully confirms and powerfully preserves them therein, even to the end”, V, article 3), it is something Dort holds. Both senses are checked against Wesley’s position, because he has other grounds besides the condition, his three texts, and his conclusion does not stand or fall with how the condition is taken (\texttt{what\allowbreak{}The\allowbreak{}Condition\allowbreak{}Means}). Wesley parts from Dort either way. In Dort’s sense of the condition he parts at any of his three readings; the check names Calvin’s reading of Hebrews 6. In his own sense he parts at the step from John 10, the step he himself disputes.

\subsubsection*{The Remonstrants left it open}
The Remonstrants, the Dutch followers of Arminius whose Remonstrance of 1610 the Synod of Dort answered, did not deny perseverance. Their fifth article says that Christ, “if only they are ready for the conflict, and desire his help, and are not inactive, keeps them from falling”, so that they cannot be “plucked out of Christ’s hands” (John 10:28). Then it stops: “But whether they are capable, through negligence, of forsaking again the first beginnings of their life in Christ … that must be more particularly determined out of the Holy Scripture, before we ourselves can teach it with the full persuasion of our minds.” Arminius wrote, “I never taught that a true believer can either totally or finally fall away from the faith, and perish; yet I will not conceal, that there are passages of Scripture which seem to me to wear this aspect” (\emph{Declaration of Sentiments}). What the article concludes from John 10:28 leaves the question open (\texttt{remonstrance\_\allowbreak{}leaves\_\allowbreak{}perseverance\_\allowbreak{}open}).

\subsubsection*{The view from eternity}
Boethius, in the sixth century, asked whether God’s foreknowledge leaves us free. He answered that God does not foresee the future as we guess at it, but sees all of time at once: “Now, eternity is the possession of endless life whole and perfect at a single moment” (\emph{Consolation of Philosophy} V, prose 6). Knowing a thing does not make it happen: “there are two necessities—one simple, as that men are necessarily mortal; the other conditioned, as that, if you know that someone is walking, he must necessarily be walking”. Things God foresees, “viewed, however, in themselves, … desist not from the absolute freedom naturally theirs”. Aquinas carried the answer from foreknowledge to predestination: “all things that are in time are present to God from eternity” (\emph{Summa} I q.14 a.13), and “it is possible for one who is predestinated considered in himself to die in mortal sin; yet it is not possible, supposed … that he is predestinated” (I q.23 a.6, ad 2).

So “could he have fallen?” can be asked of the believer considered in himself, apart from God’s decree, and of the believer as he is, under it. Considered in himself, nothing in the one who perseveres makes his perseverance necessary. Dort grants that and says more: falling, “with respect to themselves is not only possible, but would undoubtedly happen; but with respect to God, it is utterly impossible, since his counsel can not be changed” (V, article 8). Aquinas means less than Dort’s “would”, only that it is possible. The claim encoded here is the one both grant: considered in himself, the believer could fall. The other question is whether one who is truly justified can at the last fall away and be lost: by his own fault, as Trent, the Lutherans and Wesley hold, or because God does not give him the gift of perseverance, as Augustine holds. That is the claim Westminster denies.

The view from eternity, as encoded here, holds that God knows all things as present, that the elect persevere, and that each of them could, considered in himself, have fallen. It says nothing on whether the truly justified can finally fall away. So it can be added to Westminster’s position without contradiction (\texttt{from\_\allowbreak{}eternity\_\allowbreak{}holds\_\allowbreak{}with\_\allowbreak{}westminster}), and to Wesley’s, together with the claim that the elect persevere, which Wesley’s position is not encoded as holding (\texttt{from\_\allowbreak{}eternity\_\allowbreak{}holds\_\allowbreak{}with\_\allowbreak{}wesley}). Granted the view alone, it neither follows nor fails that the truly justified can finally fall away (\texttt{from\_\allowbreak{}eternity\_\allowbreak{}leaves\_\allowbreak{}perseverance\_\allowbreak{}open}). Add Owen’s reading of 1 John 2:19, that all who fall away were never true believers, and it follows that none can (\texttt{from\_\allowbreak{}eternity\_\allowbreak{}with\_\allowbreak{}owens\_\allowbreak{}reading\_\allowbreak{}settles\_\allowbreak{}it}). So what decides the question is not God’s knowledge of time but who the elect are, and that is the question of 1 John 2:19.

Wesley holds, with Boethius, that all time is present to God and that his knowing compels nothing: “Men are as free in believing or not believing as if he did not know it at all” (Sermon 58, §5). He parts from Aquinas over what that knowledge does. For Wesley God knows what we do because we do it, and chooses those he foresees believing; for Aquinas, predestination is not grounded in anything God foresees in us. This chapter does not weigh that question.

\subsubsection*{Falling for a time}
Westminster says that believers “may … fall into grievous sins; and for a time continue therein”, and “come to be deprived of some measure of their graces and comforts” (XVII.3); they do not, it says in XVII.1, fall “totally” from the state of grace. Dort names “the lamentable fall of David, Peter, and other saints” (V, article 4), and says God does not suffer them “to lose the grace of adoption”, that is, of being God’s children, “and forfeit the state of justification” (V, article 6). The Smalcald Articles, Luther’s articles of 1537, read David otherwise. Against “sectarists” who held that “all those who had once received the Spirit … even though they should afterwards sin, would still remain in the faith” (III.iii.42), they say that when holy men “happen to fall into manifest sins, as David into adultery, murder, and blasphemy, … then faith and the Holy Ghost has departed from them” (III.iii.43). Trent holds that “the received grace of Justification is lost, not only by infidelity”, that is, by losing faith, “but also by any other mortal sin whatever” (Session VI, ch. 15); a mortal sin, in Catholic teaching, is a grave sin that destroys the life of grace.

The text the Reformed answer from is David’s own prayer after his sin: “Cast me not away from your presence, and take not your Holy Spirit from me” (Psalm 51:11). Calvin took the prayer itself as evidence: the words “imply that the Spirit had not altogether been taken away from him, however much his gifts had been temporarily obscured”, and grace is “the incorruptible seed of God, which never can perish in any heart where it has been deposited”. So the Reformed and the Lutherans part over what happened to David. The point on which they divide, the \emph{crux}, is Calvin’s reading, that David still had the Spirit when he prayed. The Smalcald Articles deny it, and each reading is rated \emph{disputed}, so neither side wins here (\texttt{where\allowbreak{}They\allowbreak{}Part\allowbreak{}Over\allowbreak{}David}). This is the Reformed answer on falling for a time; the case that no one truly justified finally falls away does not rest on it. The encoding carries the Lutheran claim as “justifying grace can be lost”, without the words “for a time”.

\subsubsection*{Knowing that one will persevere}
“The Spirit himself bears witness with our spirit that we are children of God” (Romans 8:16). Dort says that “true believers for themselves may and do obtain assurance” of their perseverance, “according to the measure of their faith” (V, article 9), and that it springs “from the testimony of the Holy Spirit, witnessing with our spirit, that we are children and heirs of God (Rom. viii. 16)” (V, article 10). Westminster calls it “an infallible assurance of faith”, attained “without extraordinary revelation” (XVIII.2–3). Trent condemns anyone who says he will have “that great gift of perseverance unto the end” “of an absolute and infallible certainty”, “unless he have learned this by special revelation” (canon 16). In its chapter on perseverance (ch. 13) it quotes Paul: “let those who think themselves to stand, take heed lest they fall” (1 Corinthians 10:12).

Calvin answers the verse with “two kinds of assurance”. One “rests on the promises of God, because a pious conscience feels assured that God will never be wanting to it”, and yet, “keeping in mind its own infirmity, casts itself upon God, and with carefulness and anxiety commits itself to him”. That assurance, he says, “is inseparable from faith”. The other “arises from negligence, when men … give themselves no concern, as if they were beyond the reach of danger”, and it is that one Paul warns against. So for Calvin assurance rests on God’s promise, and the care Paul asks for belongs to that same faith, which, “keeping in mind its own infirmity, casts itself upon God”.

Trent, too, commends “a most firm hope in God’s help” (ch. 13), and parts from Calvin over whether anyone can know he is among those who will persevere: “except by special revelation, it cannot be known whom God hath chosen unto Himself” (ch. 12). For Calvin, those who have “the living root of faith” also “carry a testimony of their own adoption firmly fixed in their hearts” (\emph{Catholic Epistles}, on 1 John 2:19). So the question between them is whether the promise, received by faith, lets a believer know his own perseverance. It is tied to the question of 1 John 2:19, though the encoding does not join the two: if some who had true faith fall away, faith alone cannot tell a believer that he will not. Trent’s step from 1 Corinthians 10:12 is rated \emph{disputed}, because Calvin grants the verse and denies that it asks for doubt (\texttt{where\allowbreak{}They\allowbreak{}Part\allowbreak{}Over\allowbreak{}Assurance}).

The Reformed step from Romans 8:16 is rated \emph{well supported}. Wesley grants the verse to every believer, and he allows an assurance of perseverance too: “The full assurance of faith relates to present pardon; the full assurance of hope, to future glory”, the latter being “Divine evidence (wrought in the soul by the same immediate inspiration of the Holy Ghost) of persevering grace, and of eternal glory”, which “you cannot expect, if you abate your diligence” (\emph{Notes}, on Hebrews 6:11). This page reads his “immediate inspiration” as the ordinary witness of the Spirit, since he says the same of the assurance of faith, and so counts him as granting the step. He holds that this assurance is given to some, and lost with diligence.

\subsubsection*{Weighed together}
Weighed as one dispute, Westminster is in a tie with each of the four positions that hold the justified can fall away. One position defeats another when it contradicts something the other holds and its weakest premise or step is rated no lower than the other’s; here each defeats the other, so between them it is a tie. With the Lutherans and Trent the tie is also over falling for a time, and with Trent over knowing that one will persevere. The four do not attack one another. However the ties are settled, the one thing that must be accepted is what the Remonstrants concluded from John 10:28: no power from outside snatches believers out of Christ’s hand (\texttt{only\_\allowbreak{}the\_\allowbreak{}remonstrance\_\allowbreak{}prevails}). Westminster and each of the four can be defended, which means each is accepted on at least one way of settling the ties, and none of the five must be accepted on every way. Every tie rests on a reading or step rated \emph{disputed}, so the verdict moves only if a new, cited argument moves one of those ratings.

\subsubsection*{What is not encoded}
“No one born of God makes a practice of sinning, for God’s seed abides in him” (1 John 3:9). Calvin cites it against the opinion that the elect may lose the Spirit altogether; the Smalcald Articles cite it for the opposite, that where sin has its way “the Holy Ghost and faith are [certainly] not present”. No step from it is encoded yet. Nor are the readings that take Hebrews 6 and 10 to describe true believers but deny that they are lost forever, reading the warnings as warnings or as hypothetical; no one holding them has yet been checked at a verified page.

Dort itself grants that believers, left to themselves, “could not persevere in a state of grace” (V, article 3), and that their falling “with respect to themselves is not only possible, but would undoubtedly happen” (V, article 8). So Dort and Wesley agree about what believers could do if left to themselves. They part over whether God keeps them, and that is what the claims here are about.

\subsubsection*{The positions weighed}
Six parties are heard: the Remonstrance, with what it concludes from John 10:28; Westminster and Dort; Augustine and Aquinas; Trent; the Lutherans; and Wesley. The view from eternity is not among them: as encoded, it says nothing on whether the truly justified can finally fall away, so it attacks none of them and none attacks it. One position defeats another when it contradicts something the other holds and its weakest premise or step is rated no lower than the other’s. Who defeats whom is computed from their premises and checked by the kernel.

\subsubsection*{How this is checked}
\paragraph*{The conclusion}
\leavevmode\par\nobreak
\par\medskip\noindent\textbf{\texttt{perseverance}}\par\nobreak\noindent
\textbf{The perseverance of the saints}, as two claims: the elect persevere, and no one truly justified can finally fall away. The first, Augustine, Aquinas and the Lutherans hold as firmly as Westminster. The second is where they part.

\begin{derivation}
  \item[] \(\pvar{TestimonyArguments.SolaGratia}{41} \land \lnot \pvar{TestimonyArguments.SolaGratia}{43}\)
\end{derivation}
\paragraph*{The steps}
\leavevmode\par\nobreak
\par\medskip\noindent\textbf{\texttt{john10\allowbreak{}Step}}\par\nobreak\noindent
John 10:27–29: those the Father gave to Christ “will never perish”. So the elect persevere. Rated \texttt{disputed}: Wesley grants the text and reads its promise as conditional, quoting 2 Peter 1:10.

\begin{derivation}
  \item[] \(\pvar{TestimonyArguments.SolaGratia}{32} \rightarrow \pvar{TestimonyArguments.SolaGratia}{41}\)\gloss{\emph{would give:} Those God has chosen to save cannot finally fall away from grace}
\end{derivation}
\par\medskip\noindent\textbf{\texttt{bridge\allowbreak{}Step}}\par\nobreak\noindent
1 John 2:19, “if they had been of us, they would have continued with us”, read with Owen of every apostate, that all who fall away were never true believers: so everyone truly justified is among the elect. Rated \texttt{well\allowbreak{}Supported}; the dissent is from the reading, which Augustine denies.

\begin{derivation}
  \item[] \((\pvar{TestimonyArguments.SolaGratia}{35} \land \pvar{TestimonyArguments.SolaGratia}{55}) \rightarrow \pvar{TestimonyArguments.SolaGratia}{42}\)\gloss{\emph{would give:} Everyone who is truly justified is among the elect}
\end{derivation}
\par\medskip\noindent\textbf{\texttt{join\allowbreak{}Step}}\par\nobreak\noindent
If the elect persevere and the truly justified are elect, no one truly justified can finally fall away. Rated \texttt{consensus}: it follows from what the words mean.

\begin{derivation}
  \item[] \((\pvar{TestimonyArguments.SolaGratia}{41} \land \pvar{TestimonyArguments.SolaGratia}{42}) \rightarrow \lnot \pvar{TestimonyArguments.SolaGratia}{43}\)\gloss{\emph{would deny:} One who is truly justified can fall away from grace finally, and perish}
\end{derivation}
\par\medskip\noindent\textbf{\texttt{philippians\allowbreak{}Step}}\par\nobreak\noindent
Philippians 1:6: he who began a good work will bring it to completion. So no one truly justified can finally fall away. Rated \texttt{disputed}: Trent quotes the verse with “unless men be themselves wanting to His grace”.

\begin{derivation}
  \item[] \(\pvar{TestimonyArguments.SolaGratia}{33} \rightarrow \lnot \pvar{TestimonyArguments.SolaGratia}{43}\)\gloss{\emph{would deny:} One who is truly justified can fall away from grace finally, and perish}
\end{derivation}
\par\medskip\noindent\textbf{\texttt{conditional\allowbreak{}Step}}\par\nobreak\noindent
If God’s promise holds unless men are wanting to his grace, one truly justified can finally fall away. Rated \texttt{well\allowbreak{}Supported}.

\begin{derivation}
  \item[] \(\pvar{TestimonyArguments.SolaGratia}{38} \rightarrow \pvar{TestimonyArguments.SolaGratia}{43}\)\gloss{\emph{would give:} One who is truly justified can fall away from grace finally, and perish}
\end{derivation}
\par\medskip\noindent\textbf{\texttt{hebrews\allowbreak{}Step}}\par\nobreak\noindent
Hebrews 6:4–6, read as describing true believers who fall away: so one truly justified can finally fall away. Rated \texttt{well\allowbreak{}Supported}; the dissent is from the reading.

\begin{derivation}
  \item[] \((\pvar{TestimonyArguments.SolaGratia}{34} \land \pvar{TestimonyArguments.SolaGratia}{40}) \rightarrow \pvar{TestimonyArguments.SolaGratia}{43}\)\gloss{\emph{would give:} One who is truly justified can fall away from grace finally, and perish}
\end{derivation}
\par\medskip\noindent\textbf{\texttt{augustine\allowbreak{}Step}}\par\nobreak\noindent
If some whom God has regenerated are not given perseverance, one truly justified can finally fall away. Rated \texttt{well\allowbreak{}Supported}.

\begin{derivation}
  \item[] \(\pvar{TestimonyArguments.SolaGratia}{37} \rightarrow \pvar{TestimonyArguments.SolaGratia}{43}\)\gloss{\emph{would give:} One who is truly justified can fall away from grace finally, and perish}
\end{derivation}
\par\medskip\noindent\textbf{\texttt{remonstrance\allowbreak{}Step}}\par\nobreak\noindent
John 10:28: no one will snatch them out of Christ’s hand. So no power from outside snatches believers out of it. Rated \texttt{consensus}.

\begin{derivation}
  \item[] \(\pvar{TestimonyArguments.SolaGratia}{32} \rightarrow \pvar{TestimonyArguments.SolaGratia}{36}\)\gloss{\emph{would give:} No power from outside can snatch true believers out of Christ’s hand}
\end{derivation}
\par\medskip\noindent\textbf{\texttt{romans8\allowbreak{}Step}}\par\nobreak\noindent
Romans 8:30, “those whom he justified he also glorified”, read with Dort as a chain from which no one drops: so no one truly justified can finally fall away. Rated \texttt{well\allowbreak{}Supported}; the dissent is from the reading, which Wesley denies.

\begin{derivation}
  \item[] \((\pvar{TestimonyArguments.SolaGratia}{44} \land \pvar{TestimonyArguments.SolaGratia}{54}) \rightarrow \lnot \pvar{TestimonyArguments.SolaGratia}{43}\)\gloss{\emph{would deny:} One who is truly justified can fall away from grace finally, and perish}
\end{derivation}
\par\medskip\noindent\textbf{\texttt{timeless\allowbreak{}Step}}\par\nobreak\noindent
God knows all things as present to his eternity, and his knowing imposes no necessity: so the one who perseveres could, considered in himself, have fallen away. Rated \texttt{well\allowbreak{}Supported}.

\begin{derivation}
  \item[] \(\pvar{TestimonyArguments.SolaGratia}{57} \rightarrow \pvar{TestimonyArguments.SolaGratia}{56}\)\gloss{\emph{would give:} Considered in himself, apart from God’s decree, the one who perseveres could have fallen}
\end{derivation}
\par\medskip\noindent\textbf{\texttt{hebrews10\allowbreak{}Step}}\par\nobreak\noindent
Hebrews 10:29, read as describing one truly sanctified by Christ’s blood who falls away: so one truly justified can finally fall away. Rated \texttt{well\allowbreak{}Supported}; the dissent is from the reading.

\begin{derivation}
  \item[] \((\pvar{TestimonyArguments.SolaGratia}{45} \land \pvar{TestimonyArguments.SolaGratia}{50}) \rightarrow \pvar{TestimonyArguments.SolaGratia}{43}\)\gloss{\emph{would give:} One who is truly justified can fall away from grace finally, and perish}
\end{derivation}
\par\medskip\noindent\textbf{\texttt{second\allowbreak{}Peter\allowbreak{}Step}}\par\nobreak\noindent
2 Peter 2:20–22, read as describing true believers who fall back into the world: so one truly justified can finally fall away. Rated \texttt{well\allowbreak{}Supported}; the dissent is from the reading.

\begin{derivation}
  \item[] \((\pvar{TestimonyArguments.SolaGratia}{46} \land \pvar{TestimonyArguments.SolaGratia}{51}) \rightarrow \pvar{TestimonyArguments.SolaGratia}{43}\)\gloss{\emph{would give:} One who is truly justified can fall away from grace finally, and perish}
\end{derivation}
\par\medskip\noindent\textbf{\texttt{david\allowbreak{}Step}}\par\nobreak\noindent
If David lost the Spirit in his sin, justifying grace can be lost. Rated \texttt{well\allowbreak{}Supported}; Calvin denies the ground.

\begin{derivation}
  \item[] \(\pvar{TestimonyArguments.SolaGratia}{52} \rightarrow \pvar{TestimonyArguments.SolaGratia}{20}\)\gloss{\emph{would give:} Justifying grace, once given, can be lost}
\end{derivation}
\par\medskip\noindent\textbf{\texttt{psalm\allowbreak{}Step}}\par\nobreak\noindent
Psalm 51:11, read as Calvin reads it, that David prayed as one who had not lost the Spirit: so justifying grace, once given, is never lost. Rated \texttt{well\allowbreak{}Supported}.

\begin{derivation}
  \item[] \((\pvar{TestimonyArguments.SolaGratia}{47} \land \lnot \pvar{TestimonyArguments.SolaGratia}{52}) \rightarrow \lnot \pvar{TestimonyArguments.SolaGratia}{20}\)\gloss{\emph{would deny:} Justifying grace, once given, can be lost}
\end{derivation}
\par\medskip\noindent\textbf{\texttt{assurance\allowbreak{}Step}}\par\nobreak\noindent
Romans 8:16: the Spirit “bears witness with our spirit that we are children of God”. So a believer can know, without a special revelation, that he will persevere. Rated \texttt{well\allowbreak{}Supported}.

\begin{derivation}
  \item[] \(\pvar{TestimonyArguments.SolaGratia}{48} \rightarrow \pvar{TestimonyArguments.SolaGratia}{53}\)\gloss{\emph{would give:} A believer can know, without a special revelation, that he will persevere}
\end{derivation}
\par\medskip\noindent\textbf{\texttt{caution\allowbreak{}Step}}\par\nobreak\noindent
1 Corinthians 10:12: “let anyone who thinks that he stands take heed lest he fall”. So no believer can know, without a special revelation, that he will persevere. Rated \texttt{disputed}: Calvin grants the verse and denies that it means doubt.

\begin{derivation}
  \item[] \(\pvar{TestimonyArguments.SolaGratia}{49} \rightarrow \lnot \pvar{TestimonyArguments.SolaGratia}{53}\)\gloss{\emph{would deny:} A believer can know, without a special revelation, that he will persevere}
\end{derivation}
\paragraph*{The lines}
\leavevmode\par\nobreak
\par\medskip\noindent\textbf{\texttt{john10\allowbreak{}Line}} --- John 10:27–29 (“no one will snatch them out of my hand”)\par\nobreak\noindent
\textbf{John 10:27–29}: the elect persevere.

\begin{derivation}
  \item[(1)] \(\pvar{TestimonyArguments.SolaGratia}{32}\)\gloss{John 10:27–29 — they will never perish, and no one will snatch them out of my hand}
  \item[(2)] \(\pvar{TestimonyArguments.SolaGratia}{32} \rightarrow \pvar{TestimonyArguments.SolaGratia}{41}\)
  \item[\(\vdash\)] \(\pvar{TestimonyArguments.SolaGratia}{41}\)\gloss{\emph{at issue:} Those God has chosen to save cannot finally fall away from grace}
\end{derivation}
\par\medskip\noindent\textbf{\texttt{bridge\allowbreak{}Line}} --- 1 John 2:19 (“they went out from us, but they were not of us”)\par\nobreak\noindent
\textbf{1 John 2:19}, read of every apostate: the truly justified are elect.

\begin{derivation}
  \item[(1)] \(\pvar{TestimonyArguments.SolaGratia}{35}\)\gloss{1 John 2:19 — they went out from us, but they were not of us}
  \item[(2)] \(\pvar{TestimonyArguments.SolaGratia}{55}\)\gloss{1 John 2:19, read of every apostate: all who fall away were never true believers}
  \item[(3)] \((\pvar{TestimonyArguments.SolaGratia}{35} \land \pvar{TestimonyArguments.SolaGratia}{55}) \rightarrow \pvar{TestimonyArguments.SolaGratia}{42}\)
  \item[\(\vdash\)] \(\pvar{TestimonyArguments.SolaGratia}{42}\)\gloss{\emph{at issue:} Everyone who is truly justified is among the elect}
\end{derivation}
\par\medskip\noindent\textbf{\texttt{philippians\allowbreak{}Line}} --- Philippians 1:6 (“he who began a good work in you will bring it to completion”)\par\nobreak\noindent
\textbf{Philippians 1:6}: no one truly justified can finally fall away.

\begin{derivation}
  \item[(1)] \(\pvar{TestimonyArguments.SolaGratia}{33}\)\gloss{Philippians 1:6 — he who began a good work in you will bring it to completion}
  \item[(2)] \(\pvar{TestimonyArguments.SolaGratia}{33} \rightarrow \lnot \pvar{TestimonyArguments.SolaGratia}{43}\)
  \item[\(\vdash\)] \(\lnot \pvar{TestimonyArguments.SolaGratia}{43}\)\gloss{\emph{at issue, the denial of:} One who is truly justified can fall away from grace finally, and perish}
\end{derivation}
\par\medskip\noindent\textbf{\texttt{conditional\allowbreak{}Line}} --- The promise, kept “unless men be themselves wanting to His grace”\par\nobreak\noindent
\textbf{The condition}: the promise holds unless men are wanting to God’s grace.

\begin{derivation}
  \item[(1)] \(\pvar{TestimonyArguments.SolaGratia}{38}\)\gloss{God’s promise to keep those he has begun to save holds unless they fall away by their own fault}
  \item[(2)] \(\pvar{TestimonyArguments.SolaGratia}{38} \rightarrow \pvar{TestimonyArguments.SolaGratia}{43}\)
  \item[\(\vdash\)] \(\pvar{TestimonyArguments.SolaGratia}{43}\)\gloss{\emph{at issue:} One who is truly justified can fall away from grace finally, and perish}
\end{derivation}
\par\medskip\noindent\textbf{\texttt{hebrews\allowbreak{}Line}} --- Hebrews 6:4–6, as describing true believers who fall away\par\nobreak\noindent
\textbf{Hebrews 6:4–6}, read as describing true believers who fall away.

\begin{derivation}
  \item[(1)] \(\pvar{TestimonyArguments.SolaGratia}{34}\)\gloss{Hebrews 6:4–6 — it is impossible to restore again to repentance those once enlightened who have fallen away}
  \item[(2)] \(\pvar{TestimonyArguments.SolaGratia}{40}\)\gloss{Hebrews 6:4–6 describes true believers who actually fall away}
  \item[(3)] \((\pvar{TestimonyArguments.SolaGratia}{34} \land \pvar{TestimonyArguments.SolaGratia}{40}) \rightarrow \pvar{TestimonyArguments.SolaGratia}{43}\)
  \item[\(\vdash\)] \(\pvar{TestimonyArguments.SolaGratia}{43}\)\gloss{\emph{at issue:} One who is truly justified can fall away from grace finally, and perish}
\end{derivation}
\par\medskip\noindent\textbf{\texttt{augustine\allowbreak{}Line}} --- Some regenerate are not given perseverance (Augustine)\par\nobreak\noindent
\textbf{Augustine}: some regenerate are not given perseverance.

\begin{derivation}
  \item[(1)] \(\pvar{TestimonyArguments.SolaGratia}{37}\)\gloss{To some whom God has regenerated, and given faith, hope and love, he does not give perseverance}
  \item[(2)] \(\pvar{TestimonyArguments.SolaGratia}{37} \rightarrow \pvar{TestimonyArguments.SolaGratia}{43}\)
  \item[\(\vdash\)] \(\pvar{TestimonyArguments.SolaGratia}{43}\)\gloss{\emph{at issue:} One who is truly justified can fall away from grace finally, and perish}
\end{derivation}
\par\medskip\noindent\textbf{\texttt{remonstrance\allowbreak{}Line}} --- Remonstrance, article 5 (John 10:28)\par\nobreak\noindent
\textbf{The Remonstrance, article 5}: no one plucks believers out of Christ’s hand.

\begin{derivation}
  \item[(1)] \(\pvar{TestimonyArguments.SolaGratia}{32}\)\gloss{John 10:27–29 — they will never perish, and no one will snatch them out of my hand}
  \item[(2)] \(\pvar{TestimonyArguments.SolaGratia}{32} \rightarrow \pvar{TestimonyArguments.SolaGratia}{36}\)
  \item[\(\vdash\)] \(\pvar{TestimonyArguments.SolaGratia}{36}\)\gloss{\emph{at issue:} No power from outside can snatch true believers out of Christ’s hand}
\end{derivation}
\par\medskip\noindent\textbf{\texttt{romans8\allowbreak{}Line}} --- Romans 8:29–30 (“those whom he justified he also glorified”)\par\nobreak\noindent
\textbf{Romans 8:29–30}, read as an unbroken chain: no one truly justified can finally fall away.

\begin{derivation}
  \item[(1)] \(\pvar{TestimonyArguments.SolaGratia}{44}\)\gloss{Romans 8:29–30 — those whom he justified he also glorified}
  \item[(2)] \(\pvar{TestimonyArguments.SolaGratia}{54}\)\gloss{Romans 8:29–30 is a chain from which no one drops: every one justified is glorified}
  \item[(3)] \((\pvar{TestimonyArguments.SolaGratia}{44} \land \pvar{TestimonyArguments.SolaGratia}{54}) \rightarrow \lnot \pvar{TestimonyArguments.SolaGratia}{43}\)
  \item[\(\vdash\)] \(\lnot \pvar{TestimonyArguments.SolaGratia}{43}\)\gloss{\emph{at issue, the denial of:} One who is truly justified can fall away from grace finally, and perish}
\end{derivation}
\par\medskip\noindent\textbf{\texttt{hebrews10\allowbreak{}Line}} --- Hebrews 10:29, as describing one truly sanctified who falls away\par\nobreak\noindent
\textbf{Hebrews 10:26–29}, read as describing one truly sanctified who falls.

\begin{derivation}
  \item[(1)] \(\pvar{TestimonyArguments.SolaGratia}{45}\)\gloss{Hebrews 10:26–29 — the one who has profaned the blood of the covenant by which he was sanctified}
  \item[(2)] \(\pvar{TestimonyArguments.SolaGratia}{50}\)\gloss{Hebrews 10:29 describes one truly sanctified by Christ’s blood who falls away}
  \item[(3)] \((\pvar{TestimonyArguments.SolaGratia}{45} \land \pvar{TestimonyArguments.SolaGratia}{50}) \rightarrow \pvar{TestimonyArguments.SolaGratia}{43}\)
  \item[\(\vdash\)] \(\pvar{TestimonyArguments.SolaGratia}{43}\)\gloss{\emph{at issue:} One who is truly justified can fall away from grace finally, and perish}
\end{derivation}
\par\medskip\noindent\textbf{\texttt{second\allowbreak{}Peter\allowbreak{}Line}} --- 2 Peter 2:20–22, as describing true believers who fall back\par\nobreak\noindent
\textbf{2 Peter 2:20–22}, read as describing true believers who fall back.

\begin{derivation}
  \item[(1)] \(\pvar{TestimonyArguments.SolaGratia}{46}\)\gloss{2 Peter 2:20–22 — having escaped the defilements of the world, they are again entangled in them}
  \item[(2)] \(\pvar{TestimonyArguments.SolaGratia}{51}\)\gloss{2 Peter 2:20–22 describes true believers who fall back into the world}
  \item[(3)] \((\pvar{TestimonyArguments.SolaGratia}{46} \land \pvar{TestimonyArguments.SolaGratia}{51}) \rightarrow \pvar{TestimonyArguments.SolaGratia}{43}\)
  \item[\(\vdash\)] \(\pvar{TestimonyArguments.SolaGratia}{43}\)\gloss{\emph{at issue:} One who is truly justified can fall away from grace finally, and perish}
\end{derivation}
\par\medskip\noindent\textbf{\texttt{david\allowbreak{}Line}} --- David’s fall, as losing the Spirit (Smalcald Articles III.iii.43)\par\nobreak\noindent
\textbf{David’s fall}, read with the Smalcald Articles: he lost the Spirit.

\begin{derivation}
  \item[(1)] \(\pvar{TestimonyArguments.SolaGratia}{52}\)\gloss{In his sin David lost the Holy Spirit and his faith}
  \item[(2)] \(\pvar{TestimonyArguments.SolaGratia}{52} \rightarrow \pvar{TestimonyArguments.SolaGratia}{20}\)
  \item[\(\vdash\)] \(\pvar{TestimonyArguments.SolaGratia}{20}\)\gloss{\emph{at issue:} Justifying grace, once given, can be lost}
\end{derivation}
\par\medskip\noindent\textbf{\texttt{psalm\allowbreak{}Line}} --- Psalm 51:11 (“take not your Holy Spirit from me”), as Calvin reads it\par\nobreak\noindent
\textbf{Psalm 51:11}, read with Calvin: David had not lost the Spirit.

\begin{derivation}
  \item[(1)] \(\pvar{TestimonyArguments.SolaGratia}{47}\)\gloss{Psalm 51:11 — take not your Holy Spirit from me}
  \item[(2)] \(\lnot \pvar{TestimonyArguments.SolaGratia}{52}\)\gloss{\emph{denied:} In his sin David lost the Holy Spirit and his faith}
  \item[(3)] \((\pvar{TestimonyArguments.SolaGratia}{47} \land \lnot \pvar{TestimonyArguments.SolaGratia}{52}) \rightarrow \lnot \pvar{TestimonyArguments.SolaGratia}{20}\)
  \item[\(\vdash\)] \(\lnot \pvar{TestimonyArguments.SolaGratia}{20}\)\gloss{\emph{at issue, the denial of:} Justifying grace, once given, can be lost}
\end{derivation}
\par\medskip\noindent\textbf{\texttt{assurance\allowbreak{}Line}} --- Romans 8:16 (“the Spirit himself bears witness with our spirit”)\par\nobreak\noindent
\textbf{Romans 8:16}: a believer can know that he will persevere.

\begin{derivation}
  \item[(1)] \(\pvar{TestimonyArguments.SolaGratia}{48}\)\gloss{Romans 8:16 — the Spirit himself bears witness with our spirit}
  \item[(2)] \(\pvar{TestimonyArguments.SolaGratia}{48} \rightarrow \pvar{TestimonyArguments.SolaGratia}{53}\)
  \item[\(\vdash\)] \(\pvar{TestimonyArguments.SolaGratia}{53}\)\gloss{\emph{at issue:} A believer can know, without a special revelation, that he will persevere}
\end{derivation}
\par\medskip\noindent\textbf{\texttt{caution\allowbreak{}Line}} --- 1 Corinthians 10:12 (“let anyone who thinks that he stands take heed”)\par\nobreak\noindent
\textbf{1 Corinthians 10:12}: no believer can know that he will persevere.

\begin{derivation}
  \item[(1)] \(\pvar{TestimonyArguments.SolaGratia}{49}\)\gloss{1 Corinthians 10:12 — let anyone who thinks that he stands take heed lest he fall}
  \item[(2)] \(\pvar{TestimonyArguments.SolaGratia}{49} \rightarrow \lnot \pvar{TestimonyArguments.SolaGratia}{53}\)
  \item[\(\vdash\)] \(\lnot \pvar{TestimonyArguments.SolaGratia}{53}\)\gloss{\emph{at issue, the denial of:} A believer can know, without a special revelation, that he will persevere}
\end{derivation}
\par\medskip\noindent\textbf{\texttt{timeless\allowbreak{}Line}} --- God knows all things as present, and imposes no necessity (Boethius)\par\nobreak\noindent
\textbf{The view from eternity} (Boethius; Aquinas): God knows all things as present, so the one who perseveres could, considered in himself, have fallen away.

\begin{derivation}
  \item[(1)] \(\pvar{TestimonyArguments.SolaGratia}{57}\)\gloss{God knows all things as present to his eternity, which imposes no necessity on them}
  \item[(2)] \(\pvar{TestimonyArguments.SolaGratia}{57} \rightarrow \pvar{TestimonyArguments.SolaGratia}{56}\)
  \item[\(\vdash\)] \(\pvar{TestimonyArguments.SolaGratia}{56}\)\gloss{\emph{at issue:} Considered in himself, apart from God’s decree, the one who perseveres could have fallen}
\end{derivation}
\paragraph*{The positions}
\leavevmode\par\nobreak
\par\medskip\noindent\textbf{\texttt{remonstrance}} --- Remonstrance (article 5, 1610)\par\nobreak\noindent
\textbf{The Remonstrants of 1610}: what their fifth article concludes from John 10:28, that no one plucks believers out of Christ’s hand. The article puts a condition on Christ’s keeping, “if only they are ready for the conflict, and desire his help, and are not inactive”. That is the condition the dilemma below takes two ways, and the article does not say which way it means, so it is not carried here as a premise. It leaves whether believers can fall away by negligence to be “more particularly determined out of the Holy Scripture”.

\begin{derivation}
  \item[(1)] \(\pvar{TestimonyArguments.SolaGratia}{32}\)\gloss{John 10:27–29 — they will never perish, and no one will snatch them out of my hand}
  \item[(2)] \(\pvar{TestimonyArguments.SolaGratia}{32} \rightarrow \pvar{TestimonyArguments.SolaGratia}{36}\)
  \item[\(\vdash\)] \(\pvar{TestimonyArguments.SolaGratia}{36}\)\gloss{\emph{at issue:} No power from outside can snatch true believers out of Christ’s hand}
  \keepwithnext
\end{derivation}
\emph{Grounded in scripture alone:} John 10:27–29 — they will never perish, and no one will snatch them out of my hand.
\par\medskip\noindent\textbf{\texttt{augustinian}} --- Augustine and Aquinas (On Rebuke and Grace; Summa I q.23)\par\nobreak\noindent
\textbf{Augustine and Aquinas.} The elect persevere; and, on Augustine’s ground, to some whom God has regenerated he does not give perseverance, so one truly justified can finally fall away. They read 1 John 2:19 as Augustine does, of those who were in the faith and were not elect, and deny that every apostate never had true faith.

\begin{derivation}
  \item[(1)] \(\pvar{TestimonyArguments.SolaGratia}{37}\)\gloss{To some whom God has regenerated, and given faith, hope and love, he does not give perseverance}
  \item[(2)] \(\pvar{TestimonyArguments.SolaGratia}{41}\)\gloss{Those God has chosen to save cannot finally fall away from grace}
  \item[(3)] \(\pvar{TestimonyArguments.SolaGratia}{35}\)\gloss{1 John 2:19 — they went out from us, but they were not of us}
  \breakrow{(4)} \(\lnot \pvar{TestimonyArguments.SolaGratia}{55}\)\gloss{\emph{denied:} 1 John 2:19, read of every apostate: all who fall away were never true believers}
  \item[(5)] \(\pvar{TestimonyArguments.SolaGratia}{37} \rightarrow \pvar{TestimonyArguments.SolaGratia}{43}\)\gloss{\emph{would give:} One who is truly justified can fall away from grace finally, and perish}
  \item[\(\vdash\)] \(\pvar{TestimonyArguments.SolaGratia}{41} \land \pvar{TestimonyArguments.SolaGratia}{43}\)\gloss{\emph{at issue:} \(\pvar{TestimonyArguments.SolaGratia}{41}\) Those God has chosen to save cannot finally fall away from grace; \(\pvar{TestimonyArguments.SolaGratia}{43}\) One who is truly justified can fall away from grace finally, and perish}
  \keepwithnext
\end{derivation}
\emph{Grounded in scripture alone:} 1 John 2:19 — they went out from us, but they were not of us.
\par\medskip\noindent\textbf{\texttt{trent\allowbreak{}On\allowbreak{}Perseverance}} --- Trent (Session VI, chapters 13 and 15, canons 16 and 23)\par\nobreak\noindent
\textbf{Trent}, from the sixth session (1547): Philippians 1:6, as its thirteenth chapter reads it, with the condition “unless men be themselves wanting to His grace”; 1 Corinthians 10:12, which the same chapter quotes, against certainty of persevering (canon 16); and its fifteenth chapter, that grace “is lost … by any other mortal sin whatever”.

\begin{derivation}
  \item[(1)] \(\pvar{TestimonyArguments.SolaGratia}{38}\)\gloss{God’s promise to keep those he has begun to save holds unless they fall away by their own fault}
  \item[(2)] \(\pvar{TestimonyArguments.SolaGratia}{49}\)\gloss{1 Corinthians 10:12 — let anyone who thinks that he stands take heed lest he fall}
  \item[(3)] \(\pvar{TestimonyArguments.SolaGratia}{33}\)\gloss{Philippians 1:6 — he who began a good work in you will bring it to completion}
  \breakrow{(4)} \(\pvar{TestimonyArguments.SolaGratia}{20}\)\gloss{Justifying grace, once given, can be lost}
  \breakrow{(5)} \(\pvar{TestimonyArguments.SolaGratia}{38} \rightarrow \pvar{TestimonyArguments.SolaGratia}{43}\)
  \item[(6)] \(\pvar{TestimonyArguments.SolaGratia}{49} \rightarrow \lnot \pvar{TestimonyArguments.SolaGratia}{53}\)\gloss{\emph{would deny:} A believer can know, without a special revelation, that he will persevere}
  \item[\(\vdash\)] \(\pvar{TestimonyArguments.SolaGratia}{43}\)\gloss{\emph{at issue:} One who is truly justified can fall away from grace finally, and perish}
  \keepwithnext
\end{derivation}
\emph{Grounded in scripture alone:} 1 Corinthians 10:12 — let anyone who thinks that he stands take heed lest he fall, Philippians 1:6 — he who began a good work in you will bring it to completion.
\par\medskip\noindent\textbf{\texttt{lutheran\allowbreak{}On\allowbreak{}Perseverance}} --- Lutheran (the Formula of Concord, XI; the Smalcald Articles)\par\nobreak\noindent
\textbf{The Lutherans}, from the Formula of Concord (1577). The elect persevere, “from which no one can pluck us, John 10:28” (Solid Declaration XI.45); those in whom God began the good work can fall away by their own fault (XI.42); and, from the Smalcald Articles (1537), when holy men fall into manifest sins, as David did, “faith and the Holy Ghost has departed from them” (III.iii.43).

\begin{derivation}
  \item[(1)] \(\pvar{TestimonyArguments.SolaGratia}{38}\)\gloss{God’s promise to keep those he has begun to save holds unless they fall away by their own fault}
  \item[(2)] \(\pvar{TestimonyArguments.SolaGratia}{52}\)\gloss{In his sin David lost the Holy Spirit and his faith}
  \item[(3)] \(\pvar{TestimonyArguments.SolaGratia}{32}\)\gloss{John 10:27–29 — they will never perish, and no one will snatch them out of my hand}
  \breakrow{(4)} \(\pvar{TestimonyArguments.SolaGratia}{41}\)\gloss{Those God has chosen to save cannot finally fall away from grace}
  \breakrow{(5)} \(\pvar{TestimonyArguments.SolaGratia}{38} \rightarrow \pvar{TestimonyArguments.SolaGratia}{43}\)\gloss{\emph{would give:} One who is truly justified can fall away from grace finally, and perish}
  \item[(6)] \(\pvar{TestimonyArguments.SolaGratia}{52} \rightarrow \pvar{TestimonyArguments.SolaGratia}{20}\)\gloss{\emph{would give:} Justifying grace, once given, can be lost}
  \item[\(\vdash\)] \(\pvar{TestimonyArguments.SolaGratia}{41} \land \pvar{TestimonyArguments.SolaGratia}{43}\)\gloss{\emph{at issue:} \(\pvar{TestimonyArguments.SolaGratia}{41}\) Those God has chosen to save cannot finally fall away from grace; \(\pvar{TestimonyArguments.SolaGratia}{43}\) One who is truly justified can fall away from grace finally, and perish}
  \keepwithnext
\end{derivation}
\emph{Grounded in scripture alone:} John 10:27–29 — they will never perish, and no one will snatch them out of my hand.
\par\medskip\noindent\textbf{\texttt{wesley\allowbreak{}On\allowbreak{}Perseverance}} --- Wesley (Serious Thoughts upon the Perseverance of the Saints)\par\nobreak\noindent
\textbf{Wesley}, from \emph{Serious Thoughts upon the Perseverance of the Saints} (1751). John 10:27–29, with the condition he reads in it, “if ye do these things, ye shall never fall” (2 Peter 1:10); and three texts read as describing true believers who fell: Hebrews 6:4–6, Hebrews 10:26–29 and 2 Peter 2:20–22. The condition is carried as Wesley states it, with no step drawn from it, because what it commits him to is what the dilemma below asks. His conclusion comes from the three texts.

\begin{derivation}
  \item[(1)] \(\pvar{TestimonyArguments.SolaGratia}{34}\)\gloss{Hebrews 6:4–6 — it is impossible to restore again to repentance those once enlightened who have fallen away}
  \item[(2)] \(\pvar{TestimonyArguments.SolaGratia}{40}\)\gloss{Hebrews 6:4–6 describes true believers who actually fall away}
  \item[(3)] \(\pvar{TestimonyArguments.SolaGratia}{45}\)\gloss{Hebrews 10:26–29 — the one who has profaned the blood of the covenant by which he was sanctified}
  \breakrow{(4)} \(\pvar{TestimonyArguments.SolaGratia}{50}\)\gloss{Hebrews 10:29 describes one truly sanctified by Christ’s blood who falls away}
  \breakrow{(5)} \(\pvar{TestimonyArguments.SolaGratia}{46}\)\gloss{2 Peter 2:20–22 — having escaped the defilements of the world, they are again entangled in them}
  \breakrow{(6)} \(\pvar{TestimonyArguments.SolaGratia}{51}\)\gloss{2 Peter 2:20–22 describes true believers who fall back into the world}
  \breakrow{(7)} \(\pvar{TestimonyArguments.SolaGratia}{32}\)\gloss{John 10:27–29 — they will never perish, and no one will snatch them out of my hand}
  \breakrow{(8)} \(\pvar{TestimonyArguments.SolaGratia}{38}\)\gloss{God’s promise to keep those he has begun to save holds unless they fall away by their own fault}
  \breakrow{(9)} \(\pvar{TestimonyArguments.SolaGratia}{44}\)\gloss{Romans 8:29–30 — those whom he justified he also glorified}
  \breakrow{(10)} \(\lnot \pvar{TestimonyArguments.SolaGratia}{54}\)\gloss{\emph{denied:} Romans 8:29–30 is a chain from which no one drops: every one justified is glorified}
  \breakrow{(11)} \((\pvar{TestimonyArguments.SolaGratia}{34} \land \pvar{TestimonyArguments.SolaGratia}{40}) \rightarrow \pvar{TestimonyArguments.SolaGratia}{43}\)
  \breakrow{(12)} \((\pvar{TestimonyArguments.SolaGratia}{45} \land \pvar{TestimonyArguments.SolaGratia}{50}) \rightarrow \pvar{TestimonyArguments.SolaGratia}{43}\)
  \item[(13)] \((\pvar{TestimonyArguments.SolaGratia}{46} \land \pvar{TestimonyArguments.SolaGratia}{51}) \rightarrow \pvar{TestimonyArguments.SolaGratia}{43}\)
  \item[\(\vdash\)] \(\pvar{TestimonyArguments.SolaGratia}{43}\)\gloss{\emph{at issue:} One who is truly justified can fall away from grace finally, and perish}
  \keepwithnext
\end{derivation}
\emph{Grounded in scripture alone:} Hebrews 6:4–6 — it is impossible to restore again to repentance those once enlightened who have fallen away, Hebrews 10:26–29 — the one who has profaned the blood of the covenant by which he was sanctified, 2 Peter 2:20–22 — having escaped the defilements of the world, they are again entangled in them, John 10:27–29 — they will never perish, and no one will snatch them out of my hand, Romans 8:29–30 — those whom he justified he also glorified.
\paragraph*{The position argued for}
\leavevmode\par\nobreak
\par\medskip\noindent\textbf{\texttt{perseverance\allowbreak{}Shared}}\par\nobreak\noindent
The Reformed answers to the texts Wesley reads of true believers who fell. Hebrews 6, with Calvin: those described had only “some taste” of grace. Hebrews 10:29, with Owen: “It is not real or internal sanctification that is here intended”. 2 Peter 2, with Calvin: “The faithful also do indeed sin; but … they do not fall away from the grace of God”.

\begin{derivation}
  \item[(1)] \(\pvar{TestimonyArguments.SolaGratia}{34}\)\gloss{Hebrews 6:4–6 — it is impossible to restore again to repentance those once enlightened who have fallen away}
  \item[(2)] \(\lnot \pvar{TestimonyArguments.SolaGratia}{40}\)\gloss{\emph{denied:} Hebrews 6:4–6 describes true believers who actually fall away}
  \item[(3)] \(\pvar{TestimonyArguments.SolaGratia}{45}\)\gloss{Hebrews 10:26–29 — the one who has profaned the blood of the covenant by which he was sanctified}
  \breakrow{(4)} \(\lnot \pvar{TestimonyArguments.SolaGratia}{50}\)\gloss{\emph{denied:} Hebrews 10:29 describes one truly sanctified by Christ’s blood who falls away}
  \item[(5)] \(\pvar{TestimonyArguments.SolaGratia}{46}\)\gloss{2 Peter 2:20–22 — having escaped the defilements of the world, they are again entangled in them}
  \item[(6)] \(\lnot \pvar{TestimonyArguments.SolaGratia}{51}\)\gloss{\emph{denied:} 2 Peter 2:20–22 describes true believers who fall back into the world}
\end{derivation}
\par\medskip\noindent\textbf{\texttt{perseverance\allowbreak{}Reformed}} --- Reformed (Westminster XVII; Dort, Fifth Head)\par\nobreak\noindent
\textbf{The Reformed position}, as Westminster (1647) and Dort (1619) state it. Three routes to the claim that no one truly justified can finally fall away: John 10:27–29 with 1 John 2:19, joined; Philippians 1:6; and Romans 8:29–30. With the answers to the texts Wesley reads of believers who fell; Psalm 51:11, read with Calvin, against losing grace for a time; and Romans 8:16, for assurance.

\begin{derivation}
  \item[(1)] \(\pvar{TestimonyArguments.SolaGratia}{32}\)\gloss{John 10:27–29 — they will never perish, and no one will snatch them out of my hand}
  \item[(2)] \(\pvar{TestimonyArguments.SolaGratia}{35}\)\gloss{1 John 2:19 — they went out from us, but they were not of us}
  \item[(3)] \(\pvar{TestimonyArguments.SolaGratia}{55}\)\gloss{1 John 2:19, read of every apostate: all who fall away were never true believers}
  \breakrow{(4)} \(\pvar{TestimonyArguments.SolaGratia}{33}\)\gloss{Philippians 1:6 — he who began a good work in you will bring it to completion}
  \breakrow{(5)} \(\pvar{TestimonyArguments.SolaGratia}{44}\)\gloss{Romans 8:29–30 — those whom he justified he also glorified}
  \breakrow{(6)} \(\pvar{TestimonyArguments.SolaGratia}{54}\)\gloss{Romans 8:29–30 is a chain from which no one drops: every one justified is glorified}
  \breakrow{(7)} \(\pvar{TestimonyArguments.SolaGratia}{47}\)\gloss{Psalm 51:11 — take not your Holy Spirit from me}
  \breakrow{(8)} \(\lnot \pvar{TestimonyArguments.SolaGratia}{52}\)\gloss{\emph{denied:} In his sin David lost the Holy Spirit and his faith}
  \breakrow{(9)} \(\pvar{TestimonyArguments.SolaGratia}{48}\)\gloss{Romans 8:16 — the Spirit himself bears witness with our spirit}
  \breakrow{(10)} \(\pvar{TestimonyArguments.SolaGratia}{34}\)\gloss{Hebrews 6:4–6 — it is impossible to restore again to repentance those once enlightened who have fallen away}
  \breakrow{(11)} \(\lnot \pvar{TestimonyArguments.SolaGratia}{40}\)\gloss{\emph{denied:} Hebrews 6:4–6 describes true believers who actually fall away}
  \breakrow{(12)} \(\pvar{TestimonyArguments.SolaGratia}{45}\)\gloss{Hebrews 10:26–29 — the one who has profaned the blood of the covenant by which he was sanctified}
  \breakrow{(13)} \(\lnot \pvar{TestimonyArguments.SolaGratia}{50}\)\gloss{\emph{denied:} Hebrews 10:29 describes one truly sanctified by Christ’s blood who falls away}
  \breakrow{(14)} \(\pvar{TestimonyArguments.SolaGratia}{46}\)\gloss{2 Peter 2:20–22 — having escaped the defilements of the world, they are again entangled in them}
  \breakrow{(15)} \(\lnot \pvar{TestimonyArguments.SolaGratia}{51}\)\gloss{\emph{denied:} 2 Peter 2:20–22 describes true believers who fall back into the world}
  \breakrow{(16)} \(\pvar{TestimonyArguments.SolaGratia}{32} \rightarrow \pvar{TestimonyArguments.SolaGratia}{41}\)\gloss{\emph{would give:} Those God has chosen to save cannot finally fall away from grace}
  \breakrow{(17)} \((\pvar{TestimonyArguments.SolaGratia}{35} \land \pvar{TestimonyArguments.SolaGratia}{55}) \rightarrow \pvar{TestimonyArguments.SolaGratia}{42}\)\gloss{\emph{would give:} Everyone who is truly justified is among the elect}
  \breakrow{(18)} \(\pvar{TestimonyArguments.SolaGratia}{33} \rightarrow \lnot \pvar{TestimonyArguments.SolaGratia}{43}\)\gloss{\emph{would deny:} One who is truly justified can fall away from grace finally, and perish}
  \breakrow{(19)} \((\pvar{TestimonyArguments.SolaGratia}{44} \land \pvar{TestimonyArguments.SolaGratia}{54}) \rightarrow \lnot \pvar{TestimonyArguments.SolaGratia}{43}\)\gloss{\emph{would deny:} One who is truly justified can fall away from grace finally, and perish}
  \breakrow{(20)} \((\pvar{TestimonyArguments.SolaGratia}{47} \land \lnot \pvar{TestimonyArguments.SolaGratia}{52}) \rightarrow \lnot \pvar{TestimonyArguments.SolaGratia}{20}\)\gloss{\emph{would deny:} Justifying grace, once given, can be lost}
  \breakrow{(21)} \(\pvar{TestimonyArguments.SolaGratia}{48} \rightarrow \pvar{TestimonyArguments.SolaGratia}{53}\)\gloss{\emph{would give:} A believer can know, without a special revelation, that he will persevere}
  \item[(22)] \((\pvar{TestimonyArguments.SolaGratia}{41} \land \pvar{TestimonyArguments.SolaGratia}{42}) \rightarrow \lnot \pvar{TestimonyArguments.SolaGratia}{43}\)\gloss{\emph{would deny:} One who is truly justified can fall away from grace finally, and perish}
  \item[\(\vdash\)] \(\pvar{TestimonyArguments.SolaGratia}{41} \land \lnot \pvar{TestimonyArguments.SolaGratia}{43}\)\gloss{\emph{at issue:} \(\pvar{TestimonyArguments.SolaGratia}{41}\) Those God has chosen to save cannot finally fall away from grace; \(\pvar{TestimonyArguments.SolaGratia}{43}\) \emph{denied:} One who is truly justified can fall away from grace finally, and perish}
  \keepwithnext
\end{derivation}
\emph{Grounded in scripture alone:} John 10:27–29 — they will never perish, and no one will snatch them out of my hand, 1 John 2:19 — they went out from us, but they were not of us, Philippians 1:6 — he who began a good work in you will bring it to completion, Romans 8:29–30 — those whom he justified he also glorified, Psalm 51:11 — take not your Holy Spirit from me, Romans 8:16 — the Spirit himself bears witness with our spirit, Hebrews 6:4–6 — it is impossible to restore again to repentance those once enlightened who have fallen away, Hebrews 10:26–29 — the one who has profaned the blood of the covenant by which he was sanctified, 2 Peter 2:20–22 — having escaped the defilements of the world, they are again entangled in them.
\par\medskip\noindent\textbf{\texttt{john\allowbreak{}Strand}} --- Reformed, the John strand alone\par\nobreak\noindent
The John strand alone: John 10:27–29 and 1 John 2:19, joined.

\begin{derivation}
  \item[(1)] \(\pvar{TestimonyArguments.SolaGratia}{32}\)\gloss{John 10:27–29 — they will never perish, and no one will snatch them out of my hand}
  \item[(2)] \(\pvar{TestimonyArguments.SolaGratia}{35}\)\gloss{1 John 2:19 — they went out from us, but they were not of us}
  \item[(3)] \(\pvar{TestimonyArguments.SolaGratia}{55}\)\gloss{1 John 2:19, read of every apostate: all who fall away were never true believers}
  \breakrow{(4)} \(\pvar{TestimonyArguments.SolaGratia}{34}\)\gloss{Hebrews 6:4–6 — it is impossible to restore again to repentance those once enlightened who have fallen away}
  \breakrow{(5)} \(\lnot \pvar{TestimonyArguments.SolaGratia}{40}\)\gloss{\emph{denied:} Hebrews 6:4–6 describes true believers who actually fall away}
  \breakrow{(6)} \(\pvar{TestimonyArguments.SolaGratia}{45}\)\gloss{Hebrews 10:26–29 — the one who has profaned the blood of the covenant by which he was sanctified}
  \breakrow{(7)} \(\lnot \pvar{TestimonyArguments.SolaGratia}{50}\)\gloss{\emph{denied:} Hebrews 10:29 describes one truly sanctified by Christ’s blood who falls away}
  \breakrow{(8)} \(\pvar{TestimonyArguments.SolaGratia}{46}\)\gloss{2 Peter 2:20–22 — having escaped the defilements of the world, they are again entangled in them}
  \breakrow{(9)} \(\lnot \pvar{TestimonyArguments.SolaGratia}{51}\)\gloss{\emph{denied:} 2 Peter 2:20–22 describes true believers who fall back into the world}
  \breakrow{(10)} \(\pvar{TestimonyArguments.SolaGratia}{32} \rightarrow \pvar{TestimonyArguments.SolaGratia}{41}\)\gloss{\emph{would give:} Those God has chosen to save cannot finally fall away from grace}
  \breakrow{(11)} \((\pvar{TestimonyArguments.SolaGratia}{35} \land \pvar{TestimonyArguments.SolaGratia}{55}) \rightarrow \pvar{TestimonyArguments.SolaGratia}{42}\)\gloss{\emph{would give:} Everyone who is truly justified is among the elect}
  \item[(12)] \((\pvar{TestimonyArguments.SolaGratia}{41} \land \pvar{TestimonyArguments.SolaGratia}{42}) \rightarrow \lnot \pvar{TestimonyArguments.SolaGratia}{43}\)\gloss{\emph{would deny:} One who is truly justified can fall away from grace finally, and perish}
  \item[\(\vdash\)] \(\pvar{TestimonyArguments.SolaGratia}{41} \land \lnot \pvar{TestimonyArguments.SolaGratia}{43}\)\gloss{\emph{at issue:} \(\pvar{TestimonyArguments.SolaGratia}{41}\) Those God has chosen to save cannot finally fall away from grace; \(\pvar{TestimonyArguments.SolaGratia}{43}\) \emph{denied:} One who is truly justified can fall away from grace finally, and perish}
  \keepwithnext
\end{derivation}
\emph{Grounded in scripture alone:} John 10:27–29 — they will never perish, and no one will snatch them out of my hand, 1 John 2:19 — they went out from us, but they were not of us, Hebrews 6:4–6 — it is impossible to restore again to repentance those once enlightened who have fallen away, Hebrews 10:26–29 — the one who has profaned the blood of the covenant by which he was sanctified, 2 Peter 2:20–22 — having escaped the defilements of the world, they are again entangled in them.
\par\medskip\noindent\textbf{\texttt{philippians\allowbreak{}Strand}} --- Reformed, the Philippians strand alone\par\nobreak\noindent
The Philippians strand alone, with John 10:27–29 for the elect.

\begin{derivation}
  \item[(1)] \(\pvar{TestimonyArguments.SolaGratia}{32}\)\gloss{John 10:27–29 — they will never perish, and no one will snatch them out of my hand}
  \item[(2)] \(\pvar{TestimonyArguments.SolaGratia}{33}\)\gloss{Philippians 1:6 — he who began a good work in you will bring it to completion}
  \item[(3)] \(\pvar{TestimonyArguments.SolaGratia}{34}\)\gloss{Hebrews 6:4–6 — it is impossible to restore again to repentance those once enlightened who have fallen away}
  \breakrow{(4)} \(\lnot \pvar{TestimonyArguments.SolaGratia}{40}\)\gloss{\emph{denied:} Hebrews 6:4–6 describes true believers who actually fall away}
  \breakrow{(5)} \(\pvar{TestimonyArguments.SolaGratia}{45}\)\gloss{Hebrews 10:26–29 — the one who has profaned the blood of the covenant by which he was sanctified}
  \breakrow{(6)} \(\lnot \pvar{TestimonyArguments.SolaGratia}{50}\)\gloss{\emph{denied:} Hebrews 10:29 describes one truly sanctified by Christ’s blood who falls away}
  \breakrow{(7)} \(\pvar{TestimonyArguments.SolaGratia}{46}\)\gloss{2 Peter 2:20–22 — having escaped the defilements of the world, they are again entangled in them}
  \breakrow{(8)} \(\lnot \pvar{TestimonyArguments.SolaGratia}{51}\)\gloss{\emph{denied:} 2 Peter 2:20–22 describes true believers who fall back into the world}
  \breakrow{(9)} \(\pvar{TestimonyArguments.SolaGratia}{32} \rightarrow \pvar{TestimonyArguments.SolaGratia}{41}\)\gloss{\emph{would give:} Those God has chosen to save cannot finally fall away from grace}
  \item[(10)] \(\pvar{TestimonyArguments.SolaGratia}{33} \rightarrow \lnot \pvar{TestimonyArguments.SolaGratia}{43}\)\gloss{\emph{would deny:} One who is truly justified can fall away from grace finally, and perish}
  \item[\(\vdash\)] \(\pvar{TestimonyArguments.SolaGratia}{41} \land \lnot \pvar{TestimonyArguments.SolaGratia}{43}\)\gloss{\emph{at issue:} \(\pvar{TestimonyArguments.SolaGratia}{41}\) Those God has chosen to save cannot finally fall away from grace; \(\pvar{TestimonyArguments.SolaGratia}{43}\) \emph{denied:} One who is truly justified can fall away from grace finally, and perish}
  \keepwithnext
\end{derivation}
\emph{Grounded in scripture alone:} John 10:27–29 — they will never perish, and no one will snatch them out of my hand, Philippians 1:6 — he who began a good work in you will bring it to completion, Hebrews 6:4–6 — it is impossible to restore again to repentance those once enlightened who have fallen away, Hebrews 10:26–29 — the one who has profaned the blood of the covenant by which he was sanctified, 2 Peter 2:20–22 — having escaped the defilements of the world, they are again entangled in them.
\paragraph*{Readings}
\leavevmode\par\nobreak
\par\medskip\noindent\textbf{\texttt{persevering\allowbreak{}Reading}}\par\nobreak\noindent
The Reformed reading: no one truly justified falls away finally; Hebrews 6 describes those who had only a taste of grace; and no regenerate person is left without perseverance. The promise is not conditional on anything God does not himself keep.

{\small\begin{verbatim}
def perseveringReading : Valuation Claim :=
  fun a =>
    match a with
    | Claim.justifiedCanFinallyFallAway => False
    | Claim.hebrews6DescribesBelieversWhoFall => False
    | Claim.hebrews10DescribesTheSanctifiedWhoFall => False
    | Claim.secondPeter2DescribesBelieversWhoFall => False
    | Claim.someRegenerateNotGivenPerseverance => False
    | Claim.promiseIsConditional => False
    | Claim.davidLostTheSpirit => False
    | Claim.justifyingGraceCanBeLost => False
    | x => True
\end{verbatim}}
\par\medskip\noindent\textbf{\texttt{falling\allowbreak{}Away\allowbreak{}Reading}}\par\nobreak\noindent
A reading on which every rival’s premises hold: the elect persevere, some truly justified, not being elect, finally fall away, and no believer can know without a special revelation that he will persevere. Not a position the rivals share: Wesley, for one, is not encoded as holding that the elect persevere.

{\small\begin{verbatim}
def fallingAwayReading : Valuation Claim :=
  fun a =>
    match a with
    | Claim.justifiedAreElect => False
    | Claim.believersCanBeSureOfPersevering => False
    | Claim.apostatesNeverHadTrueFaith => False
    | Claim.romans8ChainIsUnbroken => False
    | x => True
\end{verbatim}}
\paragraph*{What follows}
\leavevmode\par\nobreak
\par\medskip\noindent\textbf{\texttt{perseverance\allowbreak{}Reformed\_\allowbreak{}establishes}}\par\nobreak\noindent
\textbf{The Reformed premises establish the perseverance of the saints}: the elect persevere, and no one truly justified can finally fall away. That the second follows rests on readings and steps rated \texttt{disputed}: the step from John 10 with the reading of 1 John 2:19, the step from Philippians 1:6, or the reading of Romans 8:29–30.

{\small\begin{verbatim}
theorem perseveranceReformed_establishes : Establishes perseveranceReformed
-- axioms: propext, Quot.sound
\end{verbatim}}
\par\medskip\noindent\textbf{\texttt{perseverance\allowbreak{}Reformed\_\allowbreak{}is\_\allowbreak{}satisfiable}}\par\nobreak\noindent
The Reformed premises can be held without contradiction.

{\small\begin{verbatim}
theorem perseveranceReformed_is_satisfiable : Satisfiable
    perseveranceReformed.premises
-- axioms: propext, Quot.sound
\end{verbatim}}
\par\medskip\noindent\textbf{\texttt{john\_\allowbreak{}strand\_\allowbreak{}establishes}}\par\nobreak\noindent
\textbf{The John strand reaches the conclusion alone.} John 10:27–29 for the elect, 1 John 2:19 for the truly justified being elect, and the two joined.

{\small\begin{verbatim}
theorem john_strand_establishes : Establishes johnStrand
-- axioms: propext, Quot.sound
\end{verbatim}}
\par\medskip\noindent\textbf{\texttt{john\allowbreak{}Strand\_\allowbreak{}is\_\allowbreak{}satisfiable}}\par\nobreak\noindent
The John strand can be held without contradiction.

{\small\begin{verbatim}
theorem johnStrand_is_satisfiable : Satisfiable johnStrand.premises
-- axioms: propext, Quot.sound
\end{verbatim}}
\par\medskip\noindent\textbf{\texttt{philippians\_\allowbreak{}strand\_\allowbreak{}establishes}}\par\nobreak\noindent
\textbf{The Philippians strand reaches the conclusion alone}, without 1 John 2:19: Philippians 1:6 for the truly justified, and John 10:27–29 for the elect.

{\small\begin{verbatim}
theorem philippians_strand_establishes : Establishes philippiansStrand
-- axioms: propext, Quot.sound
\end{verbatim}}
\par\medskip\noindent\textbf{\texttt{philippians\allowbreak{}Strand\_\allowbreak{}is\_\allowbreak{}satisfiable}}\par\nobreak\noindent
The Philippians strand can be held without contradiction.

{\small\begin{verbatim}
theorem philippiansStrand_is_satisfiable : Satisfiable
    philippiansStrand.premises
-- axioms: propext, Quot.sound
\end{verbatim}}
\par\medskip\noindent\textbf{\texttt{remonstrance\_\allowbreak{}establishes}}\par\nobreak\noindent
The Remonstrance’s conclusion follows from its premises.

{\small\begin{verbatim}
theorem remonstrance_establishes : Establishes remonstrance
-- axioms: propext, Quot.sound
\end{verbatim}}
\par\medskip\noindent\textbf{\texttt{remonstrance\_\allowbreak{}is\_\allowbreak{}satisfiable}}\par\nobreak\noindent
The Remonstrance can be held without contradiction.

{\small\begin{verbatim}
theorem remonstrance_is_satisfiable : Satisfiable remonstrance.premises
-- axioms: propext, Quot.sound
\end{verbatim}}
\par\medskip\noindent\textbf{\texttt{remonstrance\_\allowbreak{}leaves\_\allowbreak{}perseverance\_\allowbreak{}open}}\par\nobreak\noindent
\textbf{The Remonstrants’ premises leave perseverance open.} Their fifth article grants from John 10:28 that no one plucks believers out of Christ’s hand, and leaves whether they can fall away by negligence to be “more particularly determined out of the Holy Scripture”. The Reformed reading satisfies it, and so does a reading on which every rival’s premises hold. So the article neither teaches the perseverance of the saints nor denies it, as Arminius said of himself: “I never taught that a true believer can either totally or finally fall away from the faith, and perish”.

{\small\begin{verbatim}
theorem remonstrance_leaves_perseverance_open : Independent
    remonstrance.premises (p Claim.justifiedCanFinallyFallAway)
-- axioms: propext, Quot.sound
\end{verbatim}}
\par\medskip\noindent\textbf{\texttt{augustinian\_\allowbreak{}establishes}}\par\nobreak\noindent
Augustine and Aquinas’s conclusion follows from their premises.

{\small\begin{verbatim}
theorem augustinian_establishes : Establishes augustinian
-- axioms: propext, Quot.sound
\end{verbatim}}
\par\medskip\noindent\textbf{\texttt{augustinian\_\allowbreak{}is\_\allowbreak{}satisfiable}}\par\nobreak\noindent
Augustine and Aquinas can be held without contradiction.

{\small\begin{verbatim}
theorem augustinian_is_satisfiable : Satisfiable augustinian.premises
-- axioms: propext, Quot.sound
\end{verbatim}}
\par\medskip\noindent\textbf{\texttt{trent\allowbreak{}On\allowbreak{}Perseverance\_\allowbreak{}establishes}}\par\nobreak\noindent
Trent’s conclusion follows from its premises.

{\small\begin{verbatim}
theorem trentOnPerseverance_establishes : Establishes trentOnPerseverance
-- axioms: propext, Quot.sound
\end{verbatim}}
\par\medskip\noindent\textbf{\texttt{trent\allowbreak{}On\allowbreak{}Perseverance\_\allowbreak{}is\_\allowbreak{}satisfiable}}\par\nobreak\noindent
Trent can be held without contradiction.

{\small\begin{verbatim}
theorem trentOnPerseverance_is_satisfiable : Satisfiable
    trentOnPerseverance.premises
-- axioms: propext, Quot.sound
\end{verbatim}}
\par\medskip\noindent\textbf{\texttt{lutheran\allowbreak{}On\allowbreak{}Perseverance\_\allowbreak{}establishes}}\par\nobreak\noindent
The Lutherans’ conclusion follows from their premises.

{\small\begin{verbatim}
theorem lutheranOnPerseverance_establishes : Establishes
    lutheranOnPerseverance
-- axioms: propext, Quot.sound
\end{verbatim}}
\par\medskip\noindent\textbf{\texttt{lutheran\allowbreak{}On\allowbreak{}Perseverance\_\allowbreak{}is\_\allowbreak{}satisfiable}}\par\nobreak\noindent
The Lutherans can be held without contradiction.

{\small\begin{verbatim}
theorem lutheranOnPerseverance_is_satisfiable : Satisfiable
    lutheranOnPerseverance.premises
-- axioms: propext, Quot.sound
\end{verbatim}}
\par\medskip\noindent\textbf{\texttt{wesley\allowbreak{}On\allowbreak{}Perseverance\_\allowbreak{}establishes}}\par\nobreak\noindent
Wesley’s conclusion follows from his premises.

{\small\begin{verbatim}
theorem wesleyOnPerseverance_establishes : Establishes wesleyOnPerseverance
-- axioms: propext, Quot.sound
\end{verbatim}}
\par\medskip\noindent\textbf{\texttt{wesley\allowbreak{}On\allowbreak{}Perseverance\_\allowbreak{}is\_\allowbreak{}satisfiable}}\par\nobreak\noindent
Wesley can be held without contradiction.

{\small\begin{verbatim}
theorem wesleyOnPerseverance_is_satisfiable : Satisfiable
    wesleyOnPerseverance.premises
-- axioms: propext, Quot.sound
\end{verbatim}}
\paragraph*{Where Augustine parts from Westminster}
\leavevmode\par\nobreak
\par\medskip\noindent\textbf{\texttt{where\allowbreak{}Augustine\allowbreak{}Parts\allowbreak{}From\allowbreak{}Johns\allowbreak{}Strand}}\par\nobreak\noindent
\textbf{Where Augustine parts from the John route: over what 1 John 2:19 says of those who left.} Both read “they went out from us, but they were not of us”. The Reformed, with Owen and Calvin, read it of every apostate: those who fall away “never … have been true believers”. Augustine reads it otherwise: those who left “were in goodness”, “in the faith of children”, and were not elect. So the question is whether one who leaves the faith never had it or had it and gave it up, and the crux is the Reformed reading, rated \texttt{disputed}.

\paragraph{Why \emph{Reformed, the John strand alone} stands against \emph{Augustine and Aquinas (On Rebuke and Grace; Summa I q.23)}.}
\begin{itemize}
\item \textbf{The crux:} \(\pvar{TestimonyArguments.SolaGratia}{55}\), a premise of \emph{Reformed, the John strand alone}.
\item \textbf{What it does:} its conclusion does not follow without it.
\item \textbf{Where the rival breaks:} \(\lnot \pvar{TestimonyArguments.SolaGratia}{55}\) cannot be held together with \(\pvar{TestimonyArguments.SolaGratia}{55}\); each is needed for the break, and without the crux they stand.
\item \textbf{What this rests on:}
\begin{itemize}
\item \(\pvar{TestimonyArguments.SolaGratia}{55}\) 1 John 2:19, read of every apostate: all who fall away were never true believers --- \emph{disputed}: \cite[594--595]{owen-works-11-perseverance}; \cite[191--192]{calvin-commentaries-catholic-epistles}; \cite[479--480]{npnf1-05-1887}
\end{itemize}
\end{itemize}
\par\medskip\noindent\textbf{\texttt{where\allowbreak{}Augustine\allowbreak{}Parts\allowbreak{}From\allowbreak{}Philippians}}\par\nobreak\noindent
\textbf{Where Augustine’s position meets the Philippians route: at its one step.} Augustine is not cited on Philippians 1:6. The step from the verse to “no one truly justified can finally fall away” is rated \texttt{disputed} because Trent quotes the verse with its condition, and the Formula of Concord cites it while explaining why some fall away.

\paragraph{Why \emph{Reformed, the Philippians strand alone} stands against \emph{Augustine and Aquinas (On Rebuke and Grace; Summa I q.23)}.}
\begin{itemize}
\item \textbf{The crux:} \(\pvar{TestimonyArguments.SolaGratia}{33} \rightarrow \lnot \pvar{TestimonyArguments.SolaGratia}{43}\), a premise of \emph{Reformed, the Philippians strand alone}.
\item \textbf{What it does:} its conclusion does not follow without it.
\item \textbf{Granted:} \(\pvar{TestimonyArguments.SolaGratia}{33}\), which the break also needs.
\item \textbf{Where the rival breaks:} \(\pvar{TestimonyArguments.SolaGratia}{37}\), \(\pvar{TestimonyArguments.SolaGratia}{37} \rightarrow \pvar{TestimonyArguments.SolaGratia}{43}\) cannot be held together with \(\pvar{TestimonyArguments.SolaGratia}{33} \rightarrow \lnot \pvar{TestimonyArguments.SolaGratia}{43}\) and \(\pvar{TestimonyArguments.SolaGratia}{33}\); each is needed for the break, and without the crux they stand.
\item \textbf{What this rests on:}
\begin{itemize}
\item \(\pvar{TestimonyArguments.SolaGratia}{33}\) Philippians 1:6 — he who began a good work in you will bring it to completion --- \emph{consensus}: Phil 1:6
\item the step itself --- \emph{disputed}: \cite[444--445]{wesley-works-9-1827}; \cite[415--416]{calvin-commentary-john-1847}; \cite[Solid Declaration XI.45]{concordia-triglotta-1921}; \emph{disputed}: \cite[Trent, Session VI (1547), ch. 13]{tanner-decrees-1990}; \cite[40]{waterworth-trent-1848}; \cite[Solid Declaration XI.42]{concordia-triglotta-1921}
\end{itemize}
\end{itemize}
\par\medskip\noindent\textbf{\texttt{where\allowbreak{}Wesley\allowbreak{}Parts\allowbreak{}Over\allowbreak{}Romans8}}\par\nobreak\noindent
\textbf{Where Wesley parts from the Romans route: over what Romans 8:30 says.} Dort reads the verse as a chain from which no one drops, so that every one justified is glorified. Wesley reads it backwards, so that every one glorified was first justified: Paul “does not affirm … that precisely the same number of men are called, justified, and glorified”. The crux is Dort’s reading, rated \texttt{disputed}; it is not needed for Westminster’s conclusion, which other routes reach, so it is Westminster’s answer to Wesley on this text.

\paragraph{Why \emph{Reformed (Westminster XVII; Dort, Fifth Head)} stands against \emph{Wesley (Serious Thoughts upon the Perseverance of the Saints)}.}
\begin{itemize}
\item \textbf{The crux:} \(\pvar{TestimonyArguments.SolaGratia}{54}\), a premise of \emph{Reformed (Westminster XVII; Dort, Fifth Head)}.
\item \textbf{What it does:} its conclusion follows without it --- the crux is its answer to \emph{Wesley (Serious Thoughts upon the Perseverance of the Saints)}.
\item \textbf{Where the rival breaks:} \(\lnot \pvar{TestimonyArguments.SolaGratia}{54}\) cannot be held together with \(\pvar{TestimonyArguments.SolaGratia}{54}\); each is needed for the break, and without the crux they stand.
\item \textbf{What this rests on:}
\begin{itemize}
\item \(\pvar{TestimonyArguments.SolaGratia}{54}\) Romans 8:29–30 is a chain from which no one drops: every one justified is glorified --- \emph{disputed}: \cite[582--583]{schaff-creeds-christendom-3-1977}; \cite[557]{schaff-creeds-christendom-3-1977}; \cite[384]{wesley-explanatory-notes-nt}; \cite[413]{wesley-works-9-1827}
\end{itemize}
\end{itemize}
\paragraph*{The view from eternity}
\leavevmode\par\nobreak
\par\medskip\noindent\textbf{\texttt{from\allowbreak{}Eternity}} --- The view from eternity (Boethius; Aquinas)\par\nobreak\noindent
\textbf{The view from eternity}, from Boethius (\emph{Consolation} V, prose 6) and Aquinas (\emph{Summa} I q.14 a.13; I q.23 a.6). God knows all things, the future included, as present to his eternity, and his knowing them imposes no necessity on them. So the one who perseveres is predestined, and still, considered in himself, could have fallen away.

\begin{derivation}
  \item[(1)] \(\pvar{TestimonyArguments.SolaGratia}{57}\)\gloss{God knows all things as present to his eternity, which imposes no necessity on them}
  \item[(2)] \(\pvar{TestimonyArguments.SolaGratia}{41}\)\gloss{Those God has chosen to save cannot finally fall away from grace}
  \item[(3)] \(\pvar{TestimonyArguments.SolaGratia}{57} \rightarrow \pvar{TestimonyArguments.SolaGratia}{56}\)\gloss{\emph{would give:} Considered in himself, apart from God’s decree, the one who perseveres could have fallen}
  \item[\(\vdash\)] \(\pvar{TestimonyArguments.SolaGratia}{41} \land \pvar{TestimonyArguments.SolaGratia}{56}\)\gloss{\emph{at issue:} \(\pvar{TestimonyArguments.SolaGratia}{41}\) Those God has chosen to save cannot finally fall away from grace; \(\pvar{TestimonyArguments.SolaGratia}{56}\) Considered in himself, apart from God’s decree, the one who perseveres could have fallen}
  \keepwithnext
\end{derivation}
\emph{No premise here rests on scripture alone.}
\par\medskip\noindent\textbf{\texttt{from\allowbreak{}Eternity\_\allowbreak{}establishes}}\par\nobreak\noindent
The view from eternity delivers its conclusion.

{\small\begin{verbatim}
theorem fromEternity_establishes : Establishes fromEternity
-- axioms: propext, Quot.sound
\end{verbatim}}
\par\medskip\noindent\textbf{\texttt{from\allowbreak{}Eternity\_\allowbreak{}is\_\allowbreak{}satisfiable}}\par\nobreak\noindent
The view from eternity can be held without contradiction.

{\small\begin{verbatim}
theorem fromEternity_is_satisfiable : Satisfiable fromEternity.premises
-- axioms: propext, Quot.sound
\end{verbatim}}
\par\medskip\noindent\textbf{\texttt{from\_\allowbreak{}eternity\_\allowbreak{}holds\_\allowbreak{}with\_\allowbreak{}westminster}}\par\nobreak\noindent
\textbf{The view from eternity can be added to Westminster’s position without contradiction.} One reading satisfies every premise of both. Dort itself grants that the believers’ falling, “with respect to themselves is not only possible, but would undoubtedly happen; but with respect to God, it is utterly impossible” (V, article 8).

What this does not claim: that Westminster holds the view from eternity, or that the view supports Westminster. As encoded, the view says nothing on whether the truly justified can finally fall away, so it can be added to either side.

{\small\begin{verbatim}
theorem from_eternity_holds_with_westminster : Satisfiable
    (perseveranceReformed.premises ++ fromEternity.premises)
-- axioms: propext, Quot.sound
\end{verbatim}}
\par\medskip\noindent\textbf{\texttt{from\_\allowbreak{}eternity\_\allowbreak{}holds\_\allowbreak{}with\_\allowbreak{}wesley}}\par\nobreak\noindent
\textbf{The view from eternity can be added to Wesley’s position without contradiction}, with its claim that the elect persevere, which Wesley’s position is not encoded as holding. One reading satisfies every premise of both. Wesley holds, with Boethius, that all time is present to God and that his knowing compels nothing: “Men are as free in believing or not believing as if he did not know it at all” (Sermon 58, §5).

What this does not claim: that Wesley holds the view as Aquinas does. For Wesley God knows what we do because we do it; for Aquinas predestination is not grounded in anything God foresees.

{\small\begin{verbatim}
theorem from_eternity_holds_with_wesley : Satisfiable
    (wesleyOnPerseverance.premises ++ fromEternity.premises)
-- axioms: propext, Quot.sound
\end{verbatim}}
\par\medskip\noindent\textbf{\texttt{from\_\allowbreak{}eternity\_\allowbreak{}leaves\_\allowbreak{}perseverance\_\allowbreak{}open}}\par\nobreak\noindent
\textbf{The view from eternity, by itself, leaves perseverance open.} Granted that God knows all things as present, that the elect persevere, and that each could, considered in himself, have fallen, it neither follows nor fails that one truly justified can finally fall away. This holds because the view, as encoded, says nothing about final falling; the next result shows what does decide it.

{\small\begin{verbatim}
theorem from_eternity_leaves_perseverance_open : Independent
    fromEternity.premises (p Claim.justifiedCanFinallyFallAway)
-- axioms: propext, Quot.sound
\end{verbatim}}
\par\medskip\noindent\textbf{\texttt{from\allowbreak{}Eternity\allowbreak{}With\allowbreak{}Owens\allowbreak{}Reading}} --- The view from eternity, with Owen’s reading of 1 John 2:19\par\nobreak\noindent
The view from eternity with Owen’s reading of 1 John 2:19, that all who fall away were never true believers, and the step from it to the claim that the truly justified are elect. The one difference from \texttt{from\allowbreak{}Eternity} is that line, with the step that joins it to the elect’s perseverance.

\begin{derivation}
  \item[(1)] \(\pvar{TestimonyArguments.SolaGratia}{57}\)\gloss{God knows all things as present to his eternity, which imposes no necessity on them}
  \item[(2)] \(\pvar{TestimonyArguments.SolaGratia}{35}\)\gloss{1 John 2:19 — they went out from us, but they were not of us}
  \item[(3)] \(\pvar{TestimonyArguments.SolaGratia}{55}\)\gloss{1 John 2:19, read of every apostate: all who fall away were never true believers}
  \breakrow{(4)} \(\pvar{TestimonyArguments.SolaGratia}{41}\)\gloss{Those God has chosen to save cannot finally fall away from grace}
  \breakrow{(5)} \(\pvar{TestimonyArguments.SolaGratia}{57} \rightarrow \pvar{TestimonyArguments.SolaGratia}{56}\)\gloss{\emph{would give:} Considered in himself, apart from God’s decree, the one who perseveres could have fallen}
  \breakrow{(6)} \((\pvar{TestimonyArguments.SolaGratia}{35} \land \pvar{TestimonyArguments.SolaGratia}{55}) \rightarrow \pvar{TestimonyArguments.SolaGratia}{42}\)\gloss{\emph{would give:} Everyone who is truly justified is among the elect}
  \item[(7)] \((\pvar{TestimonyArguments.SolaGratia}{41} \land \pvar{TestimonyArguments.SolaGratia}{42}) \rightarrow \lnot \pvar{TestimonyArguments.SolaGratia}{43}\)\gloss{\emph{would deny:} One who is truly justified can fall away from grace finally, and perish}
  \item[\(\vdash\)] \(\pvar{TestimonyArguments.SolaGratia}{41} \land \pvar{TestimonyArguments.SolaGratia}{56}\)\gloss{\emph{at issue:} \(\pvar{TestimonyArguments.SolaGratia}{41}\) Those God has chosen to save cannot finally fall away from grace; \(\pvar{TestimonyArguments.SolaGratia}{56}\) Considered in himself, apart from God’s decree, the one who perseveres could have fallen}
  \keepwithnext
\end{derivation}
\emph{Grounded in scripture alone:} 1 John 2:19 — they went out from us, but they were not of us.
\par\medskip\noindent\textbf{\texttt{from\allowbreak{}Eternity\allowbreak{}With\allowbreak{}Owens\allowbreak{}Reading\_\allowbreak{}is\_\allowbreak{}satisfiable}}\par\nobreak\noindent
The view from eternity with Owen’s reading can be held without contradiction.

{\small\begin{verbatim}
theorem fromEternityWithOwensReading_is_satisfiable : Satisfiable
    fromEternityWithOwensReading.premises
-- axioms: propext, Quot.sound
\end{verbatim}}
\par\medskip\noindent\textbf{\texttt{from\_\allowbreak{}eternity\_\allowbreak{}with\_\allowbreak{}owens\_\allowbreak{}reading\_\allowbreak{}settles\_\allowbreak{}it}}\par\nobreak\noindent
\textbf{With Owen’s reading of 1 John 2:19, the view from eternity denies that the truly justified can finally fall away.} Add to the view the reading that all who fall away were never true believers, and the step from it to “the truly justified are elect”, and the elect’s perseverance does the rest. So what decides the question is not God’s knowledge of time but who the elect are: the reading of 1 John 2:19, which Augustine denies and which is rated \texttt{disputed}.

{\small\begin{verbatim}
theorem from_eternity_with_owens_reading_settles_it : Entails
    fromEternityWithOwensReading.premises (notP
    Claim.justifiedCanFinallyFallAway)
-- axioms: propext, Quot.sound
\end{verbatim}}
\paragraph*{Falling for a time, and knowing that one will persevere}
\leavevmode\par\nobreak
\par\medskip\noindent\textbf{\texttt{where\allowbreak{}They\allowbreak{}Part\allowbreak{}Over\allowbreak{}David}}\par\nobreak\noindent
\textbf{Where the Smalcald Articles and Calvin part over David.} The Lutherans read David’s fall as losing the Spirit and his faith, and so hold that justifying grace can be lost. The Reformed answer, in Calvin’s words, is his reading of David’s own prayer, “take not your Holy Spirit from me” (Psalm 51:11): the words “imply that the Spirit had not altogether been taken away from him”. That reading is the crux, and it is rated \texttt{disputed}. It is not part of the case for the conclusion; it is Westminster’s answer to the Lutherans on falling for a time.

\paragraph{Why \emph{Reformed (Westminster XVII; Dort, Fifth Head)} stands against \emph{Lutheran (the Formula of Concord, XI; the Smalcald Articles)}.}
\begin{itemize}
\item \textbf{The crux:} \(\lnot \pvar{TestimonyArguments.SolaGratia}{52}\), a premise of \emph{Reformed (Westminster XVII; Dort, Fifth Head)}.
\item \textbf{What it does:} its conclusion follows without it --- the crux is its answer to \emph{Lutheran (the Formula of Concord, XI; the Smalcald Articles)}.
\item \textbf{Where the rival breaks:} \(\pvar{TestimonyArguments.SolaGratia}{52}\) cannot be held together with \(\lnot \pvar{TestimonyArguments.SolaGratia}{52}\); each is needed for the break, and without the crux they stand.
\item \textbf{What this rests on:}
\begin{itemize}
\item \(\pvar{TestimonyArguments.SolaGratia}{52}\) In his sin David lost the Holy Spirit and his faith --- \emph{disputed}: \cite[299--300]{calvin-psalms-2-1846}
\end{itemize}
\end{itemize}
\par\medskip\noindent\textbf{\texttt{where\allowbreak{}They\allowbreak{}Part\allowbreak{}Over\allowbreak{}Assurance}}\par\nobreak\noindent
\textbf{Where Trent and Dort part over knowing that one will persevere.} Dort holds that believers “may and do obtain assurance” of their perseverance, from the Spirit “witnessing with our spirit, that we are children and heirs of God” (Romans 8:16). Trent condemns anyone who says he will have “that great gift of perseverance unto the end” “of an absolute and infallible certainty”, “unless he have learned this by special revelation” (canon 16), and in its chapter on perseverance quotes “let those who think themselves to stand, take heed lest they fall” (1 Corinthians 10:12; ch. 13). The crux is Trent’s step from that verse, rated \texttt{disputed}: Calvin grants the verse and answers that “there are two kinds of assurance”, and that Paul does not tell believers “to tremble from uncertainty as to their salvation”: the assurance that “rests on the promises of God” is “inseparable from faith”, and Paul warns only against the one that “arises from negligence”.

\paragraph{Why \emph{Trent (Session VI, chapters 13 and 15, canons 16 and 23)} stands against \emph{Reformed (Westminster XVII; Dort, Fifth Head)}.}
\begin{itemize}
\item \textbf{The crux:} \(\pvar{TestimonyArguments.SolaGratia}{49} \rightarrow \lnot \pvar{TestimonyArguments.SolaGratia}{53}\), a premise of \emph{Trent (Session VI, chapters 13 and 15, canons 16 and 23)}.
\item \textbf{What it does:} its conclusion follows without it --- the crux is its answer to \emph{Reformed (Westminster XVII; Dort, Fifth Head)}.
\item \textbf{Granted:} \(\pvar{TestimonyArguments.SolaGratia}{49}\), which the break also needs.
\item \textbf{Where the rival breaks:} \(\pvar{TestimonyArguments.SolaGratia}{48}\), \(\pvar{TestimonyArguments.SolaGratia}{48} \rightarrow \pvar{TestimonyArguments.SolaGratia}{53}\) cannot be held together with \(\pvar{TestimonyArguments.SolaGratia}{49} \rightarrow \lnot \pvar{TestimonyArguments.SolaGratia}{53}\) and \(\pvar{TestimonyArguments.SolaGratia}{49}\); each is needed for the break, and without the crux they stand.
\item \textbf{What this rests on:}
\begin{itemize}
\item \(\pvar{TestimonyArguments.SolaGratia}{49}\) 1 Corinthians 10:12 — let anyone who thinks that he stands take heed lest he fall --- \emph{consensus}: 1 Cor 10:12
\item the step itself --- \emph{well supported}: \cite[Trent, Session VI (1547), canon 23]{tanner-decrees-1990}; \cite[47]{waterworth-trent-1848}; \cite[439]{wesley-works-9-1827}; \emph{disputed}: \cite[329--330]{calvin-corinthians-1-1848}; \cite[Trent, Session VI (1547), ch. 13]{tanner-decrees-1990}; \cite[40--41]{waterworth-trent-1848}
\end{itemize}
\end{itemize}
\paragraph*{What the condition means}
\leavevmode\par\nobreak
\par\medskip\noindent\textbf{\texttt{as\allowbreak{}Condition\allowbreak{}God\allowbreak{}Keeps}}\par\nobreak\noindent
The condition read as one God himself keeps: those who are Christ’s sheep do hear and follow, because God keeps them in faith. Dort’s reading.

{\small\begin{verbatim}
def asConditionGodKeeps : Reading Claim :=
  { name := "as a condition God himself keeps in his sheep, Dort's sense",
    commits := p Claim.godKeepsHisSheepInFaith,
    source := asConditionGodKeepsSource }
\end{verbatim}}
\par\medskip\noindent\textbf{\texttt{as\allowbreak{}Condition\allowbreak{}That\allowbreak{}Can\allowbreak{}Fail}}\par\nobreak\noindent
The condition read as one that can fail: one truly justified can fall away finally. Wesley’s sense.

{\small\begin{verbatim}
def asConditionThatCanFail : Reading Claim :=
  { name := "as a condition that can fail, Wesley's sense",
    commits := p Claim.justifiedCanFinallyFallAway,
    source := asConditionThatCanFailSource }
\end{verbatim}}
\par\medskip\noindent\textbf{\texttt{wesley\allowbreak{}As\allowbreak{}God\allowbreak{}Keeps}} --- Wesley (Serious Thoughts upon the Perseverance of the Saints), read as a condition God himself keeps in his sheep, Dort’s sense\par\nobreak\noindent
Wesley, with the condition read in Dort’s sense.

\begin{derivation}
  \item[(1)] \(\pvar{TestimonyArguments.SolaGratia}{34}\)\gloss{Hebrews 6:4–6 — it is impossible to restore again to repentance those once enlightened who have fallen away}
  \item[(2)] \(\pvar{TestimonyArguments.SolaGratia}{40}\)\gloss{Hebrews 6:4–6 describes true believers who actually fall away}
  \item[(3)] \(\pvar{TestimonyArguments.SolaGratia}{45}\)\gloss{Hebrews 10:26–29 — the one who has profaned the blood of the covenant by which he was sanctified}
  \breakrow{(4)} \(\pvar{TestimonyArguments.SolaGratia}{50}\)\gloss{Hebrews 10:29 describes one truly sanctified by Christ’s blood who falls away}
  \breakrow{(5)} \(\pvar{TestimonyArguments.SolaGratia}{46}\)\gloss{2 Peter 2:20–22 — having escaped the defilements of the world, they are again entangled in them}
  \breakrow{(6)} \(\pvar{TestimonyArguments.SolaGratia}{51}\)\gloss{2 Peter 2:20–22 describes true believers who fall back into the world}
  \breakrow{(7)} \(\pvar{TestimonyArguments.SolaGratia}{32}\)\gloss{John 10:27–29 — they will never perish, and no one will snatch them out of my hand}
  \breakrow{(8)} \(\pvar{TestimonyArguments.SolaGratia}{38}\)\gloss{God’s promise to keep those he has begun to save holds unless they fall away by their own fault}
  \breakrow{(9)} \(\pvar{TestimonyArguments.SolaGratia}{44}\)\gloss{Romans 8:29–30 — those whom he justified he also glorified}
  \breakrow{(10)} \(\lnot \pvar{TestimonyArguments.SolaGratia}{54}\)\gloss{\emph{denied:} Romans 8:29–30 is a chain from which no one drops: every one justified is glorified}
  \breakrow{(11)} \((\pvar{TestimonyArguments.SolaGratia}{34} \land \pvar{TestimonyArguments.SolaGratia}{40}) \rightarrow \pvar{TestimonyArguments.SolaGratia}{43}\)
  \breakrow{(12)} \((\pvar{TestimonyArguments.SolaGratia}{45} \land \pvar{TestimonyArguments.SolaGratia}{50}) \rightarrow \pvar{TestimonyArguments.SolaGratia}{43}\)
  \breakrow{(13)} \((\pvar{TestimonyArguments.SolaGratia}{46} \land \pvar{TestimonyArguments.SolaGratia}{51}) \rightarrow \pvar{TestimonyArguments.SolaGratia}{43}\)
  \item[(14)] \(\pvar{TestimonyArguments.SolaGratia}{38} \rightarrow \pvar{TestimonyArguments.SolaGratia}{39}\)\gloss{\emph{would give:} Those who are Christ’s sheep do hear and follow, because God keeps them in faith}
  \item[\(\vdash\)] \(\pvar{TestimonyArguments.SolaGratia}{43}\)\gloss{\emph{at issue:} One who is truly justified can fall away from grace finally, and perish}
  \keepwithnext
\end{derivation}
\emph{Grounded in scripture alone:} Hebrews 6:4–6 — it is impossible to restore again to repentance those once enlightened who have fallen away, Hebrews 10:26–29 — the one who has profaned the blood of the covenant by which he was sanctified, 2 Peter 2:20–22 — having escaped the defilements of the world, they are again entangled in them, John 10:27–29 — they will never perish, and no one will snatch them out of my hand, Romans 8:29–30 — those whom he justified he also glorified.
\par\medskip\noindent\textbf{\texttt{wesley\allowbreak{}As\allowbreak{}Can\allowbreak{}Fail}} --- Wesley (Serious Thoughts upon the Perseverance of the Saints), read as a condition that can fail, Wesley’s sense\par\nobreak\noindent
Wesley, with the condition read in his own sense.

\begin{derivation}
  \item[(1)] \(\pvar{TestimonyArguments.SolaGratia}{34}\)\gloss{Hebrews 6:4–6 — it is impossible to restore again to repentance those once enlightened who have fallen away}
  \item[(2)] \(\pvar{TestimonyArguments.SolaGratia}{40}\)\gloss{Hebrews 6:4–6 describes true believers who actually fall away}
  \item[(3)] \(\pvar{TestimonyArguments.SolaGratia}{45}\)\gloss{Hebrews 10:26–29 — the one who has profaned the blood of the covenant by which he was sanctified}
  \breakrow{(4)} \(\pvar{TestimonyArguments.SolaGratia}{50}\)\gloss{Hebrews 10:29 describes one truly sanctified by Christ’s blood who falls away}
  \breakrow{(5)} \(\pvar{TestimonyArguments.SolaGratia}{46}\)\gloss{2 Peter 2:20–22 — having escaped the defilements of the world, they are again entangled in them}
  \breakrow{(6)} \(\pvar{TestimonyArguments.SolaGratia}{51}\)\gloss{2 Peter 2:20–22 describes true believers who fall back into the world}
  \breakrow{(7)} \(\pvar{TestimonyArguments.SolaGratia}{32}\)\gloss{John 10:27–29 — they will never perish, and no one will snatch them out of my hand}
  \breakrow{(8)} \(\pvar{TestimonyArguments.SolaGratia}{38}\)\gloss{God’s promise to keep those he has begun to save holds unless they fall away by their own fault}
  \breakrow{(9)} \(\pvar{TestimonyArguments.SolaGratia}{44}\)\gloss{Romans 8:29–30 — those whom he justified he also glorified}
  \breakrow{(10)} \(\lnot \pvar{TestimonyArguments.SolaGratia}{54}\)\gloss{\emph{denied:} Romans 8:29–30 is a chain from which no one drops: every one justified is glorified}
  \breakrow{(11)} \((\pvar{TestimonyArguments.SolaGratia}{34} \land \pvar{TestimonyArguments.SolaGratia}{40}) \rightarrow \pvar{TestimonyArguments.SolaGratia}{43}\)
  \breakrow{(12)} \((\pvar{TestimonyArguments.SolaGratia}{45} \land \pvar{TestimonyArguments.SolaGratia}{50}) \rightarrow \pvar{TestimonyArguments.SolaGratia}{43}\)
  \breakrow{(13)} \((\pvar{TestimonyArguments.SolaGratia}{46} \land \pvar{TestimonyArguments.SolaGratia}{51}) \rightarrow \pvar{TestimonyArguments.SolaGratia}{43}\)
  \item[(14)] \(\pvar{TestimonyArguments.SolaGratia}{38} \rightarrow \pvar{TestimonyArguments.SolaGratia}{43}\)
  \item[\(\vdash\)] \(\pvar{TestimonyArguments.SolaGratia}{43}\)\gloss{\emph{at issue:} One who is truly justified can fall away from grace finally, and perish}
  \keepwithnext
\end{derivation}
\emph{Grounded in scripture alone:} Hebrews 6:4–6 — it is impossible to restore again to repentance those once enlightened who have fallen away, Hebrews 10:26–29 — the one who has profaned the blood of the covenant by which he was sanctified, 2 Peter 2:20–22 — having escaped the defilements of the world, they are again entangled in them, John 10:27–29 — they will never perish, and no one will snatch them out of my hand, Romans 8:29–30 — those whom he justified he also glorified.
\par\medskip\noindent\textbf{\texttt{wesley\allowbreak{}As\allowbreak{}God\allowbreak{}Keeps\_\allowbreak{}is\_\allowbreak{}satisfiable}}\par\nobreak\noindent
Wesley read the first way can be held without contradiction.

{\small\begin{verbatim}
theorem wesleyAsGodKeeps_is_satisfiable : Satisfiable
    wesleyAsGodKeeps.premises
-- axioms: propext, Quot.sound
\end{verbatim}}
\par\medskip\noindent\textbf{\texttt{wesley\allowbreak{}As\allowbreak{}Can\allowbreak{}Fail\_\allowbreak{}is\_\allowbreak{}satisfiable}}\par\nobreak\noindent
Wesley read the second way can be held without contradiction.

{\small\begin{verbatim}
theorem wesleyAsCanFail_is_satisfiable : Satisfiable wesleyAsCanFail.premises
-- axioms: propext, Quot.sound
\end{verbatim}}
\par\medskip\noindent\textbf{\texttt{where\allowbreak{}Wesleys\allowbreak{}Hebrews\allowbreak{}Falls}}\par\nobreak\noindent
\textbf{Why Westminster stands against Wesley, even with the condition read in Dort’s sense.} Wesley’s reading of Hebrews 6:4–6 still delivers that one truly justified can finally fall away, as do his readings of Hebrews 10 and 2 Peter 2. The crux named here is Calvin’s reading of Hebrews 6, that those described had only “some taste of his grace”. It is Westminster’s answer to Wesley, rated \texttt{disputed}.

\paragraph{Why \emph{Reformed (Westminster XVII; Dort, Fifth Head)} stands against \emph{Wesley (Serious Thoughts upon the Perseverance of the Saints), read as a condition God himself keeps in his sheep, Dort’s sense}.}
\begin{itemize}
\item \textbf{The crux:} \(\lnot \pvar{TestimonyArguments.SolaGratia}{40}\), a premise of \emph{Reformed (Westminster XVII; Dort, Fifth Head)}.
\item \textbf{What it does:} its conclusion follows without it --- the crux is its answer to \emph{Wesley (Serious Thoughts upon the Perseverance of the Saints), read as a condition God himself keeps in his sheep, Dort’s sense}.
\item \textbf{Where the rival breaks:} \(\pvar{TestimonyArguments.SolaGratia}{40}\) cannot be held together with \(\lnot \pvar{TestimonyArguments.SolaGratia}{40}\); each is needed for the break, and without the crux they stand.
\item \textbf{What this rests on:}
\begin{itemize}
\item \(\pvar{TestimonyArguments.SolaGratia}{40}\) Hebrews 6:4–6 describes true believers who actually fall away --- \emph{disputed}: \cite[137--138]{calvin-hebrews-1853}
\end{itemize}
\end{itemize}
\par\medskip\noindent\textbf{\texttt{where\allowbreak{}Wesleys\allowbreak{}Condition\allowbreak{}Falls}}\par\nobreak\noindent
\textbf{Why the John route stands against Wesley, with the condition taken as he means it.} Then the condition says that one truly justified can finally fall away. The John route needs both its steps to deny that, and the crux is the step from John 10:27–29, the one Wesley himself disputes (§18).

\paragraph{Why \emph{Reformed, the John strand alone} stands against \emph{Wesley (Serious Thoughts upon the Perseverance of the Saints), read as a condition that can fail, Wesley’s sense}.}
\begin{itemize}
\item \textbf{The crux:} \(\pvar{TestimonyArguments.SolaGratia}{32} \rightarrow \pvar{TestimonyArguments.SolaGratia}{41}\), a premise of \emph{Reformed, the John strand alone}.
\item \textbf{What it does:} its conclusion does not follow without it.
\item \textbf{Granted:} \(\pvar{TestimonyArguments.SolaGratia}{32}\), \(\pvar{TestimonyArguments.SolaGratia}{35}\), \(\pvar{TestimonyArguments.SolaGratia}{55}\), \((\pvar{TestimonyArguments.SolaGratia}{35} \land \pvar{TestimonyArguments.SolaGratia}{55}) \rightarrow \pvar{TestimonyArguments.SolaGratia}{42}\), \((\pvar{TestimonyArguments.SolaGratia}{41} \land \pvar{TestimonyArguments.SolaGratia}{42}) \rightarrow \lnot \pvar{TestimonyArguments.SolaGratia}{43}\), which the break also needs.
\item \textbf{Where the rival breaks:} \(\pvar{TestimonyArguments.SolaGratia}{38}\), \(\pvar{TestimonyArguments.SolaGratia}{38} \rightarrow \pvar{TestimonyArguments.SolaGratia}{43}\) cannot be held together with \(\pvar{TestimonyArguments.SolaGratia}{32} \rightarrow \pvar{TestimonyArguments.SolaGratia}{41}\) and \(\pvar{TestimonyArguments.SolaGratia}{32}\), \(\pvar{TestimonyArguments.SolaGratia}{35}\), \(\pvar{TestimonyArguments.SolaGratia}{55}\), \((\pvar{TestimonyArguments.SolaGratia}{35} \land \pvar{TestimonyArguments.SolaGratia}{55}) \rightarrow \pvar{TestimonyArguments.SolaGratia}{42}\), \((\pvar{TestimonyArguments.SolaGratia}{41} \land \pvar{TestimonyArguments.SolaGratia}{42}) \rightarrow \lnot \pvar{TestimonyArguments.SolaGratia}{43}\); each is needed for the break, and without the crux they stand.
\item \textbf{What this rests on:}
\begin{itemize}
\item \(\pvar{TestimonyArguments.SolaGratia}{32}\) John 10:27–29 — they will never perish, and no one will snatch them out of my hand --- \emph{consensus}: John 10:27-29
\item \(\pvar{TestimonyArguments.SolaGratia}{35}\) 1 John 2:19 — they went out from us, but they were not of us --- \emph{consensus}: 1 John 2:19
\item \(\pvar{TestimonyArguments.SolaGratia}{55}\) 1 John 2:19, read of every apostate: all who fall away were never true believers --- \emph{disputed}: \cite[594--595]{owen-works-11-perseverance}; \cite[191--192]{calvin-commentaries-catholic-epistles}; \cite[479--480]{npnf1-05-1887}
\item \(\pvar{TestimonyArguments.SolaGratia}{42}\) Everyone who is truly justified is among the elect --- \emph{disputed}: \cite[X.4]{westminster-confession-1647}; \cite[137--138]{calvin-hebrews-1853}
\item \(\pvar{TestimonyArguments.SolaGratia}{41}\) Those God has chosen to save cannot finally fall away from grace --- \emph{well supported}: \cite[XVII.1]{westminster-confession-1647}; \cite[636]{schaff-creeds-christendom-3-1977}; \cite[593]{schaff-creeds-christendom-3-1977}
\item \(\pvar{TestimonyArguments.SolaGratia}{43}\) One who is truly justified can fall away from grace finally, and perish --- \emph{disputed}: \cite[Trent, Session VI (1547), canon 23]{tanner-decrees-1990}; \cite[47]{waterworth-trent-1848}; \cite[439]{wesley-works-9-1827}
\item the step itself --- \emph{disputed}: \cite[444--445]{wesley-works-9-1827}; \cite[415--416]{calvin-commentary-john-1847}; \cite[Solid Declaration XI.45]{concordia-triglotta-1921}; \emph{well supported}: \cite[595]{owen-works-11-perseverance}; \cite[191--192]{calvin-commentaries-catholic-epistles}; \cite[137--138]{calvin-hebrews-1853}; \cite[Trent, Session VI (1547), canon 23]{tanner-decrees-1990}; \cite[47]{waterworth-trent-1848}; \emph{consensus}: \cite[XVII.1]{westminster-confession-1647}; \cite[636]{schaff-creeds-christendom-3-1977}
\end{itemize}
\end{itemize}
\par\medskip\noindent\textbf{\texttt{what\allowbreak{}The\allowbreak{}Condition\allowbreak{}Means}}\par\nobreak\noindent
\textbf{What the condition means, read both ways.} The claim is the one Trent, the Lutherans and Wesley share: God’s promise to keep those he has begun to save holds unless the believer turns away. Both senses are checked against Wesley’s position, because he has three grounds besides the condition, his readings of Hebrews 6, Hebrews 10 and 2 Peter 2, so his conclusion does not stand or fall with how the condition is taken.

Taken as a condition God himself keeps, in the sense of Dort’s Fifth Head (article 3), the condition is something Dort holds. Wesley still cannot be held with Westminster, because each of his three readings, of Hebrews 6, Hebrews 10 and 2 Peter 2, delivers the falling away by itself; the check names Calvin’s reading of Hebrews 6.

Taken as a condition that can fail, which is how Wesley means it (“a Saint may fall away”), the condition is the falling away. It cannot be held with the John route, at the step from John 10:27–29, which Wesley disputes.

Each reading is consistent: Wesley read either way can be held without contradiction. Both breaks come at readings and steps rated \texttt{disputed}, so neither shows the other side right.

\paragraph{Every reading of \(P_{38}\), answered.} \emph{Wesley (Serious Thoughts upon the Perseverance of the Saints)} holds \(\pvar{TestimonyArguments.SolaGratia}{38}\). It is read 2 ways here, and each reading is checked; none can be left out.

\paragraph{1. Read as a condition God himself keeps in his sheep, Dort’s sense.} The claim commits \emph{Wesley (Serious Thoughts upon the Perseverance of the Saints)} to \(\pvar{TestimonyArguments.SolaGratia}{39}\) --- so read by \cite[593]{schaff-creeds-christendom-3-1977}; \cite[415--416]{calvin-commentary-john-1847} (\emph{disputed}).

Against \emph{Reformed (Westminster XVII; Dort, Fifth Head)}, it cannot be held. Where it breaks:
\paragraph{Why \emph{Reformed (Westminster XVII; Dort, Fifth Head)} stands against \emph{Wesley (Serious Thoughts upon the Perseverance of the Saints), read as a condition God himself keeps in his sheep, Dort’s sense}.}
\begin{itemize}
\item \textbf{The crux:} \(\lnot \pvar{TestimonyArguments.SolaGratia}{40}\), a premise of \emph{Reformed (Westminster XVII; Dort, Fifth Head)}.
\item \textbf{What it does:} its conclusion follows without it --- the crux is its answer to \emph{Wesley (Serious Thoughts upon the Perseverance of the Saints), read as a condition God himself keeps in his sheep, Dort’s sense}.
\item \textbf{Where the rival breaks:} \(\pvar{TestimonyArguments.SolaGratia}{40}\) cannot be held together with \(\lnot \pvar{TestimonyArguments.SolaGratia}{40}\); each is needed for the break, and without the crux they stand.
\item \textbf{What this rests on:}
\begin{itemize}
\item \(\pvar{TestimonyArguments.SolaGratia}{40}\) Hebrews 6:4–6 describes true believers who actually fall away --- \emph{disputed}: \cite[137--138]{calvin-hebrews-1853}
\end{itemize}
\end{itemize}
\paragraph{2. Read as a condition that can fail, Wesley’s sense.} The claim commits \emph{Wesley (Serious Thoughts upon the Perseverance of the Saints)} to \(\pvar{TestimonyArguments.SolaGratia}{43}\) --- so read by \cite[439]{wesley-works-9-1827}; \cite[444--445]{wesley-works-9-1827} (\emph{disputed}).

Against \emph{Reformed, the John strand alone}, it cannot be held. Where it breaks:
\paragraph{Why \emph{Reformed, the John strand alone} stands against \emph{Wesley (Serious Thoughts upon the Perseverance of the Saints), read as a condition that can fail, Wesley’s sense}.}
\begin{itemize}
\item \textbf{The crux:} \(\pvar{TestimonyArguments.SolaGratia}{32} \rightarrow \pvar{TestimonyArguments.SolaGratia}{41}\), a premise of \emph{Reformed, the John strand alone}.
\item \textbf{What it does:} its conclusion does not follow without it.
\item \textbf{Granted:} \(\pvar{TestimonyArguments.SolaGratia}{32}\), \(\pvar{TestimonyArguments.SolaGratia}{35}\), \(\pvar{TestimonyArguments.SolaGratia}{55}\), \((\pvar{TestimonyArguments.SolaGratia}{35} \land \pvar{TestimonyArguments.SolaGratia}{55}) \rightarrow \pvar{TestimonyArguments.SolaGratia}{42}\), \((\pvar{TestimonyArguments.SolaGratia}{41} \land \pvar{TestimonyArguments.SolaGratia}{42}) \rightarrow \lnot \pvar{TestimonyArguments.SolaGratia}{43}\), which the break also needs.
\item \textbf{Where the rival breaks:} \(\pvar{TestimonyArguments.SolaGratia}{38}\), \(\pvar{TestimonyArguments.SolaGratia}{38} \rightarrow \pvar{TestimonyArguments.SolaGratia}{43}\) cannot be held together with \(\pvar{TestimonyArguments.SolaGratia}{32} \rightarrow \pvar{TestimonyArguments.SolaGratia}{41}\) and \(\pvar{TestimonyArguments.SolaGratia}{32}\), \(\pvar{TestimonyArguments.SolaGratia}{35}\), \(\pvar{TestimonyArguments.SolaGratia}{55}\), \((\pvar{TestimonyArguments.SolaGratia}{35} \land \pvar{TestimonyArguments.SolaGratia}{55}) \rightarrow \pvar{TestimonyArguments.SolaGratia}{42}\), \((\pvar{TestimonyArguments.SolaGratia}{41} \land \pvar{TestimonyArguments.SolaGratia}{42}) \rightarrow \lnot \pvar{TestimonyArguments.SolaGratia}{43}\); each is needed for the break, and without the crux they stand.
\item \textbf{What this rests on:}
\begin{itemize}
\item \(\pvar{TestimonyArguments.SolaGratia}{32}\) John 10:27–29 — they will never perish, and no one will snatch them out of my hand --- \emph{consensus}: John 10:27-29
\item \(\pvar{TestimonyArguments.SolaGratia}{35}\) 1 John 2:19 — they went out from us, but they were not of us --- \emph{consensus}: 1 John 2:19
\item \(\pvar{TestimonyArguments.SolaGratia}{55}\) 1 John 2:19, read of every apostate: all who fall away were never true believers --- \emph{disputed}: \cite[594--595]{owen-works-11-perseverance}; \cite[191--192]{calvin-commentaries-catholic-epistles}; \cite[479--480]{npnf1-05-1887}
\item \(\pvar{TestimonyArguments.SolaGratia}{42}\) Everyone who is truly justified is among the elect --- \emph{disputed}: \cite[X.4]{westminster-confession-1647}; \cite[137--138]{calvin-hebrews-1853}
\item \(\pvar{TestimonyArguments.SolaGratia}{41}\) Those God has chosen to save cannot finally fall away from grace --- \emph{well supported}: \cite[XVII.1]{westminster-confession-1647}; \cite[636]{schaff-creeds-christendom-3-1977}; \cite[593]{schaff-creeds-christendom-3-1977}
\item \(\pvar{TestimonyArguments.SolaGratia}{43}\) One who is truly justified can fall away from grace finally, and perish --- \emph{disputed}: \cite[Trent, Session VI (1547), canon 23]{tanner-decrees-1990}; \cite[47]{waterworth-trent-1848}; \cite[439]{wesley-works-9-1827}
\item the step itself --- \emph{disputed}: \cite[444--445]{wesley-works-9-1827}; \cite[415--416]{calvin-commentary-john-1847}; \cite[Solid Declaration XI.45]{concordia-triglotta-1921}; \emph{well supported}: \cite[595]{owen-works-11-perseverance}; \cite[191--192]{calvin-commentaries-catholic-epistles}; \cite[137--138]{calvin-hebrews-1853}; \cite[Trent, Session VI (1547), canon 23]{tanner-decrees-1990}; \cite[47]{waterworth-trent-1848}; \emph{consensus}: \cite[XVII.1]{westminster-confession-1647}; \cite[636]{schaff-creeds-christendom-3-1977}
\end{itemize}
\end{itemize}
\par\medskip\noindent\textbf{\texttt{perseverance\allowbreak{}Node}}\par\nobreak\noindent
The package each party argues from.

{\small\begin{verbatim}
def perseveranceNode : PerseveranceParty → ArgumentPackage Claim :=
  fun x =>
    match x with
    | PerseveranceParty.remonstrance => remonstrance
    | PerseveranceParty.reformed => perseveranceReformed
    | PerseveranceParty.augustinian => augustinian
    | PerseveranceParty.trent => trentOnPerseverance
    | PerseveranceParty.lutheran => lutheranOnPerseverance
    | PerseveranceParty.wesleyan => wesleyOnPerseverance
\end{verbatim}}
\par\medskip\noindent\textbf{\texttt{perseverance\allowbreak{}Dispute}}\par\nobreak\noindent
The dispute over the perseverance of the saints: every party can be held without contradiction, delivers its conclusion, and has its steps rated.

{\small\begin{verbatim}
def perseveranceDispute : Dispute Claim PerseveranceParty :=
  { node := perseveranceNode,
    consistent := perseveranceDispute._proof_1,
    sound := perseveranceDispute._proof_2,
    rated := perseveranceDispute._proof_9 }
\end{verbatim}}
\par\medskip\noindent\textbf{\texttt{perseverance\allowbreak{}Dispute\_\allowbreak{}defeats}}\par\nobreak\noindent
\textbf{Who defeats whom}, all 36 pairs: eight defeats, four standoffs. Westminster stands off against Augustine and Aquinas, Trent, the Lutherans and Wesley, over whether one truly justified can finally fall away; with the Lutherans and Trent also over whether grace can be lost for a time, and with Trent over whether a believer can know that he will persevere. The four do not attack one another, and nothing attacks the Remonstrance. So each of the four can be defended alongside the other three, and none must be accepted, because Westminster defeats each of them. Every cell is checked by the kernel.

{\small\begin{verbatim}
theorem perseveranceDispute_defeats : ∀ (i j : PerseveranceParty),
    perseveranceDispute.defeats i j ↔ (i, j) ∈ [(PerseveranceParty.reformed,
    PerseveranceParty.augustinian), (PerseveranceParty.reformed,
    PerseveranceParty.trent), (PerseveranceParty.reformed,
    PerseveranceParty.lutheran), (PerseveranceParty.reformed,
    PerseveranceParty.wesleyan), (PerseveranceParty.augustinian,
    PerseveranceParty.reformed), (PerseveranceParty.trent,
    PerseveranceParty.reformed), (PerseveranceParty.lutheran,
    PerseveranceParty.reformed), (PerseveranceParty.wesleyan,
    PerseveranceParty.reformed)]
-- axioms: propext, Classical.choice, Quot.sound
\end{verbatim}}
\par\medskip\noindent\textbf{\texttt{remonstrance\allowbreak{}Unanswered}}\par\nobreak\noindent
How far the dispute is settled before any choice. Nothing attacks the Remonstrance; every other party has a defeater that the Remonstrance does not answer.

\paragraph{What the dispute forces is exactly \emph{Remonstrance (article 5, 1610)}.}
\begin{itemize}
\item Stage 1: \emph{Remonstrance (article 5, 1610)}.
\begin{itemize}
\item \emph{Remonstrance (article 5, 1610)} is defeated by nothing.
\end{itemize}
\item Nothing else is forced:
\begin{itemize}
\item \emph{Augustine and Aquinas (On Rebuke and Grace; Summa I q.23)} defeats \emph{Reformed (Westminster XVII; Dort, Fifth Head)}, and nothing forced defeats \emph{Augustine and Aquinas (On Rebuke and Grace; Summa I q.23)}.
\item \emph{Reformed (Westminster XVII; Dort, Fifth Head)} defeats \emph{Augustine and Aquinas (On Rebuke and Grace; Summa I q.23)}, and nothing forced defeats \emph{Reformed (Westminster XVII; Dort, Fifth Head)}.
\item \emph{Reformed (Westminster XVII; Dort, Fifth Head)} defeats \emph{Trent (Session VI, chapters 13 and 15, canons 16 and 23)}, and nothing forced defeats \emph{Reformed (Westminster XVII; Dort, Fifth Head)}.
\item \emph{Reformed (Westminster XVII; Dort, Fifth Head)} defeats \emph{Lutheran (the Formula of Concord, XI; the Smalcald Articles)} (\texttt{where\allowbreak{}They\allowbreak{}Part\allowbreak{}Over\allowbreak{}David}), and nothing forced defeats \emph{Reformed (Westminster XVII; Dort, Fifth Head)}.
\item \emph{Reformed (Westminster XVII; Dort, Fifth Head)} defeats \emph{Wesley (Serious Thoughts upon the Perseverance of the Saints)} (\texttt{where\allowbreak{}Wesley\allowbreak{}Parts\allowbreak{}Over\allowbreak{}Romans8}), and nothing forced defeats \emph{Reformed (Westminster XVII; Dort, Fifth Head)}.
\end{itemize}
\end{itemize}
\par\noindent\textbf{What this rests on.} The reasons state 5 defeats and 8 absences of defeat, each a cell of the defeat table, \texttt{perseverance\allowbreak{}Dispute\_\allowbreak{}defeats}, computed from the two parties' premises and checked by the kernel. Whether an attack survives turns on the attacker's weakest link and on the rating of what it attacks, so a changed rating can change the verdict. The weakest links:
\begin{itemize}
\item \emph{Remonstrance (article 5, 1610)}, weakest at \emph{consensus}:
\begin{itemize}
\item John 10:27–29 — they will never perish, and no one will snatch them out of my hand --- John 10:27-29
\item an inference step --- \cite[548--549]{schaff-creeds-christendom-3-1977}; \cite[603]{arminius-works-1-1825}; \cite[Solid Declaration XI.45]{concordia-triglotta-1921}; \cite[444--445]{wesley-works-9-1827}
\end{itemize}
\item \emph{Reformed (Westminster XVII; Dort, Fifth Head)}, weakest at \emph{disputed}:
\begin{itemize}
\item 1 John 2:19, read of every apostate: all who fall away were never true believers --- \cite[594--595]{owen-works-11-perseverance}; \cite[191--192]{calvin-commentaries-catholic-epistles}; \cite[479--480]{npnf1-05-1887}
\item Romans 8:29–30 is a chain from which no one drops: every one justified is glorified --- \cite[582--583]{schaff-creeds-christendom-3-1977}; \cite[557]{schaff-creeds-christendom-3-1977}; \cite[384]{wesley-explanatory-notes-nt}; \cite[413]{wesley-works-9-1827}
\item In his sin David lost the Holy Spirit and his faith (encoded as a denial, so it ranks disputed) --- \cite[299--300]{calvin-psalms-2-1846}
\item Hebrews 6:4–6 describes true believers who actually fall away (encoded as a denial, so it ranks disputed) --- \cite[137--138]{calvin-hebrews-1853}
\item Hebrews 10:29 describes one truly sanctified by Christ’s blood who falls away (encoded as a denial, so it ranks disputed) --- \cite[545--546]{owen-works-23-1855}
\item 2 Peter 2:20–22 describes true believers who fall back into the world (encoded as a denial, so it ranks disputed) --- \cite[410--412]{calvin-commentaries-catholic-epistles}
\item an inference step --- \cite[444--445]{wesley-works-9-1827}; \cite[415--416]{calvin-commentary-john-1847}; \cite[Solid Declaration XI.45]{concordia-triglotta-1921}
\item an inference step --- \cite[Trent, Session VI (1547), ch. 13]{tanner-decrees-1990}; \cite[40]{waterworth-trent-1848}; \cite[Solid Declaration XI.42]{concordia-triglotta-1921}
\end{itemize}
\item \emph{Augustine and Aquinas (On Rebuke and Grace; Summa I q.23)}, weakest at \emph{disputed}:
\begin{itemize}
\item To some whom God has regenerated, and given faith, hope and love, he does not give perseverance --- \cite[478]{npnf1-05-1887}; \cite[On Rebuke and Grace, ch. 18]{npnf1-05-1887}
\item 1 John 2:19, read of every apostate: all who fall away were never true believers (encoded as a denial, so it ranks disputed) --- \cite[479--480]{npnf1-05-1887}
\end{itemize}
\item \emph{Trent (Session VI, chapters 13 and 15, canons 16 and 23)}, weakest at \emph{disputed}:
\begin{itemize}
\item God’s promise to keep those he has begun to save holds unless they fall away by their own fault --- \cite[Trent, Session VI (1547), ch. 13]{tanner-decrees-1990}; \cite[40]{waterworth-trent-1848}
\item Justifying grace, once given, can be lost --- \cite[Trent, Session VI (1547), ch. 15]{tanner-decrees-1990}; \cite[42]{waterworth-trent-1848}; \cite[Trent, Session VI (1547), canon 23]{tanner-decrees-1990}; \cite[47]{waterworth-trent-1848}
\item an inference step --- \cite[329--330]{calvin-corinthians-1-1848}; \cite[Trent, Session VI (1547), ch. 13]{tanner-decrees-1990}; \cite[40--41]{waterworth-trent-1848}
\end{itemize}
\item \emph{Lutheran (the Formula of Concord, XI; the Smalcald Articles)}, weakest at \emph{disputed}:
\begin{itemize}
\item God’s promise to keep those he has begun to save holds unless they fall away by their own fault --- \cite[Solid Declaration XI.42]{concordia-triglotta-1921}
\item In his sin David lost the Holy Spirit and his faith --- \cite[Smalcald Articles III.iii.43]{concordia-triglotta-1921}
\end{itemize}
\item \emph{Wesley (Serious Thoughts upon the Perseverance of the Saints)}, weakest at \emph{disputed}:
\begin{itemize}
\item Hebrews 6:4–6 describes true believers who actually fall away --- \cite[449]{wesley-works-9-1827}
\item Hebrews 10:29 describes one truly sanctified by Christ’s blood who falls away --- \cite[448--449]{wesley-works-9-1827}; \cite[585]{wesley-explanatory-notes-nt}
\item 2 Peter 2:20–22 describes true believers who fall back into the world --- \cite[445--446]{wesley-works-9-1827}; \cite[625]{wesley-explanatory-notes-nt}
\item God’s promise to keep those he has begun to save holds unless they fall away by their own fault --- \cite[444--445]{wesley-works-9-1827}
\item Romans 8:29–30 is a chain from which no one drops: every one justified is glorified (encoded as a denial, so it ranks disputed) --- \cite[384]{wesley-explanatory-notes-nt}; \cite[413]{wesley-works-9-1827}
\end{itemize}
\end{itemize}
\noindent Why \emph{Reformed (Westminster XVII; Dort, Fifth Head)} defeats \emph{Lutheran (the Formula of Concord, XI; the Smalcald Articles)}, at the level of the claims: \texttt{where\allowbreak{}They\allowbreak{}Part\allowbreak{}Over\allowbreak{}David}.
\noindent Why \emph{Reformed (Westminster XVII; Dort, Fifth Head)} defeats \emph{Wesley (Serious Thoughts upon the Perseverance of the Saints)}, at the level of the claims: \texttt{where\allowbreak{}Wesley\allowbreak{}Parts\allowbreak{}Over\allowbreak{}Romans8}.
\par\medskip\noindent\textbf{\texttt{only\_\allowbreak{}the\_\allowbreak{}remonstrance\_\allowbreak{}prevails}}\par\nobreak\noindent
\textbf{Only what the Remonstrants concluded from John 10:28 prevails: no power from outside snatches believers out of Christ’s hand.} Every other party is in a tie with another, and this is all that must be accepted however the ties are settled.

What this does not claim: that the perseverance of the saints is false, or that one truly justified can finally fall away. Both are in a standoff, and neither is forced.

{\small\begin{verbatim}
theorem only_the_remonstrance_prevails : Framework.grounded
    perseveranceDispute.defeats = Solver.toSet
    [PerseveranceParty.remonstrance]
-- axioms: propext, Classical.choice, Quot.sound
\end{verbatim}}
\par\medskip\noindent\textbf{\texttt{westminster\allowbreak{}Defended}}\par\nobreak\noindent
Why Westminster can be defended: with the Remonstrance, it answers every attacker itself.

\paragraph{\emph{Reformed (Westminster XVII; Dort, Fifth Head)} is accepted on some resolution.}
\begin{itemize}
\item A position holding \emph{Remonstrance (article 5, 1610)}, \emph{Reformed (Westminster XVII; Dort, Fifth Head)} can be held.
\item None of \emph{Remonstrance (article 5, 1610)}, \emph{Reformed (Westminster XVII; Dort, Fifth Head)} defeats another, and each attack on them is answered from among them:
\begin{itemize}
\item \emph{Augustine and Aquinas (On Rebuke and Grace; Summa I q.23)} defeats \emph{Reformed (Westminster XVII; Dort, Fifth Head)}, and \emph{Reformed (Westminster XVII; Dort, Fifth Head)} defeats \emph{Augustine and Aquinas (On Rebuke and Grace; Summa I q.23)}.
\item \emph{Trent (Session VI, chapters 13 and 15, canons 16 and 23)} defeats \emph{Reformed (Westminster XVII; Dort, Fifth Head)} (\texttt{where\allowbreak{}They\allowbreak{}Part\allowbreak{}Over\allowbreak{}Assurance}), and \emph{Reformed (Westminster XVII; Dort, Fifth Head)} defeats \emph{Trent (Session VI, chapters 13 and 15, canons 16 and 23)}.
\item \emph{Lutheran (the Formula of Concord, XI; the Smalcald Articles)} defeats \emph{Reformed (Westminster XVII; Dort, Fifth Head)}, and \emph{Reformed (Westminster XVII; Dort, Fifth Head)} defeats \emph{Lutheran (the Formula of Concord, XI; the Smalcald Articles)} (\texttt{where\allowbreak{}They\allowbreak{}Part\allowbreak{}Over\allowbreak{}David}).
\item \emph{Wesley (Serious Thoughts upon the Perseverance of the Saints)} defeats \emph{Reformed (Westminster XVII; Dort, Fifth Head)}, and \emph{Reformed (Westminster XVII; Dort, Fifth Head)} defeats \emph{Wesley (Serious Thoughts upon the Perseverance of the Saints)} (\texttt{where\allowbreak{}Wesley\allowbreak{}Parts\allowbreak{}Over\allowbreak{}Romans8}).
\end{itemize}
\end{itemize}
\par\noindent\textbf{What this rests on.} The reasons state 8 defeats and 8 absences of defeat, each a cell of the defeat table, \texttt{perseverance\allowbreak{}Dispute\_\allowbreak{}defeats}, computed from the two parties' premises and checked by the kernel. Whether an attack survives turns on the attacker's weakest link and on the rating of what it attacks, so a changed rating can change the verdict. The weakest links:
\begin{itemize}
\item \emph{Remonstrance (article 5, 1610)}, weakest at \emph{consensus}:
\begin{itemize}
\item John 10:27–29 — they will never perish, and no one will snatch them out of my hand --- John 10:27-29
\item an inference step --- \cite[548--549]{schaff-creeds-christendom-3-1977}; \cite[603]{arminius-works-1-1825}; \cite[Solid Declaration XI.45]{concordia-triglotta-1921}; \cite[444--445]{wesley-works-9-1827}
\end{itemize}
\item \emph{Reformed (Westminster XVII; Dort, Fifth Head)}, weakest at \emph{disputed}:
\begin{itemize}
\item 1 John 2:19, read of every apostate: all who fall away were never true believers --- \cite[594--595]{owen-works-11-perseverance}; \cite[191--192]{calvin-commentaries-catholic-epistles}; \cite[479--480]{npnf1-05-1887}
\item Romans 8:29–30 is a chain from which no one drops: every one justified is glorified --- \cite[582--583]{schaff-creeds-christendom-3-1977}; \cite[557]{schaff-creeds-christendom-3-1977}; \cite[384]{wesley-explanatory-notes-nt}; \cite[413]{wesley-works-9-1827}
\item In his sin David lost the Holy Spirit and his faith (encoded as a denial, so it ranks disputed) --- \cite[299--300]{calvin-psalms-2-1846}
\item Hebrews 6:4–6 describes true believers who actually fall away (encoded as a denial, so it ranks disputed) --- \cite[137--138]{calvin-hebrews-1853}
\item Hebrews 10:29 describes one truly sanctified by Christ’s blood who falls away (encoded as a denial, so it ranks disputed) --- \cite[545--546]{owen-works-23-1855}
\item 2 Peter 2:20–22 describes true believers who fall back into the world (encoded as a denial, so it ranks disputed) --- \cite[410--412]{calvin-commentaries-catholic-epistles}
\item an inference step --- \cite[444--445]{wesley-works-9-1827}; \cite[415--416]{calvin-commentary-john-1847}; \cite[Solid Declaration XI.45]{concordia-triglotta-1921}
\item an inference step --- \cite[Trent, Session VI (1547), ch. 13]{tanner-decrees-1990}; \cite[40]{waterworth-trent-1848}; \cite[Solid Declaration XI.42]{concordia-triglotta-1921}
\end{itemize}
\item \emph{Augustine and Aquinas (On Rebuke and Grace; Summa I q.23)}, weakest at \emph{disputed}:
\begin{itemize}
\item To some whom God has regenerated, and given faith, hope and love, he does not give perseverance --- \cite[478]{npnf1-05-1887}; \cite[On Rebuke and Grace, ch. 18]{npnf1-05-1887}
\item 1 John 2:19, read of every apostate: all who fall away were never true believers (encoded as a denial, so it ranks disputed) --- \cite[479--480]{npnf1-05-1887}
\end{itemize}
\item \emph{Trent (Session VI, chapters 13 and 15, canons 16 and 23)}, weakest at \emph{disputed}:
\begin{itemize}
\item God’s promise to keep those he has begun to save holds unless they fall away by their own fault --- \cite[Trent, Session VI (1547), ch. 13]{tanner-decrees-1990}; \cite[40]{waterworth-trent-1848}
\item Justifying grace, once given, can be lost --- \cite[Trent, Session VI (1547), ch. 15]{tanner-decrees-1990}; \cite[42]{waterworth-trent-1848}; \cite[Trent, Session VI (1547), canon 23]{tanner-decrees-1990}; \cite[47]{waterworth-trent-1848}
\item an inference step --- \cite[329--330]{calvin-corinthians-1-1848}; \cite[Trent, Session VI (1547), ch. 13]{tanner-decrees-1990}; \cite[40--41]{waterworth-trent-1848}
\end{itemize}
\item \emph{Lutheran (the Formula of Concord, XI; the Smalcald Articles)}, weakest at \emph{disputed}:
\begin{itemize}
\item God’s promise to keep those he has begun to save holds unless they fall away by their own fault --- \cite[Solid Declaration XI.42]{concordia-triglotta-1921}
\item In his sin David lost the Holy Spirit and his faith --- \cite[Smalcald Articles III.iii.43]{concordia-triglotta-1921}
\end{itemize}
\item \emph{Wesley (Serious Thoughts upon the Perseverance of the Saints)}, weakest at \emph{disputed}:
\begin{itemize}
\item Hebrews 6:4–6 describes true believers who actually fall away --- \cite[449]{wesley-works-9-1827}
\item Hebrews 10:29 describes one truly sanctified by Christ’s blood who falls away --- \cite[448--449]{wesley-works-9-1827}; \cite[585]{wesley-explanatory-notes-nt}
\item 2 Peter 2:20–22 describes true believers who fall back into the world --- \cite[445--446]{wesley-works-9-1827}; \cite[625]{wesley-explanatory-notes-nt}
\item God’s promise to keep those he has begun to save holds unless they fall away by their own fault --- \cite[444--445]{wesley-works-9-1827}
\item Romans 8:29–30 is a chain from which no one drops: every one justified is glorified (encoded as a denial, so it ranks disputed) --- \cite[384]{wesley-explanatory-notes-nt}; \cite[413]{wesley-works-9-1827}
\end{itemize}
\end{itemize}
\noindent Why \emph{Trent (Session VI, chapters 13 and 15, canons 16 and 23)} defeats \emph{Reformed (Westminster XVII; Dort, Fifth Head)}, at the level of the claims: \texttt{where\allowbreak{}They\allowbreak{}Part\allowbreak{}Over\allowbreak{}Assurance}.
\noindent Why \emph{Reformed (Westminster XVII; Dort, Fifth Head)} defeats \emph{Lutheran (the Formula of Concord, XI; the Smalcald Articles)}, at the level of the claims: \texttt{where\allowbreak{}They\allowbreak{}Part\allowbreak{}Over\allowbreak{}David}.
\noindent Why \emph{Reformed (Westminster XVII; Dort, Fifth Head)} defeats \emph{Wesley (Serious Thoughts upon the Perseverance of the Saints)}, at the level of the claims: \texttt{where\allowbreak{}Wesley\allowbreak{}Parts\allowbreak{}Over\allowbreak{}Romans8}.
\par\medskip\noindent\textbf{\texttt{perseverance\_\allowbreak{}defensible}}\par\nobreak\noindent
\textbf{Westminster can be defended}: on at least one way of settling the ties, it is accepted, with the Remonstrance, because it defeats each of the four that attack it.

{\small\begin{verbatim}
theorem perseverance_defensible : Framework.CredulouslyAccepted
    perseveranceDispute.defeats PerseveranceParty.reformed
-- axioms: propext, Classical.choice, Quot.sound
\end{verbatim}}
\par\medskip\noindent\textbf{\texttt{westminster\allowbreak{}Answered\allowbreak{}By\allowbreak{}Augustine}}\par\nobreak\noindent
Why Westminster is not forced: a maximal defensible position holds Augustine and Aquinas, and they defeat Westminster.

\paragraph{\emph{Reformed (Westminster XVII; Dort, Fifth Head)} is not accepted on every resolution.}
\begin{itemize}
\item A position holding \emph{Remonstrance (article 5, 1610)}, \emph{Augustine and Aquinas (On Rebuke and Grace; Summa I q.23)}, \emph{Trent (Session VI, chapters 13 and 15, canons 16 and 23)}, \emph{Lutheran (the Formula of Concord, XI; the Smalcald Articles)}, \emph{Wesley (Serious Thoughts upon the Perseverance of the Saints)} can be held, and it holds \emph{Augustine and Aquinas (On Rebuke and Grace; Summa I q.23)}; \emph{Augustine and Aquinas (On Rebuke and Grace; Summa I q.23)} conflicts with \emph{Reformed (Westminster XVII; Dort, Fifth Head)}: \emph{Augustine and Aquinas (On Rebuke and Grace; Summa I q.23)} defeats \emph{Reformed (Westminster XVII; Dort, Fifth Head)}, and \emph{Reformed (Westminster XVII; Dort, Fifth Head)} defeats \emph{Augustine and Aquinas (On Rebuke and Grace; Summa I q.23)}.
\item None of \emph{Remonstrance (article 5, 1610)}, \emph{Augustine and Aquinas (On Rebuke and Grace; Summa I q.23)}, \emph{Trent (Session VI, chapters 13 and 15, canons 16 and 23)}, \emph{Lutheran (the Formula of Concord, XI; the Smalcald Articles)}, \emph{Wesley (Serious Thoughts upon the Perseverance of the Saints)} defeats another, and each attack on them is answered from among them:
\begin{itemize}
\item \emph{Reformed (Westminster XVII; Dort, Fifth Head)} defeats \emph{Augustine and Aquinas (On Rebuke and Grace; Summa I q.23)}, and \emph{Augustine and Aquinas (On Rebuke and Grace; Summa I q.23)} defeats \emph{Reformed (Westminster XVII; Dort, Fifth Head)}.
\item \emph{Reformed (Westminster XVII; Dort, Fifth Head)} defeats \emph{Trent (Session VI, chapters 13 and 15, canons 16 and 23)}, and \emph{Augustine and Aquinas (On Rebuke and Grace; Summa I q.23)} defeats \emph{Reformed (Westminster XVII; Dort, Fifth Head)}.
\item \emph{Reformed (Westminster XVII; Dort, Fifth Head)} defeats \emph{Lutheran (the Formula of Concord, XI; the Smalcald Articles)} (\texttt{where\allowbreak{}They\allowbreak{}Part\allowbreak{}Over\allowbreak{}David}), and \emph{Augustine and Aquinas (On Rebuke and Grace; Summa I q.23)} defeats \emph{Reformed (Westminster XVII; Dort, Fifth Head)}.
\item \emph{Reformed (Westminster XVII; Dort, Fifth Head)} defeats \emph{Wesley (Serious Thoughts upon the Perseverance of the Saints)} (\texttt{where\allowbreak{}Wesley\allowbreak{}Parts\allowbreak{}Over\allowbreak{}Romans8}), and \emph{Augustine and Aquinas (On Rebuke and Grace; Summa I q.23)} defeats \emph{Reformed (Westminster XVII; Dort, Fifth Head)}.
\end{itemize}
\end{itemize}
\par\noindent\textbf{What this rests on.} The reasons state 5 defeats and 26 absences of defeat, each a cell of the defeat table, \texttt{perseverance\allowbreak{}Dispute\_\allowbreak{}defeats}, computed from the two parties' premises and checked by the kernel. Whether an attack survives turns on the attacker's weakest link and on the rating of what it attacks, so a changed rating can change the verdict. The weakest links:
\begin{itemize}
\item \emph{Augustine and Aquinas (On Rebuke and Grace; Summa I q.23)}, weakest at \emph{disputed}:
\begin{itemize}
\item To some whom God has regenerated, and given faith, hope and love, he does not give perseverance --- \cite[478]{npnf1-05-1887}; \cite[On Rebuke and Grace, ch. 18]{npnf1-05-1887}
\item 1 John 2:19, read of every apostate: all who fall away were never true believers (encoded as a denial, so it ranks disputed) --- \cite[479--480]{npnf1-05-1887}
\end{itemize}
\item \emph{Reformed (Westminster XVII; Dort, Fifth Head)}, weakest at \emph{disputed}:
\begin{itemize}
\item 1 John 2:19, read of every apostate: all who fall away were never true believers --- \cite[594--595]{owen-works-11-perseverance}; \cite[191--192]{calvin-commentaries-catholic-epistles}; \cite[479--480]{npnf1-05-1887}
\item Romans 8:29–30 is a chain from which no one drops: every one justified is glorified --- \cite[582--583]{schaff-creeds-christendom-3-1977}; \cite[557]{schaff-creeds-christendom-3-1977}; \cite[384]{wesley-explanatory-notes-nt}; \cite[413]{wesley-works-9-1827}
\item In his sin David lost the Holy Spirit and his faith (encoded as a denial, so it ranks disputed) --- \cite[299--300]{calvin-psalms-2-1846}
\item Hebrews 6:4–6 describes true believers who actually fall away (encoded as a denial, so it ranks disputed) --- \cite[137--138]{calvin-hebrews-1853}
\item Hebrews 10:29 describes one truly sanctified by Christ’s blood who falls away (encoded as a denial, so it ranks disputed) --- \cite[545--546]{owen-works-23-1855}
\item 2 Peter 2:20–22 describes true believers who fall back into the world (encoded as a denial, so it ranks disputed) --- \cite[410--412]{calvin-commentaries-catholic-epistles}
\item an inference step --- \cite[444--445]{wesley-works-9-1827}; \cite[415--416]{calvin-commentary-john-1847}; \cite[Solid Declaration XI.45]{concordia-triglotta-1921}
\item an inference step --- \cite[Trent, Session VI (1547), ch. 13]{tanner-decrees-1990}; \cite[40]{waterworth-trent-1848}; \cite[Solid Declaration XI.42]{concordia-triglotta-1921}
\end{itemize}
\item \emph{Remonstrance (article 5, 1610)}, weakest at \emph{consensus}:
\begin{itemize}
\item John 10:27–29 — they will never perish, and no one will snatch them out of my hand --- John 10:27-29
\item an inference step --- \cite[548--549]{schaff-creeds-christendom-3-1977}; \cite[603]{arminius-works-1-1825}; \cite[Solid Declaration XI.45]{concordia-triglotta-1921}; \cite[444--445]{wesley-works-9-1827}
\end{itemize}
\item \emph{Trent (Session VI, chapters 13 and 15, canons 16 and 23)}, weakest at \emph{disputed}:
\begin{itemize}
\item God’s promise to keep those he has begun to save holds unless they fall away by their own fault --- \cite[Trent, Session VI (1547), ch. 13]{tanner-decrees-1990}; \cite[40]{waterworth-trent-1848}
\item Justifying grace, once given, can be lost --- \cite[Trent, Session VI (1547), ch. 15]{tanner-decrees-1990}; \cite[42]{waterworth-trent-1848}; \cite[Trent, Session VI (1547), canon 23]{tanner-decrees-1990}; \cite[47]{waterworth-trent-1848}
\item an inference step --- \cite[329--330]{calvin-corinthians-1-1848}; \cite[Trent, Session VI (1547), ch. 13]{tanner-decrees-1990}; \cite[40--41]{waterworth-trent-1848}
\end{itemize}
\item \emph{Lutheran (the Formula of Concord, XI; the Smalcald Articles)}, weakest at \emph{disputed}:
\begin{itemize}
\item God’s promise to keep those he has begun to save holds unless they fall away by their own fault --- \cite[Solid Declaration XI.42]{concordia-triglotta-1921}
\item In his sin David lost the Holy Spirit and his faith --- \cite[Smalcald Articles III.iii.43]{concordia-triglotta-1921}
\end{itemize}
\item \emph{Wesley (Serious Thoughts upon the Perseverance of the Saints)}, weakest at \emph{disputed}:
\begin{itemize}
\item Hebrews 6:4–6 describes true believers who actually fall away --- \cite[449]{wesley-works-9-1827}
\item Hebrews 10:29 describes one truly sanctified by Christ’s blood who falls away --- \cite[448--449]{wesley-works-9-1827}; \cite[585]{wesley-explanatory-notes-nt}
\item 2 Peter 2:20–22 describes true believers who fall back into the world --- \cite[445--446]{wesley-works-9-1827}; \cite[625]{wesley-explanatory-notes-nt}
\item God’s promise to keep those he has begun to save holds unless they fall away by their own fault --- \cite[444--445]{wesley-works-9-1827}
\item Romans 8:29–30 is a chain from which no one drops: every one justified is glorified (encoded as a denial, so it ranks disputed) --- \cite[384]{wesley-explanatory-notes-nt}; \cite[413]{wesley-works-9-1827}
\end{itemize}
\end{itemize}
\noindent Why \emph{Reformed (Westminster XVII; Dort, Fifth Head)} defeats \emph{Lutheran (the Formula of Concord, XI; the Smalcald Articles)}, at the level of the claims: \texttt{where\allowbreak{}They\allowbreak{}Part\allowbreak{}Over\allowbreak{}David}.
\noindent Why \emph{Reformed (Westminster XVII; Dort, Fifth Head)} defeats \emph{Wesley (Serious Thoughts upon the Perseverance of the Saints)}, at the level of the claims: \texttt{where\allowbreak{}Wesley\allowbreak{}Parts\allowbreak{}Over\allowbreak{}Romans8}.
\par\medskip\noindent\textbf{\texttt{perseverance\_\allowbreak{}not\_\allowbreak{}forced}}\par\nobreak\noindent
\textbf{Nor must Westminster be accepted.} On another way of settling the ties, Augustine and Aquinas, Trent, the Lutherans and Wesley are accepted together, and Westminster cannot be accepted with them.

{\small\begin{verbatim}
theorem perseverance_not_forced : ¬Framework.SkepticallyAccepted
    perseveranceDispute.defeats PerseveranceParty.reformed
-- axioms: propext, Classical.choice, Quot.sound
\end{verbatim}}
\par\medskip\noindent\textbf{\texttt{augustine\allowbreak{}Defended}}\par\nobreak\noindent
Why Augustine and Aquinas can be defended: with Trent, the Lutherans and Wesley, each of whom defeats Westminster.

\paragraph{\emph{Augustine and Aquinas (On Rebuke and Grace; Summa I q.23)} is accepted on some resolution.}
\begin{itemize}
\item A position holding \emph{Remonstrance (article 5, 1610)}, \emph{Augustine and Aquinas (On Rebuke and Grace; Summa I q.23)}, \emph{Trent (Session VI, chapters 13 and 15, canons 16 and 23)}, \emph{Lutheran (the Formula of Concord, XI; the Smalcald Articles)}, \emph{Wesley (Serious Thoughts upon the Perseverance of the Saints)} can be held.
\item None of \emph{Remonstrance (article 5, 1610)}, \emph{Augustine and Aquinas (On Rebuke and Grace; Summa I q.23)}, \emph{Trent (Session VI, chapters 13 and 15, canons 16 and 23)}, \emph{Lutheran (the Formula of Concord, XI; the Smalcald Articles)}, \emph{Wesley (Serious Thoughts upon the Perseverance of the Saints)} defeats another, and each attack on them is answered from among them:
\begin{itemize}
\item \emph{Reformed (Westminster XVII; Dort, Fifth Head)} defeats \emph{Augustine and Aquinas (On Rebuke and Grace; Summa I q.23)}, and \emph{Augustine and Aquinas (On Rebuke and Grace; Summa I q.23)} defeats \emph{Reformed (Westminster XVII; Dort, Fifth Head)}.
\item \emph{Reformed (Westminster XVII; Dort, Fifth Head)} defeats \emph{Trent (Session VI, chapters 13 and 15, canons 16 and 23)}, and \emph{Augustine and Aquinas (On Rebuke and Grace; Summa I q.23)} defeats \emph{Reformed (Westminster XVII; Dort, Fifth Head)}.
\item \emph{Reformed (Westminster XVII; Dort, Fifth Head)} defeats \emph{Lutheran (the Formula of Concord, XI; the Smalcald Articles)} (\texttt{where\allowbreak{}They\allowbreak{}Part\allowbreak{}Over\allowbreak{}David}), and \emph{Augustine and Aquinas (On Rebuke and Grace; Summa I q.23)} defeats \emph{Reformed (Westminster XVII; Dort, Fifth Head)}.
\item \emph{Reformed (Westminster XVII; Dort, Fifth Head)} defeats \emph{Wesley (Serious Thoughts upon the Perseverance of the Saints)} (\texttt{where\allowbreak{}Wesley\allowbreak{}Parts\allowbreak{}Over\allowbreak{}Romans8}), and \emph{Augustine and Aquinas (On Rebuke and Grace; Summa I q.23)} defeats \emph{Reformed (Westminster XVII; Dort, Fifth Head)}.
\end{itemize}
\end{itemize}
\par\noindent\textbf{What this rests on.} The reasons state 5 defeats and 26 absences of defeat, each a cell of the defeat table, \texttt{perseverance\allowbreak{}Dispute\_\allowbreak{}defeats}, computed from the two parties' premises and checked by the kernel. Whether an attack survives turns on the attacker's weakest link and on the rating of what it attacks, so a changed rating can change the verdict. The weakest links:
\begin{itemize}
\item \emph{Remonstrance (article 5, 1610)}, weakest at \emph{consensus}:
\begin{itemize}
\item John 10:27–29 — they will never perish, and no one will snatch them out of my hand --- John 10:27-29
\item an inference step --- \cite[548--549]{schaff-creeds-christendom-3-1977}; \cite[603]{arminius-works-1-1825}; \cite[Solid Declaration XI.45]{concordia-triglotta-1921}; \cite[444--445]{wesley-works-9-1827}
\end{itemize}
\item \emph{Augustine and Aquinas (On Rebuke and Grace; Summa I q.23)}, weakest at \emph{disputed}:
\begin{itemize}
\item To some whom God has regenerated, and given faith, hope and love, he does not give perseverance --- \cite[478]{npnf1-05-1887}; \cite[On Rebuke and Grace, ch. 18]{npnf1-05-1887}
\item 1 John 2:19, read of every apostate: all who fall away were never true believers (encoded as a denial, so it ranks disputed) --- \cite[479--480]{npnf1-05-1887}
\end{itemize}
\item \emph{Trent (Session VI, chapters 13 and 15, canons 16 and 23)}, weakest at \emph{disputed}:
\begin{itemize}
\item God’s promise to keep those he has begun to save holds unless they fall away by their own fault --- \cite[Trent, Session VI (1547), ch. 13]{tanner-decrees-1990}; \cite[40]{waterworth-trent-1848}
\item Justifying grace, once given, can be lost --- \cite[Trent, Session VI (1547), ch. 15]{tanner-decrees-1990}; \cite[42]{waterworth-trent-1848}; \cite[Trent, Session VI (1547), canon 23]{tanner-decrees-1990}; \cite[47]{waterworth-trent-1848}
\item an inference step --- \cite[329--330]{calvin-corinthians-1-1848}; \cite[Trent, Session VI (1547), ch. 13]{tanner-decrees-1990}; \cite[40--41]{waterworth-trent-1848}
\end{itemize}
\item \emph{Lutheran (the Formula of Concord, XI; the Smalcald Articles)}, weakest at \emph{disputed}:
\begin{itemize}
\item God’s promise to keep those he has begun to save holds unless they fall away by their own fault --- \cite[Solid Declaration XI.42]{concordia-triglotta-1921}
\item In his sin David lost the Holy Spirit and his faith --- \cite[Smalcald Articles III.iii.43]{concordia-triglotta-1921}
\end{itemize}
\item \emph{Wesley (Serious Thoughts upon the Perseverance of the Saints)}, weakest at \emph{disputed}:
\begin{itemize}
\item Hebrews 6:4–6 describes true believers who actually fall away --- \cite[449]{wesley-works-9-1827}
\item Hebrews 10:29 describes one truly sanctified by Christ’s blood who falls away --- \cite[448--449]{wesley-works-9-1827}; \cite[585]{wesley-explanatory-notes-nt}
\item 2 Peter 2:20–22 describes true believers who fall back into the world --- \cite[445--446]{wesley-works-9-1827}; \cite[625]{wesley-explanatory-notes-nt}
\item God’s promise to keep those he has begun to save holds unless they fall away by their own fault --- \cite[444--445]{wesley-works-9-1827}
\item Romans 8:29–30 is a chain from which no one drops: every one justified is glorified (encoded as a denial, so it ranks disputed) --- \cite[384]{wesley-explanatory-notes-nt}; \cite[413]{wesley-works-9-1827}
\end{itemize}
\item \emph{Reformed (Westminster XVII; Dort, Fifth Head)}, weakest at \emph{disputed}:
\begin{itemize}
\item 1 John 2:19, read of every apostate: all who fall away were never true believers --- \cite[594--595]{owen-works-11-perseverance}; \cite[191--192]{calvin-commentaries-catholic-epistles}; \cite[479--480]{npnf1-05-1887}
\item Romans 8:29–30 is a chain from which no one drops: every one justified is glorified --- \cite[582--583]{schaff-creeds-christendom-3-1977}; \cite[557]{schaff-creeds-christendom-3-1977}; \cite[384]{wesley-explanatory-notes-nt}; \cite[413]{wesley-works-9-1827}
\item In his sin David lost the Holy Spirit and his faith (encoded as a denial, so it ranks disputed) --- \cite[299--300]{calvin-psalms-2-1846}
\item Hebrews 6:4–6 describes true believers who actually fall away (encoded as a denial, so it ranks disputed) --- \cite[137--138]{calvin-hebrews-1853}
\item Hebrews 10:29 describes one truly sanctified by Christ’s blood who falls away (encoded as a denial, so it ranks disputed) --- \cite[545--546]{owen-works-23-1855}
\item 2 Peter 2:20–22 describes true believers who fall back into the world (encoded as a denial, so it ranks disputed) --- \cite[410--412]{calvin-commentaries-catholic-epistles}
\item an inference step --- \cite[444--445]{wesley-works-9-1827}; \cite[415--416]{calvin-commentary-john-1847}; \cite[Solid Declaration XI.45]{concordia-triglotta-1921}
\item an inference step --- \cite[Trent, Session VI (1547), ch. 13]{tanner-decrees-1990}; \cite[40]{waterworth-trent-1848}; \cite[Solid Declaration XI.42]{concordia-triglotta-1921}
\end{itemize}
\end{itemize}
\noindent Why \emph{Reformed (Westminster XVII; Dort, Fifth Head)} defeats \emph{Lutheran (the Formula of Concord, XI; the Smalcald Articles)}, at the level of the claims: \texttt{where\allowbreak{}They\allowbreak{}Part\allowbreak{}Over\allowbreak{}David}.
\noindent Why \emph{Reformed (Westminster XVII; Dort, Fifth Head)} defeats \emph{Wesley (Serious Thoughts upon the Perseverance of the Saints)}, at the level of the claims: \texttt{where\allowbreak{}Wesley\allowbreak{}Parts\allowbreak{}Over\allowbreak{}Romans8}.
\par\medskip\noindent\textbf{\texttt{augustinian\_\allowbreak{}defensible}}\par\nobreak\noindent
\textbf{Augustine and Aquinas can be defended}: accepted together with Trent, the Lutherans and Wesley, who agree that one truly justified can finally fall away, and each of whom answers Westminster.

{\small\begin{verbatim}
theorem augustinian_defensible : Framework.CredulouslyAccepted
    perseveranceDispute.defeats PerseveranceParty.augustinian
-- axioms: propext, Classical.choice, Quot.sound
\end{verbatim}}
\par\medskip\noindent\textbf{\texttt{augustine\allowbreak{}Answered\allowbreak{}By\allowbreak{}Westminster}}\par\nobreak\noindent
Why Augustine and Aquinas are not forced: Westminster defeats them.

\paragraph{\emph{Augustine and Aquinas (On Rebuke and Grace; Summa I q.23)} is not accepted on every resolution.}
\begin{itemize}
\item A position holding \emph{Remonstrance (article 5, 1610)}, \emph{Reformed (Westminster XVII; Dort, Fifth Head)} can be held, and it holds \emph{Reformed (Westminster XVII; Dort, Fifth Head)}; \emph{Reformed (Westminster XVII; Dort, Fifth Head)} conflicts with \emph{Augustine and Aquinas (On Rebuke and Grace; Summa I q.23)}: \emph{Reformed (Westminster XVII; Dort, Fifth Head)} defeats \emph{Augustine and Aquinas (On Rebuke and Grace; Summa I q.23)}, and \emph{Augustine and Aquinas (On Rebuke and Grace; Summa I q.23)} defeats \emph{Reformed (Westminster XVII; Dort, Fifth Head)}.
\item None of \emph{Remonstrance (article 5, 1610)}, \emph{Reformed (Westminster XVII; Dort, Fifth Head)} defeats another, and each attack on them is answered from among them:
\begin{itemize}
\item \emph{Augustine and Aquinas (On Rebuke and Grace; Summa I q.23)} defeats \emph{Reformed (Westminster XVII; Dort, Fifth Head)}, and \emph{Reformed (Westminster XVII; Dort, Fifth Head)} defeats \emph{Augustine and Aquinas (On Rebuke and Grace; Summa I q.23)}.
\item \emph{Trent (Session VI, chapters 13 and 15, canons 16 and 23)} defeats \emph{Reformed (Westminster XVII; Dort, Fifth Head)} (\texttt{where\allowbreak{}They\allowbreak{}Part\allowbreak{}Over\allowbreak{}Assurance}), and \emph{Reformed (Westminster XVII; Dort, Fifth Head)} defeats \emph{Trent (Session VI, chapters 13 and 15, canons 16 and 23)}.
\item \emph{Lutheran (the Formula of Concord, XI; the Smalcald Articles)} defeats \emph{Reformed (Westminster XVII; Dort, Fifth Head)}, and \emph{Reformed (Westminster XVII; Dort, Fifth Head)} defeats \emph{Lutheran (the Formula of Concord, XI; the Smalcald Articles)} (\texttt{where\allowbreak{}They\allowbreak{}Part\allowbreak{}Over\allowbreak{}David}).
\item \emph{Wesley (Serious Thoughts upon the Perseverance of the Saints)} defeats \emph{Reformed (Westminster XVII; Dort, Fifth Head)}, and \emph{Reformed (Westminster XVII; Dort, Fifth Head)} defeats \emph{Wesley (Serious Thoughts upon the Perseverance of the Saints)} (\texttt{where\allowbreak{}Wesley\allowbreak{}Parts\allowbreak{}Over\allowbreak{}Romans8}).
\end{itemize}
\end{itemize}
\par\noindent\textbf{What this rests on.} The reasons state 8 defeats and 8 absences of defeat, each a cell of the defeat table, \texttt{perseverance\allowbreak{}Dispute\_\allowbreak{}defeats}, computed from the two parties' premises and checked by the kernel. Whether an attack survives turns on the attacker's weakest link and on the rating of what it attacks, so a changed rating can change the verdict. The weakest links:
\begin{itemize}
\item \emph{Reformed (Westminster XVII; Dort, Fifth Head)}, weakest at \emph{disputed}:
\begin{itemize}
\item 1 John 2:19, read of every apostate: all who fall away were never true believers --- \cite[594--595]{owen-works-11-perseverance}; \cite[191--192]{calvin-commentaries-catholic-epistles}; \cite[479--480]{npnf1-05-1887}
\item Romans 8:29–30 is a chain from which no one drops: every one justified is glorified --- \cite[582--583]{schaff-creeds-christendom-3-1977}; \cite[557]{schaff-creeds-christendom-3-1977}; \cite[384]{wesley-explanatory-notes-nt}; \cite[413]{wesley-works-9-1827}
\item In his sin David lost the Holy Spirit and his faith (encoded as a denial, so it ranks disputed) --- \cite[299--300]{calvin-psalms-2-1846}
\item Hebrews 6:4–6 describes true believers who actually fall away (encoded as a denial, so it ranks disputed) --- \cite[137--138]{calvin-hebrews-1853}
\item Hebrews 10:29 describes one truly sanctified by Christ’s blood who falls away (encoded as a denial, so it ranks disputed) --- \cite[545--546]{owen-works-23-1855}
\item 2 Peter 2:20–22 describes true believers who fall back into the world (encoded as a denial, so it ranks disputed) --- \cite[410--412]{calvin-commentaries-catholic-epistles}
\item an inference step --- \cite[444--445]{wesley-works-9-1827}; \cite[415--416]{calvin-commentary-john-1847}; \cite[Solid Declaration XI.45]{concordia-triglotta-1921}
\item an inference step --- \cite[Trent, Session VI (1547), ch. 13]{tanner-decrees-1990}; \cite[40]{waterworth-trent-1848}; \cite[Solid Declaration XI.42]{concordia-triglotta-1921}
\end{itemize}
\item \emph{Augustine and Aquinas (On Rebuke and Grace; Summa I q.23)}, weakest at \emph{disputed}:
\begin{itemize}
\item To some whom God has regenerated, and given faith, hope and love, he does not give perseverance --- \cite[478]{npnf1-05-1887}; \cite[On Rebuke and Grace, ch. 18]{npnf1-05-1887}
\item 1 John 2:19, read of every apostate: all who fall away were never true believers (encoded as a denial, so it ranks disputed) --- \cite[479--480]{npnf1-05-1887}
\end{itemize}
\item \emph{Remonstrance (article 5, 1610)}, weakest at \emph{consensus}:
\begin{itemize}
\item John 10:27–29 — they will never perish, and no one will snatch them out of my hand --- John 10:27-29
\item an inference step --- \cite[548--549]{schaff-creeds-christendom-3-1977}; \cite[603]{arminius-works-1-1825}; \cite[Solid Declaration XI.45]{concordia-triglotta-1921}; \cite[444--445]{wesley-works-9-1827}
\end{itemize}
\item \emph{Trent (Session VI, chapters 13 and 15, canons 16 and 23)}, weakest at \emph{disputed}:
\begin{itemize}
\item God’s promise to keep those he has begun to save holds unless they fall away by their own fault --- \cite[Trent, Session VI (1547), ch. 13]{tanner-decrees-1990}; \cite[40]{waterworth-trent-1848}
\item Justifying grace, once given, can be lost --- \cite[Trent, Session VI (1547), ch. 15]{tanner-decrees-1990}; \cite[42]{waterworth-trent-1848}; \cite[Trent, Session VI (1547), canon 23]{tanner-decrees-1990}; \cite[47]{waterworth-trent-1848}
\item an inference step --- \cite[329--330]{calvin-corinthians-1-1848}; \cite[Trent, Session VI (1547), ch. 13]{tanner-decrees-1990}; \cite[40--41]{waterworth-trent-1848}
\end{itemize}
\item \emph{Lutheran (the Formula of Concord, XI; the Smalcald Articles)}, weakest at \emph{disputed}:
\begin{itemize}
\item God’s promise to keep those he has begun to save holds unless they fall away by their own fault --- \cite[Solid Declaration XI.42]{concordia-triglotta-1921}
\item In his sin David lost the Holy Spirit and his faith --- \cite[Smalcald Articles III.iii.43]{concordia-triglotta-1921}
\end{itemize}
\item \emph{Wesley (Serious Thoughts upon the Perseverance of the Saints)}, weakest at \emph{disputed}:
\begin{itemize}
\item Hebrews 6:4–6 describes true believers who actually fall away --- \cite[449]{wesley-works-9-1827}
\item Hebrews 10:29 describes one truly sanctified by Christ’s blood who falls away --- \cite[448--449]{wesley-works-9-1827}; \cite[585]{wesley-explanatory-notes-nt}
\item 2 Peter 2:20–22 describes true believers who fall back into the world --- \cite[445--446]{wesley-works-9-1827}; \cite[625]{wesley-explanatory-notes-nt}
\item God’s promise to keep those he has begun to save holds unless they fall away by their own fault --- \cite[444--445]{wesley-works-9-1827}
\item Romans 8:29–30 is a chain from which no one drops: every one justified is glorified (encoded as a denial, so it ranks disputed) --- \cite[384]{wesley-explanatory-notes-nt}; \cite[413]{wesley-works-9-1827}
\end{itemize}
\end{itemize}
\noindent Why \emph{Trent (Session VI, chapters 13 and 15, canons 16 and 23)} defeats \emph{Reformed (Westminster XVII; Dort, Fifth Head)}, at the level of the claims: \texttt{where\allowbreak{}They\allowbreak{}Part\allowbreak{}Over\allowbreak{}Assurance}.
\noindent Why \emph{Reformed (Westminster XVII; Dort, Fifth Head)} defeats \emph{Lutheran (the Formula of Concord, XI; the Smalcald Articles)}, at the level of the claims: \texttt{where\allowbreak{}They\allowbreak{}Part\allowbreak{}Over\allowbreak{}David}.
\noindent Why \emph{Reformed (Westminster XVII; Dort, Fifth Head)} defeats \emph{Wesley (Serious Thoughts upon the Perseverance of the Saints)}, at the level of the claims: \texttt{where\allowbreak{}Wesley\allowbreak{}Parts\allowbreak{}Over\allowbreak{}Romans8}.
\par\medskip\noindent\textbf{\texttt{augustinian\_\allowbreak{}not\_\allowbreak{}forced}}\par\nobreak\noindent
Nor must Augustine and Aquinas be accepted: the way of settling the ties that accepts Westminster cannot accept them.

{\small\begin{verbatim}
theorem augustinian_not_forced : ¬Framework.SkepticallyAccepted
    perseveranceDispute.defeats PerseveranceParty.augustinian
-- axioms: propext, Classical.choice, Quot.sound
\end{verbatim}}
\par\medskip\noindent\textbf{\texttt{trent\allowbreak{}Defended\allowbreak{}On\allowbreak{}Perseverance}}\par\nobreak\noindent
Why Trent can be defended on perseverance.

\paragraph{\emph{Trent (Session VI, chapters 13 and 15, canons 16 and 23)} is accepted on some resolution.}
\begin{itemize}
\item A position holding \emph{Remonstrance (article 5, 1610)}, \emph{Augustine and Aquinas (On Rebuke and Grace; Summa I q.23)}, \emph{Trent (Session VI, chapters 13 and 15, canons 16 and 23)}, \emph{Lutheran (the Formula of Concord, XI; the Smalcald Articles)}, \emph{Wesley (Serious Thoughts upon the Perseverance of the Saints)} can be held.
\item None of \emph{Remonstrance (article 5, 1610)}, \emph{Augustine and Aquinas (On Rebuke and Grace; Summa I q.23)}, \emph{Trent (Session VI, chapters 13 and 15, canons 16 and 23)}, \emph{Lutheran (the Formula of Concord, XI; the Smalcald Articles)}, \emph{Wesley (Serious Thoughts upon the Perseverance of the Saints)} defeats another, and each attack on them is answered from among them:
\begin{itemize}
\item \emph{Reformed (Westminster XVII; Dort, Fifth Head)} defeats \emph{Augustine and Aquinas (On Rebuke and Grace; Summa I q.23)}, and \emph{Augustine and Aquinas (On Rebuke and Grace; Summa I q.23)} defeats \emph{Reformed (Westminster XVII; Dort, Fifth Head)}.
\item \emph{Reformed (Westminster XVII; Dort, Fifth Head)} defeats \emph{Trent (Session VI, chapters 13 and 15, canons 16 and 23)}, and \emph{Augustine and Aquinas (On Rebuke and Grace; Summa I q.23)} defeats \emph{Reformed (Westminster XVII; Dort, Fifth Head)}.
\item \emph{Reformed (Westminster XVII; Dort, Fifth Head)} defeats \emph{Lutheran (the Formula of Concord, XI; the Smalcald Articles)} (\texttt{where\allowbreak{}They\allowbreak{}Part\allowbreak{}Over\allowbreak{}David}), and \emph{Augustine and Aquinas (On Rebuke and Grace; Summa I q.23)} defeats \emph{Reformed (Westminster XVII; Dort, Fifth Head)}.
\item \emph{Reformed (Westminster XVII; Dort, Fifth Head)} defeats \emph{Wesley (Serious Thoughts upon the Perseverance of the Saints)} (\texttt{where\allowbreak{}Wesley\allowbreak{}Parts\allowbreak{}Over\allowbreak{}Romans8}), and \emph{Augustine and Aquinas (On Rebuke and Grace; Summa I q.23)} defeats \emph{Reformed (Westminster XVII; Dort, Fifth Head)}.
\end{itemize}
\end{itemize}
\par\noindent\textbf{What this rests on.} The reasons state 5 defeats and 26 absences of defeat, each a cell of the defeat table, \texttt{perseverance\allowbreak{}Dispute\_\allowbreak{}defeats}, computed from the two parties' premises and checked by the kernel. Whether an attack survives turns on the attacker's weakest link and on the rating of what it attacks, so a changed rating can change the verdict. The weakest links:
\begin{itemize}
\item \emph{Remonstrance (article 5, 1610)}, weakest at \emph{consensus}:
\begin{itemize}
\item John 10:27–29 — they will never perish, and no one will snatch them out of my hand --- John 10:27-29
\item an inference step --- \cite[548--549]{schaff-creeds-christendom-3-1977}; \cite[603]{arminius-works-1-1825}; \cite[Solid Declaration XI.45]{concordia-triglotta-1921}; \cite[444--445]{wesley-works-9-1827}
\end{itemize}
\item \emph{Augustine and Aquinas (On Rebuke and Grace; Summa I q.23)}, weakest at \emph{disputed}:
\begin{itemize}
\item To some whom God has regenerated, and given faith, hope and love, he does not give perseverance --- \cite[478]{npnf1-05-1887}; \cite[On Rebuke and Grace, ch. 18]{npnf1-05-1887}
\item 1 John 2:19, read of every apostate: all who fall away were never true believers (encoded as a denial, so it ranks disputed) --- \cite[479--480]{npnf1-05-1887}
\end{itemize}
\item \emph{Trent (Session VI, chapters 13 and 15, canons 16 and 23)}, weakest at \emph{disputed}:
\begin{itemize}
\item God’s promise to keep those he has begun to save holds unless they fall away by their own fault --- \cite[Trent, Session VI (1547), ch. 13]{tanner-decrees-1990}; \cite[40]{waterworth-trent-1848}
\item Justifying grace, once given, can be lost --- \cite[Trent, Session VI (1547), ch. 15]{tanner-decrees-1990}; \cite[42]{waterworth-trent-1848}; \cite[Trent, Session VI (1547), canon 23]{tanner-decrees-1990}; \cite[47]{waterworth-trent-1848}
\item an inference step --- \cite[329--330]{calvin-corinthians-1-1848}; \cite[Trent, Session VI (1547), ch. 13]{tanner-decrees-1990}; \cite[40--41]{waterworth-trent-1848}
\end{itemize}
\item \emph{Lutheran (the Formula of Concord, XI; the Smalcald Articles)}, weakest at \emph{disputed}:
\begin{itemize}
\item God’s promise to keep those he has begun to save holds unless they fall away by their own fault --- \cite[Solid Declaration XI.42]{concordia-triglotta-1921}
\item In his sin David lost the Holy Spirit and his faith --- \cite[Smalcald Articles III.iii.43]{concordia-triglotta-1921}
\end{itemize}
\item \emph{Wesley (Serious Thoughts upon the Perseverance of the Saints)}, weakest at \emph{disputed}:
\begin{itemize}
\item Hebrews 6:4–6 describes true believers who actually fall away --- \cite[449]{wesley-works-9-1827}
\item Hebrews 10:29 describes one truly sanctified by Christ’s blood who falls away --- \cite[448--449]{wesley-works-9-1827}; \cite[585]{wesley-explanatory-notes-nt}
\item 2 Peter 2:20–22 describes true believers who fall back into the world --- \cite[445--446]{wesley-works-9-1827}; \cite[625]{wesley-explanatory-notes-nt}
\item God’s promise to keep those he has begun to save holds unless they fall away by their own fault --- \cite[444--445]{wesley-works-9-1827}
\item Romans 8:29–30 is a chain from which no one drops: every one justified is glorified (encoded as a denial, so it ranks disputed) --- \cite[384]{wesley-explanatory-notes-nt}; \cite[413]{wesley-works-9-1827}
\end{itemize}
\item \emph{Reformed (Westminster XVII; Dort, Fifth Head)}, weakest at \emph{disputed}:
\begin{itemize}
\item 1 John 2:19, read of every apostate: all who fall away were never true believers --- \cite[594--595]{owen-works-11-perseverance}; \cite[191--192]{calvin-commentaries-catholic-epistles}; \cite[479--480]{npnf1-05-1887}
\item Romans 8:29–30 is a chain from which no one drops: every one justified is glorified --- \cite[582--583]{schaff-creeds-christendom-3-1977}; \cite[557]{schaff-creeds-christendom-3-1977}; \cite[384]{wesley-explanatory-notes-nt}; \cite[413]{wesley-works-9-1827}
\item In his sin David lost the Holy Spirit and his faith (encoded as a denial, so it ranks disputed) --- \cite[299--300]{calvin-psalms-2-1846}
\item Hebrews 6:4–6 describes true believers who actually fall away (encoded as a denial, so it ranks disputed) --- \cite[137--138]{calvin-hebrews-1853}
\item Hebrews 10:29 describes one truly sanctified by Christ’s blood who falls away (encoded as a denial, so it ranks disputed) --- \cite[545--546]{owen-works-23-1855}
\item 2 Peter 2:20–22 describes true believers who fall back into the world (encoded as a denial, so it ranks disputed) --- \cite[410--412]{calvin-commentaries-catholic-epistles}
\item an inference step --- \cite[444--445]{wesley-works-9-1827}; \cite[415--416]{calvin-commentary-john-1847}; \cite[Solid Declaration XI.45]{concordia-triglotta-1921}
\item an inference step --- \cite[Trent, Session VI (1547), ch. 13]{tanner-decrees-1990}; \cite[40]{waterworth-trent-1848}; \cite[Solid Declaration XI.42]{concordia-triglotta-1921}
\end{itemize}
\end{itemize}
\noindent Why \emph{Reformed (Westminster XVII; Dort, Fifth Head)} defeats \emph{Lutheran (the Formula of Concord, XI; the Smalcald Articles)}, at the level of the claims: \texttt{where\allowbreak{}They\allowbreak{}Part\allowbreak{}Over\allowbreak{}David}.
\noindent Why \emph{Reformed (Westminster XVII; Dort, Fifth Head)} defeats \emph{Wesley (Serious Thoughts upon the Perseverance of the Saints)}, at the level of the claims: \texttt{where\allowbreak{}Wesley\allowbreak{}Parts\allowbreak{}Over\allowbreak{}Romans8}.
\par\medskip\noindent\textbf{\texttt{trent\_\allowbreak{}on\_\allowbreak{}perseverance\_\allowbreak{}defensible}}\par\nobreak\noindent
\textbf{Trent can be defended} on perseverance, in the same company.

{\small\begin{verbatim}
theorem trent_on_perseverance_defensible : Framework.CredulouslyAccepted
    perseveranceDispute.defeats PerseveranceParty.trent
-- axioms: propext, Classical.choice, Quot.sound
\end{verbatim}}
\par\medskip\noindent\textbf{\texttt{trent\allowbreak{}Answered\allowbreak{}On\allowbreak{}Perseverance}}\par\nobreak\noindent
Why Trent is not forced on perseverance: Westminster defeats it.

\paragraph{\emph{Trent (Session VI, chapters 13 and 15, canons 16 and 23)} is not accepted on every resolution.}
\begin{itemize}
\item A position holding \emph{Remonstrance (article 5, 1610)}, \emph{Reformed (Westminster XVII; Dort, Fifth Head)} can be held, and it holds \emph{Reformed (Westminster XVII; Dort, Fifth Head)}; \emph{Reformed (Westminster XVII; Dort, Fifth Head)} conflicts with \emph{Trent (Session VI, chapters 13 and 15, canons 16 and 23)}: \emph{Reformed (Westminster XVII; Dort, Fifth Head)} defeats \emph{Trent (Session VI, chapters 13 and 15, canons 16 and 23)}, and \emph{Trent (Session VI, chapters 13 and 15, canons 16 and 23)} defeats \emph{Reformed (Westminster XVII; Dort, Fifth Head)} (\texttt{where\allowbreak{}They\allowbreak{}Part\allowbreak{}Over\allowbreak{}Assurance}).
\item None of \emph{Remonstrance (article 5, 1610)}, \emph{Reformed (Westminster XVII; Dort, Fifth Head)} defeats another, and each attack on them is answered from among them:
\begin{itemize}
\item \emph{Augustine and Aquinas (On Rebuke and Grace; Summa I q.23)} defeats \emph{Reformed (Westminster XVII; Dort, Fifth Head)}, and \emph{Reformed (Westminster XVII; Dort, Fifth Head)} defeats \emph{Augustine and Aquinas (On Rebuke and Grace; Summa I q.23)}.
\item \emph{Trent (Session VI, chapters 13 and 15, canons 16 and 23)} defeats \emph{Reformed (Westminster XVII; Dort, Fifth Head)} (\texttt{where\allowbreak{}They\allowbreak{}Part\allowbreak{}Over\allowbreak{}Assurance}), and \emph{Reformed (Westminster XVII; Dort, Fifth Head)} defeats \emph{Trent (Session VI, chapters 13 and 15, canons 16 and 23)}.
\item \emph{Lutheran (the Formula of Concord, XI; the Smalcald Articles)} defeats \emph{Reformed (Westminster XVII; Dort, Fifth Head)}, and \emph{Reformed (Westminster XVII; Dort, Fifth Head)} defeats \emph{Lutheran (the Formula of Concord, XI; the Smalcald Articles)} (\texttt{where\allowbreak{}They\allowbreak{}Part\allowbreak{}Over\allowbreak{}David}).
\item \emph{Wesley (Serious Thoughts upon the Perseverance of the Saints)} defeats \emph{Reformed (Westminster XVII; Dort, Fifth Head)}, and \emph{Reformed (Westminster XVII; Dort, Fifth Head)} defeats \emph{Wesley (Serious Thoughts upon the Perseverance of the Saints)} (\texttt{where\allowbreak{}Wesley\allowbreak{}Parts\allowbreak{}Over\allowbreak{}Romans8}).
\end{itemize}
\end{itemize}
\par\noindent\textbf{What this rests on.} The reasons state 8 defeats and 8 absences of defeat, each a cell of the defeat table, \texttt{perseverance\allowbreak{}Dispute\_\allowbreak{}defeats}, computed from the two parties' premises and checked by the kernel. Whether an attack survives turns on the attacker's weakest link and on the rating of what it attacks, so a changed rating can change the verdict. The weakest links:
\begin{itemize}
\item \emph{Reformed (Westminster XVII; Dort, Fifth Head)}, weakest at \emph{disputed}:
\begin{itemize}
\item 1 John 2:19, read of every apostate: all who fall away were never true believers --- \cite[594--595]{owen-works-11-perseverance}; \cite[191--192]{calvin-commentaries-catholic-epistles}; \cite[479--480]{npnf1-05-1887}
\item Romans 8:29–30 is a chain from which no one drops: every one justified is glorified --- \cite[582--583]{schaff-creeds-christendom-3-1977}; \cite[557]{schaff-creeds-christendom-3-1977}; \cite[384]{wesley-explanatory-notes-nt}; \cite[413]{wesley-works-9-1827}
\item In his sin David lost the Holy Spirit and his faith (encoded as a denial, so it ranks disputed) --- \cite[299--300]{calvin-psalms-2-1846}
\item Hebrews 6:4–6 describes true believers who actually fall away (encoded as a denial, so it ranks disputed) --- \cite[137--138]{calvin-hebrews-1853}
\item Hebrews 10:29 describes one truly sanctified by Christ’s blood who falls away (encoded as a denial, so it ranks disputed) --- \cite[545--546]{owen-works-23-1855}
\item 2 Peter 2:20–22 describes true believers who fall back into the world (encoded as a denial, so it ranks disputed) --- \cite[410--412]{calvin-commentaries-catholic-epistles}
\item an inference step --- \cite[444--445]{wesley-works-9-1827}; \cite[415--416]{calvin-commentary-john-1847}; \cite[Solid Declaration XI.45]{concordia-triglotta-1921}
\item an inference step --- \cite[Trent, Session VI (1547), ch. 13]{tanner-decrees-1990}; \cite[40]{waterworth-trent-1848}; \cite[Solid Declaration XI.42]{concordia-triglotta-1921}
\end{itemize}
\item \emph{Trent (Session VI, chapters 13 and 15, canons 16 and 23)}, weakest at \emph{disputed}:
\begin{itemize}
\item God’s promise to keep those he has begun to save holds unless they fall away by their own fault --- \cite[Trent, Session VI (1547), ch. 13]{tanner-decrees-1990}; \cite[40]{waterworth-trent-1848}
\item Justifying grace, once given, can be lost --- \cite[Trent, Session VI (1547), ch. 15]{tanner-decrees-1990}; \cite[42]{waterworth-trent-1848}; \cite[Trent, Session VI (1547), canon 23]{tanner-decrees-1990}; \cite[47]{waterworth-trent-1848}
\item an inference step --- \cite[329--330]{calvin-corinthians-1-1848}; \cite[Trent, Session VI (1547), ch. 13]{tanner-decrees-1990}; \cite[40--41]{waterworth-trent-1848}
\end{itemize}
\item \emph{Remonstrance (article 5, 1610)}, weakest at \emph{consensus}:
\begin{itemize}
\item John 10:27–29 — they will never perish, and no one will snatch them out of my hand --- John 10:27-29
\item an inference step --- \cite[548--549]{schaff-creeds-christendom-3-1977}; \cite[603]{arminius-works-1-1825}; \cite[Solid Declaration XI.45]{concordia-triglotta-1921}; \cite[444--445]{wesley-works-9-1827}
\end{itemize}
\item \emph{Augustine and Aquinas (On Rebuke and Grace; Summa I q.23)}, weakest at \emph{disputed}:
\begin{itemize}
\item To some whom God has regenerated, and given faith, hope and love, he does not give perseverance --- \cite[478]{npnf1-05-1887}; \cite[On Rebuke and Grace, ch. 18]{npnf1-05-1887}
\item 1 John 2:19, read of every apostate: all who fall away were never true believers (encoded as a denial, so it ranks disputed) --- \cite[479--480]{npnf1-05-1887}
\end{itemize}
\item \emph{Lutheran (the Formula of Concord, XI; the Smalcald Articles)}, weakest at \emph{disputed}:
\begin{itemize}
\item God’s promise to keep those he has begun to save holds unless they fall away by their own fault --- \cite[Solid Declaration XI.42]{concordia-triglotta-1921}
\item In his sin David lost the Holy Spirit and his faith --- \cite[Smalcald Articles III.iii.43]{concordia-triglotta-1921}
\end{itemize}
\item \emph{Wesley (Serious Thoughts upon the Perseverance of the Saints)}, weakest at \emph{disputed}:
\begin{itemize}
\item Hebrews 6:4–6 describes true believers who actually fall away --- \cite[449]{wesley-works-9-1827}
\item Hebrews 10:29 describes one truly sanctified by Christ’s blood who falls away --- \cite[448--449]{wesley-works-9-1827}; \cite[585]{wesley-explanatory-notes-nt}
\item 2 Peter 2:20–22 describes true believers who fall back into the world --- \cite[445--446]{wesley-works-9-1827}; \cite[625]{wesley-explanatory-notes-nt}
\item God’s promise to keep those he has begun to save holds unless they fall away by their own fault --- \cite[444--445]{wesley-works-9-1827}
\item Romans 8:29–30 is a chain from which no one drops: every one justified is glorified (encoded as a denial, so it ranks disputed) --- \cite[384]{wesley-explanatory-notes-nt}; \cite[413]{wesley-works-9-1827}
\end{itemize}
\end{itemize}
\noindent Why \emph{Trent (Session VI, chapters 13 and 15, canons 16 and 23)} defeats \emph{Reformed (Westminster XVII; Dort, Fifth Head)}, at the level of the claims: \texttt{where\allowbreak{}They\allowbreak{}Part\allowbreak{}Over\allowbreak{}Assurance}.
\noindent Why \emph{Reformed (Westminster XVII; Dort, Fifth Head)} defeats \emph{Lutheran (the Formula of Concord, XI; the Smalcald Articles)}, at the level of the claims: \texttt{where\allowbreak{}They\allowbreak{}Part\allowbreak{}Over\allowbreak{}David}.
\noindent Why \emph{Reformed (Westminster XVII; Dort, Fifth Head)} defeats \emph{Wesley (Serious Thoughts upon the Perseverance of the Saints)}, at the level of the claims: \texttt{where\allowbreak{}Wesley\allowbreak{}Parts\allowbreak{}Over\allowbreak{}Romans8}.
\par\medskip\noindent\textbf{\texttt{trent\_\allowbreak{}on\_\allowbreak{}perseverance\_\allowbreak{}not\_\allowbreak{}forced}}\par\nobreak\noindent
Nor must Trent be accepted, for the same reason.

{\small\begin{verbatim}
theorem trent_on_perseverance_not_forced : ¬Framework.SkepticallyAccepted
    perseveranceDispute.defeats PerseveranceParty.trent
-- axioms: propext, Classical.choice, Quot.sound
\end{verbatim}}
\par\medskip\noindent\textbf{\texttt{lutherans\allowbreak{}Defended\allowbreak{}On\allowbreak{}Perseverance}}\par\nobreak\noindent
Why the Lutherans can be defended on perseverance.

\paragraph{\emph{Lutheran (the Formula of Concord, XI; the Smalcald Articles)} is accepted on some resolution.}
\begin{itemize}
\item A position holding \emph{Remonstrance (article 5, 1610)}, \emph{Augustine and Aquinas (On Rebuke and Grace; Summa I q.23)}, \emph{Trent (Session VI, chapters 13 and 15, canons 16 and 23)}, \emph{Lutheran (the Formula of Concord, XI; the Smalcald Articles)}, \emph{Wesley (Serious Thoughts upon the Perseverance of the Saints)} can be held.
\item None of \emph{Remonstrance (article 5, 1610)}, \emph{Augustine and Aquinas (On Rebuke and Grace; Summa I q.23)}, \emph{Trent (Session VI, chapters 13 and 15, canons 16 and 23)}, \emph{Lutheran (the Formula of Concord, XI; the Smalcald Articles)}, \emph{Wesley (Serious Thoughts upon the Perseverance of the Saints)} defeats another, and each attack on them is answered from among them:
\begin{itemize}
\item \emph{Reformed (Westminster XVII; Dort, Fifth Head)} defeats \emph{Augustine and Aquinas (On Rebuke and Grace; Summa I q.23)}, and \emph{Augustine and Aquinas (On Rebuke and Grace; Summa I q.23)} defeats \emph{Reformed (Westminster XVII; Dort, Fifth Head)}.
\item \emph{Reformed (Westminster XVII; Dort, Fifth Head)} defeats \emph{Trent (Session VI, chapters 13 and 15, canons 16 and 23)}, and \emph{Augustine and Aquinas (On Rebuke and Grace; Summa I q.23)} defeats \emph{Reformed (Westminster XVII; Dort, Fifth Head)}.
\item \emph{Reformed (Westminster XVII; Dort, Fifth Head)} defeats \emph{Lutheran (the Formula of Concord, XI; the Smalcald Articles)} (\texttt{where\allowbreak{}They\allowbreak{}Part\allowbreak{}Over\allowbreak{}David}), and \emph{Augustine and Aquinas (On Rebuke and Grace; Summa I q.23)} defeats \emph{Reformed (Westminster XVII; Dort, Fifth Head)}.
\item \emph{Reformed (Westminster XVII; Dort, Fifth Head)} defeats \emph{Wesley (Serious Thoughts upon the Perseverance of the Saints)} (\texttt{where\allowbreak{}Wesley\allowbreak{}Parts\allowbreak{}Over\allowbreak{}Romans8}), and \emph{Augustine and Aquinas (On Rebuke and Grace; Summa I q.23)} defeats \emph{Reformed (Westminster XVII; Dort, Fifth Head)}.
\end{itemize}
\end{itemize}
\par\noindent\textbf{What this rests on.} The reasons state 5 defeats and 26 absences of defeat, each a cell of the defeat table, \texttt{perseverance\allowbreak{}Dispute\_\allowbreak{}defeats}, computed from the two parties' premises and checked by the kernel. Whether an attack survives turns on the attacker's weakest link and on the rating of what it attacks, so a changed rating can change the verdict. The weakest links:
\begin{itemize}
\item \emph{Remonstrance (article 5, 1610)}, weakest at \emph{consensus}:
\begin{itemize}
\item John 10:27–29 — they will never perish, and no one will snatch them out of my hand --- John 10:27-29
\item an inference step --- \cite[548--549]{schaff-creeds-christendom-3-1977}; \cite[603]{arminius-works-1-1825}; \cite[Solid Declaration XI.45]{concordia-triglotta-1921}; \cite[444--445]{wesley-works-9-1827}
\end{itemize}
\item \emph{Augustine and Aquinas (On Rebuke and Grace; Summa I q.23)}, weakest at \emph{disputed}:
\begin{itemize}
\item To some whom God has regenerated, and given faith, hope and love, he does not give perseverance --- \cite[478]{npnf1-05-1887}; \cite[On Rebuke and Grace, ch. 18]{npnf1-05-1887}
\item 1 John 2:19, read of every apostate: all who fall away were never true believers (encoded as a denial, so it ranks disputed) --- \cite[479--480]{npnf1-05-1887}
\end{itemize}
\item \emph{Trent (Session VI, chapters 13 and 15, canons 16 and 23)}, weakest at \emph{disputed}:
\begin{itemize}
\item God’s promise to keep those he has begun to save holds unless they fall away by their own fault --- \cite[Trent, Session VI (1547), ch. 13]{tanner-decrees-1990}; \cite[40]{waterworth-trent-1848}
\item Justifying grace, once given, can be lost --- \cite[Trent, Session VI (1547), ch. 15]{tanner-decrees-1990}; \cite[42]{waterworth-trent-1848}; \cite[Trent, Session VI (1547), canon 23]{tanner-decrees-1990}; \cite[47]{waterworth-trent-1848}
\item an inference step --- \cite[329--330]{calvin-corinthians-1-1848}; \cite[Trent, Session VI (1547), ch. 13]{tanner-decrees-1990}; \cite[40--41]{waterworth-trent-1848}
\end{itemize}
\item \emph{Lutheran (the Formula of Concord, XI; the Smalcald Articles)}, weakest at \emph{disputed}:
\begin{itemize}
\item God’s promise to keep those he has begun to save holds unless they fall away by their own fault --- \cite[Solid Declaration XI.42]{concordia-triglotta-1921}
\item In his sin David lost the Holy Spirit and his faith --- \cite[Smalcald Articles III.iii.43]{concordia-triglotta-1921}
\end{itemize}
\item \emph{Wesley (Serious Thoughts upon the Perseverance of the Saints)}, weakest at \emph{disputed}:
\begin{itemize}
\item Hebrews 6:4–6 describes true believers who actually fall away --- \cite[449]{wesley-works-9-1827}
\item Hebrews 10:29 describes one truly sanctified by Christ’s blood who falls away --- \cite[448--449]{wesley-works-9-1827}; \cite[585]{wesley-explanatory-notes-nt}
\item 2 Peter 2:20–22 describes true believers who fall back into the world --- \cite[445--446]{wesley-works-9-1827}; \cite[625]{wesley-explanatory-notes-nt}
\item God’s promise to keep those he has begun to save holds unless they fall away by their own fault --- \cite[444--445]{wesley-works-9-1827}
\item Romans 8:29–30 is a chain from which no one drops: every one justified is glorified (encoded as a denial, so it ranks disputed) --- \cite[384]{wesley-explanatory-notes-nt}; \cite[413]{wesley-works-9-1827}
\end{itemize}
\item \emph{Reformed (Westminster XVII; Dort, Fifth Head)}, weakest at \emph{disputed}:
\begin{itemize}
\item 1 John 2:19, read of every apostate: all who fall away were never true believers --- \cite[594--595]{owen-works-11-perseverance}; \cite[191--192]{calvin-commentaries-catholic-epistles}; \cite[479--480]{npnf1-05-1887}
\item Romans 8:29–30 is a chain from which no one drops: every one justified is glorified --- \cite[582--583]{schaff-creeds-christendom-3-1977}; \cite[557]{schaff-creeds-christendom-3-1977}; \cite[384]{wesley-explanatory-notes-nt}; \cite[413]{wesley-works-9-1827}
\item In his sin David lost the Holy Spirit and his faith (encoded as a denial, so it ranks disputed) --- \cite[299--300]{calvin-psalms-2-1846}
\item Hebrews 6:4–6 describes true believers who actually fall away (encoded as a denial, so it ranks disputed) --- \cite[137--138]{calvin-hebrews-1853}
\item Hebrews 10:29 describes one truly sanctified by Christ’s blood who falls away (encoded as a denial, so it ranks disputed) --- \cite[545--546]{owen-works-23-1855}
\item 2 Peter 2:20–22 describes true believers who fall back into the world (encoded as a denial, so it ranks disputed) --- \cite[410--412]{calvin-commentaries-catholic-epistles}
\item an inference step --- \cite[444--445]{wesley-works-9-1827}; \cite[415--416]{calvin-commentary-john-1847}; \cite[Solid Declaration XI.45]{concordia-triglotta-1921}
\item an inference step --- \cite[Trent, Session VI (1547), ch. 13]{tanner-decrees-1990}; \cite[40]{waterworth-trent-1848}; \cite[Solid Declaration XI.42]{concordia-triglotta-1921}
\end{itemize}
\end{itemize}
\noindent Why \emph{Reformed (Westminster XVII; Dort, Fifth Head)} defeats \emph{Lutheran (the Formula of Concord, XI; the Smalcald Articles)}, at the level of the claims: \texttt{where\allowbreak{}They\allowbreak{}Part\allowbreak{}Over\allowbreak{}David}.
\noindent Why \emph{Reformed (Westminster XVII; Dort, Fifth Head)} defeats \emph{Wesley (Serious Thoughts upon the Perseverance of the Saints)}, at the level of the claims: \texttt{where\allowbreak{}Wesley\allowbreak{}Parts\allowbreak{}Over\allowbreak{}Romans8}.
\par\medskip\noindent\textbf{\texttt{lutheran\_\allowbreak{}on\_\allowbreak{}perseverance\_\allowbreak{}defensible}}\par\nobreak\noindent
\textbf{The Lutherans can be defended} on perseverance, in the same company.

{\small\begin{verbatim}
theorem lutheran_on_perseverance_defensible : Framework.CredulouslyAccepted
    perseveranceDispute.defeats PerseveranceParty.lutheran
-- axioms: propext, Classical.choice, Quot.sound
\end{verbatim}}
\par\medskip\noindent\textbf{\texttt{lutherans\allowbreak{}Answered\allowbreak{}On\allowbreak{}Perseverance}}\par\nobreak\noindent
Why the Lutherans are not forced on perseverance: Westminster defeats them.

\paragraph{\emph{Lutheran (the Formula of Concord, XI; the Smalcald Articles)} is not accepted on every resolution.}
\begin{itemize}
\item A position holding \emph{Remonstrance (article 5, 1610)}, \emph{Reformed (Westminster XVII; Dort, Fifth Head)} can be held, and it holds \emph{Reformed (Westminster XVII; Dort, Fifth Head)}; \emph{Reformed (Westminster XVII; Dort, Fifth Head)} conflicts with \emph{Lutheran (the Formula of Concord, XI; the Smalcald Articles)}: \emph{Reformed (Westminster XVII; Dort, Fifth Head)} defeats \emph{Lutheran (the Formula of Concord, XI; the Smalcald Articles)} (\texttt{where\allowbreak{}They\allowbreak{}Part\allowbreak{}Over\allowbreak{}David}), and \emph{Lutheran (the Formula of Concord, XI; the Smalcald Articles)} defeats \emph{Reformed (Westminster XVII; Dort, Fifth Head)}.
\item None of \emph{Remonstrance (article 5, 1610)}, \emph{Reformed (Westminster XVII; Dort, Fifth Head)} defeats another, and each attack on them is answered from among them:
\begin{itemize}
\item \emph{Augustine and Aquinas (On Rebuke and Grace; Summa I q.23)} defeats \emph{Reformed (Westminster XVII; Dort, Fifth Head)}, and \emph{Reformed (Westminster XVII; Dort, Fifth Head)} defeats \emph{Augustine and Aquinas (On Rebuke and Grace; Summa I q.23)}.
\item \emph{Trent (Session VI, chapters 13 and 15, canons 16 and 23)} defeats \emph{Reformed (Westminster XVII; Dort, Fifth Head)} (\texttt{where\allowbreak{}They\allowbreak{}Part\allowbreak{}Over\allowbreak{}Assurance}), and \emph{Reformed (Westminster XVII; Dort, Fifth Head)} defeats \emph{Trent (Session VI, chapters 13 and 15, canons 16 and 23)}.
\item \emph{Lutheran (the Formula of Concord, XI; the Smalcald Articles)} defeats \emph{Reformed (Westminster XVII; Dort, Fifth Head)}, and \emph{Reformed (Westminster XVII; Dort, Fifth Head)} defeats \emph{Lutheran (the Formula of Concord, XI; the Smalcald Articles)} (\texttt{where\allowbreak{}They\allowbreak{}Part\allowbreak{}Over\allowbreak{}David}).
\item \emph{Wesley (Serious Thoughts upon the Perseverance of the Saints)} defeats \emph{Reformed (Westminster XVII; Dort, Fifth Head)}, and \emph{Reformed (Westminster XVII; Dort, Fifth Head)} defeats \emph{Wesley (Serious Thoughts upon the Perseverance of the Saints)} (\texttt{where\allowbreak{}Wesley\allowbreak{}Parts\allowbreak{}Over\allowbreak{}Romans8}).
\end{itemize}
\end{itemize}
\par\noindent\textbf{What this rests on.} The reasons state 8 defeats and 8 absences of defeat, each a cell of the defeat table, \texttt{perseverance\allowbreak{}Dispute\_\allowbreak{}defeats}, computed from the two parties' premises and checked by the kernel. Whether an attack survives turns on the attacker's weakest link and on the rating of what it attacks, so a changed rating can change the verdict. The weakest links:
\begin{itemize}
\item \emph{Reformed (Westminster XVII; Dort, Fifth Head)}, weakest at \emph{disputed}:
\begin{itemize}
\item 1 John 2:19, read of every apostate: all who fall away were never true believers --- \cite[594--595]{owen-works-11-perseverance}; \cite[191--192]{calvin-commentaries-catholic-epistles}; \cite[479--480]{npnf1-05-1887}
\item Romans 8:29–30 is a chain from which no one drops: every one justified is glorified --- \cite[582--583]{schaff-creeds-christendom-3-1977}; \cite[557]{schaff-creeds-christendom-3-1977}; \cite[384]{wesley-explanatory-notes-nt}; \cite[413]{wesley-works-9-1827}
\item In his sin David lost the Holy Spirit and his faith (encoded as a denial, so it ranks disputed) --- \cite[299--300]{calvin-psalms-2-1846}
\item Hebrews 6:4–6 describes true believers who actually fall away (encoded as a denial, so it ranks disputed) --- \cite[137--138]{calvin-hebrews-1853}
\item Hebrews 10:29 describes one truly sanctified by Christ’s blood who falls away (encoded as a denial, so it ranks disputed) --- \cite[545--546]{owen-works-23-1855}
\item 2 Peter 2:20–22 describes true believers who fall back into the world (encoded as a denial, so it ranks disputed) --- \cite[410--412]{calvin-commentaries-catholic-epistles}
\item an inference step --- \cite[444--445]{wesley-works-9-1827}; \cite[415--416]{calvin-commentary-john-1847}; \cite[Solid Declaration XI.45]{concordia-triglotta-1921}
\item an inference step --- \cite[Trent, Session VI (1547), ch. 13]{tanner-decrees-1990}; \cite[40]{waterworth-trent-1848}; \cite[Solid Declaration XI.42]{concordia-triglotta-1921}
\end{itemize}
\item \emph{Lutheran (the Formula of Concord, XI; the Smalcald Articles)}, weakest at \emph{disputed}:
\begin{itemize}
\item God’s promise to keep those he has begun to save holds unless they fall away by their own fault --- \cite[Solid Declaration XI.42]{concordia-triglotta-1921}
\item In his sin David lost the Holy Spirit and his faith --- \cite[Smalcald Articles III.iii.43]{concordia-triglotta-1921}
\end{itemize}
\item \emph{Remonstrance (article 5, 1610)}, weakest at \emph{consensus}:
\begin{itemize}
\item John 10:27–29 — they will never perish, and no one will snatch them out of my hand --- John 10:27-29
\item an inference step --- \cite[548--549]{schaff-creeds-christendom-3-1977}; \cite[603]{arminius-works-1-1825}; \cite[Solid Declaration XI.45]{concordia-triglotta-1921}; \cite[444--445]{wesley-works-9-1827}
\end{itemize}
\item \emph{Augustine and Aquinas (On Rebuke and Grace; Summa I q.23)}, weakest at \emph{disputed}:
\begin{itemize}
\item To some whom God has regenerated, and given faith, hope and love, he does not give perseverance --- \cite[478]{npnf1-05-1887}; \cite[On Rebuke and Grace, ch. 18]{npnf1-05-1887}
\item 1 John 2:19, read of every apostate: all who fall away were never true believers (encoded as a denial, so it ranks disputed) --- \cite[479--480]{npnf1-05-1887}
\end{itemize}
\item \emph{Trent (Session VI, chapters 13 and 15, canons 16 and 23)}, weakest at \emph{disputed}:
\begin{itemize}
\item God’s promise to keep those he has begun to save holds unless they fall away by their own fault --- \cite[Trent, Session VI (1547), ch. 13]{tanner-decrees-1990}; \cite[40]{waterworth-trent-1848}
\item Justifying grace, once given, can be lost --- \cite[Trent, Session VI (1547), ch. 15]{tanner-decrees-1990}; \cite[42]{waterworth-trent-1848}; \cite[Trent, Session VI (1547), canon 23]{tanner-decrees-1990}; \cite[47]{waterworth-trent-1848}
\item an inference step --- \cite[329--330]{calvin-corinthians-1-1848}; \cite[Trent, Session VI (1547), ch. 13]{tanner-decrees-1990}; \cite[40--41]{waterworth-trent-1848}
\end{itemize}
\item \emph{Wesley (Serious Thoughts upon the Perseverance of the Saints)}, weakest at \emph{disputed}:
\begin{itemize}
\item Hebrews 6:4–6 describes true believers who actually fall away --- \cite[449]{wesley-works-9-1827}
\item Hebrews 10:29 describes one truly sanctified by Christ’s blood who falls away --- \cite[448--449]{wesley-works-9-1827}; \cite[585]{wesley-explanatory-notes-nt}
\item 2 Peter 2:20–22 describes true believers who fall back into the world --- \cite[445--446]{wesley-works-9-1827}; \cite[625]{wesley-explanatory-notes-nt}
\item God’s promise to keep those he has begun to save holds unless they fall away by their own fault --- \cite[444--445]{wesley-works-9-1827}
\item Romans 8:29–30 is a chain from which no one drops: every one justified is glorified (encoded as a denial, so it ranks disputed) --- \cite[384]{wesley-explanatory-notes-nt}; \cite[413]{wesley-works-9-1827}
\end{itemize}
\end{itemize}
\noindent Why \emph{Reformed (Westminster XVII; Dort, Fifth Head)} defeats \emph{Lutheran (the Formula of Concord, XI; the Smalcald Articles)}, at the level of the claims: \texttt{where\allowbreak{}They\allowbreak{}Part\allowbreak{}Over\allowbreak{}David}.
\noindent Why \emph{Trent (Session VI, chapters 13 and 15, canons 16 and 23)} defeats \emph{Reformed (Westminster XVII; Dort, Fifth Head)}, at the level of the claims: \texttt{where\allowbreak{}They\allowbreak{}Part\allowbreak{}Over\allowbreak{}Assurance}.
\noindent Why \emph{Reformed (Westminster XVII; Dort, Fifth Head)} defeats \emph{Wesley (Serious Thoughts upon the Perseverance of the Saints)}, at the level of the claims: \texttt{where\allowbreak{}Wesley\allowbreak{}Parts\allowbreak{}Over\allowbreak{}Romans8}.
\par\medskip\noindent\textbf{\texttt{lutheran\_\allowbreak{}on\_\allowbreak{}perseverance\_\allowbreak{}not\_\allowbreak{}forced}}\par\nobreak\noindent
Nor must the Lutherans be accepted, for the same reason.

{\small\begin{verbatim}
theorem lutheran_on_perseverance_not_forced : ¬Framework.SkepticallyAccepted
    perseveranceDispute.defeats PerseveranceParty.lutheran
-- axioms: propext, Classical.choice, Quot.sound
\end{verbatim}}
\par\medskip\noindent\textbf{\texttt{wesley\allowbreak{}Defended\allowbreak{}On\allowbreak{}Perseverance}}\par\nobreak\noindent
Why Wesley can be defended on perseverance.

\paragraph{\emph{Wesley (Serious Thoughts upon the Perseverance of the Saints)} is accepted on some resolution.}
\begin{itemize}
\item A position holding \emph{Remonstrance (article 5, 1610)}, \emph{Augustine and Aquinas (On Rebuke and Grace; Summa I q.23)}, \emph{Trent (Session VI, chapters 13 and 15, canons 16 and 23)}, \emph{Lutheran (the Formula of Concord, XI; the Smalcald Articles)}, \emph{Wesley (Serious Thoughts upon the Perseverance of the Saints)} can be held.
\item None of \emph{Remonstrance (article 5, 1610)}, \emph{Augustine and Aquinas (On Rebuke and Grace; Summa I q.23)}, \emph{Trent (Session VI, chapters 13 and 15, canons 16 and 23)}, \emph{Lutheran (the Formula of Concord, XI; the Smalcald Articles)}, \emph{Wesley (Serious Thoughts upon the Perseverance of the Saints)} defeats another, and each attack on them is answered from among them:
\begin{itemize}
\item \emph{Reformed (Westminster XVII; Dort, Fifth Head)} defeats \emph{Augustine and Aquinas (On Rebuke and Grace; Summa I q.23)}, and \emph{Augustine and Aquinas (On Rebuke and Grace; Summa I q.23)} defeats \emph{Reformed (Westminster XVII; Dort, Fifth Head)}.
\item \emph{Reformed (Westminster XVII; Dort, Fifth Head)} defeats \emph{Trent (Session VI, chapters 13 and 15, canons 16 and 23)}, and \emph{Augustine and Aquinas (On Rebuke and Grace; Summa I q.23)} defeats \emph{Reformed (Westminster XVII; Dort, Fifth Head)}.
\item \emph{Reformed (Westminster XVII; Dort, Fifth Head)} defeats \emph{Lutheran (the Formula of Concord, XI; the Smalcald Articles)} (\texttt{where\allowbreak{}They\allowbreak{}Part\allowbreak{}Over\allowbreak{}David}), and \emph{Augustine and Aquinas (On Rebuke and Grace; Summa I q.23)} defeats \emph{Reformed (Westminster XVII; Dort, Fifth Head)}.
\item \emph{Reformed (Westminster XVII; Dort, Fifth Head)} defeats \emph{Wesley (Serious Thoughts upon the Perseverance of the Saints)} (\texttt{where\allowbreak{}Wesley\allowbreak{}Parts\allowbreak{}Over\allowbreak{}Romans8}), and \emph{Augustine and Aquinas (On Rebuke and Grace; Summa I q.23)} defeats \emph{Reformed (Westminster XVII; Dort, Fifth Head)}.
\end{itemize}
\end{itemize}
\par\noindent\textbf{What this rests on.} The reasons state 5 defeats and 26 absences of defeat, each a cell of the defeat table, \texttt{perseverance\allowbreak{}Dispute\_\allowbreak{}defeats}, computed from the two parties' premises and checked by the kernel. Whether an attack survives turns on the attacker's weakest link and on the rating of what it attacks, so a changed rating can change the verdict. The weakest links:
\begin{itemize}
\item \emph{Remonstrance (article 5, 1610)}, weakest at \emph{consensus}:
\begin{itemize}
\item John 10:27–29 — they will never perish, and no one will snatch them out of my hand --- John 10:27-29
\item an inference step --- \cite[548--549]{schaff-creeds-christendom-3-1977}; \cite[603]{arminius-works-1-1825}; \cite[Solid Declaration XI.45]{concordia-triglotta-1921}; \cite[444--445]{wesley-works-9-1827}
\end{itemize}
\item \emph{Augustine and Aquinas (On Rebuke and Grace; Summa I q.23)}, weakest at \emph{disputed}:
\begin{itemize}
\item To some whom God has regenerated, and given faith, hope and love, he does not give perseverance --- \cite[478]{npnf1-05-1887}; \cite[On Rebuke and Grace, ch. 18]{npnf1-05-1887}
\item 1 John 2:19, read of every apostate: all who fall away were never true believers (encoded as a denial, so it ranks disputed) --- \cite[479--480]{npnf1-05-1887}
\end{itemize}
\item \emph{Trent (Session VI, chapters 13 and 15, canons 16 and 23)}, weakest at \emph{disputed}:
\begin{itemize}
\item God’s promise to keep those he has begun to save holds unless they fall away by their own fault --- \cite[Trent, Session VI (1547), ch. 13]{tanner-decrees-1990}; \cite[40]{waterworth-trent-1848}
\item Justifying grace, once given, can be lost --- \cite[Trent, Session VI (1547), ch. 15]{tanner-decrees-1990}; \cite[42]{waterworth-trent-1848}; \cite[Trent, Session VI (1547), canon 23]{tanner-decrees-1990}; \cite[47]{waterworth-trent-1848}
\item an inference step --- \cite[329--330]{calvin-corinthians-1-1848}; \cite[Trent, Session VI (1547), ch. 13]{tanner-decrees-1990}; \cite[40--41]{waterworth-trent-1848}
\end{itemize}
\item \emph{Lutheran (the Formula of Concord, XI; the Smalcald Articles)}, weakest at \emph{disputed}:
\begin{itemize}
\item God’s promise to keep those he has begun to save holds unless they fall away by their own fault --- \cite[Solid Declaration XI.42]{concordia-triglotta-1921}
\item In his sin David lost the Holy Spirit and his faith --- \cite[Smalcald Articles III.iii.43]{concordia-triglotta-1921}
\end{itemize}
\item \emph{Wesley (Serious Thoughts upon the Perseverance of the Saints)}, weakest at \emph{disputed}:
\begin{itemize}
\item Hebrews 6:4–6 describes true believers who actually fall away --- \cite[449]{wesley-works-9-1827}
\item Hebrews 10:29 describes one truly sanctified by Christ’s blood who falls away --- \cite[448--449]{wesley-works-9-1827}; \cite[585]{wesley-explanatory-notes-nt}
\item 2 Peter 2:20–22 describes true believers who fall back into the world --- \cite[445--446]{wesley-works-9-1827}; \cite[625]{wesley-explanatory-notes-nt}
\item God’s promise to keep those he has begun to save holds unless they fall away by their own fault --- \cite[444--445]{wesley-works-9-1827}
\item Romans 8:29–30 is a chain from which no one drops: every one justified is glorified (encoded as a denial, so it ranks disputed) --- \cite[384]{wesley-explanatory-notes-nt}; \cite[413]{wesley-works-9-1827}
\end{itemize}
\item \emph{Reformed (Westminster XVII; Dort, Fifth Head)}, weakest at \emph{disputed}:
\begin{itemize}
\item 1 John 2:19, read of every apostate: all who fall away were never true believers --- \cite[594--595]{owen-works-11-perseverance}; \cite[191--192]{calvin-commentaries-catholic-epistles}; \cite[479--480]{npnf1-05-1887}
\item Romans 8:29–30 is a chain from which no one drops: every one justified is glorified --- \cite[582--583]{schaff-creeds-christendom-3-1977}; \cite[557]{schaff-creeds-christendom-3-1977}; \cite[384]{wesley-explanatory-notes-nt}; \cite[413]{wesley-works-9-1827}
\item In his sin David lost the Holy Spirit and his faith (encoded as a denial, so it ranks disputed) --- \cite[299--300]{calvin-psalms-2-1846}
\item Hebrews 6:4–6 describes true believers who actually fall away (encoded as a denial, so it ranks disputed) --- \cite[137--138]{calvin-hebrews-1853}
\item Hebrews 10:29 describes one truly sanctified by Christ’s blood who falls away (encoded as a denial, so it ranks disputed) --- \cite[545--546]{owen-works-23-1855}
\item 2 Peter 2:20–22 describes true believers who fall back into the world (encoded as a denial, so it ranks disputed) --- \cite[410--412]{calvin-commentaries-catholic-epistles}
\item an inference step --- \cite[444--445]{wesley-works-9-1827}; \cite[415--416]{calvin-commentary-john-1847}; \cite[Solid Declaration XI.45]{concordia-triglotta-1921}
\item an inference step --- \cite[Trent, Session VI (1547), ch. 13]{tanner-decrees-1990}; \cite[40]{waterworth-trent-1848}; \cite[Solid Declaration XI.42]{concordia-triglotta-1921}
\end{itemize}
\end{itemize}
\noindent Why \emph{Reformed (Westminster XVII; Dort, Fifth Head)} defeats \emph{Lutheran (the Formula of Concord, XI; the Smalcald Articles)}, at the level of the claims: \texttt{where\allowbreak{}They\allowbreak{}Part\allowbreak{}Over\allowbreak{}David}.
\noindent Why \emph{Reformed (Westminster XVII; Dort, Fifth Head)} defeats \emph{Wesley (Serious Thoughts upon the Perseverance of the Saints)}, at the level of the claims: \texttt{where\allowbreak{}Wesley\allowbreak{}Parts\allowbreak{}Over\allowbreak{}Romans8}.
\par\medskip\noindent\textbf{\texttt{wesley\_\allowbreak{}on\_\allowbreak{}perseverance\_\allowbreak{}defensible}}\par\nobreak\noindent
\textbf{Wesley can be defended} on perseverance, in the same company.

{\small\begin{verbatim}
theorem wesley_on_perseverance_defensible : Framework.CredulouslyAccepted
    perseveranceDispute.defeats PerseveranceParty.wesleyan
-- axioms: propext, Classical.choice, Quot.sound
\end{verbatim}}
\par\medskip\noindent\textbf{\texttt{wesley\allowbreak{}Answered\allowbreak{}On\allowbreak{}Perseverance}}\par\nobreak\noindent
Why Wesley is not forced on perseverance: Westminster defeats him.

\paragraph{\emph{Wesley (Serious Thoughts upon the Perseverance of the Saints)} is not accepted on every resolution.}
\begin{itemize}
\item A position holding \emph{Remonstrance (article 5, 1610)}, \emph{Reformed (Westminster XVII; Dort, Fifth Head)} can be held, and it holds \emph{Reformed (Westminster XVII; Dort, Fifth Head)}; \emph{Reformed (Westminster XVII; Dort, Fifth Head)} conflicts with \emph{Wesley (Serious Thoughts upon the Perseverance of the Saints)}: \emph{Reformed (Westminster XVII; Dort, Fifth Head)} defeats \emph{Wesley (Serious Thoughts upon the Perseverance of the Saints)} (\texttt{where\allowbreak{}Wesley\allowbreak{}Parts\allowbreak{}Over\allowbreak{}Romans8}), and \emph{Wesley (Serious Thoughts upon the Perseverance of the Saints)} defeats \emph{Reformed (Westminster XVII; Dort, Fifth Head)}.
\item None of \emph{Remonstrance (article 5, 1610)}, \emph{Reformed (Westminster XVII; Dort, Fifth Head)} defeats another, and each attack on them is answered from among them:
\begin{itemize}
\item \emph{Augustine and Aquinas (On Rebuke and Grace; Summa I q.23)} defeats \emph{Reformed (Westminster XVII; Dort, Fifth Head)}, and \emph{Reformed (Westminster XVII; Dort, Fifth Head)} defeats \emph{Augustine and Aquinas (On Rebuke and Grace; Summa I q.23)}.
\item \emph{Trent (Session VI, chapters 13 and 15, canons 16 and 23)} defeats \emph{Reformed (Westminster XVII; Dort, Fifth Head)} (\texttt{where\allowbreak{}They\allowbreak{}Part\allowbreak{}Over\allowbreak{}Assurance}), and \emph{Reformed (Westminster XVII; Dort, Fifth Head)} defeats \emph{Trent (Session VI, chapters 13 and 15, canons 16 and 23)}.
\item \emph{Lutheran (the Formula of Concord, XI; the Smalcald Articles)} defeats \emph{Reformed (Westminster XVII; Dort, Fifth Head)}, and \emph{Reformed (Westminster XVII; Dort, Fifth Head)} defeats \emph{Lutheran (the Formula of Concord, XI; the Smalcald Articles)} (\texttt{where\allowbreak{}They\allowbreak{}Part\allowbreak{}Over\allowbreak{}David}).
\item \emph{Wesley (Serious Thoughts upon the Perseverance of the Saints)} defeats \emph{Reformed (Westminster XVII; Dort, Fifth Head)}, and \emph{Reformed (Westminster XVII; Dort, Fifth Head)} defeats \emph{Wesley (Serious Thoughts upon the Perseverance of the Saints)} (\texttt{where\allowbreak{}Wesley\allowbreak{}Parts\allowbreak{}Over\allowbreak{}Romans8}).
\end{itemize}
\end{itemize}
\par\noindent\textbf{What this rests on.} The reasons state 8 defeats and 8 absences of defeat, each a cell of the defeat table, \texttt{perseverance\allowbreak{}Dispute\_\allowbreak{}defeats}, computed from the two parties' premises and checked by the kernel. Whether an attack survives turns on the attacker's weakest link and on the rating of what it attacks, so a changed rating can change the verdict. The weakest links:
\begin{itemize}
\item \emph{Reformed (Westminster XVII; Dort, Fifth Head)}, weakest at \emph{disputed}:
\begin{itemize}
\item 1 John 2:19, read of every apostate: all who fall away were never true believers --- \cite[594--595]{owen-works-11-perseverance}; \cite[191--192]{calvin-commentaries-catholic-epistles}; \cite[479--480]{npnf1-05-1887}
\item Romans 8:29–30 is a chain from which no one drops: every one justified is glorified --- \cite[582--583]{schaff-creeds-christendom-3-1977}; \cite[557]{schaff-creeds-christendom-3-1977}; \cite[384]{wesley-explanatory-notes-nt}; \cite[413]{wesley-works-9-1827}
\item In his sin David lost the Holy Spirit and his faith (encoded as a denial, so it ranks disputed) --- \cite[299--300]{calvin-psalms-2-1846}
\item Hebrews 6:4–6 describes true believers who actually fall away (encoded as a denial, so it ranks disputed) --- \cite[137--138]{calvin-hebrews-1853}
\item Hebrews 10:29 describes one truly sanctified by Christ’s blood who falls away (encoded as a denial, so it ranks disputed) --- \cite[545--546]{owen-works-23-1855}
\item 2 Peter 2:20–22 describes true believers who fall back into the world (encoded as a denial, so it ranks disputed) --- \cite[410--412]{calvin-commentaries-catholic-epistles}
\item an inference step --- \cite[444--445]{wesley-works-9-1827}; \cite[415--416]{calvin-commentary-john-1847}; \cite[Solid Declaration XI.45]{concordia-triglotta-1921}
\item an inference step --- \cite[Trent, Session VI (1547), ch. 13]{tanner-decrees-1990}; \cite[40]{waterworth-trent-1848}; \cite[Solid Declaration XI.42]{concordia-triglotta-1921}
\end{itemize}
\item \emph{Wesley (Serious Thoughts upon the Perseverance of the Saints)}, weakest at \emph{disputed}:
\begin{itemize}
\item Hebrews 6:4–6 describes true believers who actually fall away --- \cite[449]{wesley-works-9-1827}
\item Hebrews 10:29 describes one truly sanctified by Christ’s blood who falls away --- \cite[448--449]{wesley-works-9-1827}; \cite[585]{wesley-explanatory-notes-nt}
\item 2 Peter 2:20–22 describes true believers who fall back into the world --- \cite[445--446]{wesley-works-9-1827}; \cite[625]{wesley-explanatory-notes-nt}
\item God’s promise to keep those he has begun to save holds unless they fall away by their own fault --- \cite[444--445]{wesley-works-9-1827}
\item Romans 8:29–30 is a chain from which no one drops: every one justified is glorified (encoded as a denial, so it ranks disputed) --- \cite[384]{wesley-explanatory-notes-nt}; \cite[413]{wesley-works-9-1827}
\end{itemize}
\item \emph{Remonstrance (article 5, 1610)}, weakest at \emph{consensus}:
\begin{itemize}
\item John 10:27–29 — they will never perish, and no one will snatch them out of my hand --- John 10:27-29
\item an inference step --- \cite[548--549]{schaff-creeds-christendom-3-1977}; \cite[603]{arminius-works-1-1825}; \cite[Solid Declaration XI.45]{concordia-triglotta-1921}; \cite[444--445]{wesley-works-9-1827}
\end{itemize}
\item \emph{Augustine and Aquinas (On Rebuke and Grace; Summa I q.23)}, weakest at \emph{disputed}:
\begin{itemize}
\item To some whom God has regenerated, and given faith, hope and love, he does not give perseverance --- \cite[478]{npnf1-05-1887}; \cite[On Rebuke and Grace, ch. 18]{npnf1-05-1887}
\item 1 John 2:19, read of every apostate: all who fall away were never true believers (encoded as a denial, so it ranks disputed) --- \cite[479--480]{npnf1-05-1887}
\end{itemize}
\item \emph{Trent (Session VI, chapters 13 and 15, canons 16 and 23)}, weakest at \emph{disputed}:
\begin{itemize}
\item God’s promise to keep those he has begun to save holds unless they fall away by their own fault --- \cite[Trent, Session VI (1547), ch. 13]{tanner-decrees-1990}; \cite[40]{waterworth-trent-1848}
\item Justifying grace, once given, can be lost --- \cite[Trent, Session VI (1547), ch. 15]{tanner-decrees-1990}; \cite[42]{waterworth-trent-1848}; \cite[Trent, Session VI (1547), canon 23]{tanner-decrees-1990}; \cite[47]{waterworth-trent-1848}
\item an inference step --- \cite[329--330]{calvin-corinthians-1-1848}; \cite[Trent, Session VI (1547), ch. 13]{tanner-decrees-1990}; \cite[40--41]{waterworth-trent-1848}
\end{itemize}
\item \emph{Lutheran (the Formula of Concord, XI; the Smalcald Articles)}, weakest at \emph{disputed}:
\begin{itemize}
\item God’s promise to keep those he has begun to save holds unless they fall away by their own fault --- \cite[Solid Declaration XI.42]{concordia-triglotta-1921}
\item In his sin David lost the Holy Spirit and his faith --- \cite[Smalcald Articles III.iii.43]{concordia-triglotta-1921}
\end{itemize}
\end{itemize}
\noindent Why \emph{Reformed (Westminster XVII; Dort, Fifth Head)} defeats \emph{Wesley (Serious Thoughts upon the Perseverance of the Saints)}, at the level of the claims: \texttt{where\allowbreak{}Wesley\allowbreak{}Parts\allowbreak{}Over\allowbreak{}Romans8}.
\noindent Why \emph{Trent (Session VI, chapters 13 and 15, canons 16 and 23)} defeats \emph{Reformed (Westminster XVII; Dort, Fifth Head)}, at the level of the claims: \texttt{where\allowbreak{}They\allowbreak{}Part\allowbreak{}Over\allowbreak{}Assurance}.
\noindent Why \emph{Reformed (Westminster XVII; Dort, Fifth Head)} defeats \emph{Lutheran (the Formula of Concord, XI; the Smalcald Articles)}, at the level of the claims: \texttt{where\allowbreak{}They\allowbreak{}Part\allowbreak{}Over\allowbreak{}David}.
\par\medskip\noindent\textbf{\texttt{wesley\_\allowbreak{}on\_\allowbreak{}perseverance\_\allowbreak{}not\_\allowbreak{}forced}}\par\nobreak\noindent
Nor must Wesley be accepted, for the same reason.

{\small\begin{verbatim}
theorem wesley_on_perseverance_not_forced : ¬Framework.SkepticallyAccepted
    perseveranceDispute.defeats PerseveranceParty.wesleyan
-- axioms: propext, Classical.choice, Quot.sound
\end{verbatim}}
\subsection*{Appendix: the claims and their sources}
\addcontentsline{toc}{subsection}{Appendix: the claims and their sources}
The claims, in groups: the texts every party reads; the readings of them that are in dispute; the claims each rival position adds; the three parts of the conclusion; and, last, the claims of the chapter on the perseverance of the saints.

The conclusion has three parts so that a rival can hold one part and deny another, and both rivals do. Everyone in this dispute says that salvation is by grace and not earned. Only the Reformed say that grace brings about the consent it asks for, and that in being made alive the will contributes nothing.

\subsubsection*{Citation and classification for every atom}
Every claim has a source, and the build fails if one does not. The three positions agree about the texts and disagree about what they mean, so each position’s list of sources is a common list with the entries it cites differently replaced. The Reformed position cites its own account of Acts 7:51, and the Remonstrant position cites Arminius for faith as God’s gift.

Dort and the Remonstrance are cited from Schaff’s \emph{Creeds of Christendom}, which prints both. Trent is cited by session and canon from Tanner’s \emph{Decrees}, the critical edition, and quoted in Waterworth’s translation of 1848, at its page. Arminius is cited from Nichols’s translation of his works, and Aquinas from the English Dominicans’ translation of the \emph{Summa}.

Every claim and every step carries one of four ratings: \emph{consensus}, accepted across traditions; \emph{well supported}, though not universally held; \emph{plausible}; and \emph{disputed}, actively contested. A step is rated \emph{disputed} when a reader the library cites grants its grounds and denies its conclusion.

\subsubsection*{Ratings for inference steps}
A step’s rating is its own, cited to whoever grants its grounds and denies its conclusion, or, where no one was found who does, to whoever grants both.

\subsubsection*{Two readings of the ability God gives}
Augustine distinguishes two kinds of help. “The aid without which a thing does not come to pass is one thing, and the aid by which a thing comes to pass is another.” Adam was given the first, “not that by which it should be brought about that he should persevere, but that without which he could not of free will persevere”; the saints are given the second, “so that by means of this gift they cannot help persevering” (\emph{On Rebuke and Grace}, ch. 34). The third way’s claim, that God gives the ability to believe as Adam was made able, can be read as either.

\subsubsection*{The claims, numbered}
\begin{longtable}{@{}r >{\raggedright\arraybackslash}p{0.30\textwidth} l >{\raggedright\arraybackslash}p{0.19\textwidth} >{\raggedright\arraybackslash}p{0.21\textwidth}@{}}
\toprule
 & claim & kind & held & source \\
\midrule
\endhead
\pvartarget{TestimonyArguments.SolaGratia}{1}$P_{1}$ & Ephesians 2:1–5 — dead in trespasses, God made us alive together with Christ\textsuperscript{\dag} & textual & Christian, hist.-gramm., consensus & Eph 2:1-5 \\
\pvartarget{TestimonyArguments.SolaGratia}{2}$P_{2}$ & Ephesians 2:8–9 — by grace you have been saved through faith, and this is not your own doing; it is the gift of God\textsuperscript{\dag} & textual & Christian, hist.-gramm., consensus & Eph 2:8-9 \\
\pvartarget{TestimonyArguments.SolaGratia}{3}$P_{3}$ & Romans 9:16 — it depends not on human will or exertion, but on God, who has mercy\textsuperscript{\dag} & textual & Christian, hist.-gramm., consensus & Rom 9:16 \\
\pvartarget{TestimonyArguments.SolaGratia}{4}$P_{4}$ & Romans 11:6 — if it is by grace, it is no longer on the basis of works\textsuperscript{\dag} & textual & Christian, hist.-gramm., consensus & Rom 11:6 \\
\pvartarget{TestimonyArguments.SolaGratia}{5}$P_{5}$ & Acts 7:51 — you always resist the Holy Spirit\textsuperscript{\dag} & textual & Christian, hist.-gramm., consensus & Acts 7:51 \\
\pvartarget{TestimonyArguments.SolaGratia}{6}$P_{6}$ & Ephesians 4:30 — do not grieve the Holy Spirit of God, by whom you were sealed\textsuperscript{\dag} & textual & Christian, hist.-gramm., consensus & Eph 4:30 \\
\pvartarget{TestimonyArguments.SolaGratia}{7}$P_{7}$ & 1 Thessalonians 5:19 — do not quench the Spirit\textsuperscript{\dag} & textual & Christian, hist.-gramm., consensus & 1 Thess 5:19 \\
\pvartarget{TestimonyArguments.SolaGratia}{8}$P_{8}$ & Matthew 12:31–32 — the blasphemy against the Spirit will not be forgiven\textsuperscript{\dag} & textual & Christian, hist.-gramm., consensus & Matt 12:31-32 \\
\pvartarget{TestimonyArguments.SolaGratia}{9}$P_{9}$ & Genesis 2:16–17 — of the tree of the knowledge of good and evil you shall not eat\textsuperscript{\dag} & textual & Christian, hist.-gramm., consensus & Gen 2:16-17 \\
\pvartarget{TestimonyArguments.SolaGratia}{10}$P_{10}$ & The Holy Spirit’s work is resisted, by unbelievers and by believers & theological & Christian, hist.-gramm., consensus & \cite[593]{schaff-creeds-christendom-3-1977}; \cite[Solid Declaration II.60]{concordia-triglotta-1921} \\
\pvartarget{TestimonyArguments.SolaGratia}{11}$P_{11}$ & Those who obstinately resist the Spirit can grieve and lose him & interpretive & Christian, hist.-gramm., disputed & \cite[Solid Declaration II.83]{concordia-triglotta-1921}; \cite[Solid Declaration XI.42]{concordia-triglotta-1921} \\
\pvartarget{TestimonyArguments.SolaGratia}{12}$P_{12}$ & The blasphemy against the Spirit is impenitence to the end: grace refused till death & interpretive & Christian, hist.-gramm., disputed & \cite[325]{npnf1-06-1888}; \cite[II-II q.14 a.1]{aquinas-summa-1920}; \cite[74]{calvin-harmony-evangelists-2-1845} \\
\pvartarget{TestimonyArguments.SolaGratia}{13}$P_{13}$ & The grace that saves can be finally refused & theological & Arminian, disputed & \cite[Five Arminian Articles (1610), art. 4]{schaff-creeds-christendom-3-1977}; \cite[Solid Declaration XI.41]{concordia-triglotta-1921}; \cite[397]{wesley-works-9-1827} \\
\pvartarget{TestimonyArguments.SolaGratia}{14}$P_{14}$ & God gives the ability to believe, as Adam was made able to stand or fall & theological & Christian, hist.-gramm., well supported & \cite[484]{npnf1-05-1887}; \cite[Sermon 85, III.4]{wesley-sermons-1872} \\
\pvartarget{TestimonyArguments.SolaGratia}{15}$P_{15}$ & Grace is received by faith alone & theological & Christian, hist.-gramm., disputed & \cite[Augsburg Confession IV.1–3]{concordia-triglotta-1921}; \cite[Sermon 1, III.1]{wesley-sermons-1872}; \cite[Trent, Session VI (1547), canon 9]{tanner-decrees-1990} \\
\pvartarget{TestimonyArguments.SolaGratia}{16}$P_{16}$ & In converting grace God is the sole mover: the mind is moved and does not move & theological & Roman Catholic, disputed & \cite[I-II q.111 a.2]{aquinas-summa-1920}; \cite[ch. IV]{garrigou-lagrange-grace-1952} \\
\pvartarget{TestimonyArguments.SolaGratia}{17}$P_{17}$ & God moves the will intending that it attain grace & theological & Roman Catholic, well supported & \cite[I-II q.112 a.3]{aquinas-summa-1920}; \cite{pohle-molinism-1911} \\
\pvartarget{TestimonyArguments.SolaGratia}{18}$P_{18}$ & Under operating grace the will is not merely passive: it freely consents & theological & Roman Catholic, well supported & \cite[ch. IV]{garrigou-lagrange-grace-1952}; \cite[Trent, Session VI (1547), canon 4]{tanner-decrees-1990}; \cite[45]{waterworth-trent-1848} \\
\pvartarget{TestimonyArguments.SolaGratia}{19}$P_{19}$ & Grace truly sufficient, giving the power to do good, is resisted & theological & Roman Catholic, disputed & \cite[ch. VI]{garrigou-lagrange-grace-1952} \\
\pvartarget{TestimonyArguments.SolaGratia}{20}$P_{20}$ & Justifying grace, once given, can be lost & theological & Roman Catholic, disputed & \cite[Trent, Session VI (1547), ch. 15]{tanner-decrees-1990}; \cite[42]{waterworth-trent-1848}; \cite[Trent, Session VI (1547), canon 23]{tanner-decrees-1990}; \cite[47]{waterworth-trent-1848} \\
\pvartarget{TestimonyArguments.SolaGratia}{21}$P_{21}$ & Zechariah 1:3 (“return to me, and I will return to you”) with Lamentations 5:21 (“restore us to yourself, O Lord”) & textual & Christian, hist.-gramm., consensus & Zech 1:3; Lam 5:21; \cite[Trent, Session VI (1547), ch. 5]{tanner-decrees-1990}; \cite[33]{waterworth-trent-1848} \\
\pvartarget{TestimonyArguments.SolaGratia}{22}$P_{22}$ & In Ephesians 2:8, “this” (τοῦτο) points to the whole of being saved by grace through faith & linguistic & Christian, hist.-gramm., well supported & \cite[343]{hoehner-ephesians-2002}; \cite[ad loc. Eph 2:8]{robertson-word-pictures-4-1931} \\
\pvartarget{TestimonyArguments.SolaGratia}{23}$P_{23}$ & Romans 9:16 is about God’s mercy in saving, not only in choosing a people for history & interpretive & Christian, hist.-gramm., well supported & \cite[527--565]{arminius-works-3-1853}; \cite[351]{abasciano-corporate-election-romans-9-2006} \\
\pvartarget{TestimonyArguments.SolaGratia}{24}$P_{24}$ & Faith itself is God’s gift & theological & Christian, hist.-gramm., well supported & \cite[Canons of Dort, III/IV, art. 14]{schaff-creeds-christendom-3-1977}; \cite[681]{arminius-works-1-1825} \\
\pvartarget{TestimonyArguments.SolaGratia}{25}$P_{25}$ & What Stephen’s hearers resisted in Acts 7:51 was the grace that saves & interpretive & Arminian, disputed & \cite[Five Arminian Articles (1610), art. 4]{schaff-creeds-christendom-3-1977} \\
\pvartarget{TestimonyArguments.SolaGratia}{26}$P_{26}$ & The free will, moved by God, cooperates by assenting to God’s call & theological & Roman Catholic, disputed & \cite[Trent, Session VI (1547), canon 4]{tanner-decrees-1990}; \cite[45]{waterworth-trent-1848} \\
\pvartarget{TestimonyArguments.SolaGratia}{27}$P_{27}$ & The free will, moved by God, can refuse its assent if it wishes & theological & Roman Catholic, disputed & \cite[Trent, Session VI (1547), canon 4]{tanner-decrees-1990}; \cite[45]{waterworth-trent-1848} \\
\pvartarget{TestimonyArguments.SolaGratia}{28}$P_{28}$ & Conversion is God’s work alone: God regenerates, and the will does nothing to it & theological & Reformed Protestant, disputed & \cite[Canons of Dort, III/IV, arts. 10–12]{schaff-creeds-christendom-3-1977} \\
\pvartarget{TestimonyArguments.SolaGratia}{29}$P_{29}$ & Salvation is by grace and not earned: grace goes before every saving act & theological & Christian, hist.-gramm., consensus & \cite[Trent, Session VI (1547), chs. 5, 8]{tanner-decrees-1990}; \cite[32]{waterworth-trent-1848}; \cite[36]{waterworth-trent-1848}; \cite[Five Arminian Articles (1610), arts. 3–4]{schaff-creeds-christendom-3-1977}; \cite[Canons of Dort, III/IV, art. 10]{schaff-creeds-christendom-3-1977} \\
\pvartarget{TestimonyArguments.SolaGratia}{30}$P_{30}$ & The grace that saves brings about the consent it seeks & theological & Reformed Protestant, disputed & \cite[Canons of Dort, III/IV, art. 12]{schaff-creeds-christendom-3-1977} \\
\pvartarget{TestimonyArguments.SolaGratia}{31}$P_{31}$ & In being made alive the will is passive, and contributes nothing to it & theological & Reformed Protestant, disputed & \cite[X.2]{westminster-confession-1647}; \cite[Canons of Dort, III/IV, art. 12]{schaff-creeds-christendom-3-1977} \\
\pvartarget{TestimonyArguments.SolaGratia}{32}$P_{32}$ & John 10:27–29 — they will never perish, and no one will snatch them out of my hand\textsuperscript{\dag} & textual & Christian, hist.-gramm., consensus & John 10:27-29 \\
\pvartarget{TestimonyArguments.SolaGratia}{33}$P_{33}$ & Philippians 1:6 — he who began a good work in you will bring it to completion\textsuperscript{\dag} & textual & Christian, hist.-gramm., consensus & Phil 1:6 \\
\pvartarget{TestimonyArguments.SolaGratia}{34}$P_{34}$ & Hebrews 6:4–6 — it is impossible to restore again to repentance those once enlightened who have fallen away\textsuperscript{\dag} & textual & Christian, hist.-gramm., consensus & Heb 6:4-6 \\
\pvartarget{TestimonyArguments.SolaGratia}{35}$P_{35}$ & 1 John 2:19 — they went out from us, but they were not of us\textsuperscript{\dag} & textual & Christian, hist.-gramm., consensus & 1 John 2:19 \\
\pvartarget{TestimonyArguments.SolaGratia}{36}$P_{36}$ & No power from outside can snatch true believers out of Christ’s hand & theological & Christian, hist.-gramm., consensus & \cite[548]{schaff-creeds-christendom-3-1977}; \cite[Solid Declaration XI.45]{concordia-triglotta-1921}; \cite[444--445]{wesley-works-9-1827} \\
\pvartarget{TestimonyArguments.SolaGratia}{37}$P_{37}$ & To some whom God has regenerated, and given faith, hope and love, he does not give perseverance & theological & Christian, hist.-gramm., disputed & \cite[478]{npnf1-05-1887}; \cite[On Rebuke and Grace, ch. 18]{npnf1-05-1887} \\
\pvartarget{TestimonyArguments.SolaGratia}{38}$P_{38}$ & God’s promise to keep those he has begun to save holds unless they fall away by their own fault & theological & Roman Catholic, disputed & \cite[Trent, Session VI (1547), ch. 13]{tanner-decrees-1990}; \cite[40]{waterworth-trent-1848} \\
\pvartarget{TestimonyArguments.SolaGratia}{39}$P_{39}$ & Those who are Christ’s sheep do hear and follow, because God keeps them in faith & theological & Reformed Protestant, disputed & \cite[593]{schaff-creeds-christendom-3-1977}; \cite[415--416]{calvin-commentary-john-1847} \\
\pvartarget{TestimonyArguments.SolaGratia}{40}$P_{40}$ & Hebrews 6:4–6 describes true believers who actually fall away & interpretive & Arminian, disputed & \cite[449]{wesley-works-9-1827} \\
\pvartarget{TestimonyArguments.SolaGratia}{41}$P_{41}$ & Those God has chosen to save cannot finally fall away from grace & theological & Reformed Protestant, well supported & \cite[XVII.1]{westminster-confession-1647}; \cite[636]{schaff-creeds-christendom-3-1977}; \cite[593]{schaff-creeds-christendom-3-1977} \\
\pvartarget{TestimonyArguments.SolaGratia}{42}$P_{42}$ & Everyone who is truly justified is among the elect & theological & Reformed Protestant, disputed & \cite[X.4]{westminster-confession-1647}; \cite[137--138]{calvin-hebrews-1853} \\
\pvartarget{TestimonyArguments.SolaGratia}{43}$P_{43}$ & One who is truly justified can fall away from grace finally, and perish & theological & Roman Catholic, disputed & \cite[Trent, Session VI (1547), canon 23]{tanner-decrees-1990}; \cite[47]{waterworth-trent-1848}; \cite[439]{wesley-works-9-1827} \\
\pvartarget{TestimonyArguments.SolaGratia}{44}$P_{44}$ & Romans 8:29–30 — those whom he justified he also glorified\textsuperscript{\dag} & textual & Christian, hist.-gramm., consensus & Rom 8:29-30 \\
\pvartarget{TestimonyArguments.SolaGratia}{45}$P_{45}$ & Hebrews 10:26–29 — the one who has profaned the blood of the covenant by which he was sanctified\textsuperscript{\dag} & textual & Christian, hist.-gramm., consensus & Heb 10:26-29 \\
\pvartarget{TestimonyArguments.SolaGratia}{46}$P_{46}$ & 2 Peter 2:20–22 — having escaped the defilements of the world, they are again entangled in them\textsuperscript{\dag} & textual & Christian, hist.-gramm., consensus & 2 Pet 2:20-22 \\
\pvartarget{TestimonyArguments.SolaGratia}{47}$P_{47}$ & Psalm 51:11 — take not your Holy Spirit from me\textsuperscript{\dag} & textual & Christian, hist.-gramm., consensus & Ps 51:11 \\
\pvartarget{TestimonyArguments.SolaGratia}{48}$P_{48}$ & Romans 8:16 — the Spirit himself bears witness with our spirit\textsuperscript{\dag} & textual & Christian, hist.-gramm., consensus & Rom 8:16 \\
\pvartarget{TestimonyArguments.SolaGratia}{49}$P_{49}$ & 1 Corinthians 10:12 — let anyone who thinks that he stands take heed lest he fall\textsuperscript{\dag} & textual & Christian, hist.-gramm., consensus & 1 Cor 10:12 \\
\pvartarget{TestimonyArguments.SolaGratia}{50}$P_{50}$ & Hebrews 10:29 describes one truly sanctified by Christ’s blood who falls away & interpretive & Arminian, disputed & \cite[448--449]{wesley-works-9-1827}; \cite[585]{wesley-explanatory-notes-nt} \\
\pvartarget{TestimonyArguments.SolaGratia}{51}$P_{51}$ & 2 Peter 2:20–22 describes true believers who fall back into the world & interpretive & Arminian, disputed & \cite[445--446]{wesley-works-9-1827}; \cite[625]{wesley-explanatory-notes-nt} \\
\pvartarget{TestimonyArguments.SolaGratia}{52}$P_{52}$ & In his sin David lost the Holy Spirit and his faith & interpretive & Christian, hist.-gramm., disputed & \cite[Smalcald Articles III.iii.43]{concordia-triglotta-1921} \\
\pvartarget{TestimonyArguments.SolaGratia}{53}$P_{53}$ & A believer can know, without a special revelation, that he will persevere & theological & Roman Catholic, disputed & \cite[Trent, Session VI (1547), canon 16]{tanner-decrees-1990}; \cite[46]{waterworth-trent-1848}; \cite[Trent, Session VI (1547), ch. 12]{tanner-decrees-1990}; \cite[39--40]{waterworth-trent-1848} \\
\pvartarget{TestimonyArguments.SolaGratia}{54}$P_{54}$ & Romans 8:29–30 is a chain from which no one drops: every one justified is glorified & interpretive & Reformed Protestant, disputed & \cite[582--583]{schaff-creeds-christendom-3-1977}; \cite[557]{schaff-creeds-christendom-3-1977}; \cite[384]{wesley-explanatory-notes-nt}; \cite[413]{wesley-works-9-1827} \\
\pvartarget{TestimonyArguments.SolaGratia}{55}$P_{55}$ & 1 John 2:19, read of every apostate: all who fall away were never true believers & interpretive & Reformed Protestant, disputed & \cite[594--595]{owen-works-11-perseverance}; \cite[191--192]{calvin-commentaries-catholic-epistles}; \cite[479--480]{npnf1-05-1887} \\
\pvartarget{TestimonyArguments.SolaGratia}{56}$P_{56}$ & Considered in himself, apart from God’s decree, the one who perseveres could have fallen & theological & Christian, hist.-gramm., consensus & \cite[594]{schaff-creeds-christendom-3-1977}; \cite[I q.23 a.6 ad 2]{aquinas-summa-1920}; \cite[263--264]{boethius-consolation-james-1897} \\
\pvartarget{TestimonyArguments.SolaGratia}{57}$P_{57}$ & God knows all things as present to his eternity, which imposes no necessity on them & theological & Christian, hist.-gramm., well supported & \cite[257--262]{boethius-consolation-james-1897}; \cite[I q.14 a.13]{aquinas-summa-1920}; \cite[Sermon 58, §5]{wesley-sermons-1872} \\
\bottomrule
\end{longtable}

\noindent\footnotesize \textsuperscript{\dag} grounded in scripture alone: the reading is assumed, not argued for.
\normalsize
\subsubsection*{How this is checked}
\par\medskip\noindent\textbf{\texttt{grace\allowbreak{}Inference}}\par\nobreak\noindent
The rating of the step from Romans 11:6 to “salvation is by grace”: \texttt{consensus}. Trent quotes the verse to that effect, “if it be a grace, it is not now by works” (Session VI, chapter 8); the Remonstrants’ fourth article calls grace “the beginning, continuance, and accomplishment of all good”; Dort says that conversion “must be wholly ascribed to God” (III/IV, article 10).

{\small\begin{verbatim}
def graceInference : Source :=
  { primary := trent✝ "ch. 8",
    supporting :=
      [waterworth✝ 36,
        Reference.work schaffCreeds3
          (Bib.Locus.sectionRef "Five Arminian Articles (1610), art. 4"),
        Reference.work schaffCreeds3
          (Bib.Locus.sectionRef "Canons of Dort, III/IV, art. 10")],
    tradition := Tradition.christianHistoricalGrammatical,
    confidence := Confidence.consensus }
\end{verbatim}}
\par\medskip\noindent\textbf{\texttt{ephesians\allowbreak{}Inference}}\par\nobreak\noindent
\textbf{The rating of the step from “this” pointing to the whole, to “faith is God’s gift”: \texttt{disputed}.} Robertson grants the ground and, in the same note, denies the conclusion. “This” refers “not to [pistis] (feminine) or to [charis] (feminine also), but to the act of being saved by grace conditioned on faith on our part”, and of the two words he says that grace is “God’s part, ‘faith’ ours”.

Calvin is cited for less. He is Reformed, and holds elsewhere that faith is God’s gift, so he does not deny the conclusion. But he denies that this verse teaches it: Paul’s “meaning is, not that faith is the gift of God, but that salvation is given to us by God”. So a reader may grant the reading of “this” and still deny that the verse delivers faith as a gift.

{\small\begin{verbatim}
def ephesiansInference : Source :=
  {
    primary :=
      Reference.work robertsonWordPictures4
        (Bib.Locus.adLoc { book := Book.ephesians, chapter := 2, verse := 8 }),
    supporting :=
      [Reference.work calvinGalatiansEphesians (Bib.Locus.pages 228 229)],
    tradition := Tradition.christianHistoricalGrammatical,
    confidence := Confidence.disputed }
\end{verbatim}}
\par\medskip\noindent\textbf{\texttt{gift\allowbreak{}Inference}}\par\nobreak\noindent
\textbf{The rating of the step from the gift to God’s work alone: \texttt{disputed}.} Arminius grants both grounds. Faith is God’s gift: he speaks of God, “whose benefit and gift I acknowledge faith to be” (\emph{Apology}, article 4). And a man cannot turn to God by himself: he ascribes to grace “the commencement, the continuance and the consummation of all good”. He denies the conclusion: “many persons resist the Holy Spirit and reject the grace that is offered” (\emph{Declaration of Sentiments}).

{\small\begin{verbatim}
def giftInference : Source :=
  { primary := Reference.work arminiusWorks1 (Bib.Locus.page 600),
    supporting :=
      [Reference.work arminiusWorks1 (Bib.Locus.page 681),
        Reference.work schaffCreeds3
          (Bib.Locus.sectionRef "Five Arminian Articles (1610), art. 4")],
    tradition := Tradition.arminian, confidence := Confidence.disputed }
\end{verbatim}}
\par\medskip\noindent\textbf{\texttt{romans\allowbreak{}Inference}}\par\nobreak\noindent
\textbf{The rating of the step from Romans 9:16, read as being about salvation, to God’s work alone: \texttt{disputed}.} Arminius grants the ground: the chapter’s theme is “That the Gospel, not the law, is the power of God unto salvation, not to him that worketh, but to him that believeth”, and the willing and running of 9:16 is the effort by which a man “hopes that he will obtain justification and salvation apart from the mercy of God”. He denies the conclusion: the mercy goes to those “who seek righteousness and salvation by faith in Christ”, and grace can be refused.

{\small\begin{verbatim}
def romansInference : Source :=
  { primary := Reference.work arminiusWorks3 (Bib.Locus.pages 527 565),
    tradition := Tradition.arminian, confidence := Confidence.disputed }
\end{verbatim}}
\par\medskip\noindent\textbf{\texttt{cooperation\allowbreak{}Inference}}\par\nobreak\noindent
\textbf{The rating of Trent’s step from cooperation to “the will is not passive”: \texttt{disputed}, and the dissent rests on a distinction.} Canon 4 speaks of the will “disposing and preparing itself” for justification. Aquinas grants that grace works with a will that moves, which he calls cooperating grace, and still holds that in grace’s first effect “our mind is moved and does not move, but … God is the sole mover” (I-II q.111 a.2). Dort likewise grants that the renewed will “becomes itself active” (III/IV, article 12). Both place the will’s activity after a first moment in which it is only moved; the canon does not say which moment it means.

{\small\begin{verbatim}
def cooperationInference : Source :=
  {
    primary :=
      Reference.work aquinasSumma (Bib.Locus.sectionRef "I-II q.111 a.2"),
    supporting :=
      [Reference.work schaffCreeds3
          (Bib.Locus.sectionRef "Canons of Dort, III/IV, art. 12"),
        trent✝ "canon 4",
        waterworth✝ 45],
    tradition := Tradition.romanCatholic, confidence := Confidence.disputed }
\end{verbatim}}
\par\medskip\noindent\textbf{\texttt{dissent\allowbreak{}Inference}}\par\nobreak\noindent
\textbf{The rating of Trent’s step from “the will can refuse” to “grace is not effectual of itself”: \texttt{disputed}.} Aquinas grants the ground and denies the conclusion. “If God moves the will to anything, it is incompatible with this supposition, that the will be not moved thereto. But it is not impossible simply” (I-II q.10 a.4 ad 3). And “if God intends, while moving, that the one whose heart He moves should attain to grace, he will infallibly attain to it” (I-II q.112 a.3). On that reading the will can refuse, taken simply, and grace is effectual all the same.

{\small\begin{verbatim}
def dissentInference : Source :=
  {
    primary :=
      Reference.work aquinasSumma (Bib.Locus.sectionRef "I-II q.10 a.4 ad 3"),
    supporting :=
      [Reference.work aquinasSumma (Bib.Locus.sectionRef "I-II q.112 a.3"),
        trent✝ "canon 4",
        waterworth✝ 45],
    tradition := Tradition.romanCatholic, confidence := Confidence.disputed }
\end{verbatim}}
\par\medskip\noindent\textbf{\texttt{texts\allowbreak{}Of\allowbreak{}Resistance\allowbreak{}Inference}}\par\nobreak\noindent
The rating of the step from the texts of resistance to “the Spirit’s work is resisted”: \texttt{consensus}. Every party grants it. Dort says that believers “grieve the Holy Spirit” (V, article 5); Calvin, on Ephesians 4:30, that “we grieve him when we do not follow his guidance”; the Formula of Concord, citing Acts 7:51, that “those who always resist the Holy Ghost … are not converted” (Solid Declaration II.60).

{\small\begin{verbatim}
def textsOfResistanceInference : Source :=
  { primary := dort✝ 593,
    supporting :=
      [Reference.work calvinGalatiansEphesians (Bib.Locus.page 301),
        concord✝ "Solid Declaration II.60"],
    tradition := Tradition.christianHistoricalGrammatical,
    confidence := Confidence.consensus }
\end{verbatim}}
\par\medskip\noindent\textbf{\texttt{resistance\allowbreak{}Inference}}\par\nobreak\noindent
The rating of the Remonstrants’ step from resistance, read as resistance to the grace that saves, to “that grace can be finally refused”: \texttt{well\allowbreak{}Supported}. Grant that what Stephen’s hearers resisted was the grace that saves, and no reader was found who denies that it can be finally refused. The Reformed deny the ground instead, reading the resistance as made to the preached word and not to the inward work of grace. The word “resist” carries two senses here, resisting for a time and refusing to the end, and the step needs the second.

{\small\begin{verbatim}
def resistanceInference : Source :=
  {
    primary :=
      Reference.work schaffCreeds3
        (Bib.Locus.sectionRef "Five Arminian Articles (1610), art. 4"),
    supporting := [Reference.work arminiusWorks1 (Bib.Locus.page 600)],
    tradition := Tradition.arminian, confidence := Confidence.wellSupported }
\end{verbatim}}
\par\medskip\noindent\textbf{\texttt{grief\allowbreak{}Inference}}\par\nobreak\noindent
The rating of the step from grieving the Spirit, read as able to end in losing him, to “the grace that saves can be finally refused”: \texttt{well\allowbreak{}Supported}. The Formula of Concord holds both: the obstinate “do not receive, but grieve and lose, the Holy Ghost” (Solid Declaration II.83). Dort denies the ground instead: God “does not wholly withdraw the Holy Spirit from his own people” (V, article 6).

{\small\begin{verbatim}
def griefInference : Source :=
  { primary := concord✝ "Solid Declaration II.83",
    supporting :=
      [concord✝ "Solid Declaration XI.42",
        dort✝ 593],
    tradition := Tradition.christianHistoricalGrammatical,
    confidence := Confidence.wellSupported }
\end{verbatim}}
\par\medskip\noindent\textbf{\texttt{blasphemy\allowbreak{}Inference}}\par\nobreak\noindent
\textbf{The rating of the step from the unforgivable sin to final refusal: \texttt{disputed}.} Augustine grants both grounds. The sin that is never forgiven is impenitence: “This impenitence then is ‘the blasphemy of the Spirit, which shall not be forgiven’” (Sermon 71, §20). And he denies the conclusion: God is the one “whom no man’s will resists when He wills to give salvation” (\emph{On Rebuke and Grace}, ch. 43). The grace refused to the end is, for him, not the grace that saves.

{\small\begin{verbatim}
def blasphemyInference : Source :=
  { primary := Reference.work npnfAugustineSermons (Bib.Locus.page 325),
    supporting :=
      [Reference.work npnfAugustineAntiPelagian
          (Bib.Locus.sectionRef "On Rebuke and Grace, ch. 43")],
    tradition := Tradition.christianHistoricalGrammatical,
    confidence := Confidence.disputed }
\end{verbatim}}
\par\medskip\noindent\textbf{\texttt{passive\allowbreak{}Inference}}\par\nobreak\noindent
The rating of the step from God’s work alone to the will’s passivity: \texttt{well\allowbreak{}Supported}. The Lutherans, who part from Dort at the next step, grant this one: man “does nothing whatever towards it, but only suffers what God works in him” (Formula of Concord, Solid Declaration II.89). No reader was found who grants that God alone regenerates and denies that the will is passive in it. Aquinas holds that in the justification of one who has the use of reason, God’s motion “does not take place without a movement of the free-will” (I-II q.113 a.3); that is passivity in another sense, the will moved and so moving, which is not what this step denies.

{\small\begin{verbatim}
def passiveInference : Source :=
  { primary := concord✝ "Solid Declaration II.89",
    supporting :=
      [Reference.work schaffCreeds3
          (Bib.Locus.sectionRef "Canons of Dort, III/IV, art. 12")],
    tradition := Tradition.christianHistoricalGrammatical,
    confidence := Confidence.wellSupported }
\end{verbatim}}
\par\medskip\noindent\textbf{\texttt{effectual\allowbreak{}Inference}}\par\nobreak\noindent
\textbf{The rating of the step from God’s work alone to “grace cannot be finally refused”: \texttt{disputed}.} The Formula of Concord grants the ground: conversion “is not only in part, but altogether an operation, gift, and present, and work of the Holy Ghost alone” (Solid Declaration II.89). And it denies the conclusion: “God does not force man to become godly (for those who always resist the Holy Ghost … as Stephen says of the hardened Jews, Acts 7:51, are not converted)” (II.60); the perverse will “resists the Holy Ghost, who wishes to be efficacious” (XI.41).

{\small\begin{verbatim}
def effectualInference : Source :=
  { primary := concord✝ "Solid Declaration II.89",
    supporting :=
      [concord✝ "Solid Declaration II.60",
        concord✝ "Solid Declaration XI.41"],
    tradition := Tradition.christianHistoricalGrammatical,
    confidence := Confidence.disputed }
\end{verbatim}}
\par\medskip\noindent\textbf{\texttt{as\allowbreak{}Aid\allowbreak{}Without\allowbreak{}Which\allowbreak{}Source}}\par\nobreak\noindent
Who reads the ability as Adam’s kind of help, one that makes believing possible and can be refused: Wesley, for whom “a measure of free-will” is “supernaturally restored to every man”, with “power either to resist or yield to the grace of God” (\emph{Predestination Calmly Considered}, §§XLVI–XLVII).

{\small\begin{verbatim}
def asAidWithoutWhichSource : Source :=
  { primary := Reference.work wesleyWorks9 (Bib.Locus.page 397),
    supporting :=
      [Reference.work npnfAugustineAntiPelagian (Bib.Locus.page 485)],
    tradition := Tradition.arminian, confidence := Confidence.disputed }
\end{verbatim}}
\par\medskip\noindent\textbf{\texttt{as\allowbreak{}Aid\allowbreak{}By\allowbreak{}Which\allowbreak{}Source}}\par\nobreak\noindent
Who reads the ability as the saints’ kind of help, one that brings believing about: Augustine, of the grace given to the saints (\emph{On Rebuke and Grace}, ch. 34), and Dort, for which faith is God’s gift “not on account of its being offered by God to man, to be accepted or rejected at his pleasure, but because it is in reality conferred” (III/IV, article 14).

{\small\begin{verbatim}
def asAidByWhichSource : Source :=
  { primary := Reference.work npnfAugustineAntiPelagian (Bib.Locus.page 485),
    supporting := [dort✝ 591],
    tradition := Tradition.reformedProtestant,
    confidence := Confidence.disputed }
\end{verbatim}}
\paragraph*{Aquinas and the Thomists}
\leavevmode\par\nobreak
\par\medskip\noindent\textbf{\texttt{operating\allowbreak{}Grace\allowbreak{}Inference}}\par\nobreak\noindent
The rating of the step from operating grace to God’s work alone: \texttt{well\allowbreak{}Supported}. Grant that in converting grace God is the sole mover of the will, and no reader was found who denies that God alone is the one who regenerates. The Molinists deny the ground instead, placing God as sole mover only in the will’s first, indeliberate motions; for them the converting act is “due to” the will’s consent, “in no way” to “the Thomistic predetermination” (Pohle, “Molinism”).

{\small\begin{verbatim}
def operatingGraceInference : Source :=
  {
    primary :=
      Reference.work aquinasSumma (Bib.Locus.sectionRef "I-II q.111 a.2"),
    supporting :=
      [garrigou✝ "ch. IV",
        Reference.work pohleMolinism],
    tradition := Tradition.romanCatholic,
    confidence := Confidence.wellSupported }
\end{verbatim}}
\par\medskip\noindent\textbf{\texttt{efficacy\allowbreak{}Inference}}\par\nobreak\noindent
\textbf{The rating of the step from God’s intention in moving to grace effectual of itself: \texttt{disputed}.} The Molinists grant the ground and deny the conclusion. They grant that “the decree of God to bestow an efficacious grace upon a man involves with metaphysical certainty the free co-operation of the will”, and deny that the grace brings it about: “the infallible success of efficacious grace is not to be ascribed to its own intrinsic nature, but to the Divine scientia media”, God’s knowledge of what a person would freely do (Pohle, “Molinism”). Garrigou-Lagrange holds the conclusion: “grace is intrinsically efficacious” (\emph{Grace}, ch. VII).

{\small\begin{verbatim}
def efficacyInference : Source :=
  { primary := Reference.work pohleMolinism,
    supporting :=
      [Reference.work aquinasSumma (Bib.Locus.sectionRef "I-II q.112 a.3"),
        garrigou✝ "ch. VII"],
    tradition := Tradition.romanCatholic, confidence := Confidence.disputed }
\end{verbatim}}
\par\medskip\noindent\textbf{\texttt{as\allowbreak{}Agent\allowbreak{}Of\allowbreak{}New\allowbreak{}Life\allowbreak{}Source}}\par\nobreak\noindent
Who holds of the will’s part what this reading takes “not merely passive” to mean: the Molinists, for whom “it is rather the will itself which by its consent” renders grace efficacious (Pohle, “Molinism”). Rated \texttt{disputed}: no Thomist reads his own claim so, and Garrigou-Lagrange says Molina “deviates from the teaching of St. Thomas” here.

{\small\begin{verbatim}
def asAgentOfNewLifeSource : Source :=
  { primary := Reference.work pohleMolinism,
    tradition := Tradition.romanCatholic, confidence := Confidence.disputed }
\end{verbatim}}
\par\medskip\noindent\textbf{\texttt{as\allowbreak{}Disposition\allowbreak{}Source}}\par\nobreak\noindent
Who reads “not merely passive” as the will acting while God alone is the agent of the new life: Garrigou-Lagrange, for whom the will’s consent, “although as a disposition it precedes it [grace] in the genus or order of material cause, it nevertheless follows it in the genus or order of formal, efficient, and final cause” (\emph{Grace}, ch. X); and Aquinas, for whom the movement of the free will “is not the cause of grace, but the effect” (I-II q.111 a.2 ad 2). Rated \texttt{disputed}: Garrigou-Lagrange himself lists Calvin’s “merely passive” among the errors (ch. IV), and a Reformed reader may take that as a denial of the third part in Dort’s sense, not only in another.

{\small\begin{verbatim}
def asDispositionSource : Source :=
  { primary := garrigou✝ "ch. X",
    supporting :=
      [Reference.work aquinasSumma
          (Bib.Locus.sectionRef "I-II q.111 a.2 ad 2")],
    tradition := Tradition.romanCatholic, confidence := Confidence.disputed }
\end{verbatim}}
\paragraph*{Ratings for the steps on perseverance}
\leavevmode\par\nobreak
\par\medskip\noindent\textbf{\texttt{john10\allowbreak{}Inference}}\par\nobreak\noindent
\textbf{The rating of the step from John 10:27–29 to “the elect persevere”: \texttt{disputed}.} Calvin reads it so: “the salvation of all the elect is not less certain than the power of God is invincible” (\emph{John}, vol. 1, p. 416), and the Formula of Concord cites the verse for God’s purpose, “from which no one can pluck us” (Solid Declaration XI.45). Wesley grants the text and denies that it promises anyone perseverance without condition: “In the preceding text, the condition is only implied. In this it is plainly expressed. They are \emph{my sheep}, that \emph{hear my voice}, that \emph{follow me} in all holiness. And ‘if ye do these things, ye shall never fall’” (\emph{Serious Thoughts upon the Perseverance of the Saints}, §18), quoting 2 Peter 1:10.

{\small\begin{verbatim}
def john10Inference : Source :=
  { primary := Reference.work wesleyWorks9 (Bib.Locus.pages 444 445),
    supporting :=
      [Reference.work calvinJohn (Bib.Locus.pages 415 416),
        concord✝ "Solid Declaration XI.45"],
    tradition := Tradition.arminian, confidence := Confidence.disputed }
\end{verbatim}}
\par\medskip\noindent\textbf{\texttt{bridge\allowbreak{}Inference}}\par\nobreak\noindent
The rating of the step from 1 John 2:19, read of every apostate, to “everyone truly justified is among the elect”: \texttt{well\allowbreak{}Supported}. Owen draws it: “If all they who fall away totally … were never true believers … then those who are true believers cannot fall away; but the first is true, therefore the latter” (\emph{Works}, vol. 11, p. 595). No reader was found who grants the reading and denies the conclusion. The dissent is from the reading: Augustine holds that those who left “were in goodness”, “in the faith of children”, and not elect (\emph{On Rebuke and Grace}, ch. 20); Trent condemns the conclusion outright, anathema to anyone who says “that he that falls and sins was never truly justified” (Session VI, canon 23).

{\small\begin{verbatim}
def bridgeInference : Source :=
  { primary := Reference.work owenPerseverance (Bib.Locus.page 595),
    supporting :=
      [Reference.work calvinCatholicEpistles (Bib.Locus.pages 191 192),
        Reference.work calvinHebrews (Bib.Locus.pages 137 138),
        trent✝ "canon 23",
        waterworth✝ 47],
    tradition := Tradition.reformedProtestant,
    confidence := Confidence.wellSupported }
\end{verbatim}}
\par\medskip\noindent\textbf{\texttt{join\allowbreak{}Inference}}\par\nobreak\noindent
The rating of the step from “the elect persevere” and “the truly justified are elect” to “no one truly justified can finally fall away”: \texttt{consensus}. It follows from what the words mean, and no reader was found who grants both grounds and denies the conclusion. The readers who deny the conclusion, Augustine among them, deny the second ground.

{\small\begin{verbatim}
def joinInference : Source :=
  {
    primary :=
      Reference.work westminsterConfession (Bib.Locus.sectionRef "XVII.1"),
    supporting := [Reference.work schaffCreeds3 (Bib.Locus.page 636)],
    tradition := Tradition.christianHistoricalGrammatical,
    confidence := Confidence.consensus }
\end{verbatim}}
\par\medskip\noindent\textbf{\texttt{philippians\allowbreak{}Inference}}\par\nobreak\noindent
\textbf{The rating of the step from Philippians 1:6 to “no one truly justified can finally fall away”: \texttt{disputed}.} Trent quotes the verse and, in the same sentence, denies that it settles the matter: “For God, unless men be themselves wanting to His grace, as he has begun the good work, so will he perfect it” (Session VI, ch. 13). The Formula of Concord cites it while explaining why some fall away: it “is not as though God were unwilling to grant grace for perseverance to those in whom He has begun the good work” (Solid Declaration XI.42).

{\small\begin{verbatim}
def philippiansInference : Source :=
  { primary := trent✝ "ch. 13",
    supporting :=
      [waterworth✝ 40,
        concord✝ "Solid Declaration XI.42"],
    tradition := Tradition.romanCatholic, confidence := Confidence.disputed }
\end{verbatim}}
\par\medskip\noindent\textbf{\texttt{conditional\allowbreak{}Inference}}\par\nobreak\noindent
The rating of the step from a promise that holds “unless men be themselves wanting to His grace” to “the justified can finally fall away”: \texttt{well\allowbreak{}Supported}. Read as a condition that can fail, the step only says what the condition means, and Trent draws it in canon 23. Read as a condition God himself keeps in those he has begun to save, the ground is a different claim, and the step does not apply to it; the dilemma over the promise sets the two readings out.

{\small\begin{verbatim}
def conditionalInference : Source :=
  { primary := trent✝ "canon 23",
    supporting :=
      [waterworth✝ 47,
        Reference.work wesleyWorks9 (Bib.Locus.page 439)],
    tradition := Tradition.romanCatholic,
    confidence := Confidence.wellSupported }
\end{verbatim}}
\par\medskip\noindent\textbf{\texttt{hebrews\allowbreak{}Inference}}\par\nobreak\noindent
The rating of the step from Hebrews 6:4–6, read as describing true believers who fall away, to “the justified can finally fall away”: \texttt{well\allowbreak{}Supported}. Wesley draws it: “those who are holy or righteous in the judgment of God himself … may nevertheless so fall from God, as to perish everlastingly” (\emph{Serious Thoughts}, §30). No reader was found who grants the reading and denies the conclusion. The dissent is from the reading: Calvin holds that those described had only “some taste of his grace”, a knowledge “which afterwards vanishes away”.

{\small\begin{verbatim}
def hebrewsInference : Source :=
  { primary := Reference.work wesleyWorks9 (Bib.Locus.page 449),
    tradition := Tradition.arminian, confidence := Confidence.wellSupported }
\end{verbatim}}
\par\medskip\noindent\textbf{\texttt{augustine\allowbreak{}Inference}}\par\nobreak\noindent
The rating of the step from “some regenerate are not given perseverance” to “the justified can finally fall away”: \texttt{well\allowbreak{}Supported}. Augustine draws it, and Aquinas: “anyone may lose his crown by mortal sin” (I q.23 a.6 ad 1). No reader was found who grants the ground and denies the conclusion; Calvin denies the ground.

{\small\begin{verbatim}
def augustineInference : Source :=
  { primary := Reference.work npnfAugustineAntiPelagian (Bib.Locus.page 478),
    supporting :=
      [Reference.work aquinasSumma (Bib.Locus.sectionRef "I q.23 a.6 ad 1")],
    tradition := Tradition.christianHistoricalGrammatical,
    confidence := Confidence.wellSupported }
\end{verbatim}}
\par\medskip\noindent\textbf{\texttt{remonstrance\allowbreak{}Inference}}\par\nobreak\noindent
The rating of the step from John 10:28 to “no power from outside snatches believers out of Christ’s hand”: \texttt{consensus}. The Remonstrants draw it (fifth article), and every party here grants it.

{\small\begin{verbatim}
def remonstranceInference : Source :=
  { primary := Reference.work schaffCreeds3 (Bib.Locus.pages 548 549),
    supporting :=
      [Reference.work arminiusWorks1 (Bib.Locus.page 603),
        concord✝ "Solid Declaration XI.45",
        Reference.work wesleyWorks9 (Bib.Locus.pages 444 445)],
    tradition := Tradition.christianHistoricalGrammatical,
    confidence := Confidence.consensus }
\end{verbatim}}
\par\medskip\noindent\textbf{\texttt{as\allowbreak{}Condition\allowbreak{}God\allowbreak{}Keeps\allowbreak{}Source}}\par\nobreak\noindent
The promise read as a condition God himself keeps: Dort’s “God is faithful, who having conferred grace, mercifully confirms and powerfully preserves them therein, even to the end” (V, article 3).

{\small\begin{verbatim}
def asConditionGodKeepsSource : Source :=
  { primary := dort✝ 593,
    supporting := [Reference.work calvinJohn (Bib.Locus.pages 415 416)],
    tradition := Tradition.reformedProtestant,
    confidence := Confidence.disputed }
\end{verbatim}}
\par\medskip\noindent\textbf{\texttt{as\allowbreak{}Condition\allowbreak{}That\allowbreak{}Can\allowbreak{}Fail\allowbreak{}Source}}\par\nobreak\noindent
The promise read as a condition that can fail: Wesley’s “a Saint may fall away; … one who is holy or righteous in the judgment of God himself, may nevertheless so fall from God, as to perish everlastingly” (\emph{Serious Thoughts}, §5).

{\small\begin{verbatim}
def asConditionThatCanFailSource : Source :=
  { primary := Reference.work wesleyWorks9 (Bib.Locus.page 439),
    supporting := [Reference.work wesleyWorks9 (Bib.Locus.pages 444 445)],
    tradition := Tradition.arminian, confidence := Confidence.disputed }
\end{verbatim}}
\paragraph*{Ratings for the further texts on perseverance}
\leavevmode\par\nobreak
\par\medskip\noindent\textbf{\texttt{romans8\allowbreak{}Inference}}\par\nobreak\noindent
The rating of the step from Romans 8:29–30, read as a chain from which no one drops, to “no one truly justified can finally fall away”: \texttt{well\allowbreak{}Supported}. Dort draws it: God decreed to give the elect “true faith, justification, and sanctification; and having powerfully preserved them … finally to glorify them” (I, article 7), and its Rejection of Errors condemns those who teach election to justifying faith without election to salvation, “dissolving this golden chain of salvation” (I.2, in Latin). No reader was found who grants the reading and denies the conclusion. The dissent is from the reading: Wesley answers that the verse “proves no more than that all who are glorified, are pardoned and sanctified first” (\emph{Predestination Calmly Considered}, §76), and that Paul “does not affirm … that precisely the same number of men are called, justified, and glorified” (\emph{Notes}, on Romans 8:30). Augustine reads the chain of those “called according to His purpose”, and holds that others truly justified fall away (\emph{On Rebuke and Grace}, chs. 14, 23).

{\small\begin{verbatim}
def romans8Inference : Source :=
  { primary := Reference.work schaffCreeds3 (Bib.Locus.pages 582 583),
    supporting :=
      [Reference.work schaffCreeds3 (Bib.Locus.page 557),
        Reference.work wesleyWorks9 (Bib.Locus.page 413),
        Reference.work wesleyNotes (Bib.Locus.page 384),
        Reference.work npnfAugustineAntiPelagian (Bib.Locus.pages 480 481),
        Reference.work calvinRomans (Bib.Locus.pages 319 320)],
    tradition := Tradition.reformedProtestant,
    confidence := Confidence.wellSupported }
\end{verbatim}}
\par\medskip\noindent\textbf{\texttt{timeless\allowbreak{}Inference}}\par\nobreak\noindent
The rating of the step from God’s knowing all things as present to “the one who perseveres could, considered in himself, have fallen away”: \texttt{well\allowbreak{}Supported}. Boethius draws it: things foreseen, “viewed, however, in themselves, … desist not from the absolute freedom naturally theirs” (\emph{Consolation} V, prose 6, pp. 263–264); and Aquinas: “it is possible for one who is predestinated considered in himself to die in mortal sin; yet it is not possible, supposed … that he is predestinated” (I q.23 a.6, ad 2). No reader was found who grants the premise and denies the conclusion.

{\small\begin{verbatim}
def timelessInference : Source :=
  { primary := Reference.work boethiusConsolation (Bib.Locus.pages 263 264),
    supporting :=
      [Reference.work aquinasSumma (Bib.Locus.sectionRef "I q.23 a.6 ad 2"),
        Reference.work aquinasSumma (Bib.Locus.sectionRef "I q.14 a.13")],
    tradition := Tradition.christianHistoricalGrammatical,
    confidence := Confidence.wellSupported }
\end{verbatim}}
\par\medskip\noindent\textbf{\texttt{hebrews10\allowbreak{}Inference}}\par\nobreak\noindent
The rating of the step from Hebrews 10:29, read as describing one truly sanctified who falls away, to “the justified can finally fall away”: \texttt{well\allowbreak{}Supported}. Wesley draws it: “It is undeniably plain, 1. That the person mentioned here, was once sanctified by the blood of the covenant” (\emph{Serious Thoughts}, §27). No reader was found who grants the reading and denies the conclusion; the dissent, Owen’s, is from the reading.

{\small\begin{verbatim}
def hebrews10Inference : Source :=
  { primary := Reference.work wesleyWorks9 (Bib.Locus.pages 448 449),
    supporting :=
      [Reference.work wesleyWorks9 (Bib.Locus.pages 414 415),
        Reference.work wesleyNotes (Bib.Locus.page 585)],
    tradition := Tradition.arminian, confidence := Confidence.wellSupported }
\end{verbatim}}
\par\medskip\noindent\textbf{\texttt{second\allowbreak{}Peter\allowbreak{}Inference}}\par\nobreak\noindent
The rating of the step from 2 Peter 2:20–22, read as describing true believers who fall back, to “the justified can finally fall away”: \texttt{well\allowbreak{}Supported}. Wesley draws it: those who “so effectually know Christ, as by that knowledge to have escaped the pollutions of the world, may yet fall back into those pollutions, and perish everlastingly” (\emph{Serious Thoughts}, §21). No reader was found who grants the reading and denies the conclusion; the dissent, Calvin’s, is from the reading.

{\small\begin{verbatim}
def secondPeterInference : Source :=
  { primary := Reference.work wesleyWorks9 (Bib.Locus.pages 445 446),
    supporting := [Reference.work wesleyNotes (Bib.Locus.page 625)],
    tradition := Tradition.arminian, confidence := Confidence.wellSupported }
\end{verbatim}}
\par\medskip\noindent\textbf{\texttt{david\allowbreak{}Inference}}\par\nobreak\noindent
The rating of the step from David’s losing the Spirit to “justifying grace can be lost”: \texttt{well\allowbreak{}Supported}. The Smalcald Articles draw it, and Trent’s fifteenth chapter holds the conclusion: grace “is lost … by any other mortal sin whatever”. No reader was found who grants that David lost the Spirit and denies that grace can be lost; Calvin denies the ground.

{\small\begin{verbatim}
def davidInference : Source :=
  {
    primary :=
      concord✝ "Smalcald Articles III.iii.43",
    supporting :=
      [trent✝ "ch. 15",
        waterworth✝ 42],
    tradition := Tradition.christianHistoricalGrammatical,
    confidence := Confidence.wellSupported }
\end{verbatim}}
\par\medskip\noindent\textbf{\texttt{psalm\allowbreak{}Inference}}\par\nobreak\noindent
The rating of the step from Psalm 51:11, read as Calvin reads it, to “justifying grace is never lost”: \texttt{well\allowbreak{}Supported}. Calvin draws it: “it is their duty to hold fast the truth that grace is the incorruptible seed of God, which never can perish in any heart where it has been deposited” (\emph{Psalms}, vol. 2, p. 300), and Dort says God does not suffer his own “to lose the grace of adoption and forfeit the state of justification” (V, article 6). No reader was found who grants that David kept the Spirit and denies the conclusion.

{\small\begin{verbatim}
def psalmInference : Source :=
  { primary := Reference.work calvinPsalms2 (Bib.Locus.pages 299 300),
    supporting := [dort✝ 593],
    tradition := Tradition.reformedProtestant,
    confidence := Confidence.wellSupported }
\end{verbatim}}
\par\medskip\noindent\textbf{\texttt{assurance\allowbreak{}Inference}}\par\nobreak\noindent
The rating of the step from Romans 8:16 to “a believer can know, without a special revelation, that he will persevere”: \texttt{well\allowbreak{}Supported}. Dort draws it: the assurance “springs from faith in God’s promises … from the testimony of the Holy Spirit, witnessing with our spirit, that we are children and heirs of God” (V, article 10), and Westminster cites the verse for assurance (XVIII.2). No reader was found who grants the verse and denies the conclusion. Wesley grants the verse to every believer and allows “the full assurance of hope”, “Divine evidence (wrought in the soul by the same immediate inspiration of the Holy Ghost) of persevering grace, and of eternal glory”, which “you cannot expect, if you abate your diligence” (\emph{Notes}, on Hebrews 6:11). He says the same of the assurance of faith, so the “immediate inspiration” is the ordinary witness of the Spirit, and that is not a denial. Trent denies the conclusion from another text, 1 Corinthians 10:12.

{\small\begin{verbatim}
def assuranceInference : Source :=
  { primary := dort✝ 594,
    supporting :=
      [Reference.work westminsterConfession (Bib.Locus.sectionRef "XVIII.2"),
        Reference.work wesleyNotes
          (Bib.Locus.adLoc
            { book := Book.hebrews, chapter := 6, verse := 11 })],
    tradition := Tradition.reformedProtestant,
    confidence := Confidence.wellSupported }
\end{verbatim}}
\par\medskip\noindent\textbf{\texttt{caution\allowbreak{}Inference}}\par\nobreak\noindent
\textbf{The rating of the step from 1 Corinthians 10:12 to “no believer can know, without a special revelation, that he will persevere”: \texttt{disputed}.} Trent draws it, quoting the verse in its chapter on perseverance: “let those who think themselves to stand, take heed lest they fall” (Session VI, ch. 13). Calvin grants the verse and denies the conclusion: “let us observe that there are two kinds of assurance”, and Paul “does not, however, exhort them to be always anxiously in doubt as to the will of God, or to tremble from uncertainty as to their salvation” (\emph{1 Corinthians}, vol. 1, pp. 329–330).

{\small\begin{verbatim}
def cautionInference : Source :=
  { primary := Reference.work calvinFirstCorinthians (Bib.Locus.pages 329 330),
    supporting :=
      [trent✝ "ch. 13",
        Reference.work waterworthTrent (Bib.Locus.pages 40 41)],
    tradition := Tradition.reformedProtestant,
    confidence := Confidence.disputed }
\end{verbatim}}
\paragraph*{The citations}
\leavevmode\par\nobreak
\par\medskip\noindent\textbf{\texttt{base\allowbreak{}Cite}}\par\nobreak\noindent
The common list of sources: the texts, the readings, the claims of each position, and the parts of the conclusion. Each position replaces only the entries it cites differently.

{\small\begin{verbatim}
def baseCite : Claim → AtomMeta
\end{verbatim}}
\par\medskip\noindent\textbf{\texttt{reformed\allowbreak{}Cite}}\par\nobreak\noindent
The Reformed citations. Dort does not cite Acts 7:51. What it gives is an account of those who refuse the call: “It is not the fault of the gospel, nor of Christ offered therein, nor of God, who calls men by the gospel … that those who are called by the ministry of the Word refuse to come and be converted. The fault lies in themselves” (III/IV, article 9). Read with Acts 7:51, that is the Reformed answer: what Stephen’s hearers resisted was the outward call of the preached word, not the grace that saves.

{\small\begin{verbatim}
def reformedCite : Claim → AtomMeta
\end{verbatim}}
\par\medskip\noindent\textbf{\texttt{tridentine\allowbreak{}Cite}}\par\nobreak\noindent
Trent’s citations: the texts and canon 4, as \texttt{base\allowbreak{}Cite} has them; on perseverance, its fifteenth chapter, that grace “is lost … by any other mortal sin whatever”, and its sixteenth canon, against certainty of persevering.

{\small\begin{verbatim}
def tridentineCite : Claim → AtomMeta
\end{verbatim}}
\par\medskip\noindent\textbf{\texttt{remonstrant\allowbreak{}Cite}}\par\nobreak\noindent
The Remonstrant citations. That faith is God’s gift is cited to Arminius, who affirms it on grounds of his own, and not to Dort, which affirms it in a sense he rejects.

{\small\begin{verbatim}
def remonstrantCite : Claim → AtomMeta :=
  fun x =>
    match x with
    | Claim.faithIsGodsGift =>
      { label := "Faith itself is God's gift", kind := PremiseKind.theological,
        source :=
          { primary := Reference.work arminiusWorks1 (Bib.Locus.page 681),
            tradition := Tradition.arminian,
            confidence := Confidence.wellSupported } }
    | a => baseCite a
\end{verbatim}}
\par\medskip\noindent\textbf{\texttt{lutheran\allowbreak{}Cite}}\par\nobreak\noindent
The Lutheran citations: the Formula of Concord for what the Lutherans hold, in its own words.

{\small\begin{verbatim}
def lutheranCite : Claim → AtomMeta
\end{verbatim}}
\par\medskip\noindent\textbf{\texttt{wesleyan\allowbreak{}Cite}}\par\nobreak\noindent
The Wesleyan citations: Wesley for what Wesley holds.

{\small\begin{verbatim}
def wesleyanCite : Claim → AtomMeta
\end{verbatim}}
\par\medskip\noindent\textbf{\texttt{augustinian\allowbreak{}Cite}}\par\nobreak\noindent
The citations of Augustine and Aquinas: that the elect persevere, from Augustine’s “they who are not to persevere … are not elected” (\emph{On Rebuke and Grace}, ch. 16) and from Aquinas, for whom predestination “most certainly and infallibly takes effect” (\emph{Summa} I q.23 a.6). That the justified can fall finally, from Aquinas: “anyone may lose his crown by mortal sin” (ad 1).

{\small\begin{verbatim}
def augustinianCite : Claim → AtomMeta
\end{verbatim}}

\bibliographystyle{plain}
\bibliography{references}
\end{document}
